% 本文件由 LaTeX 排版 Agent 根据有效 translation.json 自动生成。
% author=Augustine of Hippo; work_id=1406; title=Contra Faustum

\chapter{\WorkTitle{Contra Faustum}}

% structure: p0001 source=h1 target=omitted reason=page-title-already-rendered-by-parent

第一卷\newline{}第二卷\newline{}第三卷\newline{}第四卷\newline{}第五卷\newline{}第六卷\newline{}第七卷\newline{}第八卷\newline{}第九卷\newline{}第十卷\newline{}第十一卷\newline{}第十二卷\newline{}第十三卷\newline{}第十四卷\newline{}第十五卷\newline{}第十六卷\newline{}第十七卷\newline{}第十八卷\newline{}第十九卷\newline{}第二十卷\newline{}第二十一卷\newline{}第二十二卷\newline{}第二十三卷\newline{}第二十四卷\newline{}第二十五卷\newline{}第二十六卷\newline{}第二十七卷\newline{}第二十八卷\newline{}第二十九卷\newline{}第三十卷\newline{}第三十一卷\newline{}第三十二卷\newline{}第三十三卷\par

\SourceRecord{Translated by Richard Stothert.；Nicene and Post-Nicene Fathers, First Series；Vol. 4.；Edited by Philip Schaff.；Buffalo, NY: Christian Literature Publishing Co.,；1887.；<http://www.newadvent.org/fathers/1406.htm>.}

% structure: p0001 source=h1 target=section reason=page-title

\section{\WorkTitle{驳福斯图斯，卷一}}

\EditorialNote{写于大约公元400年。福斯图斯无疑是奥古斯丁时代正统基督教最敏锐、最坚决、最不择手段的反对者。奥古斯丁写作这部大作反对他的契机，是福斯图斯出版了对旧约圣经以及新约圣经（凡与摩尼教谬误不符的部分）的攻击性著作。福斯图斯似乎追随了阿迪曼图斯的脚步，奥古斯丁多年前曾撰文反对阿迪曼图斯，但福斯图斯在陈述的大胆程度上大大超过了阿迪曼图斯。基督的降生成人，涉及祂由女人所生，是攻击的重点之一。他利用福音谱系记载的差异，作为拒绝全部福音为伪造的理由。他认为福音书现存的形态并非宗徒的作品，而是后来犹太化伪造者的作品。他以极度的轻蔑对待整个旧约体系，以圣祖、梅瑟、先知等的个人生活和教导为由，亵渎他们。大部分目前流行的对旧约伦理的反对意见，在奥古斯丁时代已为人熟知。奥古斯丁的回答仅部分令人满意，因为他未能完美地理解旧约与新约的关系；但在他们写作的年代，这些回答无疑非常有效。这部著作从圣经批评的角度，以及从反对摩尼教的论战角度来看，都富有启发性。—— A.H.N.}\par

\Emph{福斯图斯是何许人。福斯图斯撰写构成奥古斯丁答复之基础的论战文章的目的。奥古斯丁对此的评论。}\par

1. \Person{福斯图斯}是非洲裔，\Place{米勒乌姆}公民；他口齿伶俐、才智敏捷，但信奉了摩尼教异端骇人听闻的教义。他在我的\WorkTitle{忏悔录}中被提及，书中记述了我与他的相识。此人出版了一本书，反对真确的基督信仰和公教真理。我们得到一本，弟兄们读了之后，要求我作出答复，作为我因爱德而欠他们的服务。因此，现今我奉我们的主、救主耶稣基督之名，并依靠祂的帮助，承担这项任务，好让所有读者明白，若非主引导人的脚步，才智的敏锐和文体的雅致于人毫无益处。在神圣慈悲奥秘的公义中，天主常常帮助迟钝和软弱的人；而缺乏这种帮助，最敏锐和最雄辩的人只会更快、更任性地陷入错误。我将陈述\Person{福斯图斯}的观点，如同他本人所言，以及我的观点，如同答复他。\par

2. \Emph{\Person{福斯图斯}说：}博学的\Person{阿迪曼图斯}，自圣\Person{摩尼}以来唯一值得关注的导师，已充分揭露并彻底驳斥了犹太教和半基督信仰的错误，因此我认为，以书面形式向你们提供简短有力的答复，以应对我们对手吹毛求疵的反对意见，并非不妥；这样，当他们像狡猾的蛇的子女，试图以诡辩迷惑你们时，你们就能准备给出明智的答复。这样，他们就能被限制在主题上，而不是东拉西扯。我已尽量简明地并列我们和对手的观点，以免长篇大论、错综复杂的论述使读者困惑。\par

3. \Emph{\Person{奥古斯丁}答复：}你警告提防半基督徒，你说我们就是；但我们警告提防假基督徒，我们已经证明你们就是。半基督信仰可能不够完善，却未必虚假。那么，如果那些你们试图误导之人的信仰不够完善，与其把他们所有的夺走，岂不更好去补充他们所缺少的？正是对不够完善的基督徒，宗徒写道：\Quote{欢喜看到你们的品行}，和\Quote{你们对基督的信德有所欠缺}。宗徒所着眼的是属灵的建筑，正如他在别处所说：\Quote{你们是天主的建筑物；} \BibleRef{格前 3:9}。在这建筑中，他既找到了喜乐的理由，也找到了努力的理由。他喜乐看到一部分已经完成；而将建筑完善的需要则要求努力。我们作为不够完善的基督徒，你们追逐我们，企图以假道学来扭曲你们所谓的我们的半基督信仰；而即使那些信德欠缺到无法回答你们所有诡辩的人，也至少足够明智，知道绝不能与你们有任何瓜葛。你们寻找半基督徒来欺骗：我们希望证明你们是假基督徒，好让基督徒从驳斥你们中学到东西，让较不成熟的人学会躲避你们。你们称我们为蛇的子女？你们一定忘了你们多少次挑剔伊甸园中的禁令，并赞美蛇开了\Person{亚当}的眼。你们更有资格使用你们给我们的称号。无论你们谴责蛇还是赞美蛇，蛇都认你们为子。\par

\SourceRecord{Translated by Richard Stothert.；Nicene and Post-Nicene Fathers, First Series；Vol. 4.；Edited by Philip Schaff.；Buffalo, NY: Christian Literature Publishing Co.,；1887.；<http://www.newadvent.org/fathers/140601.htm>.}

% structure: p0001 source=h1 target=section reason=page-title

\section{\WorkTitle{驳斥福斯托斯} 第二卷}

\Emph{\Person{福斯托斯}宣称相信福音，却拒绝接受族谱表，并引多项理由，\Person{奥斯定}则试图予以驳斥。}\par

1. \Person{福斯托斯}说：我难道不相信福音吗？当然相信。那么我因此就相信基督曾受生吗？当然不信。我虽相信福音，却不必因此相信基督曾有受生。我不信这事，因为基督并未说祂曾由人所生，而福音无论在名称上或事实上，都以基督的宣讲为开始。至于族谱，作者本人也不敢称之为福音。因为他写了什么？\BibleQuote{亚巴郎之子，达味之子耶稣基督的族谱。}\BibleRef{玛 1:1} 族谱书并不是福音书。它更像是出生登记，星辰证实了此事。另一方面，\Person{马尔谷}记载了天主圣子的宣讲，没有族谱，他以最恰当的方式开头：\BibleQuote{天主之子耶稣基督的福音。}显然，族谱并非福音。\Person{玛窦}本人也说，\Person{若翰}被监禁后，\Person{耶稣}开始宣讲天国的福音；因此，此前所述的是族谱，而非福音。为什么\Person{玛窦}不以\BibleQuote{天主之子耶稣基督的福音}开头？岂不是因为他认为称族谱为福音是罪过？所以，请明白你们至今所忽略的——族谱与福音之间的区别。那么，我承认福音的真实性吗？是的；我理解的福音是基督的宣讲。关于族谱，我也有许多话要说，如果你们愿意的话。但在我看来，你们现在想知道的是否不是我所接受的福音，而是我所接受的族谱。\par

2. \Person{奥斯定}答曰：那么，针对你们自己的问题，你们首先告诉我们你们相信福音，其次你们不信基督的受生；而你们的理由是，基督的受生不在福音中。那么，当宗徒说\BibleQuote{你要记得：按照我所传的福音，达味的后裔耶稣基督从死者中复活了，}\BibleRef{弟后 2:8}你们怎么回答他？你们必定是不知道，或假装不知道什么是福音。你们使用这个词，不是按宗徒所教导的，而是按照你们自己的谬误。你们偏离了宗徒所称的福音，因为你们不相信基督是达味的后裔。这是保禄的福音，也是其他宗徒及所有忠信管家对此伟大奥迹的福音。因为保禄在别处说：\BibleQuote{所以，不拘是我，或是他们，我们都这样传，你们也这样信了。}\BibleRef{格前 15:11} 他们并非都写了福音，但他们都在宣讲福音。传福音者的名称恰当地归于叙述我们的主耶稣基督的诞生、言行、受难的人。福音这个词意思是好消息，可以用于任何好消息，但恰当地适用于救主的叙述。如果你们教导不同的东西，你们必定离开了福音。毫无疑问，那些被你们轻视为半基督信徒的婴孩，当他们听到他们的母亲爱德借宗徒的口说：\BibleQuote{即使有人传给你们另一种福音，与我们以前所传的福音不同，当受诅咒。}\BibleRef{迦 1:8}时，必会反对你们。既然保禄按照他的福音宣讲基督是达味的后裔，而你们否认这一点并宣讲别的，愿你们受诅咒！你们说基督从未声明祂自己由人所生，但祂却处处自称为人子，这又是什么意思呢？\par

3. 你们这些博学之士，诚然为我们的益处编造了一个奇妙的原初人，他从光明族类降下与黑暗族类争战，以他的水对抗敌人的水，以他的火对抗他们的火，以他的风对抗他们的风。为何不也用他的烟对抗他们的烟，用他的黑暗对抗他们的黑暗？照你们所说，他用空气对抗烟，用光明对抗黑暗。这样看来，烟和黑暗是邪恶的，因为它们不能属于他的善。而另外三种——水、风、火——是善的。那么，它们怎能属于敌人的邪恶呢？你们回答说，黑暗族类的水是邪恶的，而原初人带来的水是善的；同样，他的善风与善火对抗敌人的恶风与恶火。但他为何不能用善烟对抗恶烟呢？你们的谎言似乎消散在烟雾中。好吧，你们的原初人曾与一个相反的本性争战。然而他所带来的五样东西中，只有一样是敌对族类所拥有的对立面：光明与黑暗相对，而其他四样彼此并不对立。空气不是烟的对立面，更不用说水与水、风与风、火与火的对立面了。\par

4. 人们对你这关于\OriginalTerm{原人}{First Man}改变他所带来的元素，以便通过取悦敌人来征服他们的狂想感到震惊。因此，你们所谓虚假的王国诚实持守其本性，而真理却为了欺骗而变化。按你们所说，耶稣基督就是这原人之子。真理竟从你们的虚构中产生。你们称赞这原人在争斗中使用多变和欺骗的形式。那么，如果你们讲真话，你们就没有效法他。如果你们效法他，你们就像他一样欺骗。然而我们的主和救主耶稣基督，真实而真诚的天主之子，真实而真诚的人子，祂亲自见证了这两者，祂的天主性来自真天主，祂的降生成人来自真人。你们的原人并不是宗徒所说的第一个人。宗徒说：\BibleQuote{第一个人是出于地，属于土；第二个人是出于天，属于天。正如那属于土的怎样，凡属于土的也怎样；那属于天的怎样，凡属于天的也怎样。我们既然带着属土的形象，将来也必带着属天的形象。}\BibleRef{格前 15:47} 那出于地、属于土的第一个人就是亚当，由尘土造成。那出于天、属于天的第二个人就是主耶稣基督；因为祂是天主之子，成为血肉，使祂外表是个人，而内里仍然是天主；使祂既是真天主之子，由祂我们被造成，又是真人子，由祂我们被重塑。为什么你们要捏造这个虚构的原人，却拒绝承认宗徒所说的第一个人？这不正应验了宗徒的话：\BibleQuote{他们转离真理的耳朵，将去听信故事}\BibleRef{弟后 4:4} 吗？按\Person{保禄}所说，第一个人是出于地、属于土；按\Person{摩尼}所说，他并不是属土的，而是由某种虚幻不可理解的五种元素装备而成。\Person{保禄}说：\BibleQuote{如果有人向你们宣讲与我们传给你们的不同，当受诅咒。}因此，如果\Person{保禄}不是说谎者，那么\Person{摩尼}就当受诅咒。\par

5. 其次，你们指责那颗引导贤士朝拜婴孩基督的星，你们这样做本应感到羞愧，因为你们描述你们虚构的基督——你们虚构的\OriginalTerm{原人}{First Man}之子——不是由一颗星宣布，而是被捆绑在所有的星宿中。因为你们说，他在与黑暗种族争斗时与黑暗的本原混合，以便通过俘获这些本原，世界可以从混合中形成。因此，按照你们亵渎的幻想，基督不仅与天堂及所有星辰混合，而且与大地及其所有产物结合并混杂——不再是救主，反而需要被你们拯救，通过你们的吃和吐。\par

这种让你们的门徒给你们带来食物的愚蠢习俗，以便你们的牙齿和胃成为解救被捆绑在其中之基督的工具，是你们亵渎幻想的结果。你们宣称基督是这样被释放的——但并非完全；因为你们认为一些微不足道的无用微粒仍留在排泄物中，以便一次又一次地在各种物质形态中混合和结合，并且无论如何要在焚烧世界的火中被释放和净化，如果不是在此之前的话。不，即使那时，你们说基督也没有完全被释放；但祂的善良神性的一些极端微粒被玷污得无法洁净，被定罪永远留在黑暗的恐怖团块中。这些自称被我们关于一颗星宣布天主之子诞生的说法冒犯的人，仿佛这是将祂的诞生置于一个星座的影响之下；而他们不仅将祂置于星宿之下，而且置于与所有物质事物、所有蔬菜的汁液、所有肉体的腐朽以及所有食物的分解这样污秽的接触中——在其中祂被捆绑——以至于释放祂的唯一方式，至少是一个重大方式，就是人，即摩尼教的选民，能够成功消化他们的晚餐。\par

我们也否认星宿对任何人生育的影响；因为我们坚持，按照天主的公义律法，人的自由意志选择善恶，不受任何必然性的约束。我们更不会将永恒创造者及万有之主的降生成人置于任何星座之下！当基督按照肉身诞生时，贤士所见的星并没有统治的能力，而是作为证人相伴。它并非控制祂，而是通过向祂致敬来承认祂。此外，这颗星并非从世界之初就延续在造物主所规定的运行轨迹中的星之一。伴随着童贞女的新的诞生，出现了一颗新的星，引导那些自己寻找基督的贤士；因为它在他们前面行走，直到他们到达他们找到天主圣言以婴孩形式出现的地方。但有什么占星家曾想过让一颗星离开它的轨道，降临到被视为在该星之下出生的婴孩那里呢？他们认为星影响生育，而不是生育改变星的运行；所以，如果福音中的星是那些天体之一，它怎么能决定基督的行动，既然它被迫在基督诞生时改变自己的行动？但如果，更可能的是，一颗以前不存在的星出现以指明基督，那么这是基督诞生之果，而非其因。基督并非因为那星在那里而诞生，而是那星在那里因为基督诞生。如果有什么命运，它在诞生中，而非在星中。命运这个词来源于一个意为\Quote{说话}的词；既然基督是天主圣言，万物在存在之前都由祂说出，那么星宿的会合不是基督的命运，而是基督是星宿的命运。那创造诸天的同一意志取了我们的尘世本性。那统治星宿的同一能力舍了祂的性命，又取了回来。\par

6. 那么，为什么诞生的叙述就不是福音呢？既然它传递如此好的消息，治愈我们的疾病？难道是因为\Person{玛窦}不像\Person{马尔谷}那样以\Quote{耶稣基督福音的开始}开头，而是以\Quote{耶稣基督的族谱}开头？同样，\Person{若望}也可能被认为没有写福音，因为他没有说\Quote{福音的开始}或\Quote{福音书}，而是说\Quote{在起初已有圣言}。或许巧舌如簧的\Person{法乌斯托}会把若望福音的开场称为一个\Emph{Verbidium}，正如他把玛窦福音的开场称为\Emph{Genesidium}一样。奇怪的是，你们竟如此厚颜无耻，把你们那些愚蠢的故事称作福音。你们告诉我们，在与某个陌生的敌对民族作战时，天主为了保护自己的王国，只得将自身的一部分本性送入敌人贪婪的口中，并受到如此玷污，以至于历经那些劳苦和折磨后，仍不能完全炼净——这有什么好消息可言？这坏消息是福音吗？任何对希腊文稍有认识的人都知道，福音意为好消息。但你们的好消息在哪里？当你们的天主自己被描述为如同在黑暗中哭泣，直到黑暗和玷污从他的肢体上除去时？而当祂停止哭泣时，祂似乎又变得残忍。因为那将要陷入物质的那部分祂，做了什么该受这谴责的事？这部分必须永远哭泣。但是不；任何审视这消息的人都不会因它是坏消息而哭泣，反会因它不是真的而发笑。\par

\SourceRecord{Translated by Richard Stothert.；Nicene and Post-Nicene Fathers, First Series；Vol. 4.；Edited by Philip Schaff.；Buffalo, NY: Christian Literature Publishing Co.,；1887.；<http://www.newadvent.org/fathers/140602.htm>.}

% structure: p0001 source=h1 target=section reason=page-title

\section{\WorkTitle{驳斥福斯图斯，第三卷}}

\Emph{\Person{福斯图斯}反对天主降生成人，理由是福音作者彼此矛盾，并且降生成人不适合天主性。\Person{奥斯定}尝试消除批判性和神学性的困难。}\par

\Emph{\Person{福斯图斯}}说：我是否相信降生成人？就我而言，这正是我长期努力说服自己的事，即天主被生了下来；但\Person{路加}和\Person{玛窦}家谱的不一致使我不安，因我不知道该遵循哪一个。因为我想可能发生这样的事：由于不是全知，我可能会将真的当作假的，将假的当作真的。因此，在对解决这个争端绝望之后，我转向了\Person{马尔谷}和\Person{若望}，这两位仍然权威，且与其他两位一样是福音作者。我很有理由赞同\Person{马尔谷}和\Person{若望}的开端，因为他们完全没有提及\Person{达味}、\Person{玛利亚}或\Person{若瑟}。\Person{若望}说：\BibleQuote{在起初已有圣言，圣言与天主同在，圣言就是天主}，意指\Person{基督}。\Person{马尔谷}说：\BibleQuote{耶稣基督，天主子，福音的开始}，仿佛在纠正\Person{玛窦}，因\Person{玛窦}称呼他为达味之子。或许，\Person{玛窦}的\Person{耶稣}与\Person{马尔谷}的\Person{耶稣}是不同的人。这是我不相信\Person{基督}诞生的理由。\par

如果你能通过调和这些叙述来消除这个困难，我准备降服。然而，无论如何，相信天主，基督徒的天主，从母腹中诞生，是不相称的。\par

\Emph{\Person{奥斯定}}回答说：如果你仔细阅读福音，并查考那些你发现矛盾的地方，而不是轻率地谴责它们，你就会看到，全世界如此多博学之士承认福音作者的权威，尽管有这最明显的差异，这证明其中隐藏着比表面更多的东西。任何人，和你一样，都能看到\Person{玛窦}和\Person{路加}中\Person{基督}的祖先不同；而\Person{若瑟}出现在两者中，\Person{玛窦}的末了和\Person{路加}的开头。显然，\Person{若瑟}之所以可被称为\Person{基督}的父亲，是因为他是\Person{基督}母亲的丈夫；因此他的名字作为男性代表出现在家谱的开头或末尾。任何人，和你一样，都能看到\Person{若瑟}在\Person{玛窦}中有一个父亲，在\Person{路加}中却有另一个，祖父和所有其他直到\Person{达味}也是如此。难道所有那些认真研读圣经的有才学之士，拉丁作家确实不多，但希腊作家却无数，都忽视了这一明显的差异吗？他们当然看到了。没有人看不到。但是，由于对圣经崇高权威的应有尊重，他们相信这里有一些东西，会给那些求问的人，拒绝那些咆哮的人；会被那些寻找的人找到，却被那些批评的人夺去；会对那些叩门的人敞开，却向那些反对的人关闭。他们求问、寻找、叩门；他们领受、找到、进入。\par

整个问题是如何\Person{若瑟}有两个父亲。假设这是可能的，两份家谱都可能正确。有两个父亲，为什么不能有两个祖父，两个曾祖父，以此类推，直到\Person{达味}，他是\Person{撒罗满}和\Person{纳堂}的父亲？\Person{撒罗满}出现在\Person{玛窦}的列表中，\Person{纳堂}出现在\Person{路加}中。这是许多人的困难，他们以为两个人不可能有同一个儿子，却忘记了一个很明显的道理：一个人既可以称为生他的人的儿子，也可以称为收养他的人的儿子。\par

我们知道，收养在古代是司空见惯的；因为甚至妇女也收养其他妇女的孩子，如\Person{撒辣}收养\Person{依市玛耳}，\Person{肋阿}收养她婢女的儿子，\Person{法郎的女儿}收养\Person{梅瑟}。\Person{雅各伯}也收养了他的孙子，\Person{若瑟}的儿子。此外，\OriginalTerm{收养}{Adoption}一词在我们的信德体系中非常重要，正如宗徒著作所显示的那样。因为宗徒\Person{保禄}在谈论犹太人的优越时说道：\BibleRef{罗 9:4}\BibleQuote{他们是以色列人，义子的名分、光荣、盟约、法律的颁布、礼仪、恩许，都是他们的；圣祖也是他们的，并且按肉身说，基督也是从他们出来的，基督就是天主，在万有之上，永远受赞美。}他又说：\BibleRef{罗 8:23}\BibleQuote{我们自己也叹息，等待义子的名分，就是我们身体的救赎。}另外，其他地方：\BibleRef{迦 4:4}\BibleQuote{但是时期一满，天主就派遣了自己的儿子，由女人所生，生于法律之下，为要把在法律之下的人赎出来，使我们获得义子的名分。}这些章节清楚地表明，收养是一个重要的象征。天主有一个独生子，是从他自己的本体所生，关于他，经上说：\BibleRef{斐 2:6}\BibleQuote{他具有天主的形体，不以自己与天主同等为应当把持的。}我们不是从天主自己的本体所生，因为我们属于受造物，不是生出的，而是被造的；但为了使我们成为\Person{基督}的兄弟，他收养了我们。因此，当我们不是由他而生，而是被创造、被形成时，天主借他的圣言和恩宠生我们的行为，被称为收养。所以\Person{若望}说：\BibleRef{若 1:12}\BibleQuote{他赐给他们权力成为天主的子女。}\par

既然义子的惯例在我们祖先和圣经中都很常见，那么因谱系不同而轻率地谴责福音作者为虚假，仿佛两者不能同时为真，岂不是非理性的亵渎？为什么不冷静地思考一个简单的事实：在人类生活中，一个人常常可以有两位父亲，一位是他血肉所生，另一位是因义子关系后来成为他意志所生的儿子？如果第二位不能正确地被称为父亲，那么我们也无权称呼\BibleQuote{我们在天的父}（参见\BibleRef{玛 6:9}），因为我们并非由祂的实体所生，而是藉着祂的恩宠与最仁慈的旨意被义子接纳，这是宗徒的教导和确凿的真理。因为对我们而言，祂是唯一的天主、主和父：是天主，因我们由祂所造，虽出自人类父母；是主，因我们属于祂；是父，因藉着祂的义子我们重生。认真地研究圣经的人，稍加思考便容易看出，两位福音作者的不同谱系中，若瑟有两位父亲，因此有两个祖先列表。你们本来也能看到这一点，若不是被好辩之心蒙蔽了眼睛。其他远非你们所能理解的事物，也已在对这些叙述各部分的仔细探究中被发现。一人生子、另一人收养他，以致一人有两位父亲，这种常见的事例，你们本来可以想到作为解释——尽管你们身处摩尼教谬误之中——假如不是怀着敌意来阅读的话。\par

4. 然而，为何玛窦从亚巴郎开始往下至若瑟，而路加从若瑟开始往上，不是到亚巴郎，而是到那位造人的天主，并且藉着赐予诫命，使人因信德能成为天主的儿子；为何玛窦在他的福音开头记载族谱，路加则在救主受若翰洗礼之后记载；还有玛窦所载族谱数目的意义——他将族谱分成每段十四代的三个部分，尽管总和似乎缺少一个；而路加在洗礼之后所记载的世代数达到七十七，这数目正是主亲自在赦免罪债时所吩咐的：\BibleQuote{不止七次，而是七十七次}（\BibleRef{玛 18:22}）；又说到：\BibleQuote{祈求，就给你们；寻找，就要找到；敲门，就给你们开门}（\BibleRef{玛 7:7}）——这些事你们永远无法明白，除非你们受到某位杰出的、研究过圣经并在可能范围内取得所有进步的公教徒教导，或者你们自己转离谬误，以基督徒精神祈求，好得到；寻找，好找到；敲门，好给你们开门。\par

5. 既然这种本性与义子的双重父道消除了由谱系差异带来的困难，浮斯图就没有理由离开两位福音作者而投奔另外两位——那将是对他投奔之人的更大冒犯，而非对被他离弃之人的冒犯。因为圣作者们不希望以牺牲弟兄为代价来获得青睐。他们的喜乐在于合一，他们在基督内是一体；若是某人说了一件事，另一人说了另一件事，或者某人以某种方式，另一人以另一种方式言说，他们仍然都是讲真理，丝毫不相互矛盾；只要读者怀着虔敬谦卑的心，不以异端精神寻求争端的借口，而是以信德之心渴求建树。如今，认为福音作者给出不同父亲的祖先，这是完全可能的事，因为人可以有两位父亲，这观点与真理并无矛盾。于是福音作者们得以和谐一致，而你们呢，由于浮斯图的承诺，理应立即屈服。\par

6. 你们或许会为他附加的评论感到困扰：\Quote{无论如何，相信天主——基督徒的天主——从母胎中出生，这很难令人信服。}仿佛我们相信神性本性是从女子的子宫中出来似的。难道我刚才没有引用宗徒关于犹太人的见证吗：\BibleQuote{圣祖是他们的，并且按血统说，基督也是从他们来的，祂是万有之上的天主，永受称颂}（\BibleRef{罗 9:5}）？因此，我们的主和救主耶稣基督，按神性是天主的真子，按肉体是人的真子，并非以祂那超乎万有的天主身份由女子所生，而是以祂从我们取来的脆弱本性，为的是在其中为我们而死，并治愈我们这脆弱本性：祂并非以天主的形体——在祂内祂认为与天主同等不是僭越——由女子所生，而是以仆人的形体，在祂空虚自己时所取的形体。因此，祂被称为空虚自己，是因为祂取了仆人的形体，而非因为祂失去了天主的形体。因为祂不变地拥有那使祂在天主形体中与圣父同等的本性，同时祂取了我们有变化的本性，藉此祂可由童贞女生出。你们呢，一面反对将基督的血肉置于童贞女的子宫中，一面却将天主的神性本身不仅置于人类的子宫中，还置于狗和猪的子宫中。你们拒绝相信基督的血肉是在童贞女的子宫中受孕——在此过程中天主并未被发现，甚至未被改变——而你们却断言，在所有人兽中，在男子的精液和女子的子宫中，在陆地或水中的所有受孕中，天主的实际部分和神性本性不断地被束缚、封闭、污染，永不能完全获得释放。\par

\SourceRecord{Translated by Richard Stothert.；Nicene and Post-Nicene Fathers, First Series；Vol. 4.；Edited by Philip Schaff.；Buffalo, NY: Christian Literature Publishing Co.,；1887.；<http://www.newadvent.org/fathers/140603.htm>.}

% structure: p0001 source=h1 target=section reason=page-title

\section{\WorkTitle{驳斥福斯图斯，第四卷}}

\Emph{\Person{福斯图斯}拒绝旧约的理由，以及\Person{奥古斯丁}对此的评论。}\par

1. \Person{福斯图斯}说：\Quote{我是否相信旧约？如果它遗赠任何东西给我，我就相信；否则，我就拒绝。拿着宣告我被剥夺继承权的别人的文件，将是过分鲁莽。请记住，旧约中\Place{迦南}的应许是给犹太人，即受割礼者，他们祭献，禁食猪肉以及\Person{梅瑟}宣布不洁的其他动物，并守安息日、无酵节以及立约者所规定的其他同类事物。基督徒没有采纳这些礼仪，也没有人遵守；因此，如果我们不愿接受产业，就应该交出文件。这是我拒绝旧约的第一个理由，除非你教导我更好的。我的第二个理由是，这份产业是如此可怜、属血肉的东西，没有任何属灵的祝福，以致于有了新约及其光荣的天国和永生应许之后，我认为它不值得接受。}\par

2. \Person{奥古斯丁}回答说：没有人怀疑旧约中包含对现世事物的应许，因此被称为旧约；也没有人怀疑天国和永生的应许属于新约。然而，在这些现世事物中，是未来事物的预象，这些预象应在我们这些万世终结已经来到的人身上应验，这不是我个人的想象，而是宗徒的判断，当他说到这类事情时：\BibleQuote{这些事是我們的预像}；又说：\BibleQuote{这些事发生在他们身上，作为预像，并且写在经上，是为警戒我们这些末世的人。}\BibleRef{格前 10:6} 因此，我们接受旧约，不是为了获得这些应许的实现，而是在其中看到对新约的预言；因为旧约为新约作证。因此，主从死者中复活后，让祂的门徒不仅看见祂，而且触摸祂，但仍然担心他们怀疑自己必死和血肉的感官，便进一步从古书的见证中证实他们，说：\BibleQuote{凡在\Person{梅瑟}法律、先知和诗篇上所记载关于我的事，都必须应验。}\BibleRef{路 24:44} 因此，我们的希望不在于现世事物的应许。我们也不相信那些时期的圣善和属灵的人——圣祖和先知——专注于地上的事物。因为他们借着天主圣神的启示，明白了什么适合那个时代，以及天主如何安排所有这些言行作为未来的预像和预言。他们深切渴望新约；但在那些预言中，他们有一个个人的使命需要履行，通过这些预言，未来的新事被预告。所以这些人的生活和言语都是预言性的。属血肉的百姓确实只想到当前的祝福，尽管就连在百姓身上也有关于未来的预言。\par

这些事你们不明白，因为正如先知所说：\BibleQuote{除非你们相信，你们必不会明白。}\BibleRef{依 7:9} 因为你们没有受过天国——即基督的真实的天主教会——的教诲。如果受过，你们就会从圣经的宝库里拿出新旧的东西来。因为主亲自说：\BibleQuote{凡受过天国教诲的经师，就好像一个家主，从他的宝库里拿出新旧的东西来。}\BibleRef{玛 13:52} 所以，你们一方面声称只接受天主的新应许，另一方面却保留了血肉的旧性，只增加了错误的标新立异；关于这种标新立异，宗徒说：\BibleQuote{要躲避亵渎的新奇言论，因为它们只会加增不虔敬，他们的言论如同毒疮蔓延。其中有\Person{希默纳约斯}和\Person{斐肋托斯}，他们在信德上偏离了，说复活已经过去，并且颠覆了一些人的信德。}\BibleRef{弟后 2:16} 从这里你们可以看出你们错误教义的来源，即教导复活只是借着宣讲真理使灵魂复活，而不会有身体的复活。但你们这些保留血肉旧性的人，如果你们不拥有现世事物，又如何能理解内心之人属灵的事物——内心之人是在认识天主上更新的？你们编造关于现世事物的幻想，这些幻想由你们全部错误教义所组成的那些血肉事物的形象构成。你们夸耀自己鄙视无价值的\Place{迦南}地，那是一个实际存在的地方，并实际赐给了犹太人；然而你们却讲述一个光明之地，被黑暗种族之地像狭窄的楔子一样从一侧切开——这是一个不存在的东西，你们从你们心智的迷惑中相信它；因此你们的生活没有因拥有它而得到支撑，你们的心智却在渴望它中消耗殆尽。\par

\SourceRecord{Translated by Richard Stothert.；Nicene and Post-Nicene Fathers, First Series；Vol. 4.；Edited by Philip Schaff.；Buffalo, NY: Christian Literature Publishing Co.,；1887.；<http://www.newadvent.org/fathers/140604.htm>.}

% structure: p0001 source=h1 target=section reason=page-title

\section{\WorkTitle{驳斥福斯图斯，第五卷}}

\Emph{福斯图斯声称摩尼教徒而非天主教徒才是福音前后一致的信仰者，并通过比较摩尼教徒与天主教徒对福音诫命的遵从试图确立这一主张。奥古斯丁揭露摩尼教徒的伪善，并赞扬天主教徒的苦修。}\par

1. 福斯图斯说：\Quote{我是否相信福音？你问我是否相信，而我对其诫命的遵从表明我相信。我倒要问你是否相信，因为你没有给出相信的证明。我舍弃了父亲、母亲、妻子、儿女，以及福音所要求的一切；\BibleRef{玛 19:29} 你还问我是否相信福音？也许你不知道什么叫做福音。福音无非是基督的宣讲和诫命。我舍弃了所有金银，不再把钱带在钱囊里；以日用的食物为满足；不为明天忧虑；不为如何糊口、穿什么而挂虑：你还问我是否相信福音？你在我身上看到福音的祝福；\BibleRef{玛 5:3} 你还问我是否相信福音？你看到我贫穷、温良、缔造和平、心地纯洁、哀恸、饥渴、为义而受迫害和仇恨；你难道还怀疑我对福音的信仰吗？如今可以明白若翰洗者，在见到耶稣并听到祂的作为之后，为何还问祂是否是基督。耶稣恰当而公正地不屑于直接回答祂就是，而是提醒他已经听到的作为：\BibleQuote{瞎子看见，聋子听见，死人复活。} \BibleRef{玛 11:2} 同样地，我完全可以这样回答你问我是否相信福音的问题，即说：我舍弃了一切——父亲、母亲、妻子、儿女、金银、吃喝、享乐、逸乐；把这当作对你问题的充分答复，并相信你若在我身上不跌倒，你便是有福的。}\par

2. \Quote{但是，按照你们的说法，相信福音不仅意味着遵从它的诫命，也意味着相信其中所写的一切；而首先，就是相信天主诞生了。然而，相信福音不只是相信耶稣诞生了，还要遵行祂所命令的。所以，如果你说我不相信福音，因为我不信降生成人，那么你更不相信，因为你无视诫命。无论如何，在这些问题解决之前，我们是相等的。如果你无视诫命并不妨碍你宣称信仰福音，那么我拒绝族谱又为何妨碍我呢？并且，如果如你所说，相信福音既包括相信族谱，也包括遵从诫命，那么你为何谴责我，既然我们都不完全？你所缺的，我拥有；我所缺的，你拥有。但如果我们无疑地相信福音只在于遵行天主的诫命，那么你的罪过是双重的。正如俗语所说，逃兵指控士兵。但假设，既然你坚持如此，完美的信仰由两部分组成：一部分在于言语，即宣认基督诞生了；另一部分在于行为，即遵守诫命；那么显然，我的部分艰难而痛苦，你的部分轻松而容易。群众自然蜂拥到你那里，远离我，因为他们不知道天主的国不在于言语，而在于能力。那么，你为何责备我承担了更艰难的部分，而把容易的部分留给你，如同留给软弱的弟兄？你心中以为信仰的那一部分，即宣认基督诞生了，比其它部分更有能力拯救灵魂。}\par

3. 那么，让我们亲自询问基督，从祂自己的口中学习我们救恩的主要方式。\Quote{谁能进入祢的王国，哦基督？}祂回答：\BibleQuote{那承行我在天之父的旨意的人}\BibleRef{玛 7:21}，而不是\Quote{那宣认我出生的人}。祂又对门徒说：\BibleQuote{你们去教导万民，因父及子及圣神之名给他们授洗，教导他们遵守我所吩咐你们的一切}\BibleRef{玛 28:19}。不是\Quote{教导他们我出生了}，而是\BibleQuote{遵守我的诫命}。又说：\BibleQuote{你们如果实行我所命令你们的，就是我的朋友}\BibleRef{若 15:14}，不是\Quote{如果你们相信我出生了}。又说：\BibleQuote{如果你们遵守我的诫命，你们将存留在我的爱内}\BibleRef{若 15:10}，并且在许多其他地方。还有在山中圣训中，当祂教导说：\BibleQuote{神贫的人是有福的，温良的人是有福的，缔造和平的人是有福的，心里洁净的人是有福的，哀恸的人是有福的，饥饿的人是有福的，为义而受迫害的人是有福的}\BibleRef{玛 5:3}，祂从未说：\Quote{宣认我出生的人是有福的}。在审判时绵羊与山羊的分开中，祂说祂将对右边的人说：\BibleQuote{我饿了，你们给了我吃的；我渴了，你们给了我喝的}\BibleRef{玛 25:35}等等；因此\BibleQuote{承受王国吧}。不是\Quote{因为你们相信我出生了，承受王国吧}。又对寻求永生的富人，祂说：\BibleQuote{去，变卖你所有的一切，跟随我}\BibleRef{玛 19:21}；不是\Quote{相信我出生了，好使你获得永生}。你看，王国、生命、幸福处处应许给我所选择的、你所谓信德的两部分中的那一部分，而从未应许给你的部分。如果你能，指出来，哪里有记载说：谁宣认基督由女人所生，便是有福的，或将承受王国，或获得永生？即使假设信德有两个部分，你的部分也没有祝福。但如果我们证明你的部分根本不是信德的一部分呢？结论便是你们是愚蠢的，这甚至将毫无疑问地被证明。目前，足以表明我们的部分被真福所加冕。此外，我们也有言语宣认的真福：因为我们宣认耶稣基督是永生天主之子；耶稣以自己的口亲口宣布这宣认带有祝福，当祂对伯多禄说：\BibleQuote{约纳的儿子西满，你是有福的，因为不是血和肉启示了你，而是我在天之父}\BibleRef{玛 16:7}。因此，我们拥有信德的这两部分，不仅仅一部分，并且在这两部分中，基督都称我们为有福；因为在一部分中我们将信德付诸实践，而在另一部分中我们的宣认不掺杂亵渎。\par

4. \Emph{奥古斯丁}回答说：我已经说过，主耶稣基督屡次称自己为人子，而摩尼教人编造了一个关于某个神话般的\Quote{元人}的愚蠢故事，这人出现在他们不虔敬的异端中，不是属地的，而是混杂着虚假的元素，与宗徒所说的相反：\BibleQuote{第一个人属于地，属于土}\BibleRef{格前 15:47}；并且宗徒仔细警告我们：\BibleQuote{如果有人给你们宣讲与我们以前所宣讲的不同福音，当受诅咒}\BibleRef{迦 1:8}。因此，我们必须根据宗徒的真理，而不是根据摩尼教的谬误，相信基督为人子。既然福音作者断言基督是由女人所生，出于达味的后裔，而保禄在写给弟茂德的信中说：\BibleQuote{你要记住：根据我的福音，耶稣基督出于达味的后裔，从死者中复活了}\BibleRef{弟后 2:8}，那么很清楚，我们应该相信基督为人子是什么意思；因为祂是天主子，我们是由祂而被造的，祂也藉着降生成人成为人子，好为我们的罪而死，并为了我们的成义而复活。\BibleRef{罗 4:25} 因此祂同时称自己为天主子和人子。仅举众多例子中的一个，在\BibleBook{若望}{福音}中记载：\BibleQuote{我实实在在告诉你们：时候将到，且现在就是，死人要听见天主子的声音，听见的人必要生活。因为父在自己内有生命，祂也赐予子在自己内有生命，并且赐予祂执行审判的权柄，因为祂是人子}\BibleRef{若 5:25}。祂说：\BibleQuote{他们要听见天主子的声音}，又说：\BibleQuote{因为祂是人子}。作为人子，祂领受了执行审判的权柄，因为祂将以人的形态来审判，好让善人与恶人都能看见祂。祂以这形态升了天，门徒们听见那声音说：\BibleQuote{祂怎样升了天，也要怎样降来}\BibleRef{宗 1:14}。作为天主子，作为与父平等且同一天主，祂不会被恶人看见，因为\BibleQuote{心里洁净的人是有福的，因为他们要看见天主}。既然祂向信祂的人应许永生，而信祂就是信真实的基督，即祂自己所宣告和祂的宗徒们所宣告的那样，既是真实的天主子，又是真实的人子；那么你们摩尼教人，信一个虚假且假冒的人子的虚假且假冒的人，并教导说天主自己出于对敌对种族的攻击的恐惧，把自己的肢体交出受折磨，并且终究不能完全得解放，显然离基督向信祂的人所应许的永生甚远。诚然，当伯多禄宣认祂是天主子时，祂对伯多禄说：\BibleQuote{约纳的儿子西满，你是有福的}。但祂没有向那些相信祂是人子的人应许什么吗？既然天主子与人子是同一位，那么永生明确地应许给那些相信人子的人。祂说：\BibleQuote{梅瑟怎样在旷野里举起了蛇，人子也必照样被举起来，使凡信祂的人不至丧亡，反而获得永生}\BibleRef{若 3:14}。你们还想要什么？那么，信人子吧，好获得永生；因为祂也是天主子，能赐予永生：因为祂就是\BibleQuote{真天主和永生}，同一位若望在他的书信中这样说。若望还补充说，谁否认基督降生成人，谁就是假基督。\par

5. 既然如此，你就不必过分夸耀基督诫命的完美，因为你遵守了福音的规条。因为规条，即使你真的遵守了，若没有真信德，也不会使你得益。难道你不知道宗徒说：\Quote{我若把我所有的财产全施舍，并舍弃我的身体被火焚，但我若没有爱，我毫无益处吗？}\BibleRef{格前 13:3} 你为何夸耀自己拥有基督徒的贫穷，而你却缺乏基督徒的爱德？强盗之间也有某种彼此的爱，这出于对罪与恶的相互意识；但这并不是宗徒所称赞的爱。他在另一处区分了真爱与所有卑劣邪恶的情感，说：\Quote{诫命的归结就是爱，出自纯洁的心、良善的良心和无伪的信德。}\BibleRef{弟前 1:5} 那么，你如何能从虚假的信德中获得真爱？你坚持一种被谎言败坏的信德：因为按你的说法，你的原人在争斗中通过改变形貌使用了诡诈，而他的敌人却停留在自己的本性中；此外，你还主张基督——祂说\Quote{我是真理}——假装了祂的降生成人、十字架上的死、受难的伤痕和复活后显示的痕迹。如果你说真话，而你的基督说假话，那么你一定比祂更好。但如果你确实追随你自己的基督，那么你的诚实就可能受到质疑，而你对所谈论的规条的遵守可能只是一种伪装。福斯图斯说你的钱袋里没有钱，这是真的吗？他的意思大概是，你的钱放在盒子和袋子里；我们不会因此责备你，只要你不说一套做一套。君士坦提乌斯如今仍然健在，现在是我在天主教会中的弟兄，他曾将你教派的许多人聚集到他在罗马的家里，为了遵守摩尼的这些规条——你对这些规条如此看重，尽管它们非常愚蠢幼稚。这些规条对于你们的软弱来说太过沉重，聚集最终完全瓦解。那些坚持的人从你们的共融中分离出来，被称为马塔里亚人，因为他们睡在草席上——这与福斯图斯的羽毛床和山羊皮褥子，以及使他不仅鄙视马塔里亚人，还鄙视他在米莱维的穷苦父亲之家的所有奢华截然不同。那么，从你们的写作中摒弃这可咒的伪善吧，如果无法从行为中摒弃；否则，你们的言语会因你们欺骗性的言辞而与生命相冲突，正如你们的原人因其欺骗性的元素而与黑暗种族相冲突一样。\par

6. 然而，我并不是仅仅对未能遵行所命令之事的人说话，而是对一受骗教派的成员说话。因为摩尼的规条是：如果你们不遵守，你们就是欺骗者；如果你们遵守，你们就是受骗者。基督从未教导你们：为了害怕犯杀人罪而不采摘蔬菜；因为当祂的门徒经过麦田感到饥饿时，祂并没有禁止他们在安息日掐麦穗。这是对当时犹太人的斥责，因为那行为发生在安息日；同时那行为本身也是对未来的摩尼教徒的斥责。然而，摩尼的规条只要求你们什么都不做，而让别人为你们行杀人；实际上真正的杀人是通过这些魔鬼的教义毁灭可怜的灵魂。\par

7. 福斯图斯的话带有异端的\OriginalTerm{狂傲}{typhus}，是极度傲慢的话。他说：\Quote{你在我身上看到福音的真福；你还问我是否相信福音？你看见我贫穷、温良、缔造和平、心地纯洁、哀恸、饥渴、为义忍受迫害和仇视；你还怀疑我对福音的信仰吗？}如果自我称义就是正义，那么福斯图斯说这些话时早已飞升天堂了。关于福斯图斯的奢侈习惯，摩尼教徒中尽人皆知，尤其是在罗马的人，我这里就不提了。我将设想一个摩尼教徒，如同君士坦提乌斯所寻找的那样，他真诚地希望看到这些规条被遵守，并强制实行了它们。我怎能看见他神贫，当他如此骄傲，相信自己的灵魂就是天主，并且不羞于谈论天主受束缚？我怎能看见他温良，当他宁愿冒犯所有福音作者们的权威也不相信？怎样缔造和平，当他主张天主本性本身——天主所是的一切，唯一真实的存在——不能持久的平安？怎样心地纯洁，当他的心中充满如此多不敬的观念？怎样哀恸，除非是为他的天主被俘虏和束缚，直到被释放并逃脱，却失去一部分——这部分被圣父与黑暗之众结合，不应为之哀恸？怎样饥渴慕义——福斯图斯在他的著作中省略了这一项，无疑是怕被认为缺乏义德？但他们怎能饥渴慕义，他们的完美义德在于因他们的弟兄被定罪于黑暗而欢欣，这并非由于他们自己的过失，而是因为他们被圣父派去争战时所对抗的污染已无可救药地玷污了他们？\par

8. 你们怎能因义而受迫害和仇视，当你们认为宣讲和教导这些不敬之事是义的呢？令人惊奇的是，基督徒时代的温和竟允许如此悖逆的邪恶完全或几乎不受惩罚。然而，你们却告诉我们，你们遭受侮辱和迫害是你们正义的伟大证明，好像我们是瞎眼或愚蠢的。如果人们根据所受的苦难程度称义，那么各种穷凶极恶的罪犯受的苦远比你们多。但无论如何，如果我们承认因某种基督徒信仰的告白而受的苦难证明受难者拥有真信仰和正义，那么你们必须承认，我们所能展示的任何更大的苦难案例都证明拥有更真的信仰和更大的正义。这样的案例你们在我们的殉道者中知道很多，尤其是\Person{西彼廉}本人，他的著作也见证了他相信基督由童贞玛利亚所生。为了这个你们憎恶的信仰，他受难而死，同许多当时的基督徒信徒一起，他们受了同样多或更多的苦。\Person{福斯图斯}，当通过证据或他自己的供认被证明是\OriginalTerm{摩尼教徒}{Manichaean}时，在基督徒们自己的求情下——他们把他带到总督面前——同其他人一起，仅被流放到一个岛上，这几乎不能算作惩罚，因为天主的仆人们每天自愿这样做，当他们希望从世界的喧嚣中退隐时。此外，地上的君王常常通过公开的法令仁慈地解除这种流放。这样后来所有人都立即被释放了。那么，承认吧，那些拥有更真的信仰和更正义的生活的人，他们被认为配得上为此受比你们所受多得多的苦。否则，就停止夸耀许多人对你们的憎恶，并学会区分因亵渎受苦和因正义受苦。你们是为着什么受苦，你们自己的书会以你们最需要注意的方式表明。\par

9. 那些你们让不明真相的人相信你们遵守的卓越崇高的福音训诫，在我们团体中确实被许多人遵守。在我们中间，不是有许多男女完全禁绝了性关系，也有许多以前结婚的人现在守节吗？不是有许多人大量施舍自己的财产，或完全舍弃，还有许多人以往往频繁或每日或延长到难以置信的斋戒来克制身体吗？还有这样的兄弟会，其成员没有私有财产，一切公用，仅包括食物和衣物必需品，他们一心一意归于天主，充满共同的愛德情感。在这样的所有修持中，许多人成了欺骗者和被弃绝者，而许多这样的人从未被发现；也有许多人起初走得好，但很快因任性堕落。许多人在考验时期被发现怀着与他们声称的不同的意图过着这种生活；又有许多人在谦逊和坚定中，恒心到底，获得救恩。在这些团体中有明显的差异；但同一愛德联合了一切人，他们出于某种必要，服从宗徒的吩咐，有妻子却像没有一样，置买产业却像不留存，享用世界却像不享用。在天主丰盛的慈悲中，次等的阶层也加入他们，就是那些被告知的人：\BibleQuote{彼此不可亏负，除非两相情愿，暂时分房，为专心祈祷；以后还要同房，免得撒殚因你们的不能节制而诱惑你们。但我说这话是出于许可，不是出于命令。}\BibleRef{格前 7:5} 对同样的人，同一宗徒也说：\BibleQuote{你们彼此告状，这已经是你们的大错。}但因着他们的软弱，他补充说：\BibleQuote{在今生的事上，你们若须判断，就派教会中最不被尊敬的人来审判。}\BibleRef{格前 6:7} 因为在天国里，不仅有那些为成全而变卖或舍弃一切跟随主的人；还有那些在愛德的合伙中如同雇佣军加入基督徒军队的人，他们最后将听到：\BibleQuote{我饿了，你们给我吃，等等。}否则，对那些宗徒给予许多关切详尽指示关于他们家庭的人，就没有救恩了：他告诉妻子要顺服丈夫，丈夫要爱妻子；儿女要听从父母，父母要按主的教导和劝诫养育儿女；仆人要畏惧听从自己肉身的主人，主人要给予仆人公义公平的对待。宗徒远没有谴责这些人，认为他们漠视福音训诫或不配得永生。因为主劝勉强壮的人达到成全，说：\BibleQuote{谁若不背起自己的十字架跟随我，不能做我的门徒。}\BibleRef{玛 10:38} 祂立即为了安慰软弱的人加上：\BibleQuote{谁因义人的名接纳一个义人，必得义人的赏报；谁因先知的名接纳一个先知，必得先知的赏报。}\BibleRef{玛 10:41} 因此，不仅那为弟茂德的胃和屡次软弱给他一点酒的人，而且那仅因门徒的名给强壮的人一杯凉水的人，都不会失去他的赏报。\BibleRef{玛 10:42}\par

10. 倘若一个人不放弃一切就不能领受福音，你们为何欺骗自己的追随者，允许他们在侍奉你们时保留妻子、儿女、家人、房屋和田产？其实，你们完全可以让他们忽视福音的规诫：因为你们许诺给他们的并非复活，而是转向另一种有死的存在，在其中他们将度过你们所谓的\OriginalTerm{选民}{Elect}那种愚蠢、幼稚、不敬的生活——就是你们自己所过并备受称赞的生活；或者，如果他们拥有更大的功德，就会进入甜瓜或黄瓜或某些可食之物中，这些将被你们咀嚼，以便通过你们的消化迅速得到净化。最不应该的是，你们这些教导如此邪说的人还自称尊重福音。因为如果福音的信德与这类胡言有任何关联，主就该说：\Quote{我饿了，你们给了我吃的}，而应该说：\Quote{你们饿了，你们吃了我}，或\Quote{我饿了，我吃了你们}。因为按照你们的荒谬说法，一个人不会因给圣人们提供食物的服务而进入天国，反而因为他吃了他们并打嗝排出他们，或者他自己被吃并被嗝入天堂。义人不再说：\Quote{主啊，我们什么时候见你饿了，给你吃的？}而必须说：\Quote{我们什么时候见你饿了，被你吃了？}祂的回答也不该是：\Quote{你们对我这些最小兄弟中的一个所做的，就是对我做的}，而该是：\Quote{当你们被我这些最小兄弟中的一个吃了，就是被我吃了。}\par

11. 你们相信并教导这样的怪物，并照此生活，竟还敢声称你们遵守了福音的规诫，并诋毁天主教会——这教会包含许多软弱的和坚强的信徒，两者都蒙主祝福，因为两者都按照各自的程度遵守福音的规诫并寄希望于其应许。敌意的盲目使你们只看见我们庄稼中的莠子：如果你们愿意承认有麦子的话，你们本可轻易看见麦子。但在你们中间，那些假装的\OriginalTerm{摩尼教徒}{Manichæan}是邪恶的，而真正的摩尼教徒是愚蠢的。因为既然信仰本身是虚假的，假装信奉它的人行事诡诈，而真心相信它的人则受欺骗。这样的信仰不能产生善的生活，因为每个人的生活善恶取决于其心思的倾向。如果你们的情感专注于精神与理智之善，而非物质形态，你们就不会将物质太阳当作神圣实体和智慧之光来崇拜——人人都知道你们这样做，虽然我现在只是顺带提及。\par

\SourceRecord{Translated by Richard Stothert.；Nicene and Post-Nicene Fathers, First Series；Vol. 4.；Edited by Philip Schaff.；Buffalo, NY: Christian Literature Publishing Co.,；1887.；<http://www.newadvent.org/fathers/140605.htm>.}

% structure: p0001 source=h1 target=section reason=page-title

\section{驳斥福斯图斯，第六卷}

\Emph{\Person{福斯图斯}声称他不相信旧约，也不遵守其诫命，并指责天主教徒自相矛盾，因为他们在接受旧约为权威的同时却忽视其规条。\Person{奥古斯丁}解释天主教关于旧约与新约关系的观点。}\par

1. \Emph{\Person{福斯图斯}}说：\Quote{你问我是否相信旧约。当然不信，因为我不遵守其诫命。我想，你也不遵守。我拒绝割损礼，认为它可憎；如果我没错，你也一样。我拒绝遵守安息日，认为它多余；我想你也一样。我拒绝祭献，认为它是偶像崇拜；毫无悬念你也一样。猪肉不是我戒除的唯一肉类；你也不是只吃猪肉。我认为一切肉类都不洁；你认为没有肉是不洁的。在这一点上，我们都抛弃了旧约。我们都认为无酵饼周和帐棚节是不必要且无用的。不用紫色补麻布衣服；用麻和羊毛做衣服视为奸淫；在必要时将牛和驴同轭视为亵渎；不任命秃头、红发或有类似特征的人为司铎，认为他们在天主眼中不洁——这些事我们都鄙视并嘲笑，视之为无关紧要；然而它们都是旧约的诫命和典章。你不能因我拒绝旧约而责备我；因为不管这样做是对是错，你和我一样做了。至于你的信仰和我的信仰之间的区别，在于你选择欺诈行为，卑鄙地口是心非，而我，没有学会欺骗的艺术，坦白地声明我憎恨这些可憎的诫命及其作者。}\par

2. \Emph{\Person{奥古斯丁}}回答：新约的继承者如何以及为何接受旧约，已经解释过了。但由于\Person{福斯图斯}当时谈论的是旧约的许诺，而现在他谈论的是诫命，我回答说，他显示对道德诫命和象征性诫命之间的区别一无所知。例如，\BibleQuote{不可贪恋}是道德诫命；\BibleQuote{你应在第八天给每个男婴行割损}是象征性诫命。由于不作此区分，摩尼教徒和所有批评旧约著作的人，没有看到天主为前一个时代所规定的任何仪式都是未来事物的影子，因为这些仪式现在已被废止，就谴责它们，尽管无疑现在不适合的东西在当时是完全适合的，作为预表现在所启示的事物。在这一点上，他们与宗徒矛盾，宗徒说：\BibleQuote{这些事都发生在他们身上，作为鉴戒，并写下来，为教训我们这末世的人。}\BibleRef{格前 10:6}宗徒在这里解释为什么这些著作应当被接受，以及为什么不再需要继续象征性的仪式。因为当他说：\BibleQuote{写下来为教训我们}时，他清楚表明我们应当非常勤勉地阅读和发现旧约圣经的意义，并应极其尊敬它们，因为是为了我们才写下的。同样，当他说：\BibleQuote{这些事是我们的鉴戒}和\BibleQuote{这些事发生在他们身上，作为鉴戒}时，他表明，既然那些事物本身已清楚启示，用以预表这些事物的行动仪式就不再具有约束力。因此他在别处说：\BibleQuote{不要让任何人在饮食上，或关于节期、月朔、安息日之事论断你们，这些都是未来事物的影子。}\BibleRef{哥 2:16}这里，当他说：\BibleQuote{不要让任何人论断你们}在这些事上时，他表明我们不再受约束去遵守它们。而当他说：\BibleQuote{这些都是未来事物的影子}时，他解释这些仪式在什么时候具有约束力，即在那些现在完全向我们启示的事物被这些未来事物的影子所象征的时候。\par

3. 诚然，如果摩尼教徒因主的复活而成义——主复活之日是祂受难后的第三天，即第八日，紧随安息日之后，即第七日之后——那么他们属血肉的心意就必脱离笼罩其上的世俗情慾之黑暗，并因内心的割损而喜乐，就不会讥诮那在旧约中以肉身割损所预表的事，尽管他们不须在新约下遵守此规。但正如宗徒所说：\BibleQuote{为洁净的人，一切都是洁净的；但为败坏和不信的人，什么也不洁净，反而连他们的理智和良心都受了玷污。}\BibleRef{铎 1:15}这些人自视如此洁净，竟把自己身体的这些肢体视为——或假装视为——不洁，却因不信和谬误而如此玷污，以致他们厌恶肉身的割损——宗徒称之为信德之义的印证——的同时，却相信他们天主的神圣肢体被拘禁和污染在他们自身这些血肉的肢体中。因为他们说肉身是不洁的；随之，那被肉身所拘禁的部分中的天主，便成了不洁：他们声称必须洁净祂，在此事达成之前，尽其所能，祂承受肉身所遭受的一切情欲，不仅感受痛苦和困扰，也感受感官的享乐。因为他们说，正是为了祂的缘故，他们禁绝性交，以免祂更紧地被束缚于肉身的枷锁之中，或遭受更多的玷污。宗徒说：\BibleQuote{为洁净的人，一切都是洁净的。}如果这对人——他们可能因邪恶的意志而被引入恶——尚且成立，那么对那永远不变、无玷的天主，万物岂不更应是洁净的！在你们以粗暴非难所污损的那些书中，论及天主的智慧说：\BibleQuote{一切污秽的东西都不能进入她；她因自己的纯洁而通行各处。}\BibleRef{智 7:24} 凡指责那位——万物对祂皆是洁净的——天主所钦定、置于人类生殖器官上作为人类重生的记号，实在是好色荒唐；而你们却认为你们的天主——对祂来说没有什么是洁净的——在被不洁之人以自身肢体行邪淫时，其本性的一部分会受玷污和败坏。因为如果你们认为婚姻的男女交合有污染，那么纵欲者的一切行为又会何等呢？因此，如果你们常问：天主难道不能以别的方式来印证信德之义吗？回答是：为何不能用此方式？既然为洁净的人万物皆是洁净的，更何况为天主呢？我们有宗徒的权威，他说割损是亚巴郎信德之义的印证。至于你们，当被问及你们的天主是否无事可做，竟将其本性的一部分与这些你们如此辱骂的肢体纠缠在一起时，你们必须努力不脸红。这些是关于因人类繁衍而来的惩罚性败坏的话题，说起来微妙。它们考验贞洁者的端庄、不洁者的情欲，以及天主的正义。\par

4. 安息日的其余部分，我们不再认为必须遵守，因为永恒安息的希望已经启示出来。但是阅读并反思它是一件非常有益的事。在先知的时代，当现在显明的事通过行为和言语被预表和预言时，我们读到的这个标记是对我们所拥有的实体的预兆。但我想知道，为什么你们遵守一种部分的安息。犹太人在他们的安息日（他们仍然以\Quote{肉性的}方式遵守）既不在田间采摘任何果实，也不在家中和烹调。但你们在你们的安息中，等待你们的一个追随者拿刀或钩子到菜园，通过杀害蔬菜为你们获取食物，并带回——说来奇怪——\Quote{活着的尸体}。因为如果砍断植物不是谋杀，为什么你们害怕去做？然而，如果植物被杀害，那通过你们的咀嚼和消化获得释放和恢复的生命又怎样了？好吧，你们取来活的蔬菜，当然，如果可能的话，你们应该将它们整个吞下；这样，在你们的追随者采摘时所犯的那一处伤口之后（你们无疑会同意赦免他），它们可以毫无损失或伤害地到达你们私人的实验室，在那里你们的天主可以被医治他的伤口。相反，你们不仅用牙齿撕碎它们，而且如果味道合意，你们还会剁碎它们，以最罪恶的方式施加无数伤口。显然，如果你们也留在家中安息，不仅像犹太人那样每周一次，而是每天，那将是最有利的事。当你们烹调黄瓜时，它们受苦，而对它们内在的生命毫无益处：因为一个煮沸的锅无法与一个圣洁的胃相比。然而，你们嘲笑安息日的安息是多余的。与其指责先贤们遵守这一规定（对他们而言并非多余），不如即使在它已经多余的现今，你们自己遵守这种安息，而不是你们的安息，因为你们的安息没有象征意义，且被谴责为建立在谎言之上，岂不更好吗？\Quote{按照你们自己愚蠢的意见}，你们对自己安息的遵守是有缺陷的，尽管这种遵守在真理的判断下是愚蠢的。你们坚持认为水果在被从树上摘下时、在被切割、擦洗、烹调、食用时受苦。所以你们吃任何不能生吃且不受伤害地吞下的东西就是错的，这样所受的伤口可能不是来自你们，而是来自你们采摘水果的追随者。你们声称，如果只吃那些无需烹调或咀嚼就能吞下的东西，你们无法释放如此大量的生命。但如果这种释放补偿了你们施加的所有痛苦，为什么你们采摘水果就不合法呢？水果可以生吃，正如你们教派中的一些人坚持生吃各种蔬菜。但在它能被吃之前，它必须被摘下或掉落，或以某种方式从地上或树上取走。你们采摘它本可得到宽恕，因为不做这事便无法进行，但你们在准备食物时折磨你们天主的肢体到那种程度，则不可宽恕。你们的一个愚蠢观念是，当水果被摘下时树会哭泣。无疑树中的生命知道一切，并察觉到是谁来到它面前。如果选民来采摘水果，难道树不会为摆脱被他人采摘水果的痛苦而高兴，并因忍受一点暂时的痛苦而获得幸福吗？然而，当你们在水果被摘下后增加它的痛苦和麻烦时，你们却不去采摘它。如果你们能，请解释这一点！禁食本身对你们来说是一个错误。在从精神黄金中清除排泄物的渣滓以及从禁锢中释放神圣肢体的工作中，不应有间歇。\Quote{你们中最仁慈的人}是那位总是保持健康、吃生食并吃很多的人。但你们在吃的时候是残忍的，因为让食物遭受如此多的痛苦；你们在禁食的时候也是残忍的，因为停止了释放神圣肢体的工作。\par

5. 尽管如此，你们仍胆敢谴责\WorkTitle{旧约}的祭献，称之为偶像崇拜，并将同样的不敬思想归咎于我们。首先为自己辩护：我们虽认为祭献不再是应尽的义务，却承认祭献是启示奥秘的一部分，预表了预言之事。因为它们是我们的预像，以多种多样方式指向我们现在所纪念的那唯一祭献。如今这祭献已经启示出来，并在适当的时候被献上，祭献不再作为敬拜行为具有约束力，却保留其象征权威。因为这些事情\BibleQuote{写在我们的教训上，就是为着这些末世的人。}\BibleRef{格前 10:11}你们反对祭献的地方在于宰杀牲畜，尽管整个动物界在某种程度上是为人的使用而有条件存在的。你们怜悯牲畜，相信它们含有人的灵魂，却不肯将一块面包给饥饿的乞丐。主\Person{耶稣}却残酷对待猪群，答应了魔鬼请求允许它们进入猪身的要求。\BibleRef{玛 8:32}同一主\Person{耶稣}，在祂受难祭献之前，对一个被治愈的癞病人说：\BibleQuote{去，把你自己给司祭察看，并献上\Person{梅瑟}所命令的祭品，给他们作为证据。}\BibleRef{路 5:14}当天主通过先知一再宣称祂不需要祭品，正如理性告诉我们，祂一无所缺，不需要任何祭品时，人的心灵就被引导去探究天主藉这些祭献想要教导我们什么。因为，可以肯定的是，如果祂不需要这些祭品，祂就不会要求献祭，除非为了教导我们一些有利于我们认识的事情，而这些事情通过这些象征得到适当的表达。你们仍然受这些祭献的束缚，它们具有教导意义，尽管现在已不必要，这比要求你们的追随者向你们提供你们认为是有生命的牺牲作为食物，要好得多、荣耀得多。宗徒\Person{保禄}最恰当地论到一些为满足食欲而传福音的人说：\BibleQuote{他们的天主是肚腹。}\BibleRef{斐 3:19}而你们不敬的傲慢远甚于此；因为你们不但不以肚腹为天主，更糟的是，你们竟使肚腹成为天主的净化者。确实，假装虔诚不宰杀动物，同时又认为所有食物的灵魂都是动物的，并且让你们称为有生命的受造物在你们的手和牙齿间遭受如此折磨，这真是极大的疯狂。\par

6. 若你们不吃肉，为什么不宰杀动物祭献给你们的天主，以便使它们——你们认为不但是人类的，甚至是神圣到属于天主本身的肢体——的灵魂，从肉体的束缚中释放出来，并因你们祈祷的效力而获得拯救，不再回来呢？然而，也许你们的胃比你们的理智提供了更有效的帮助，而穿过你们肠子的那部分神性，比只得到你们祈祷的那部分更有可能得救。你们反对吃肉的理由是你们不能生吃动物，所以你们的胃的运作无助于释放它们的灵魂。多幸福的蔬菜啊，用手拔起，用刀切碎，被火折磨，被牙齿磨碎，却活着到达了你们肠子的祭台！多不幸的羊和牛啊，它们没有如此顽强的生命力，因此被拒绝进入你们身体！这就是你们观念中的荒谬。你们坚持认为我们与旧约对立，因为我们不认为任何肉是不洁的；按照宗徒的意见：\BibleQuote{为洁净的人，一切都是洁净的；}\BibleRef{铎 1:15} 以及按照主亲自说的话：\BibleQuote{不是入你口的污秽你，而是出口的。}\BibleRef{玛 16:11} 这话不只是对群众说的，正如你们的阿迪曼图斯——福斯图斯在他攻击旧约的书中称赞他仅次于摩尼——希望我们理解的那样；但主在离开群众后，更清楚、更明确地向祂的门徒重复了这一点。阿迪曼图斯引用主的话来反对旧约，旧约禁止百姓吃某些被称为不洁的动物；毫无疑问，他害怕被问：既然他从福音中引用了关于人不被入他口、入他腹、最后入茅厕之物所玷污的经文，为什么他不但认为某些肉，而是认为所有的肉都是不洁的，并戒绝吃肉？正是为了逃避这一困境，当明确的真理过于压倒他的错误时，他使主对群众说这话；仿佛主只有在人少时才说真话，而在人前就口吐谎言。用这种方式谈论主是亵渎。所有读这经文的人都可以看到，主私下里向门徒更清楚地说了一样的事。既然福斯图斯在他书的开头如此赞扬阿迪曼图斯，将他置于摩尼之后，那么请他一句话告诉我：人是否被入他口之物所玷污，这是真是假？如果是假的，为什么这位伟大的教师阿迪曼图斯引用它来反对旧约？如果是真的，为什么尽管如此，你们却相信吃任何肉都会玷污你们？如果你们这样选择，那么宗徒并没有说一切对于异端者都是洁净的，而是说：\BibleQuote{为洁净的人，一切都是洁净的。}宗徒接着解释为什么一切对于异端者不洁：\BibleQuote{为污秽不信的人，什么也不洁净，连他们的理智和良心都受玷污。}\BibleRef{铎 1:15} 所以对于摩尼派信徒来说，绝对没有任何东西是洁净的；因为他们认为天主的实质或本性不但可能被玷污，而且实际上已经被玷污，并且玷污到永远无法完全恢复和洁净的程度。当他们称动物为不洁并戒绝吃肉时，他们是什么意思？因为他们不可能认为任何东西，无论是食物还是其他什么，是洁净的。按照他们的说法，蔬菜、水果、各种作物、地和天，都已因与黑暗族的混合而被玷污。为什么他们不在其他事上，也像对待动物一样，按自己的意见行事？为什么不完全戒食，饿死自己，而不是坚持他们的亵渎？如果他们不肯悔改，这显然是他们能做的最好的事。\par

宗徒所说的\Quote{为洁净的人一切都是洁净的}，以及\Quote{天主所造的每一样都是好的}，与旧约的禁令并不冲突；其解释（如果他们能够理解的话）如下。宗徒谈论的是事物的本性，而旧约称某些动物为不洁净，并非就其本性而言，而是象征性地，因为那时期具有预象的性质。例如，猪和羊在本性上都是洁净的，因为天主所造的每一样都是好的；但象征性地，羊是洁净的，猪是不洁净的。同样，\OriginalTerm{智}{wise}和\OriginalTerm{愚}{fool}这两个词在本性上都是洁净的，因为它们都是由字母组成的词；但\OriginalTerm{愚}{fool}可能被象征性地称为不洁净，因为它表示不洁净的事物。或许猪在象征物中的地位，就如同愚人在真实事物中的地位一样。这种动物，以及组成\Quote{fool}一词的四个字母，可能意指同一事物。无疑，这种动物被法律宣布为不洁净，因为它不反刍；这不是它的过错，而是它的本性。但那些以此动物为象征的人，不洁并非出于本性，而是因自身的过错；因为他们虽然乐意聆听智慧之言，却从不事后反思。因为在宁静的安息中，将某些有益的教训从记忆的胃中召回反思的口中，是一种精神上的反刍。上述动物是那些不这样行的人的象征。而禁止食用这些动物的肉，便是对这项过错的警告。圣经的另一处经文提到智慧的宝贵宝藏，并将反刍描述为洁净，不反刍为不洁净：\BibleQuote{智慧的宝藏，在智慧人的口里；愚昧人却把它吞下。}\BibleRef{箴 21:20}这类象征，无论是言语上的还是事物上的，给聪明的头脑在探究和比较方面提供了有用而愉快的操练。但从前人们不仅被要求聆听，还要实践许多这类事物。因为在那时，必须以行为及言语来预象那些后来将要启示的事。在基督内并由基督所启示之后，信徒团体不再肩负履行这些规诫的重担，而是被劝诫要留心那预言。这就是我们依据主和宗徒的话，认为没有动物是不洁净的原因，同时我们并不反对旧约中称某些动物为不洁净。现在，让我们听听你为何认为所有动物食品是不洁净的。\par

8. 你们的一项虚假教义是，肉体因与黑暗族混合而不洁。但这不仅会使肉体不洁，而且使你们的天主本身也不洁，因为他派遣了一部分自己，使之成为被吞没和污染的对象，以便战胜并俘获敌人。此外，由于这种混合，你们所吃的一切都是不洁的。但你们特别说肉体是不洁的。要耐心听取他们为肉体这种特殊不洁所提出的所有荒谬理由。我只提足以显示这些旧约批评者根深蒂固的愚昧之点即可，他们一面指责肉体，一面却只品味属肉之事，毫无属灵知觉。就此问题的冗长讨论或许能让我们免去对其他一些点的过多论述。以下是他们在这一问题上虚幻妄想的描述：——在那场战斗中，当\Person{原人}用以欺骗的元素诱捕了黑暗族时，属于该族的男女首领都被擒获。世界是用这些首领建造的；在用于形成天体的那些首领中，有一些怀孕的雌性。当天空开始旋转时，快速的圆周运动使这些雌性流产，流产的胎儿（有男有女）落在地上，活着，长大，交合，产生后代。地上、空中和海里的一切动物生命由此产生。既然肉体的起源来自天上，就没有理由认为它特别不洁。确实，在建造世界时，他们主张这些黑暗族元素被安排得或高或低，取决于在建造世界各部分时与它们混合的善的多少。因此，肉体应该比从地里长出的蔬菜更洁净，因为它来自天上。而且，认为流产的胎儿在获得灵魂之前——尽管是流产状态——如此有活力，以至于从天上掉下后还能活下来并繁衍，这是多么不合理！相反，在获得灵魂后，如果早产就会死亡，从不太高的地方摔下就足以致命。生命王国与死亡王国争战，理应通过给予生命来改善它们，而不是使它们比以前更易腐朽。如果易腐朽性是本性变化的结果，那么说存在一个恶的本性就错了。易腐朽性只是变化所致。两种本性都是善的，尽管一种比另一种更好。那么，在这个世界中存在的肉体，据他们说来自天上，其特殊不洁又从何而来？他们确实告诉我们，这些黑暗族的原初身体像蠕虫一样从黑暗之树生出；而树木则是从五种元素产生。但假设动物的身体最初来自树木，而后来自天上，为什么它们比树木的果实更不洁？也许有人会说，死后留下的残余是不洁的，因为生命已不在其中。同样，水果和蔬菜也必须是不洁的，因为它们在采摘或切割时死了。如我们先前所见，选民让别人给他们送来食物，以免犯杀人罪。也许，因为他们说每个生物都有两个灵魂，一个来自光明族，另一个来自黑暗族，善的灵魂在死亡时离开，恶的灵魂留下。但那样的话，动物就会像它在黑暗王国时一样有生命，当时它只有本族的灵魂，曾以此反抗天主的王国。因此，既然两个灵魂在死亡时都离开了，为什么称肉体为不洁，仿佛只有善的灵魂离开了？任何残留的生命必然属于两种；因为天主的肢体的一些残余甚至存在于污秽之中。因此没有理由使肉体比水果更不洁。事实上，他们声称持守肉体不洁因它由生育而来，从而假装大贞洁。但如果神体被肉体更粗鲁地封闭，那么他们就更应该通过进食来解放它。有无数的蠕虫并非由性交产生；有些在威尼斯附近从树中生出，他们应该吃这些，因为没有同样的理由使它们不洁。此外，还有雨后从地里生出的青蛙。让他们从这些青蛙中解放天主的肢体吧。让他们责备人类的错误，人类更喜欢由雄雌所生的鸡和鸽子，而不是纯洁的青蛙，它们是天地之女。根据这一理论，从树中产生的黑暗元素的首领一定比由生育产生的\Person{摩尼}更洁净；同样理由，他的追随者一定比从他们身体的汗水生出的虱子更不洁净。但如果凡来自肉体的都是不洁的，因为肉体本身的起源是不洁的，那么水果和蔬菜也一定是不洁的，因为它们是用粪肥施肥的。此后，水果比肉更洁净的观念又成了什么？粪便是肉体最不洁的产物，也是最具肥力的肥料。他们的教义是，生命在咀嚼和消化食物时逸出，只有一粒粒子留在排泄物中。那么，这一粒生命何以对你们喜爱的食物的生长和质量有如此大的影响？肉体是靠土地的产物滋养，而不是靠它的排泄物；而土地是靠肉体的排泄物滋养，而不是靠它的产物。让他们说哪个更洁净。或者，让他们弃绝不信和不洁——对他们而言没有洁净之物——并与我们一同拥抱宗徒的教义：\BibleQuote{为洁净者，万有都是洁净的}；\BibleQuote{大地和其中的万物属于主}；\BibleQuote{天主的各样受造物都是好的}。自然界的一切事物在其秩序中都是好的；无人因使用它们而犯罪，除非因不服从天主，他逾越了自己的秩序，并因错用它们而扰乱了它们的秩序。\par

9. 那些蒙天主喜悦的长老藉着服从保持了自己的秩序，依照天主的安排，遵守了适于特定时期的规定。因此，虽然所有作为食物的动物本性都是洁净的，但他们禁食了一些在当时具有象征性不洁的动物，为给未来所预示之事的启示做准备。同样，关于无酵饼和所有类似的事物，宗徒说有未来事物的影儿，在旧约时代，当这些遵守被命令并用于预表后来要启示之事时，忽略遵守既是犯罪，正如现在新约的光明已经升起，我们却以为这些预言性的遵守能对我们有任何益处，这是愚蠢的。另一方面，既然旧约教导我们，现在启示的事是那么久以前预表的，使我们能坚定而忠信地持守它们，那么仅仅因为主如今要求我们不是以字面方式，而是以属神和明智的态度对待其中的内容，就抛弃这些书卷，便是亵渎和不敬。它们被写下来，正如宗徒所说，是为了我们的劝诫，我们就是世世代代终结的人。\BibleRef{格前 10:11} \BibleQuote{因为凡先前所写的，都是为教训我们而写的。}\BibleRef{罗 15:4} 在旧约时代，在指定的七天内不吃无酵饼是罪；在新约时代，这不是罪。但是，藉着基督——祂用正义披戴我们的灵魂，用不朽披戴我们的身体，使我们全然更新——对来世的望德使我们相信，我们原始腐化的束缚和软弱将不会胜过我们或我们的行为，这必须继续是罪，直到七个时间周期完成。在旧约时代，这在预像的伪装下被一些圣人领悟。在新约时代，它被完全宣告并公开宣讲。\par

那时圣经的命令如今成了见证。从前，不守帐棚节是罪，现在不是。但是，不参与建造天主的帐幕——即是教会——始终是罪。从前这件事以预像实践；现在这记录成了见证。古时的帐幕被称为见证的帐幕，除非它作为合宜的象征，为某些将在适当时候启示的真理作证。用紫布补缀细麻衣服，或穿羊毛和细麻混织的衣服，现在不是罪。但是，生活无节制，并希望结合相反的生活方式——例如，一位献身宗教的女子佩戴已婚妇女的装饰，或一位没有守独身的人打扮如贞女——始终是罪。因此，当任何人生活中不一致的事物结合时，就是罪。这如今是道德真理，当时则象征着服饰。那时是预像，现在是启示的真理。因此，同一部圣经当时要求象征性的行动，如今则为所指的事物作证。预表性的遵守如今成了一项记录，用以坚定我们的信德。从前，用牛和驴一同耕地是不合法的；现在则是合法的。宗徒引用关于不可笼住正在打谷的牛嘴的经文时解释了这一点。他说：\BibleQuote{天主岂关心牛吗？}那么，我们与这个过时的禁令有何关系呢？宗徒在接下来的话中教导我们：\BibleQuote{这原是为我们的缘故写的。}\BibleRef{格前 9:9} 我们若不阅读那为我们的缘故所写的，必是不敬；因为它更是为我们的缘故，我们拥有启示，而他们只有预像。在需要时，将牛和驴套在一起并无害处。但是，将智者和愚人放在一起，不是为了一个教导、另一个服从，而是为了使两者以同等权威宣讲天主圣言，这做不到而不引起冒犯。因此，同一部圣经，曾经是命令，要求那预表未来事物的影儿，如今则是权威的见证，为揭示的真理作证。\par

关于秃头或红发之人的不洁，福斯图斯说得不准确，或者他所用的抄本有误。\BibleRef{肋 21:18} 但愿福斯图斯不以在额头上佩戴基督的十字架为耻——缺少十字架就是秃头——而不是坚持认为基督——祂说\BibleQuote{我是真理}——在复活后显示了不真实的印记，不真实的伤口！福斯图斯说他未曾学会欺骗的伎俩，并且言其所思。因此，他不可能是他的基督的门徒，因为他疯狂地宣称基督在门徒怀疑时对他们显示了假的伤口印记。我们难道要相信福斯图斯，不仅相信他其他的荒谬言论，而且相信他告诉我们他没有欺骗我们，却在称基督为欺骗者时欺骗我们吗？他比基督更好吗？他不就是欺骗者，而基督不是？或者他藉着这个谎言——他夸口说未曾学会欺骗的伎俩——证明自己不是真理基督的门徒，而是欺骗者摩尼的门徒？\par

\SourceRecord{Translated by Richard Stothert.；Nicene and Post-Nicene Fathers, First Series；Vol. 4.；Edited by Philip Schaff.；Buffalo, NY: Christian Literature Publishing Co.,；1887.；<http://www.newadvent.org/fathers/140606.htm>.}

% structure: p0001 source=h1 target=section reason=page-title

\section{驳斥福斯图斯，第七卷}

\Emph{家谱问题再次被提出并在双方之间进行了辩论。}\par

1. \Person{福斯图斯}说：您问我为什么不相信耶稣的家谱。有很多原因；但主要原因是，祂从未亲口说祂有地上的父亲或家系，相反，祂说祂不属于这个世界，祂来自天主圣父，祂从天降下，祂没有母亲或兄弟，除非那些遵行祂在天之父旨意的人。此外，这些家谱的编写者似乎并不了解耶稣在出生之前或刚出生之后的情况，所以他们叙述的事情缺乏目击者的可信度。他们认识耶稣时，耶稣已经是一个约三十岁的年轻人，如果提到一个神性存在的年龄不算亵渎的话。关于见证人的问题总是他是否亲眼看到或亲耳听到他所见证的事情。但这些家谱的编写者从未断言他们从耶稣本人那里听到这个叙述，甚至没有听到祂出生的事实；直到祂受洗之后很多年，他们才认识祂，那时距离祂出生已经很久了。因此，对我以及每一个明智的人来说，相信这个叙述就像传唤一个又聋又瞎的证人上法庭一样愚蠢。\par

2. \Person{奥古斯丁}回答说：关于福斯图斯所说的他不接受耶稣基督家谱的主要理由，我们在前面引用的经文中已经找到了完全的驳斥，即基督自称为人子，以及我们所说的人子与天主子的同一性：按祂的天主性，祂没有地上的家系，而按肉身，祂是达味的后裔，正如宗徒所教导的。因此，我们应当相信祂来自圣父，祂从天降下，并且圣言成了血肉，住在人间。如果引用\BibleQuote{谁是我的母亲？谁是我的兄弟？}\BibleRef{玛 12:48}来表明基督没有地上的母亲或家系，那么我们也必须相信，祂的门徒——祂以自己的榜样教导他们不要重视地上的亲属关系，而是看重天国——也没有父亲，因为基督对他们说：\BibleQuote{不要在地上称人为父；因为你们的父只有一位，就是天上的父。}\BibleRef{玛 23:9}祂教导他们对待父亲的方式，祂自己首先对待祂的母亲和兄弟；正如在许多其他事情上，祂屈尊为我们树立榜样，走在前面，使我们能跟随祂的足迹。福斯图斯反对家谱的主要理由完全失败了；在这支无敌部队被打败之后，其余的就很容易被击溃了。他说，那些宣告基督既是人子又是天主子的宗徒不可信，因为他们没有在基督出生时在场，他们是在基督成年后才加入祂的，也没有从基督本人那里听到这件事。那么，他们为什么相信若望所说：\BibleQuote{在起初已有圣言，圣言与天主同在，圣言就是天主。他在起初就与天主同在。万物是藉着他而造成的；凡受造的，没有一样不是由他造成的，}\BibleRef{若 1:1}以及类似段落，他们同意这些却不理解它们？若望在哪里看见这些？或者他是否从主本人那里听到过？无论若望如何知道这些，那些叙述降生的人也可能知道。此外，他们怎么知道主说了\BibleQuote{谁是我的母亲？谁是我的兄弟？}？如果是根据福音书的权威，为什么他们不相信基督的母亲和兄弟当时正在寻找祂？他们相信基督说了这些他们误解的话，却否认基于同一权威的事实。再者，如果玛窦不可能知道基督的降生，因为他只在基督成年后才认识祂，那么\Person{摩尼}，这位如此之后的人，怎么知道祂没有降生？他们会说，\Person{摩尼}从在他内的圣神知道这一点。当然，圣神会使他说真话。但为什么不更相信基督自己的门徒呢？他们亲自认识祂，不仅拥有灵感的恩赐来弥补知识的不足，而且以完全自然的方式，在事件记忆犹新之际获得了基督降生和家系的信息。然而他竟敢称宗徒为又聋又瞎。为什么你不又聋又瞎，以免学到这种亵渎的胡说，并且又哑，以免说出它呢？\par

\SourceRecord{Translated by Richard Stothert.；Nicene and Post-Nicene Fathers, First Series；Vol. 4.；Edited by Philip Schaff.；Buffalo, NY: Christian Literature Publishing Co.,；1887.；<http://www.newadvent.org/fathers/140607.htm>.}

% structure: p0001 source=h1 target=section reason=page-title

\section{\WorkTitle{驳福斯图斯，卷八}}

\Emph{\Person{福斯图斯}主张：在新约颁赐之后仍持守旧约，就是将新布补在旧衣服上。\Person{奥斯定}进一步解释旧约与新约的关系，并指责摩尼教徒的肉性。}\par

1. \Emph{\Person{福斯图斯}}说：另一个不接受旧约的理由是，我既有了新约，而圣经说新旧不相合。因为\BibleQuote{没有人把新布补在旧衣服上，否则破绽就更大了}（\BibleRef{玛 9:16}）。为了避免像你们那样造成更大的破绽，我不将基督的新与希伯来的旧相混合。一个人有了新衣，不将旧衣送给下属，大家都觉得他小气。同样，即使我如宗徒们一样生来是犹太人，在接受新约之后，也应当抛弃旧约，正如宗徒们所做的一样。况且我生来就不在奴役的轭下，且幼年就进入了基督完全的自由，若我再把自己置于轭下，那是何等愚昧和忘恩！这正是\Person{保禄}责备迦拉达人的地方，因为他们回到割礼，又\BibleQuote{转向那软弱贫乏的要素，情愿再受其奴役}（\BibleRef{迦 4:9}）。我何必做那看到别人做而受责备的事呢？我去受奴役比他们回去更糟。\par

2. \Emph{\Person{奥斯定}}回答：我们已经充分说明了为何以及如何维护旧约的权威，不是为模仿犹太人的奴役，而是为确证基督徒的自由。不是我自己，而是宗徒说：\BibleQuote{这一切发生在他们身上，是为作鉴戒，并写下来是为劝戒我们这些活在末世的人}（\BibleRef{格前 10:11}）。因此，我们不是像奴仆一样遵守那些预言我们的命令，而是像自由人一样诵读那些为坚固我们而写的文字。所以任何人都可以看出，宗徒责备迦拉达人，不是因为他们虔诚地阅读圣经中关于割礼的字句，而是因为他们迷信地渴望受割礼。我们并不将新布补在旧衣服上，而是受教于天国，如同家主，主描述\BibleQuote{他从宝库里拿出新旧的东西}（\BibleRef{玛 13:52}）。那将新布补在旧衣服上的人，是那些在放弃肉性希望之前企图灵性克己的人。细查这段经文，你会发现，当主被问及守斋时，祂回答说：\BibleQuote{没有人把新布补在旧衣服上}。门徒们对主仍有肉性的情感，因为他们害怕如果主死了，他们就会失去祂。所以主称\Person{伯多禄}为撒殚，因为他劝阻主受苦，因为\BibleQuote{他不体贴天主的事，只体贴人的事}（\BibleRef{玛 16:23}）。你们对天国的幻想，以及你们对太阳之光（肉眼所见）的崇拜，仿佛它是天堂的肖像，都显明你们希望的肉性特征。因此，你们的肉性之心就是旧衣服，你们将禁食连在其上。再者，如果新布和旧衣服不相合，那么你们的天主的肢体如何不仅被连接或固定，而且通过混合与紧密接触与黑暗原则更亲密地结合呢？也许两者都是旧的，因为两者都是虚假的，都属于肉性之心。或者，你们想通过破绽变得更坏来证明一个是新的一个是旧的，即撕开那可怜的光明国度的一部分，注定永远监禁在黑暗团中。所以这个假装精通圣经章法的艺术家，竟将荒谬缝在一起，用自己发明的破布装扮自己。\par

\SourceRecord{Translated by Richard Stothert.；Nicene and Post-Nicene Fathers, First Series；Vol. 4.；Edited by Philip Schaff.；Buffalo, NY: Christian Literature Publishing Co.,；1887.；<http://www.newadvent.org/fathers/140608.htm>.}

% structure: p0001 source=h1 target=section reason=page-title

\section{\WorkTitle{驳斥福斯托斯，卷九}}

\Emph{福斯托斯主张：若生于旧盟约之下的宗徒可合法地脱离之，则他生于外邦人之中，更可如此。奥斯定阐明犹太人与外邦人同样与福音的关系。}\par

1. \Emph{\Person{福斯托斯}}说：不接受旧约的另一个理由是，若生于其下的宗徒可以弃绝之，则我未生于其下，更可免除强使自己进入其中的责任。我们外邦人不是生来为犹太人，也不是生来为基督徒。从同一外邦世界，有人被旧约引致成为犹太人，有人被新约引致成为基督徒。这犹如两棵树，一棵甜一棵苦，从同一土壤中吸取汁浆，各按本性同化。宗徒从苦树转入甜树；而我从甜树转入苦树，实属疯狂。\par

2. \Emph{\Person{奥斯定}}回答说：你说宗徒在离开犹太教时，是从苦树转入甜树。但宗徒自己说，那些不信基督的犹太人，是折下的枝子；而外邦人，野橄榄枝，被接在好橄榄上，即希伯来人的圣根上，为能共沾橄榄根的肥汁。因为在警告外邦人不要因犹太人跌倒而骄傲时，他说：\BibleQuote{我对你们外邦人说：我既是外邦人的宗徒，我光荣我的职务，若能使我的骨肉发愤，因而拯救他们一些人。若他们被弃绝是天下的和好，则他们被接纳岂不是死而复生么？若初熟之果是圣的，面团也是圣的；若根是圣的，枝条也是圣的。若有些枝条被折下了，而你这野橄榄枝接在其中，同沾橄榄根的肥汁，你就不可向枝条夸耀；你若要夸耀，该知道不是你托着根，而是根托着你。你要说：枝条被折下，是为叫我接上。不错；他们因不信被折下，你因信德站得稳。不可高傲，反要恐惧；因为天主既然没有怜惜本来的枝条，也要当心不怜惜你。所以你要看天主的恩慈与严厉：对跌倒的人，是严厉；对你，是恩慈，只要你存留在他的恩慈中；不然，你也要被砍下。他们若不执迷于不信，仍要被接上，因为天主能再把他们接上。你从野橄榄上被砍下来，逆性接在好橄榄上，何况这些本枝子接在本橄榄上呢？弟兄们，我不愿你们不知道这奥秘（免得你们自以为聪明），就是以色列有一部分硬心，直到外邦人的数目添满，那样全以色列都要得救。}\BibleRef{罗 11:16}由此可见，你这不愿接在这根上的人，虽不像属血肉不信的犹太人那样被折下，却仍留在野橄榄的苦味中。你崇拜太阳和月亮，具有真正外邦人的风味。你仍在外邦人的野橄榄中，因为你加上了新种类的荆棘，与太阳月亮一同崇拜一个假基督，这不是你手所作的，而是你邪僻的心所捏造的。所以，来吧，接在橄榄树的根上——宗徒曾因不信落在被折的枝子中，他归回这根时，便欢喜快乐。他描述自己获得释放，当他从犹太教幸福地转入基督教时。因为基督一直在橄榄树中被宣讲，那些在他来时不信他的被折下，而信的却被接上。这些人因而被警告不可骄傲：\BibleQuote{不可高傲，反要恐惧；因为天主既然没有怜惜本来的枝条，也要当心不怜惜你。}又为防止被折下的人绝望，他接着说：\BibleQuote{他们若不执迷于不信，仍要被接上，因为天主能再把他们接上。你从野橄榄上被砍下来，逆性接在好橄榄上，何况这些本枝子接在本橄榄上呢？}宗徒因脱离被折枝子的境况、恢复橄榄树的肥汁而欢喜。所以，你们这些因错误被折下的人，应当归回，再次被接上。那些仍在野橄榄中的人，应当脱离它的不结果子，而成为肥沃的分享者。\par

\SourceRecord{Translated by Richard Stothert.；Nicene and Post-Nicene Fathers, First Series；Vol. 4.；Edited by Philip Schaff.；Buffalo, NY: Christian Literature Publishing Co.,；1887.；<http://www.newadvent.org/fathers/140609.htm>.}

% structure: p0001 source=h1 target=section reason=page-title

\section{\WorkTitle{驳福斯图斯 卷十}}

\Emph{福斯图斯坚称，旧约的许诺与新约的许诺根本不同。奥斯定承认存在差异，但坚持伦理规诫在两约中都是一样的。}\par

1. \Person{福斯图斯}说：\Quote{另一个不接受旧约的原因是，旧约和新约都教导我们不要贪恋他人的财物。旧约中的一切皆是如此。它许诺财富、丰饶、子女、子孙、长寿，还有客纳罕地；但只给受割损者、守安息日者、献祭者、戒食猪肉者。如今，我如同每位基督徒一样，不理会这些事物，视之为微不足道且无益于灵魂的救恩。因此，我断定这些许诺不属于我。并且，我记着\BibleQuote{不可贪恋}的诫命，欣然将犹太人的财产留给他们，以福音和天国光明的继承为满足。倘若一个犹太人声称有分于福音，我必合理责备他妄称无权之物，因他不服从福音的规诫。同样，犹太人也可对我说同样的话，若我自称接受旧约却无视其要求。}\par

2. \Person{奥斯定}答复：\Quote{福斯图斯不厌其烦地重复同样的胡言乱语。但反复回答同样的话也令人厌倦，尽管这是重复真理。福斯图斯在这里说的，我们已经答复过了。但如果一个犹太人问我，为什么我自称相信旧约却不遵守其规诫，我的答复如下：法律的伦理规诫为基督徒所遵守；象征性规诫在预表的事物如今被启示的时代曾被恰当遵守。因此，那些我视为不再有约束力的仪式，我仍然视之为见证，正如那些肉身的许诺，旧约由此得名。因为虽然福音教导我盼望永福，我也在那些\BibleQuote{发生在他们身上作为鉴戒，并且被写下来，是为劝戒我们这些末世的人}的事中，找到了福音的证实。这就是我们给犹太人的答复。现在我们要对摩尼教人说些什么。}\par

3. 通过展示我们看待旧约权威的方式，我们答复了犹太人——福斯图斯以为犹太人的问题（我们为何不遵守规诫）会难倒我们。但你们（摩尼教人）会如何回答这个问题：你们为何欺骗单纯之人，自称相信新约，而你们不仅不信，反而全力攻击它？回答这个对你们来说，比我们回答犹太人更难。我们认为旧约中所写的一切都是真实的，并且由天主在适当的时候所命令的。但你们因为无法找到不接受新约所写内容的理由，最后只得假装那些段落不是真实的。这是一个在被真理抓住的异端者的最后挣扎；或者毋宁说是腐朽本身的气息。然而，福斯图斯承认旧约和新约都教导他不可贪恋。他自己的天主绝不可能教导他这一点。因为倘若这个天主不贪恋他人的所有，他为什么要在黑暗区域建造新世界？也许黑暗种族先贪恋他的王国。但这就模仿了他们的坏榜样。也许光明王国先前范围很小，战争是可取的，以便通过征服扩大它。那样的话，无疑就有了贪恋，尽管敌对种族被允许发动战争以证明征服的正当性。如果没有这种欲望，就没有必要将王国扩展到旧界限之外，进入被征服仇敌的区域。如果摩尼教人能从这些圣经中学习伦理规诫，其中一条是\BibleQuote{不可贪恋}，而不是因象征性规诫而反感，他们就会以温顺和真诚承认这些规诫适应了当时的时代。当我们读到旧约中那些\BibleQuote{发生在他们身上作为鉴戒，并且被写下来，是为劝戒我们这些末世的人}时，我们并没有贪恋他人的所有。当一个人读到为他利益而写的东西时，这当然不是贪恋。\par

\SourceRecord{Translated by Richard Stothert.；Nicene and Post-Nicene Fathers, First Series；Vol. 4.；Edited by Philip Schaff.；Buffalo, NY: Christian Literature Publishing Co.,；1887.；<http://www.newadvent.org/fathers/140610.htm>.}

% structure: p0001 source=h1 target=section reason=page-title

\section{驳斥福斯图斯，第十一卷}

\Emph{福斯图斯引用经文以证明宗徒保禄放弃了先前所持的降生成人信仰。奥斯定则表明宗徒在所引经文中保持了一致。}\par

1. \Emph{福斯图斯}说：我确实相信宗徒。但我并不相信天主子按肉身是出自达味的后裔\BibleRef{罗 1:3}，因为我不相信天主的宗徒会自相矛盾，对吾主耶稣时而持此见，时而持彼见。但是，假如他写了这话——既然你们不愿承认他的著作中有任何伪作——那也并不反对我们。因为这似乎是保禄对耶稣的旧有信念，当时他像众人一样，以为耶稣是达味之子。后来，当他得知这是错误的，便自行改正；在致格林多人书中，他说：\Quote{我们不再按人的看法认识谁了；纵使我们曾按人的看法认识过基督，但如今不再这样认识祂了。}\BibleRef{格后 5:16} 请留意这两节经文的差异。一节声称耶稣按肉身是达味之子；另一节则说如今不再按人的看法认识谁了。如果保禄两者都写了，那只能如我所说。接着他又说：\Quote{所以，谁若在基督内，他就是一个新受造物；旧的已成过去，看，都成了新的。}相信耶稣按肉身出自达味后裔，就是这种旧而可逝的东西；而那不再按人的看法认识谁的信德，则是新而永久的。因此，他在别处说：\Quote{当我是孩子的时候，说话像孩子，理解像孩子，思想像孩子；但一成了人，我就放下了孩子的事。}\BibleRef{格前 13:11} 所以我们有理由偏爱保禄那新的、改正过的宣认，而非他旧的、有缺陷的宣认。如果你们坚持罗马书所说，为何我们不能坚持格林多书所说？但正是由于你们坚持经文无误，才使我们以为保禄又建起他所拆毁的，而他本人却否认这种反复。如果那句话出自保禄，他已自行改正。如果保禄不应被认为写了任何需要改正的话，那么那句话就不是他的。\par

2. \Emph{奥斯定}回答说：正如我稍早所说，当这些人被圣经的明确证词围困、无法逃脱时，他们就宣称那段落是伪作。这宣告只显示他们对真理的厌恶和对错误的固执。既无法回答这些圣经的陈述，他们就否认其真实性。但如果这种回答被接纳或被赋予任何分量，那么引用任何书卷或任何段落来反对你们的谬误都将徒劳无功。拒绝书卷本身、不重视其权威是一回事——如同外教人拒绝我们的圣经、犹太人拒绝新约，以及我们拒绝你们教派或任何其他异端教派的特有书卷，以及所谓的伪经（其名称并非由于对它们的神秘尊重，而是因其来源神秘，缺乏明确证据，只靠模糊的推测）；而说\Quote{这位圣者只写了真理，这是他的书信，但有些句子是他的，有些不是}则是另一回事。然后，当你们被要求提供证据时，你们既不援引更正确或更古老的抄本、也不援引更多数量的抄本、也不援引原文，而是回答说：\Quote{这句是他的，因为它支持我；那句不是他的，因为它反对我。}那么，你们就是真理的准则吗？凡是反对你们的就不能是真理吗？但如果一个与你们一样疯狂的反对者用同样方式反驳你们说：\Quote{支持你们的那段是伪作，反对你们的那段是真实的}，你们又能如何回答？也许你们会拿出一个书卷，其全部内容都能解释为支持你们。那么，对方就不会只拒绝一段，而是谴责整个书卷是伪作。面对这样的反对者，你们毫无办法。因为你们能为书卷提出的所有来自古老或传统的见证都毫无用处。在这方面，天主教会（由宗徒原始席位延续至今的主教传承以及如此多民族的共识所支持）的见证是显著的。因此，如果对某段经文文本有疑问——正如学者们所熟知，有些经文存在不同读法——我们应首先查阅该宗教最初传授之地的抄本；如果仍有差异，则应采取数量较多或更古老的抄本。若仍有不确定，则查阅原文。凡是在关于圣经的问题上不忽视其权威、以求知而非争辩为目的而探究的人，都使用这样的方法。\par

3. 关于保禄书信中那段反对你们\Emph{异端}的教导，即天主子是由达味的后裔所生，这在所有教会的新旧手抄本和所有语言中都能找到。因此，福斯图斯自称相信宗徒，这实为虚伪。他若不是存心欺骗人，本该说\Quote{我确实不相信}，而不是\Quote{我确实相信}。他不从保禄获知第一个人，关于此人宗徒说，他出于地，属于土；又说：\BibleQuote{第一个亚当成了有生命的生灵。}福斯图斯的第一个人既不出于地、不属于土，也不是有生命的灵魂，而是出于天主的实体，与天主同体；据说这个存有将自己的肢体、外衣或武器，即五大元素，也与黑暗种族混合，这些也都是天主实体的一部分，以致它们被束缚并受污染。他也不从保禄获知第二个人，关于此人保禄说，祂来自天上，祂是末后的亚当，是赋予生命的灵；并且祂按肉身是出于达味的后裔，生于女人，生于法律之下，为要救赎那些在法律之下的人。\BibleRef{迦 4:4}关于祂，保禄对弟茂德说：\BibleQuote{你要记住：耶稣基督，达味的后裔，从死人中复活了，符合我的福音。}\BibleRef{弟后 2:8}他又用这复活作为我们复活的榜样：\BibleQuote{我当天把我所领受的，首先传给你们，就是基督照圣经记载的，为我们的罪死了，被埋葬了，又照圣经记载的，第三天复活了。}稍后，他从这教义引出推论：\BibleQuote{如果基督被传报是从死人中复活了，你们中间怎么有些人说死人没有复活呢？}我们这位自称相信保禄的人，对这些一概不信。他否认耶稣由达味的后裔所生，否认祂生于女人（按\Quote{女人}一词并非通常所指的妻子，而只是女性的性别，正如\BibleBook{}{创世纪}中所说的，天主在将女人带到亚当之前造成了女人\BibleRef{创 2:22}）；他否认祂的死亡、埋葬和复活。他主张基督没有可朽的身体，因此不能真正死亡；并且祂复活后向门徒显现时所显示给他们的伤痕的记号，保禄也提到过这些，\BibleRef{格前 11:5}都不是真实的。他也否认我们的可朽身体将复活并被改变为属神身体；正如保禄教导：\BibleQuote{播种的是自然身体，复活的是属神身体。}为说明自然身体和属神身体的区别，宗徒接着举了我已经引用的关于第一个和末后的亚当的例子。然后他继续说：\BibleQuote{弟兄们，我告诉你们，血肉不能继承天主的国。}为解释他所谓血肉是什么意思，即不是身体实体，而是朽坏，不会进入义人的复活，他立刻说：\BibleQuote{朽坏也不能继承不朽。}为了避免有人仍然以为不是埋葬的将要复活，而是像脱去一件衣服，换上一件更好的，他接着明确表明同一身体将改变得更好，正如基督在山上的衣服没有移位，而是变了形：\BibleQuote{看，我告诉你们一个奥秘：我们众人不全然改变，但我们众人必要复活。}然后他说明谁将改变：\BibleQuote{霎时间，一眨眼，在最后喇叭声中，因为喇叭一响，死人必要复活成为不朽，我们必要改变。}若有人说所改变的并非我们的可朽坏身体，而是我们的灵魂，那么应该注意，宗徒所谈的不是灵魂，而是身体，这从他开始的问题可以明显看出：\BibleQuote{但有人会说：\InnerQuote{死人怎么复活呢？他们将带着什么身体来呢？}}同样，在他辩论的结论中，他对所谈的内容毫不含糊：\BibleQuote{这必朽坏的必须穿上不朽，这必死的必须穿上不死。}\BibleRef{格前 15:35}福斯图斯否认这一点；他把保禄宣称为\BibleQuote{不朽、不可朽坏、独享光荣和尊威}\BibleRef{弟前 1:17}的天主变得可朽坏。因为在他们这骇人听闻的可怕虚构中，天主的实体和本性有被黑暗种族完全污染的威胁，并且为拯救其余部分，有一部分实际已受污染。而最糟的是，他试图欺骗那些不熟悉圣经的无知者，自称\Quote{我确实相信宗徒保禄}，而他本应说\Quote{我确实不相信}。\par

4. 但福斯图斯有一个证据，表明保禄改变了主意，并且在写给格林多人的信中纠正了他写给罗马人的内容；或者表明他从未写过那篇归于他名下的经文，即关于耶稣基督按肉身是由达味的后裔所生。这证据是什么？他说，如果\BibleBook{}{罗马书}中的那段话\BibleQuote{天主子，按肉身是由达味的后裔所成的}是真实的，那么他对格林多人所说的话就不可能是真实的：\BibleQuote{从今以后，我们不再按血肉认识人了；即使我们曾按血肉认识基督，但如今不再这样认识祂了。}因此我们必须证明这两段话都是真实的，而且并不彼此矛盾。手抄本的一致证明两者都是真作。在一些拉丁文版本中，用\Quote{生}代替了\Quote{成}，虽然这不是直译，但含义相同。因为这两种译文以及原文都教导基督按肉身是达味的后裔。我们片刻也不能以为保禄因意见改变而纠正了自己。福斯图斯本人也感到这种解释的不恰当和不虔敬，因此他宁愿说那段经文是伪造的，而不愿说保禄错了。\par

5. 至于我们的著作，它们并非信德或实践的准则，而只是对建树神益的帮助，我们可以设想，其中有些内容在晦涩难解的问题上不完全符合真理，这些错误可能会或不会在后续的论著中得到纠正。因为我们正是宗徒所说的那些人：\Quote{若你们持别有所想，天主也要将这事启示给你们。}\BibleRef{斐 3:15} 这样的著作带着判断的权利被阅读，没有任何信服的义务。为了给这些对疑难问题的有益讨论留出空间，在宗徒时代之后的所有作品与新旧约具有权威的正典经书之间有一条分明的界限。这些经书的权威是通过主教的传承和教会的扩展，从宗徒们传到我们这里，并以其崇高的至上地位，要求每个忠实和虔诚的心灵都顺服。如果我们对圣经中明显的矛盾感到困惑，不能说这本书的作者犯了错误；而应说要么是抄本有误，要么是译文错误，要么是你没有理解。在后世无数的著作中，我们有时也能找到与圣经相同的真理，但却没有同等的权威。圣经具有其独特的圣洁。在其他书中，读者可以形成自己的看法，也许由于不理解作者而与他产生分歧，可以赞成自己喜欢的东西，或反对自己不喜欢的东西。在这种情况下，人可自由地不予相信，除非有明确的论证或某些正典的权威表明该教义或陈述要么必须是真的，要么可能是真的。但由于圣经独特的特性，我们有义务将正典所显示的即使是某一位先知、宗徒或福音传道者所说的话，都当作真理接受。否则，若对正典经书健康权威的轻视导致其权威完全消失或陷入无可救药的混乱，那么就没有一页纸能用来引导人类的软弱了。\par

6. 那么，关于这段谈及天主圣子出自\Person{达味}的后裔的经文与\Quote{纵使我们曾按肉体认识过基督，但如今不再这样认识祂了}这句话之间的明显矛盾，即使这两处引文并非出自同一位宗徒的著作——一处出自\Person{保禄}，另一处出自\Person{伯多禄}、\Person{依撒意亚}或任何其他宗徒或先知——正典的权威是平等的，不容对其中任何一处有所怀疑。因为圣经的言辞，如同出自一人之口那样和谐一致，能赢得最精确、最清晰的虔敬之心的相信，并需要最冷静的理智和最巧妙的探究来发现和确认。在我们所讨论的案例中，两处引文都出自正典，即\Person{保禄}的确实书信。我们不能说抄本有误，因为最好的拉丁文译本基本一致；也不能说译文有误，因为最好的文本有相同的读法。因此，任何人若因表面的矛盾而感到困惑，唯一的结论就是他没有理解。因此，我还有责任解释这两段话如何不仅不矛盾，反而能以同一正确的信仰准则加以调和。虔诚的探究者通过仔细审视，会发现所有困惑都被消除了。\par

7. 天主圣子由达味的后裔成为人，不仅\Person{保禄}在其他地方说过，圣经的其他部分也如此教导。至于\BibleQuote{我们纵然曾按血肉认识基督，但如今不再这样认识他了}这句话，上下文表明了宗徒的意思。在此处或其他地方，他以确定的希望看待我们未来的生命，仿佛已经现在拥有，这生命如今在我们复活的元首和中保，\Person{耶稣基督}这个人身上已经实现。这生命当然不是按血肉的，正如基督的生命现在不是按血肉的。因为宗徒在此所谓\Quote{血肉}不是指我们身体的实体——主复活后曾说：\BibleQuote{你们摸摸我，看看，因为鬼神没有肉和骨头，如同你们看见我有的。}\BibleRef{路 24:39}——而是指血肉的败坏和可死性，那时在我们内已不存在，如同现在在基督内不存在一样。宗徒在前面关于复活的段落中使用\Quote{血肉}一词指败坏：\BibleQuote{血肉不能承受天主的国，败坏也不能承受不朽坏的。}因此，在下一节所描述的事件之后：\BibleQuote{看，我告诉你们一个奥秘：我们众人都要复活，但不是众人都要改变。在顷刻，在一眨眼之间，在末次号角吹响时（因为号角将吹响），死人必要复活成为不朽的，我们也要改变。因为这可朽坏的必须穿上不朽坏的，这可死的必须穿上不死的。}\BibleRef{格前 15:50}——那么，血肉，就身体的实体而言，在改变之后，将不再有血肉，即不再有可死性的败坏；但就其本性而言，它仍是同一血肉，同一复活并改变的血肉。主复活后所说的\BibleQuote{你们摸摸我，看看；因为鬼神没有肉和骨头，如同你们看见我有的。}是真实的，宗徒所说的\BibleQuote{血肉不能承受天主的国}也是真实的。前者是针对身体的实体说的，它是改变的主体；后者是针对血肉的败坏说的，它将不复存在，因为改变之后，血肉将不再败坏。因此，\BibleQuote{我们曾按血肉认识基督}，即按肉身的可死性，在祂复活之前；\BibleQuote{但如今不再这样认识他了}，因为同一宗徒说：\BibleQuote{基督既从死者中复活，就不再死，死亡不再统治他了。}\BibleRef{罗 6:9}\BibleQuote{我们曾按血肉认识基督}这句话严格来说意味着基督曾是按血肉的，因为从未有过的事物不能被认识。且说的是\Quote{认识}，而不是\Quote{猜想}。但不必执着于一个词，以免有人说\Quote{认识}被用作\Quote{猜想}的意义。令人惊讶的是（如果人们会对盲人失明感到惊讶的话），这些瞎眼的人没有看出：如果宗徒所说的不再按血肉认识基督证明基督没有血肉，那么他在同一处所说的从此不再按血肉认识任何人，就证明这里提到的所有人都没有血肉。因为当他谈到不认识任何人时，他不会只想说基督；而是在意识到那将要与那些在复活时改变的人共同享有的未来生命时，他说：\BibleQuote{从此我们不再按血肉认识任何人；}意即，我们对未来不朽和不死的希望如此确定，以至于这想法甚至现在已使我们喜乐。所以他在别处说：\BibleQuote{你们既然与基督一同复活了，就该追求上面的事，那里有基督坐在天主的右边。你们该思念上面的事，不该思念地上的事。}\BibleRef{哥 3:1}的确，我们尚未像基督那样复活，但因着我们在祂内的希望，我们被称为与祂一同复活了。同样他又说：\BibleQuote{祂按照祂的仁慈，借着重生的洗，救了我们。}\BibleRef{铎 3:5}显然，我们在重生的洗中所获得的不是救恩本身，而是它的希望。然而，因为这希望是确定的，我们被称为得救了，仿佛救恩已经赐下。别处明确地说：\BibleQuote{我们自己心中叹息，等待着义子的身份，就是我们肉身的救赎。因为我们得救是在于希望。然而看得见的希望不是希望，因为一个人看见了还希望什么呢？但如果我们希望那看不见的，我们就耐心等待。}\BibleRef{罗 8:23}宗徒不是说\Quote{我们将要得救}，而是说\Quote{我们现在得救了}，即在于希望，虽然尚未在现实中。同样，我们如今不再按血肉认识任何人，也是在于希望，虽然尚未在现实中。这希望是在基督内，在祂内我们所希望于我们的应许已经实现了。祂复活了，死亡不再统治祂。虽然我们曾按血肉认识祂，即在祂死前，当祂的身体有宗徒恰当地称为血肉的那种可死性时，但如今不再这样认识祂了；因为祂那可死的已穿上了不死，祂的血肉，即可死性，已不复存在。\par

8. 包含这段子句的上下文揭示了其真实含义，而我们的对手却对此加以滥用。我们读到：\BibleQuote{基督的爱催迫我们，因为我们如此断定：如果一个人替众人死了，那么众人就都死了；他替众人死了，是为叫活着的人不再为自己生活，而是为替他们死而复活的那位生活。因此，从今以后，我们不再按血肉认人了；虽然我们曾按血肉认识基督，但如今不再这样认识他了。} 这些话——\BibleQuote{叫活着的人不再为自己生活，而是为替他们死而复活的那位生活}——清楚地表明，\Person{保禄}宗徒的陈述是以基督的复活为根据的。不再为自己生活，而是为祂生活，这必然意味着不再按血肉生活，即不寄望于地上可朽坏的事物，而是按神生活，寄望于复活——这复活已在基督内成就。因此，对于基督为之死而复活的人们，他们从此不再为自己生活，而是为祂生活，宗徒说，他不再按血肉认识任何人，因为他们所期盼的是未来的不朽——这望德在基督内已是事实。所以，虽然他曾按血肉认识基督，在祂死之前，如今他却不再这样认识祂；因为他知道祂已复活，死亡不再统治祂。并且，由于在基督内，我们现在虽是望德中（而非现实中）拥有基督所是的一切，他又说：\BibleQuote{所以，谁若在基督内，他就是新受造物：旧的已成过去，看，都成了新的。一切都是出于天主，祂藉着基督使我们与自己和好。} \BibleRef{格后 5:14} 这新受造物——即藉着信德更新的子民——对自己所盼望的，已经在基督内拥有了；而这望德将来也会实际实现。就这望德而言，旧的已经过去，因为我们不再处在旧约的时代，不再期望一个属世、属血肉的天主之国；一切都成了新的，使天国的应许——那里没有死亡和败坏——成为我们信心的根基。但在死者复活时，不是作为望德，而是作为现实，旧的将过去，那时最后的仇敌——死亡——将被消灭；当这可朽坏的穿上不朽坏，这可死的穿上不死时，一切都成了新的。这在基督内已经实现了，因此\Person{保禄}确实不再按血肉认识祂。但对于基督为之死而复活的人们，他并非在现实中，而只是在望德中不再按血肉认识他们。因为，正如他对\Place{厄弗所}人说的：我们已经得救，是藉着恩宠。整段经文都是这个意思：\BibleQuote{但是，富于慈悲的天主，因着祂爱我们的大爱，竟在我们因过犯而死亡的时候，使我们同基督一起生活——藉着恩宠你们得救了——} \BibleQuote{使我们同基督一起生活}这些词，与他给\Place{格林多}人所说的\BibleQuote{叫活着的人不再为自己生活，而是为替他们死而复活的那位生活}相对应。而在\BibleQuote{藉着恩宠你们得救了}这些话中，他把所盼望的事说成已经成就。所以，在上面引用的段落中，他明确说：\BibleQuote{我们得救是在于望德}。然后他继续讲未来的事，仿佛已经成就：\BibleQuote{并且使我们一同复活，}他说，\BibleQuote{一同坐在天上，在基督耶稣内。} \Person{基督}当然已经坐在天上，但我们还没有。但由于那确定的望德，我们已经拥有了未来，所以他说我们坐在天上，不是在我们自己内，而是在祂内。为了表明这仍然是未来的事，免得有人认为那在望德中已经成就的事在现实中已经实现了，他又加上：\BibleQuote{为在将来的世代中，显示祂在基督耶稣内对我们所施的恩宠，超越一切的丰富。} \BibleRef{弗 2:4} 同样，我们也必须理解下面这段经文：\BibleQuote{因为我们以前属血肉的时候，那藉着法律而活动的罪恶情欲，在我们的肢体中发动，以致结出死亡的果实。} \BibleRef{罗 7:5} 他说\BibleQuote{当我们还在血肉中时}，仿佛他们已不在血肉中。他的意思是：当我们还在寄望于血肉之事的时候，这指的是法律（只有藉着属神的爱德才能成全）仍有效力的时期，为使因违命而罪过增多，好在新约启示之后，恩宠和藉恩宠的恩赐更丰富地增多。同样，他在别处说：\BibleQuote{凡属血肉的人，不能中悦天主；}然后，为了表明他指的不是尚未死去的人，又加上：\BibleQuote{但你们已不属于血肉，而是属于圣神。} \BibleRef{罗 8:8} 意思是：那些寄望于血肉之福的人不能中悦天主；但你们不是寄望于血肉之事，而是寄望于属神的事，即天国，那里如今是属生灵的身体，将因复活的变化，按本性的能力，成为属神的身体。因为 \BibleQuote{播种的是属生灵的身体，复活起来的是属神的身体。} 由此可见，如果宗徒对那些被称为不在血肉中的人（因为他们不寄望于血肉之事，尽管他们仍背负着可朽坏且必死的血肉）不再按血肉认识他们；那么，对于基督，他更能说不再按血肉认识祂，因为在基督的身体里，他们所盼望的已经成就了！诚然，与其在遇到超出我们微弱理智能力所能解决的困难时，接受一些经文为真而谴责另一些为假，不如查考圣经各段以发现它们彼此一致，更为妥善和虔敬。至于宗徒幼年时像孩子一样了解，这不过是用来作比喻。\BibleRef{格前 13:11} 当他还是孩子时，他不是属神的人，不像他后来为建立教会而写那些著作时那样——那些著作不同于其他书，不仅是可获益的研读，而且作为教会正典的一部分，具有权威性要求我们相信。\par

\SourceRecord{Translated by Richard Stothert.；Nicene and Post-Nicene Fathers, First Series；Vol. 4.；Edited by Philip Schaff.；Buffalo, NY: Christian Literature Publishing Co.,；1887.；<http://www.newadvent.org/fathers/140611.htm>.}

% structure: p0001 source=h1 target=section reason=page-title

\section{驳福斯德书，卷十二}

\Emph{福斯德否认先知预言了基督。奥斯定从新约证明这种预言，并详细阐述旧约中基督的主要预像。}\par

1. \Person{福斯德}说：为什么我不相信先知？更准确地说，为什么你们相信他们？你们会回答，因为他们关于基督的预言。就我而言，我曾以最热切的态度阅读先知书，却找不到这样的预言。而若没有证据和见证就不相信基督，这显然是一种软弱的信德。确实，你们自己惯于教导说，基督信仰如此单纯和绝对，不容许劳心费力的探究。那么，为什么你们要用证据、而且是犹太人的证据来破坏信德的单纯呢？或者，如果你们改变了对证据的看法，那么还有什么见证能比天主亲自为祂的圣子作证更可信呢？当祂派遣圣子来到世上时，并非通过先知或解释者，而是直接从天上发出的声音说：\BibleQuote{这是我的爱子，你们要听从他！}（\BibleRef{玛 3:17}）祂又为自己作证说：\BibleQuote{我从父出来，来到了世界}（\BibleRef{若 16:28}），以及许多类似的经文。当犹太人质疑这个见证，说\BibleQuote{你为自己作证，你的证据不足凭信}时，祂回答说：\BibleQuote{我即使为自己作证，我的证据也是可信的；……你们的法律上记载：两个人的见证是可凭信的。我为自己作证，并且派遣我来的父也为我作证。}（\BibleRef{若 8:13}）祂没有提及先知。祂又引用自己工程的见证说：\BibleQuote{如果你们不信我，至少该信这些工程。}（\BibleRef{若 10:38}）而不是\BibleQuote{如果你们不信我，至少该信先知}。因此，关于我们的救主，我们不需要任何见证。我们在先知身上所寻求的只是明智、德行和善表，而你们很清楚，这些在犹太人的先知身上是找不到的。毫无疑问，这解释了为什么你们立刻引述他们的预言作为相信他们的理由，而对他们的行为却只字不提。这也许是好的策略，但与圣经的声明不符，即不能从荆棘中摘葡萄，也不能从蒺藜中摘无花果。以上可以作为对这个问题的简短而充分的答复，即我们为什么不信先知。先知的著作中没有预言基督，这一点在我们父辈的著作中已有充分证明。我只补充一点：如果希伯来先知们认识并宣讲了基督，却过着如此邪恶的生活，那么保禄论及外邦智者的话就可以应用在他们身上：\BibleQuote{他们虽然认识天主，却没有以祂为天主而光荣祂，也不感谢祂；反而在妄念中变得愚昧，他们愚昧的心灵就昏暗了。}（\BibleRef{罗 1:21}）你们看，认识伟大的事物是没什么价值的，除非生活与之相称。\par

2. \Person{奥斯定}回答说：这一切的意思是，希伯来先知们没有预言基督的任何事情，而且，如果他们确实预言了，他们的预言对我们也没有用处，并且他们自己的生活与这种预言的尊严也不相称。因此，我们必须证明预言的事实；证明它们对于我们信仰的真实性和坚定性是有用的；并且证明先知们的生活与他们的言论是一致的。在这三重讨论中，如果要从所有书中引用那些可以显示基督被预言的段落，在第一个标题下会花费很长时间。\Person{福斯德}的轻浮可以被一个伟大权威的分量有效地击败。虽然\Person{福斯德}不相信先知，但他声称相信宗徒。上文他似乎为了满足某个对手的怀疑，宣称他确信\Person{保禄}宗徒。那么，让我们听听\Person{保禄}关于先知说了什么。他的话是：\BibleQuote{\Person{保禄}，\Person{耶稣基督}的仆人，蒙召作宗徒，被选拔去传天主的福音，这福音是天主先前藉着祂的先知在圣经上所预许的，论及祂的圣子，祂按肉身是出自\Person{达味}的后裔。}（\BibleRef{罗 1:1}）\Person{福斯德}还想怎样？他会坚持说，宗徒说的是别的先知，而不是希伯来先知吗？无论如何，所预许的福音是关于天主子，祂按肉身是出自\Person{达味}的后裔为祂所生的；宗徒说他就是为此被选拔的。因此，摩尼教异端与信仰福音是相对的，福音教导天主子按肉身出自\Person{达味}的后裔。此外，还有许多地方，宗徒明确地为希伯来先知作证，其权威足以折服这些骄傲的摩尼教徒的颈项。\par

3. 宗徒说：\BibleQuote{我在基督内说实话，并不说谎，我的良心在圣神内与我一同作证：我大有忧愁，心中常有伤痛。因为为我弟兄，我骨肉之亲，就是被咒诅，与基督分离，我也愿意。他们是以色列人；那里有义子的名分、荣耀、盟约、法律的颁布、礼仪和恩许；那里有圣祖，并且按肉身说，基督也是从他们出来的，祂是在万有之上，永远可赞美的天主。}\BibleRef{罗 9:1}这里有最丰富、最明确的证据和最隆重的称赞。这里所说的义子的名分，显然是通过天主子而来的；正如宗徒对迦拉达人所说的：\BibleQuote{时期一满，天主就派遣了自己的儿子来，生于女人，生于法律之下，为把在法律之下的人赎出来，使我们获得义子的地位。}\BibleRef{迦 4:4}这里所说的荣耀，主要是指他在致罗马人同书中所说的：\BibleQuote{那么，犹太人有什么长处呢？割损又有什么益处呢？很多，各方面都有：首先，因为天主的圣言交托了他们。}\BibleRef{罗 3:1}摩尼教徒能告诉我们，除了希伯来先知的圣言之外，还有什么天主的圣言交托给了犹太人？为什么盟约特别属于以色列人？不就是因为不仅旧约给了他们，而且新约也在旧约中预表了吗？我们的对手常常在攻击那赐给以色列人的法律制度时表现出无知的凶猛，他们不明白天主愿意我们不在法律之下，而在恩宠之下。宗徒本人在这里回答了这一点，他在谈论犹太人的长处时，提及他们领受了法律。如果法律是坏的，宗徒就不会在称赞犹太人时提到它。如果法律没有宣讲基督，主自己就不会说：\BibleQuote{如果你们相信梅瑟，你们必会相信我，因为他是指着我写的}\BibleRef{若 5:46}；祂复活后也不会作证说：\BibleQuote{凡梅瑟法律、先知并圣咏上指着我所记载的话，都必须应验。}\BibleRef{路 24:44}\par

4. 但是，因为摩尼教徒宣讲另一位基督，不是宗徒们所宣讲的那一位，而是他们自己虚假捏造的假基督，效法这虚假的谎言，他们自己说谎言，虽然他们或许会被相信，因为他们不羞于自称是欺骗者的追随者，宗徒对不信的犹太人所断言的事就临到了他们：\BibleQuote{几时诵读梅瑟，一片帕子就盖在他们的心上。}若他们不转向基督，这一片帕子不会从他们那里除去，以使他们明白梅瑟；不是他们自己所造的基督，而是希伯来先知的基督。因为正如宗徒所说：\BibleQuote{你们几时转向主，帕子就会除掉。}\BibleRef{格后 3:15}我们不能奇怪他们不相信那从死者中复活的基督，祂曾说：\BibleQuote{凡梅瑟法律、先知并圣咏上指着我所记载的话，都必须应验}；因为基督亲自告诉我们，亚巴郎对那在痛苦中在地狱里折磨的硬心财主说过的话，那财主请求亚巴郎派一个人到他弟兄那里教导他们，免他们也来到这痛苦的地方。亚巴郎的回答是：\BibleQuote{他们有梅瑟及先知，听听他们吧。}当财主说，除非有人从死者中复活，他们不会相信时，他得到了这最真实的回答：\BibleQuote{如果他们不听从梅瑟及先知，纵使有人从死者中复活了，他们也不会信服的。}\BibleRef{路 16:27}因此，摩尼教徒不听梅瑟及先知，所以他们不相信基督，虽然祂从死者中复活了。事实上，他们甚至不相信基督从死者中复活了。因为他们怎能相信祂复活了，当他们不相信祂死了？再者，他们怎能相信祂死了，当他们否认祂有可朽坏的身体？\par

5. 我们却拒绝那些假教师，他们的基督是假的，或者不如说，他们的基督从未存在过。因为我们有一位真实而真诚的基督，由先知预言，由宗徒宣讲，他们在无数地方引用法律和先知的见证来支持他们的宣讲。保禄在一句简短的话中，给了我们关于这主题的正确观点。他说：\BibleQuote{但如今，天主的正义在法律之外已显示出来，为法律和先知所作证。}\BibleRef{罗 3:21}什么先知？如果不是以色列的先知，正如他清楚所说，盟约、法律的颁布和恩许都属于他们？又是什么恩许？岂不是关于基督的吗？他在别处论到基督时，简洁地说：\BibleQuote{天主的恩许，在祂内都是\InnerQuote{是}。}\BibleRef{格后 1:20}保禄告诉我，法律的颁布属于以色列人。他也告诉我，基督是法律的终向，使凡相信的人都获得正义。他还告诉我，天主的全般恩许在基督内都是\InnerQuote{是}。而你却告诉我，以色列的先知们没有预言任何关于基督的事。我要相信摩尼所讲述的、与保禄相反的虚妄冗长的寓言吗？还是要相信保禄的警告：\BibleQuote{即使我们，或是从天上降下的一位天使，给你们宣讲的福音，与我们先前给你们所宣讲的福音不同，当受诅咒。}\BibleRef{迦 1:8}\par

6. 我们的对手或许会要求我们指出以色列先知预言基督的经文。人们会以为，他们应当满意于宗徒的权威，因为宗徒宣布，我们在希伯来先知书中所读到的内容已在基督身上应验；或者满意于基督自己的话，祂说这些事是论及祂写的。凡不能指出这些经文的人，应将过错归咎于自己的无知；因为宗徒、基督和圣经绝不能指控为虚假。然而，可以从众多例子中举出一个。宗徒在上述引文之后的经文中说：\BibleQuote{天主的言语不能落空。因为并非所有以色列人都属于以色列；也不因为是亚巴郎的后裔，就都是子女；而是说：\InnerQuote{由依撒格所生的，才称为你的后裔。}这就是说：肉身的子女并非天主的子女，惟有蒙恩许的子女才算是后裔。} \BibleRef{罗 9:6} 面对向亚巴郎所作的声明：\BibleQuote{地上万民都要因你的后裔蒙受祝福}，我们的对手对此能说什么呢？当时，宗徒对这一恩许作了如下阐释：\BibleQuote{恩许原是向亚巴郎和他的后裔说的，并没有说\InnerQuote{后裔们}，好似指多数，而是指一个，即你的后裔，就是基督。} \BibleRef{迦 3:16} 若在那时对此仍有怀疑，或许情有可原，因为那时万民尚未相信这位被传报为亚巴郎后裔的基督。但如今我们亲眼见到古代预言中的应验——万民确实在亚巴郎的后裔中蒙受了祝福，而这话是在数千年前对亚巴郎说的：\BibleQuote{地上万民都要因你的后裔蒙受祝福}——此时仍试图引进另一位非亚巴郎后裔的基督，或认为亚巴郎子孙的先知书中毫无基督的预言，那便是顽固的愚妄了。\par

7. 若要列举希伯来先知书中所有论及我们的主和救主耶稣基督的经文，那将超出整卷书的篇幅，更不必说这部以简短答复为体裁的论著了。这些圣经的全部内容，或直接或间接地都与基督有关。许多地方采用寓意或隐语，也许是词语上的暗示，也许是历史叙述中的预表，需要读者勤勉探究，并以发现的乐趣作为奖赏。另一些经文则明白晓畅；因为若没有清楚的部分，我们就无法理解晦涩之处。而那些比喻性的段落，若汇聚一处，将会发现它们对基督的见证如此和谐，足以使怀疑者的迟钝蒙羞。\par

8. 在 \BibleBook{创世}{纪} 中，天主用六天完成了祂的工程，并在第七天休息。世界的历史包含六个时期，以天主与人的交往为标志。第一个时期是从亚当到诺厄；第二个时期是从诺厄到亚巴郎；第三个时期是从亚巴郎到达味；第四个时期是从达味到巴比伦充军；第五个时期是从充军到我们的主耶稣基督卑微的来临；第六个时期正在进行，并将以荣耀的救主来临审判而结束。与第七天相对应的是圣徒的安息——不是在此生，而是在来世，在那里富翁看见拉匝禄安息而自己却在地狱受折磨；那里没有傍晚，因为没有朽坏。在 \BibleBook{创世}{纪} 中，人按天主的肖像被造；在世界第六时期，我们清楚地发现了我们按创造我们者的肖像更新心志而得以转变，如宗徒所说。\BibleRef{哥 3:10} 正如亚当沉睡时从他的肋旁造了妻子，教会也因她的救主死后自祂肋旁流出的血之圣事而成为祂的产业。从丈夫肋旁造出的女人被称为厄娃，即生命，众生之母；主在福音中说：\BibleQuote{你们若不吃人子的肉，不喝祂的血，在你们内便没有生命。} \BibleRef{若 6:53} \BibleBook{创世}{纪} 的全部叙述，即使是最细微的细节，都是关于基督和教会、或关于善与恶的基督徒的预言。宗徒称亚当为\BibleQuote{那要来之人的预像} \BibleRef{罗 5:14} 时，以及他说\BibleQuote{人要离开父亲和母亲，依附自己的妻子，两人成为一体。这奥秘真是伟大！但我是指基督和教会说的} \BibleRef{弗 5:31} 时，这些话具有深意。这极为明显地指向基督离开祂父亲的方式；因为\BibleQuote{祂虽具有天主的形体，却没有以自己与天主同等为应当把持不舍的，反而空虚自己，取了奴仆的形体} \BibleRef{斐 2:6}。同样，祂也离开了祂的母亲，即依恋旧约肉体性的犹太会堂，与祂神圣的新娘教会结合，使在新约的平安中两人成为一体。因为尽管祂与圣父同为天主，我们由祂所造，祂却在肉身上分享了我们的本性，好使我们成为祂的身体，祂是头。\par

9. 就如\Person{加音}献地里的出产而被拒绝，而\Person{亚伯尔}献羊群和脂油而被悦纳，同样，新约的信德在纯全的恩宠事奉中赞美天主，胜于旧约属世的规条。因为虽然犹太人遵行这些事是对的，但他们因未能分辨基督来临时新约时期与旧约时期而犯了不信的罪。天主对\Person{加音}说：\Quote{你若献得好，却分得不均，你就犯了罪。}如果加音听从天主的话，当天主说：\Quote{你要安分，因为它必归属于你，你将管辖它}，他就会归咎己罪，承担责任，并向天主认罪；如此，借恩宠的供应，他就能管辖自己的罪，而不是在杀害无辜的弟弟时作罪的奴仆。同样，犹太人（这一切都是他们的预像）若曾安分而不躁动，并承认藉恩宠赦罪的救恩时期，又听见基督说：\Quote{\BibleQuote{健康的人不需要医生，有病的人才需要；我来不是召义人，而是召罪人悔改}}（\BibleRef{玛 9:12}），以及\Quote{\BibleQuote{凡犯罪的，就是罪的奴仆}}；并说\Quote{\BibleQuote{如果儿子使你们自由，你们就真正自由了}}（\BibleRef{若 8:34}）——他们就会在认罪中归咎己罪，向良医说，正如圣咏所载：\Quote{我曾说：主，求祢怜悯我，医治我的灵魂，因为我得罪了祢。}并因恩宠的盼望而获得自由，他们就会在罪仍存于可朽身体期间管辖罪。但如今，他们不认识天主的正义，想要建立自己的正义，因法律的功业而骄傲，而不因自己的罪而谦卑，他们不曾安分；他们受制于在可朽身体内掌权的罪，以致顺从它的私欲，就绊倒在绊脚石上，并对那位他们见其作为被天主悦纳而心生嫉妒之人燃起仇恨。那个生来瞎眼却被开了眼的人对他们说：\Quote{\BibleQuote{我们知道天主不听罪人；但若有人事奉祂，遵行祂的旨意，天主才听从他}}（\BibleRef{若 9:31}）；好像他说，天主不看加音的祭献，但看中亚伯尔的祭献。年幼的弟弟亚伯尔被哥哥杀死；基督，年幼子民的首领，被年长的犹太人杀死。亚伯尔死在田野；基督死在髑髅地。\par

10. 天主问\Person{加音}他弟弟在哪里，并非因为祂不知道，而是如法官审问罪犯。\Person{加音}回答说他不知道，且他不是弟弟的看守者。如今，当我们用天主的声音，即圣经的声音，问犹太人关于基督的事时，除了说他们不知道我们所说的基督外，还能有什么回答呢？加音的无知是假装的，犹太人拒绝基督也是受蒙蔽。此外，他们如果愿意接受并持守基督徒信仰，在某种意义上就会是基督的看守者。因为凡在心里持守基督的，不会像加音一样问：\Quote{\BibleQuote{我是我弟弟的看守者吗？}}然后天主对\Person{加音}说：\Quote{\BibleQuote{你做了什么事？你弟弟的血从地里向我喊冤。}}同样，在圣经中天主的声音指控犹太人。因为基督的血在地上大有声音，当相信祂的人从万国响应\Quote{阿们}时。这就是基督的血的声音，因为被祂的血所救赎的信徒的清晰响声，就是这血本身的声音。\par

11. 然后天主对加音说：\Quote{你从土地受诅咒，这土地开口从你手接收了你兄弟的血。你耕种土地，土地不再为你出产它的力量。你在地上将成为哀伤和卑贱的人。}这并非说\Quote{土地受诅咒}，而是说\Quote{你从土地受诅咒，这土地开口从你手接收了你兄弟的血。}因此，不信的犹太民族从土地受诅咒，即从教会受诅咒；这教会因承认罪过而开口接收了为赦免罪过而流出的血，这血是由那不愿在恩宠之下、反在法律之下的民族之手流出的。这凶手被教会诅咒，即教会承认并宣告宗徒所发出的诅咒：\BibleQuote{凡出于法律的行为，都在法律的诅咒之下。}\BibleRef{迦 3:10}然后，在说了\Quote{你从土地受诅咒，这土地开口从你手接收了你兄弟的血}之后，接着的不是\Quote{你耕种它}，而是\Quote{你耕种土地，土地不再为你出产它的力量。}他所耕种的土地并不一定是那开口从他手接收他兄弟之血的土地。他从这土地受诅咒，因此他耕种的土地不再为他出产力量。也就是说，教会承认并宣告犹太民族受诅咒，因为他们在杀害基督之后，继续耕种属地的割礼、属地的安息日、属地的逾越节的土地，而耕种所隐藏的、使人认识基督的力量或德能，并不赐给犹太民族，因为他们继续处于不敬和不信中，这力量在新约中得以启示。当他们不愿归向天主时，在读旧约时蒙在他们心上的帕子就没有被除去。这帕子只被基督除去，基督并不废除旧约的诵读，而是除去掩盖其德能的遮蔽。因此，在基督被钉十字架时，圣所帐幔裂为两半，好藉基督的苦难，使那些以忏悔的口张开喝祂血的信徒得以认识隐藏的奥秘。这样，犹太民族如同加音，继续以法律的内修生活耕种土地，但土地不向他们出产力量，因为他们没有在其中认识基督的恩宠。同样，基督的肉身是土地，犹太人通过钉死祂为我们产生了救恩，因为祂为我们的过犯而死。但这土地没有向他们出产力量，因为他们没有因祂复活的德能而成为义，因为祂为我们的成义而复活。正如宗徒所说：\BibleQuote{祂因软弱被钉在十字架上，却因天主的德能而生活。}\BibleRef{格后 13:4}这就是那土地的力量，不敬和不信的人不认识它。基督复活时，没有显现给那些钉死祂的人。因此，加音不被允许看到他所耕种以撒种的土地的力量，正如天主所说：\Quote{你耕种土地，土地不再为你出产它的力量。}\par

12. \Quote{你在地上将哀叹和战栗。}这里没有人看不到，在犹太人所散居的每个地方，他们为失去王国而哀伤，并在数量远超他们的基督徒面前战栗恐惧。于是加音回答说：\Quote{我的惩罚太重，今日祢将我驱逐离开地面，我必从祢的面被隐藏，我在地上将成为哀伤和流离的人；凡遇见我的，必杀我。}在这里他确实因恐惧而哀叹，生怕失去地上的产业后还遭受身体的死亡。他认为这比土地不向他出产力量或属灵的死亡更糟。因为他的心思是属血肉的；他不以为被隐藏在天主的面之外，即处于天主的忿怒之下，是小事，只要不被遇见杀死。这就是耕种土地却不获得其力量的属血肉的心思。属血肉的心思就是死亡；但他对此无知，为失去地上的产业而哀伤，并恐惧肉身的死亡。但天主如何回答？\Quote{不然，}祂说，\Quote{凡杀死加音的，将受七倍的报复。}也就是说，并非如你所说，不敬的属血肉的犹太民族不会因身体的死亡而灭亡。因为谁若这样消灭他们，将遭受七倍的报复，即他将自己招致犹太人为钉死基督而承担七倍的惩罚。因此，直到七天时间结束，犹太人的持续保存向信道的基督徒证明，那些在王国骄傲中杀害主的人理当受到这样的束缚。\par

\Quote{上主天主给加音立了一个记号，免得任何人遇见他就杀死他。} 这是一件最显著的事实：所有被罗马征服的民族都采纳了罗马崇拜的异教仪式；而犹太民族，无论在异教或基督教君王治下，从未失去他们法律的记号，借此他们与所有其他民族人民区分开来。任何皇帝或君主，若在其治下发现了带有这记号的百姓，都不杀死他们，即，不使他们不再是犹太人，也不使他们不再在规范上与世人分离且不同。只有当犹太人皈依基督时，他就不再是加音，也不从天主面前离开，也不住在\OriginalTerm{诺得}{Nod}地。诺得据说意为\OriginalTerm{动乱}{commotion}。圣咏作者祈求反对这动乱之恶：\Quote{勿使我的脚动摇；} 又说：\Quote{不要让恶人的手挪移我；} 以及：\Quote{那些搅扰我的人，当我动摇时必欢喜；} 以及：\Quote{上主在我右边，使我不致动摇；} 等等无数处。这恶降临于那些离开天主面前的人，即离开祂的慈爱。因此圣咏作者说：\Quote{我曾在顺遂中说：我永不动摇。} 但请看下文：\Quote{上主，祢的恩宠加给了我尊荣的力量；祢掩面，我就惊惶；} 这教导我们：任何灵魂的美善、尊荣或力量，不是凭自身，而是借参与天主之光才能拥有。摩尼派应该思想这一点，以免他们犯亵渎之罪，把自己等同于天主的性体与实体。但他们不能思想，因为他们不满足。他们不认识内心的安息。如果他们满足，正如加音被告知要满足一样，他们就会把罪归给自己；即，他们会将错归咎于自己，而不是归咎于一个无人听过的黑暗种族，并藉着天主的恩宠胜过自己的罪。但现在摩尼派，以及所有以各种异端反对真理的人，都像加音和散居的犹太人一样，离开天主面前，居住在动乱之地，即肉体的不安，而不是享受天主，即不是伊甸。\OriginalTerm{伊甸}{Eden}被解释为\OriginalTerm{宴乐}{Feasting}，乐园就设在那里。但为了不过多偏离这论文的论点，我须限于在此标题下作一些简短的评论。\par

因此，省略这些书卷中许多可找到基督的段落（它们需要更长的解释和证明，尽管最隐晦的含义最甘甜），我们可以从列举下列事物中获得令人信服的证据：从亚当数起第七位的哈诺客中悦了天主，并得以转移，正如将有一个第七日的安息，所有在世界历史第六日被降生成人的圣言重新创造的人，都将在那里转移；诺厄与他的全家借水和木得了救，正如基督的家庭借代表十字架苦难的圣洗得了救；方舟是由方形木料造成的，正如教会由预备行各样善工的圣人建造：因为方形的每边都很稳固；方舟长六倍于宽，十倍于高，如同人的身体，显示基督取了人的身体；宽达到五十肘，如宗徒说：\BibleQuote{我们的心扩大了}，\BibleRef{格后 6:11} 即，以属神的爱，他又说：\BibleQuote{天主的爱藉着所赐给我们的圣神，倾注在我们心中}，\BibleRef{罗 5:5} 因为在复活后的第五十天，基督派遣了圣神来扩大祂门徒们的心；长三百肘，即六倍五十；正如世界历史有六个时期，基督从未停止被宣讲——五个时期由先知预表，第六个时期在福音中宣告；高三十肘，为长度的十分之一；因为基督是我们的高度，祂在三十岁时肯定了福音的教导，宣告祂来不是为废除法律，而是为成全法律。十诫是法律的核心；所以方舟的长度是十乘以三十。诺厄本人也是从亚当算起的第十代。方舟的木料内外涂上沥青，以象征紧凑的结合中爱的忍耐，这忍耐保持弟兄间的联系不被损害，和平的纽带不被来自教会内外的冒犯所破坏。因为沥青是一种胶粘的物质，具有极大的能量和力量，代表爱的热情，以极大的忍耐力维持属神共融中的一切。\par

各种动物都进入方舟；正如教会包含万民，这也在给伯多禄看的器皿中预表了。洁净与不洁净的动物都在方舟里；正如好和坏的人都参与教会的圣事。洁净的取七对，不洁净的取一对；不是因为坏的比好的少，而是因为善人保持心神的合一于和平的纽带；而圣经论及圣神说有七重作用，即\BibleQuote{智慧和聪敏的神，超见和刚毅的神，明达和孝爱以及敬畏上主的神}\BibleRef{依 11:2}。同样，与圣神降临相关的数五十，是由七乘七加一组成；因此经上说：\BibleQuote{努力以和平的联系，保持圣神的合一}\BibleRef{弗 4:3}。反之，坏的取一对，因为它们易于分裂，倾向于分裂。诺厄连同其全家是第八位；因为我们的复活希望已在基督身上显现，祂在第八日（即安息日之后的第二天）从死者中复活。这日是祂受难后的第三天；但在通常计算中，它既是第八日，也是第一日。\par

16. 整只方舟在上面合为一肘，正如教会——基督的身体——被聚集在合一中，得以成全。因此基督在福音中说：\Quote{不与我聚集的，就是分散的。}\BibleRef{玛 12:30} 入口在侧边，正如没有人能进入教会，除非借着从基督被刺的肋旁流出的罪赦圣事。方舟的下层分为两层和三层房间：正如教会中万民的群体分为两种，即受割礼的和未受割礼的；或分为三种，即从诺厄的三个儿子所出的。这些方舟的部分被称为下层，因为在这尘世状态中有种族的分别，而在上面我们在合一中得以完成。\Emph{上面}没有多样性；因为基督是一切，又在一切之内，仿佛以天上的合一，在一肘之上成全我们。\par

17. 诺厄进入方舟后七天洪水来临；正如我们在未来安息的希望中受圣洗，这安息由第七日所预表。方舟之外地上的所有血肉都被洪水消灭；正如在教会的共融之外，虽然圣洗的水相同，却只效力于毁灭，而非救恩。四十昼夜降雨；正如天上的圣洗圣事，洗净了违反十诫而遍及世界四极的所有罪债（四乘十为四十），无论这罪债是在顺境之日或逆境之夜所犯的。\par

18. 诺厄五百岁时天主吩咐他建造方舟，六百岁时进入方舟；这表明方舟是在一百年内建造的，这似乎对应于一个世界时代的年数。因此第六时代借着福音的宣讲忙于建造教会。接受救恩恩赐的人被制成方形的梁木，适合各样善工，并成为神圣建筑的一部分。再者，诺厄进入方舟是在六百岁的第二个月，而两个月有六十天；因此，如同在六的每一个倍数中，我们有了标示第六时代的数字。\par

19. 提到了月的第二十七日；我们已经看到了梁木中方形的意义。这里尤其重要；因为二十七是三的立方，在我们被塑成方形、即适合各样善工的方式中，有一个三一性：借着记忆，我们记念天主；借着理智，我们认识祂；借着意志，我们爱慕祂。方舟在第七个月停驻；再次提醒我们第七日的安息。而这里再次，为了标示安息者的圆满，第二次提到了月的第二十七日。所以，在希望中应许的，在经验中实现。这里有七和八的结合；因为水高过山顶十五肘，指向圣洗中的深奥奥迹——我们重生的圣事。因为第七日的安息与第八日的复活相关。当圣人们在居间状态的安息之后重新获得身体时，安息不会停止；反而，整个人——灵魂与身体结合，在不朽的健康中更新——将在享受永生中达到望德的实现。因此，圣洗圣事，如同诺厄的洪水，高过一切骄傲者的智慧。七和八也结合在一百五十这个数字中，由七十和八十构成，这是水势汹涌的天数，指出圣洗的深奥意义，即祝圣新人持守安息与复活的信仰。\par

20. 四十天后放出的乌鸦没有返回，或被水阻挡，或被漂浮的尸体吸引；正如被不洁欲望玷污、因而贪恋外界事物的人，要么受圣洗，要么被误导而加入那些在方舟之外、即在教会之外者，对他们而言圣洗是毁灭性的。鸽子被放出时找不到安息，便返回了；正如在新约中，安息并未在此世应许给圣人们。鸽子在四十天后被放出，这时期标示人的一生。再次在七天后被放出，标示圣神七重的运行，鸽子带回一根结果实的橄榄枝；正如有些甚至在教会之外受圣洗的人，如果未失去爱德的肥美，终究可能在傍晚来到，并借着鸽子之口以平安之吻被带入同一共融中。当再次在七天后被放出时，鸽子没有返回；正如在世界终结时，圣人们的安息将不再在望德的圣事中（如同现在，在教会的共融中，他们饮用从基督肋旁流出的），而是在永恒安全的圆满中，那时王国将被交付于天主圣父，而在那不变真理的无云默观中，我们将不再需要自然象征。\par

21. 还有许多其他要点，我们连这样粗略地提及都不可能。为何在\Person{诺厄}生命的第六百零一年——即满六百年后——方舟的顶盖被移开，如同隐藏的奥秘被揭露？为何大地在二月二十七日被说成已干？仿佛第五十七这个数字标志着圣洗圣事之完成。因为二月二十七日是这一年的第五十七天；而五十七是七乘八，即灵与体的数目，再加上一，以表示团结的纽带。为何他们虽分别进入方舟，却一起出来？因为经上说：\Quote{\Person{诺厄}和他的儿子、他的妻子、他儿子的妻子都与他一同进入方舟。}此处男女被分别提及，这表示情欲与精神相争、精神与情欲相争的时期。但他们出去时，是\Person{诺厄}和他的妻子、他的儿子和他们的妻子——男女一同出去。因为在世界末日、义人复活时，身体将与精神完全和谐地结合，不受必死性的需要和情欲的干扰。为何在离开方舟后，只有洁净的动物被祭献给天主，尽管方舟内既有洁净的也有不洁净的？\par

22. 再者，当天主对\Person{诺厄}说话，并好像重新开始，以各种形式的重复来引起对教会预像的注意时，\Person{诺厄}的儿子们受祝福，并被告知要充满大地，所有动物都赐给他们为食物，正如对\Person{伯多禄}所说关于那器皿的话：\Quote{宰了吃！}他们被告知吃时要倒出血来，使此前的生活不被封闭在良心中，而是如同在告解中倾倒出来。天主将彩虹——只在阳光照耀时才出现在云中——作为祂与人及一切活物立约的记号，表示祂不再以洪水毁灭他们；如同那些不在教会中分离的人没有因洪水灭亡一样，他们在天主的云彩——即在先知们和全部圣经中——辨认出\Person{基督}的光荣，而不是寻求自己的光荣。然而，崇拜太阳的人不必以此自夸；因为他们必须明白，太阳，同样狮子、羔羊和石头，都被用作基督的预像，是因为它们有些相似，而非因为它们同质。\par

23. 再者，\Person{基督}从自己民族受的苦难，显然由\Person{诺厄}喝了他所种植的葡萄园的酒而醉，并在帐幕中赤身露体所预表。因为\Person{基督}肉身的可死性被显露出来，对犹太人来说是绊脚石，对希腊人是愚妄；但对那蒙召的，无论是犹太人或希腊人，无论是\Person{闪}或\Person{雅弗}，却是天主的德能和天主的智慧。因为天主的愚妄比人明智，天主的软弱比人刚强。\BibleRef{格前 1:23}\par

此外，那两个儿子——年长的和年幼的——背着衣服倒退而行，是两个民族的预像，也是主过去的、已完成的苦难的圣事。他们没有看见父亲的赤身，因为他们没有同意\Person{基督}的死亡；然而他们用覆盖物尊敬它，因为知道他们是从那里出生的。中间的儿子是犹太民族，因为他们既没有与宗徒们同占首位，也没有后来与外邦人一同相信。他们看见了父亲的赤身，因为他们同意了\Person{基督}的死亡；他们告诉外面的弟兄们，因为那隐藏在先知们中的事由犹太人揭露出来。因此，他们成了他们弟兄的奴仆。如今这个民族除了作基督徒的书桌、携带法律和先知书、并为教会的教义作证之外，还有什么呢？以致我们在圣事中尊崇那他们在字面上所揭露的。\par

24. 再者，人人都必因那两位儿子祝福其父亲赤身的恩惠而受感动，要么得启迪，要么在信德上得坚固；他们虽转过脸去，不愿看见葡萄树所作的恶。他说：\BibleQuote{愿\Person{闪}的天主上主受赞美。}因为虽然天主是万民的天主，就连外邦人也承认祂是以一种特殊的方式是以色列的天主。这除非用\Person{雅弗}的祝福来解释，还能怎样解释呢？教会在外邦人中的普世占据，正应验了这句话：\BibleQuote{愿天主使\Person{雅弗}扩展，使他住在\Person{闪}的帐幕里。}这就是摩尼教徒要注意的。你们看看世界的实际状况。正是你们因我们而感到惊奇和忧伤的事，就是天主正在扩展\Person{雅弗}。祂不是住在\Person{闪}的帐幕里吗？——即住在由宗徒们、先知之子所建立的教会中。请听保禄对信主的外邦人所说的话：\BibleQuote{那时候，你们没有基督，与以色列的公民团体隔绝，是盟约的局外人；没有希望的应许，并在世上没有天主。}这些话描述了\Person{雅弗}住在\Person{闪}的帐幕之前的状态。但请注意下文：\BibleQuote{所以，你们不再是外方人或旅客，而是与众圣人同为国民，是天主家里的人，被建筑在宗徒和先知的根基上，而耶稣基督自己就是主要的角石。}这里我们看到了\Person{雅弗}的扩展，并住在\Person{闪}的帐幕里。这些见证取自宗徒的书信，是你们自己承认、宣读并声称跟随的。你们在一栋建筑中占据了不幸的中间位置，而基督并不是这建筑的角石。因为你们不属于那墙——就像那些属于割损的宗徒那样相信了基督；也不属于那墙——就像那些未受割损的外邦人一样，在信德的合一中，即角石的相联中相连。然而，所有接受并阅读我们正典中任何书卷，其中提到基督曾在肉身中降生和受苦，却不肯用死亡的圣事与我们一起共同遮盖，被苦难所揭露，且不凭虔敬和爱德的知识而宣认我们众人都由此而生的事——尽管他们彼此不同，或作为犹太人和异端者，或作为一种或另一种异端者——他们仍然都是对教会有益的，因为都同样为仆人，或见证某些真理，或印证某些真理。关于异端者，经上说：\BibleQuote{原来在你们中间不免有异端，好叫那些被认可的人在你们当中显明出来。}\BibleRef{格前 11:19}那么，继续你们对旧约圣经的反对吧！继续吧，你们这些\Person{含}的仆役！你们已藐视了那赤身露体时由之而生的肉身。因为除非基督如先知所预言的来到世上，并从自己的葡萄树喝了那杯不能离开祂的酒，并在祂的苦难中沉睡，如同比人智慧的愚拙之醉；并且，在天主隐藏的旨意中，那比人强壮的血肉之躯的软弱被揭露出来。因为除非天主的圣言亲自取了这软弱，你们所夸耀的基督徒的名字就不会存在于世上。那么，继续吧，如我所说。以嘲笑说出我们当以敬畏尊敬的事。让教会使用你们作仆人，以显明那些被认可的人。先知们关于教会状况和受苦的预言如此详尽，以至于我们甚至在关于毁灭性错误的预言中也能为你们找到位置，这错误使被弃绝者灭亡，而被认可者显明出来。\par

25. 你们说基督并未被以色列的先知所预言，事实上，他们的圣经充满这样的预言，只要你们仔细查考，而不是轻率对待。在\Person{亚巴郎}中，谁离开本国和家族，为在外邦人中变得富有和昌盛，难道不是那位离开犹太人的土地和国家——祂按肉身生于其中——而现在正如我们所见的，在外邦人中扩展祂的权能的那位？在\Person{依撒格}中，谁背柴为自己祭献，难道不是那位背起自己十字架的那位？那两只角被荆棘缠住的祭献用的公绵羊是谁，难道不是那位被钉在十字架上作为我们的祭品的那位？\par

26. 那与 \Person{雅各伯} 搏斗的天使，一方面被迫给他祝福，如同弱者对强者、战败者对战胜者，另一方面又使他大腿窝脱臼，这不就是那位容许 \Person{以色列} 子民胜过祂，却祝福了他们中相信的人，而群众却像 \Person{雅各伯} 的大腿一样，因肉性而跛行的那一位吗？那放在 \Person{雅各伯} 头下的石头是谁呢？不就是 \Person{基督}——人的头吗？而在傅油于石头时，\Person{基督} 的名字本身已被表达，因为众所周知，\Person{基督} 意为受傅者。\Person{基督} 在福音中提及此事，声明这是祂自己的预像，那时祂论及 \Person{纳塔乃耳} 说，他是一位真正的 \Person{以色列} 人，心中毫无诡诈；而 \Person{纳塔乃耳}，仿佛将头靠在这石头上，或靠在 \Person{基督} 上，承认祂是天主子、\Person{以色列} 的君王，以他的宣认傅油于石头，在其中他承认 \Person{耶稣} 是 \Person{基督}。在这一场合，\Person{主} 适当地提及了 \Person{雅各伯} 在梦中所见：\BibleQuote{我实实在在告诉你们：你们要看见天开了，天主的使者在人子身上上去下来。}\BibleRef{若 1:47} 这正是 \Person{雅各伯} 所看见的，他在祝福中被称为 \Person{以色列}，当他以石头为枕，看见了一个梯子从地直达天，天主的使者在上面上去下来。\BibleRef{创 28:11} 天使代表福音宣讲者，或 \Person{基督} 的宣讲者。他们上升，当他们超越受造宇宙，描述 \Person{基督} 神性本性的至高威严，即祂在起初就是与天主同在的天主，万物由祂而造。他们下降，为讲述祂由女人所生，生于法律之下，为赎出那些在法律之下的人。\Person{基督} 是从地达天的梯子，或从属肉到属神的梯子：因为藉着祂的助佑，属肉者上升到属神；而属神者可说下降，以奶喂养属肉者，因为他们不能向他们说话如同对属神的人，只能如同对属肉的人。\BibleRef{格前 3:1} 于是在 \Person{人子} 身上既有上升也有下降。因为 \Person{人子} 在上作为我们的头，祂自己就是 \Person{救主}；祂在下在祂的身体、教会中。祂就是梯子，因为祂说：\BibleQuote{我是道路。} 我们上升向祂，为在天上之处看见祂；我们下降向祂，为滋养祂软弱的肢体。这上升和下降既是藉着祂，也是向着祂。效法祂的榜样，宣讲祂的人不仅上升瞻仰被高举的祂，也让自己下降，为坦诚宣布真理。所以宗徒上升说：\BibleQuote{我们若癫狂，是为天主；} 下降说：\BibleQuote{我们若清醒，是为你们。} 他又是藉着谁上升和下降呢？\BibleQuote{因为基督的爱催迫我们，因我们如此断定：如果一人替众人死了，那么众人就都死了；祂替众人死，是为使活着的人不再为自己生活，而是为替他们死而复活的那一位生活。}\BibleRef{格后 5:13}\par

27. 那些不喜欢圣经中这些见解的人，转而走向寓言，因为他们不能忍受纯正的道理。这些寓言对各个年龄段中幼稚的心灵都有吸引力；但我们这些属于 \Person{基督} 身体的人，应同圣咏作者一起说：\Quote{上主，恶人对我说了悦耳的事，但它们不合乎你的法律。} 在这圣经的每一页中，当我作为 \Person{亚当} 之子，流汗劳苦追寻时，\Person{基督} 或公开或隐藏地与我相遇，使我振作。发现艰深时，我的热忱更加激昂，所得的掠物被热切吞食，并藏在心里作为滋养。\par

28. \Person{基督} 在 \Person{若瑟} 身上向我显现，他被兄弟们迫害出卖，在受难后在 \Place{埃及} 得到尊荣。我们看见了 \Person{基督} 在世界中的苦难，\Place{埃及} 是其预像，在殉道者的苦难中实现。现在我们看见 \Person{基督} 在同一个世界中的尊荣，祂为自己征服这世界，以交换祂所赐予的食粮。\Person{基督} 在 \Person{梅瑟} 的棍杖中向我显现，这棍杖丢在地上时变成了一条蛇，作为祂死亡的预像，这死亡来自蛇。再者，当拉住尾巴时，它又变成了棍杖，作为祂完成工作后，在复活中恢复到原先状态的预像，以祂的新生命摧毁死亡，以致不留蛇的任何痕迹。我们作为祂的身体，也在同一有死之生命中，滑过时间的褶皱；但当最后事物的尾巴被审判之手抓住，不再继续时，我们将被更新，从死亡的毁灭——最后的仇敌——中复活，成为在 \Person{天主} 右手中统治的权杖。\par

29. 关于 \Person{以色列} 从 \Place{埃及} 出走，让我们听听宗徒自己所说的话：\BibleQuote{弟兄们，我不愿意你们不知道，我们的祖先都曾在云柱下，都从海中经过，都曾在云中和海中受了圣洗归于梅瑟，并且都吃过同样的神粮，都喝过同样的神饮，因为他们所饮的，是来自伴随他们的神盘石，那盘石就是基督。}\BibleRef{格前 10:1} 对一件事的解释是其余事的钥匙。因为如果那盘石因其稳固而就是 \Person{基督}，那么玛纳不也是 \Person{基督}——那从天上降下的生命之粮，使真正食用者获得属神生命的吗？以色列人死了，因为他们只按肉性接受了预像。宗徒称它为神粮，表明它指向 \Person{基督}，正如神饮由 \Quote{那盘石就是基督} 这话所解释，这话解释了全部。那么，云柱和火柱不也是 \Person{基督} 吗？祂以自己的正直和力量支持我们的软弱；祂黑夜发光，白日不发光，使看不见的看见，使看见的瞎眼。在云中和 \Place{红海} 中有 \Person{基督} 的血所祝圣的圣洗。追来的敌人灭亡，如同过去的罪被消除。\par

以色列人经旷野被领出，一如受洗者在走向预许之地的途中处于旷野，希望并耐心等候那未见之事。旷野中有严峻考验，免得他们心中返回埃及。然而基督并未离弃他们；云柱未曾离去。苦水因木变甜，正如敌对之民因学会尊敬基督的十字架而成为友善。十二股泉浇灌七十棵棕树，是宗徒的恩宠浇灌万邦的预像。七乘十，即十诫在圣神七重运作中得以满全。敌人试图阻拦他们的道路，却被\Person{梅瑟}伸出双手（十字架的预像）所克服。致命蛇伤因铜蛇而愈，铜蛇被举起，使人仰望。主亲自解释此事：\BibleQuote{正如\Person{梅瑟}曾在旷野里举起了蛇，人子也必须被举起来，使凡信祂的人不致丧亡，反而获得永生。}\BibleRef{若 3:14}同样，在许多其他事上，我们也可找到对抗不信之心顽梗的抗议。逾越节中宰杀一只羔羊，代表基督。关于祂，福音上记载：\BibleQuote{请看，天主的羔羊，除免世罪的！}\BibleRef{若 1:29}逾越节中羔羊的骨头不可折断；十字架上主的骨头也没有折断。福音作者引用这话说：\BibleQuote{祂的一根骨头也不可折断。}\BibleRef{若 19:36}门框涂上血以避毁灭，正如人在额上划主受难的标记以得救恩。法律在逾越节后第五十天颁布；同样，圣神在主受难后第五十天降临。法律据说由天主的指头所写；主论及圣神说：\BibleQuote{我以天主的指头驱魔。}\BibleRef{路 11:20}\Person{福斯图斯}闭眼声称，在这样的圣经中他看不出任何对基督的预言。但我们不必惊奇，他有眼可读却无心领悟，因为他不以虔敬的信德叩击天上奥秘之门，反而胆敢以亵渎的敌意行事。那就这样吧，因为本该如此。愿救恩之门向骄傲者关闭。谦卑的人，天主将教诲他们祂的道路，他们必在圣经中寻见这一切；凡他所不见的，将因所见而信。\par

他将看见\Person{耶稣}带领百姓进入预许之地；因为这名字赐予以色列的领袖，并非起初或偶然，而是因他所蒙召的工程。他将看见从预许之地摘下的一挂葡萄悬于木杆上。他将看见在\Place{耶里哥}（即这败亡的世界）中，一个妓女——主曾说她们中有些人比骄傲者先进天国的那类——她从窗户垂下一条象征血的红绳，正如为得赦罪而口里承认。他将看见\Place{耶里哥}的城墙（如同世界脆弱的防御）在约柜绕行七次后倒塌；正如现今在时间的七天历程中，天主的盟约环绕全球，最终，死亡——最后的仇敌——将被毁灭，而教会如同一家，从恶人的毁灭中得救，在赦罪之血中借着告解之窗得以净化，脱离淫乱的玷污。\par

他将看见民长时期先于君王时期，如同审判先于国度。在民长与君王之下，他都看见基督和教会多次以多种方式被预表。\Person{三松}去外族娶妻时，路上遇见狮子并将它杀死，这\Person{三松}是谁？不就是那一位，当祂去从外邦人中召叫祂的教会时，曾说：\BibleQuote{你们放心，我已战胜了世界。}\BibleRef{若 16:33}那被杀之狮口中的蜂巢又是什么意思？岂不是表明，正如我们所看见的，曾经反对基督的世俗政权的法律，如今已失去其凶暴，反而成为传布福音甘甜的庇护？那一位勇敢地用木橛钉穿敌人太阳穴的妇女，岂不是教会的信德藉基督的十字架推翻魔鬼的国度？羊毛湿而地干，而后羊毛干而地湿，岂不是希伯来民族起初在它的预表性制度中独自拥有基督这奥妙的天主，而整个世界却在无知中？如今全世界已获启示这奥妙，而犹太人却不得其门。\par

只提及君王时期少数几件事：在起初，当\Person{厄里}被弃绝、\Person{撒慕尔}被拣选时，司祭职的变更；当\Person{撒乌耳}被弃绝、\Person{达味}被拣选时，王权的变更——这岂不明明预示了在我们的主\Person{耶稣基督}内将要来临的新司祭职与新国度？旧制度是新制度的影子，被弃绝了。\Person{达味}吃了只有司祭才能吃的供饼，这岂不预表了王权与司祭职在一个人——\Person{耶稣基督}——内的合一？十个支派脱离圣殿，而两个支派留下，岂不正是宗徒所断言的那个预像：\BibleQuote{按照恩宠的拣选，留下了残余。}\BibleRef{罗 11:5}？\par

在饥荒时期，\Person{厄里亚}由乌鸦喂养，早晨带饼，晚上带肉；但摩尼教徒不能在此看出基督，祂可说因我们的救恩而饥渴，罪人们在告解中来到祂面前，现今拥有圣神的初果，而在末世，即时代的晚上，他们也将拥有身体的复活。\Person{厄里亚}被派往一个外邦寡妇那里接受供养，那女人在死前去拾两根柴，象征十字架的两根木梁。她的面和油受祝福，因为爱德的果实与喜乐不会因施与而减少，因为\BibleQuote{天主喜爱乐意奉献的人。}\BibleRef{格后 9:7}\par

35. 那些呼喊\Quote{秃头}嘲弄\Person{厄里叟}的孩童，被野兽吞噬，正如那些以幼稚的愚昧嘲笑被钉在\Place{加尔瓦略}的\Person{耶稣基督}的人，被魔鬼毁灭。\Person{厄里叟}打发仆人将他的棍杖放在死尸上，死尸却没有复活；他亲自前来，完全伏在死尸上，死尸就复活了：正如\Person{天主的圣言}藉他的仆人\Person{梅瑟}颁赐了法律，对死于罪恶的人类并无益处；然而，这并非毫无目的，因为那知道必须先颁布法律的主，如此行了。然后他亲自来临，藉分受我们的死亡而与我们相似，我们便复活了。当他们用斧头砍伐树木时，铁头从木柄飞出，沉入河底，当\Person{厄里叟}将木头扔入水中时，铁头又浮了上来。同样，当\Person{耶稣基督}的肉身临在时，正如\Person{若翰}所说的：\Quote{看，斧子已经放在树根上}，\BibleRef{玛 3:10}；藉着他们所施加的死亡，\Person{基督}与他的身体分离，下降到阴间的深处；然后，当他的身体被安放在坟墓中时，如同木头在水上，他的灵魂返回，如同铁头回到木柄上，他便复活了。读者会注意到，为求简略，许多这类事例被省略了。\par

36. 至于前往\Place{巴比伦}的迁徙，\Person{天主的圣神}藉\Person{耶肋米亚}先知命令他们前往，告诉他们要为所寄居之地的人民祈祷，因为他们的平安中他们可得平安，并要建造房屋，栽种葡萄园和园圃——当我们想到真正的\Person{以色列人}（心中毫无诡诈）藉着宗徒的职务，带着福音的规条进入外邦人的国度时，比喻的意义是明显的。因此，\Person{保禄}宗徒如同\Person{耶肋米亚}的回声，对我们说：\Quote{我首先劝勉众人，要为所有的人、为在位的祈祷、恳求、转求和感谢，好使我们能怀着完全的虔敬和爱德，度宁静平安的生活；因为这在\Person{天主}我们的救主面前是美好的，可蒙悦纳的，他愿意所有的人都得救，并认识真理。}\BibleRef{弟前 2:1} 因此，信徒建造了基督徒聚会的圣堂作为平安的居所，恢复了信友的葡萄园，栽种了园圃，其中首要的植物是芥菜树，在它广展的枝叶中，外邦人的骄傲如同天上的飞鸟，在它高飞的野心中得庇护。再者，根据\Person{耶肋米亚}的预言，七十年后从被掳中回归，以及圣殿的重建，每位信友在\Person{基督}内必能看见我们作为\Person{天主}的教会从这世界的流放回归天上\Place{耶路撒冷}的预像，那时七日的时期已经完成。被掳后重建圣殿的大司祭\Person{若苏厄}，是我们重建的真正大司祭\Person{耶稣基督}的预像。\Person{匝加利亚}先知看见这位\Person{若苏厄}穿着污秽的衣服；当那站在旁边控告他的魔鬼被击败后，污秽的衣服被脱去，赐给他光荣和尊贵的衣服。同样，\Person{耶稣基督}的身体，即教会，当仇敌在世界末日的审判中被征服时，将从流放的痛苦进入永恆安全的荣耀。这是圣咏作者在奉献房屋时的诗歌：\Quote{你把我的哀痛化成了舞蹈，脱去了我的苦衣，给我披上喜乐，好使我的灵魂歌颂你，不致缄默。}\par

37. 在这样一篇插话中，要简要地提及法律和先知书中所有关于基督的比喻性预言是不可能的。难道有人会说，这些事是按常规发生的，而将它们视为基督的预像只是巧妙的想象吗？这种异议可能来自犹太人和外教人；但那些希望被视为基督徒的人必须顺从宗徒的权威，当他说：\Quote{这一切发生在他们身上，都是作为预像}，又说：\Quote{这些事是我们的预像。} 因为如果两个人，\Person{依市玛耳}和\Person{依撒格}，是两个盟约的预像，我们能否认那大量没有历史或自然价值的细节没有意义吗？假设我们在一座宏伟建筑的墙上看到一些希伯来文字，我们会愚蠢到因为不认识这些文字就得出结论说它们不打算被阅读，只是没有意义的绘画吗？同样，凡以坦诚之心阅读旧约圣经中所包含的一切的人，必会感到不得不承认它们都有意义。\par

38. 作为那些若无象征意义便毫无意义之细节的例子：假设女人必须被造为男人的助手，然而她从沉睡的男人肋旁被取出，有什么自然理由可解释？假设方舟是为了躲避洪水所需，为什么它必须有这些尺寸，为什么这些尺寸要被记录下来供后世虔敬研读？假设动物被带进方舟是为了保存各种族类，为什么洁净的有七对，不洁净的只有一对？假设方舟必须有门，为什么门在侧面，为什么这一事实要记录下来？\Person{亚巴郎}被命令献祭他的儿子：我们可以承认，这顺服的考验是必要的，以便它在所有时代显明；我们也可以承认，由儿子背负木柴而非年迈的父亲是合适的，并且最终那致命的一击被阻止，以免父亲无子嗣。但流\Person{公绵羊}的血与\Person{亚巴郎}的考验有何关系？或者，如果必须完成祭献，公绵羊被灌木丛中的角挂住是否使它更好？人心——即理性的心——通过思考这些看似多余之事如何与必要之事交织在一起，首先被引导承认它们有意义，然后尝试发现这意义。\par

39. 犹太人自己讥讽我们所信的被钉的救主，因此不愿承认旧约记载的言论和行动预报了基督，但他们不得不向我们求解释那些若不解说就显得微不足道和可笑的事。这导致一位博学的犹太人\Person{斐洛}（希腊人说他雄辩如\Person{柏拉图}）试图不以基督为视角解释一些事，而他自己并不信基督。他的尝试只表明，一切巧妙的思辨若不以基督为视角（一切预言都真正指向他），便多么低下。宗徒的话是多么真实：\BibleQuote{当他们归向主时，帕子就除去了。}\BibleRef{格后 3:16}例如，\Person{诺厄}方舟在\Person{斐洛}看来是整个人体的\OriginalTerm{预像}{type}，一部分一部分地对应：他以此观点指出比例完全吻合。因为基督的预像没有理由不同时是人体的预像，因为人类的\OriginalTerm{救主}{Saviour}取了人体显现，但人体的预像不一定是基督的预像。然而，\Person{斐洛}的解释在方舟侧面的门这一点上失败了。他实际上为了说点什么，把这门解释为人体下部的开口。他竟敢于说出并写下这话。他确实不认识那门，也不能理解这象征。如果他转向基督，帕子就会除去，他就会发现教会的\OriginalTerm{圣事}{sacraments}从基督人性身体的肋旁涌出。因为按照那宣告\BibleQuote{二人成为一体}，方舟这一基督的预像中有些事物指向基督，有些指向教会。这种以基督为视角的解释与所有其他巧妙的曲解之间的对比，在圣经一切预象的每一细节上都相同。\par

40. \OriginalTerm{外邦人}{Pagans}也不能否认我们对言语和事物赋予寓意的权利，尤其因为我们可以指出\OriginalTerm{预像}{type}和\OriginalTerm{预象}{figure}的应验。因为外邦人自己也试图在他们自己的寓言中找到自然和宗教真理的预象。有时他们给出清晰的解释，而有时他们掩盖其含义，以至于庙宇中的神圣事物在剧场中成了笑柄。他们将可耻的放荡与卑劣的迷信结合在一起。\par

41. 除了预像与所预表之事之间这种奇妙的吻合之外，对手还可能被明确的预言暗示所说服，例如：\BibleQuote{地上万民都因你的后裔而蒙祝福。}\BibleRef{创 22:18}这话对\Person{亚巴郎}说过，又对\Person{依撒格}说过（\BibleRef{创 26:4}），又对\Person{雅各伯}说过（\BibleRef{创 28:14}）。因此\BibleQuote{我是亚巴郎、依撒格、雅各伯的天主}\BibleRef{出 3:6}这句话意义重大。天主应验了他对其后裔的\OriginalTerm{恩许}{promise}，祝福了万民。同样有意义的是，\Person{亚巴郎}自己让仆人起誓时，吩咐他把手放在他的大腿下（\BibleRef{创 24:2}），因为他知道那里将产生基督的血肉，而我们现在藉这血肉得到的，不是祝福万民的恩许，而是应验了的恩许。\par

42. 我倒想知道，或者不如说，不知道才好，福斯图斯是怀着怎样的心智盲目来读这段经文：雅各伯呼唤他的儿子们，说：\BibleQuote{你们集合，我好告诉你们日后将发生的事。雅各伯的儿子们，你们集合聆听；听以色列，你们的父亲。} 这无疑是先知的话。那么，论及他的儿子犹大——按肉身说，基督出自犹大支派，达味的后裔，正如宗徒所教导的——雅各伯说了什么呢？他说：\BibleQuote{犹大，你的兄弟要赞美你；你的手要放在你仇敌的颈上；你父亲的儿子要向你俯伏。犹大是狮子的幼崽；我的儿子，我所生的：你俯伏着上去；你像狮子一样睡卧，又像母狮之子，谁能唤醒他？权杖必不离犹大，统治者的棍杖不离他两腿之间，直到那为他所存留的事物来到。他也是万民的期待：将他的驴驹拴在葡萄树上，将他的母驴的驹子拴在葡萄树上；他在酒中洗他的衣服，在葡萄的血中洗他的袍子；他的眼睛因酒而发红，他的牙齿比奶还白。} 当我们藉着基督的明光来读这些话语时，其中毫无虚假或晦涩。我们看见祂的兄弟，即宗徒们和所有与祂同为嗣子的人，赞美祂，不求自己的光荣，但求祂的光荣。我们看见祂的手放在仇敌的颈上，那些仇敌因基督团体的增长而弯曲俯伏在地，尽管他们反对。我们看见雅各伯的儿子们——按恩宠的拣选而蒙救的遗民——朝拜祂。基督，生为婴儿，就是狮子的幼崽，正如接着所说\BibleQuote{我的儿子，我所生的}，以表明为何这幼崽——有话说：\BibleQuote{狮子幼崽比群兽更强壮} \BibleRef{箴 30:30}——甚至幼年时就比年长者更强。我们看见基督登上十字架，并在交付出灵魂时俯首。我们看见祂像狮子一样睡卧，因为在死亡本身中，祂不是被征服者，而是征服者；又像母狮之子，因为祂生与死的原因相同。祂由那位没有人看见，也不能看见者从死者中复活；因为\BibleQuote{谁能唤醒祂？}这话指向一种未知的权能。权杖没有离开犹大，统治者也没有离开他的两腿之间，直到所承诺的事物适时来到。因为我们从犹太人自己的真实历史得知，基督出生时的黑落德是他们的第一个外邦君王。所以权利杖没有离开犹大的苗裔，直到为他所存留的事物来到。然后，由于这应许不仅是为相信的犹太人，还加上：\BibleQuote{他是万民的期待。}基督用麻布拴住祂的驴驹——即祂的子民——当祂穿着苦衣宣讲：\BibleQuote{你们悔改吧，因为天国临近了。} 服从于祂的外邦人由母驴的驹子所代表，祂也骑过它，领着它进入耶路撒冷，即和平的异象，教导温良的人祂的道路。我们看见祂在酒中洗祂的衣服；因为祂与荣耀的教会是一体的，祂将教会献给自己，没有污点或皱纹；依撒意亚也对教会说：\BibleQuote{你们的罪虽似朱红，我必使它洁白如雪。} \BibleRef{依 1:18} 这如何成就，除非藉着罪的赦免？而这酒正是那所说的\BibleQuote{为多人倾流，以赦免罪过}的酒。基督是挂在木杆上的葡萄串。所以接着又说：\BibleQuote{他的衣服在葡萄的血中。} 同样，关于祂的眼睛因酒而发红，可由祂身体的那些肢体来理解，他们因圣洁地脱离尘世事物的潮流，得以凝视智慧的永恒之光。保禄在先前引用的一段话中说：\BibleQuote{我们若疯狂，那是为了天主。} 那就是因酒而发红的眼睛。但他又加上：\BibleQuote{我们若清醒，那是为了你们。} 那些需要奶喂养的婴孩也没有被遗忘，正如\BibleQuote{他的牙齿比奶还白}这句话所表明的。\par

43. 我们受蒙蔽的对手对这些明明白白的例子能说什么呢？这些例子不容许乖僻的否认，甚至怀疑的不确定性。我呼吁摩尼教徒开始探究这些主题，承认这些证据的力量——我没有时间详细讨论；我也不想挑选，以免无知的读者以为没有其他证据，而基督信徒读者可能会因我遗漏许多比我此刻想到的更引人注目的要点而责备我。你们会发现许多经文不需要像这里对雅各伯预言所做的解释。例如，每一位读者都能理解这些话：\BibleQuote{祂像羔羊被牵去宰杀}，以及整个明白的预言：\BibleQuote{因祂的创伤我们得了痊愈}——\BibleQuote{祂背负了我们的罪}。\EditorialNote{依撒意亚第五十三章} 我们在诗意的福音中读到这些话：\BibleQuote{他们刺穿了我的手和脚。他们数说了我所有的骨头。他们瞪眼注视着我。他们分了我的外衣，为我的内衣拈阄。} 现在连盲人也看到这些话的应验：\BibleQuote{大地四极将记起并归向上主，万邦万国都要在祂面前朝拜。} 福音中的话：\BibleQuote{我的心灵忧闷得要死}、\BibleQuote{我的心灵烦乱}，是圣咏中话语的重复：\BibleQuote{我忧闷地躺卧}。是谁使祂躺卧？谁的声音呼喊\Quote{钉死他，钉死他}？圣咏告诉我们：\BibleQuote{人们的牙齿是矛和箭，他们的舌头是利剑。} 但他们不能阻止祂的复活，或祂升到诸天之上，或祂以祂的圣名充满大地；因为圣咏说：\BibleQuote{天主，请祢高耸于诸天之上，愿祢的光荣充满大地。} 每个人都必须将这些话应用于基督：\BibleQuote{上主对我说：祢是我的儿子，我今日生了祢。祢向我请求，我必将万民赐祢作产业，将地极赐祢作基业。} 而耶肋米亚论及智慧的话显然也适用于基督：\BibleQuote{雅各伯将它交给他的儿子，以色列将它交给他的特选者。后来它出现在大地上，与人共处。}\par

44. 同一救主在\BibleBook{达尼尔}{先知书}中也有提及，那里人子出现在万古常存者面前，领受一个永无穷尽的王国，使万民都事奉祂。\BibleRef{达7:13} 在主亲自引用的达尼尔经文\Quote{当你们看见达尼尔先知所说的\InnerQuote{招致荒凉的可憎之物}站在圣地，读这经的应当明白}\BibleRef{玛24:15}，那\Quote{七周}的数字不仅指向基督，更指向祂降临的具体时间。对于和我们一样期待基督带来救恩、却否认祂已来临并受苦的犹太人，我们可以根据实际发生的事件与他们辩论。除了外邦人的归化如今如此普遍——这已在他们自己的圣经中关于基督的预言中——还有犹太人自身历史中的事件：他们的圣地已被摧毁，祭献已停止，司铎和古老的傅油也已终止；这一切都清楚地为达尼尔所预言，当他预言了至圣者的受傅时。\BibleRef{达9:24} 如今这些事都已发生，我们向犹太人询问那受傅的至圣者，他们却无言以对。然而，犹太人关于基督及其降临的所有知识都来自旧约。他们为什么问若翰是不是基督？他们为什么对主说：\Quote{你叫我们犹疑不定到几时呢？你若是基督，就明明地告诉我们。}为什么\Person{伯多禄}、\Person{安德肋}和\Person{斐理伯}对\Person{纳塔乃耳}说：\Quote{我们找到了\OriginalTerm{默西亚}{Messias}，即\OriginalTerm{基督}{Christus}}？不正是因为这名号是他们从圣经预言中所熟知的吗？没有别的民族君王和司铎受傅并被称为\Quote{受傅者}或\Quote{基督}。这象征性的傅油不能终止，直到那借此预表的那一位来临为止。因为在所有受傅者中，犹太人期待一位要来拯救他们的。但因着天主奥秘的公义，他们被蒙蔽了；只思想默西亚的能力，却不明白祂的软弱，正是这软弱使祂为我们而死。在\BibleBook{智慧}{篇}中预言了犹太人：\Quote{我们要用耻辱的死刑定他的罪，因为他必按他的话受考验。如果他真是天主子，天主必帮助他，拯救他脱离仇人的手。他们这样思想，是出于谬误，因为他们的邪恶蒙蔽了他们。}\BibleRef{智2:18} 这些话也适用于那些面对这一切证据、面对如此一系列预言及其应验，却仍否认基督在圣经中得到预告的人。每当他们重复这种否认，我们就可在祂的帮助下提出新的证明，祂做了如此周全的预备来对抗人的执拗，以至于已提出的证明无需重复。\par

45. \Person{福斯特}有一个逃避性的反对意见，他无疑认为这是逃避最清晰预言证据的绝妙方法，但人们不愿理会它，因为回答它可能会赋予它实际上并不具备的重要性。有什么比说\Quote{没有证据就不信基督是信德薄弱}更不合理的呢？难道我们的对手相信关于基督的见证吗？福斯特希望我们相信来自天上的声音，而非人的见证。但他们听到这声音了吗？有关它的知识不正是通过人的见证传给我们吗？宗徒描述了这知识的传递，他说：\Quote{那么，人怎能呼求他们所不信从的呢？又怎能信从他们所没有听过的呢？没有宣讲者，又怎能听到呢？如果没有被派遣，又怎能宣讲呢？正如经上所载：\InnerQuote{传布和平、报道喜讯者的脚步是多么美丽！}}\BibleRef{罗10:14} 显然，在宗徒的宣讲中，提到了先知性的见证。宗徒引用先知的预言，以证明他们道理的真理性和重要性。因为尽管他们的宣讲伴有行奇迹的能力，但奇迹会被归为魔术——正如有些人至今仍斗胆暗示的那样——除非宗徒们表明先知们的权威站在他们一边。生活在如此久远以前的先知们的见证，不能归为魔术。也许福斯特不愿我们相信希伯来先知作为真基督的见证，是因为他相信关于假基督的波斯异端。\par

46. 根据天主教会训导，基督徒的心智必须首先在单纯的信仰中得到滋养，才能有理解天上永恒之事的能力。因此先知说：\Quote{除非你们相信，你们将不会明白。}单纯的信仰是指：在我们达到认识基督之爱的高峰、得以充满天主的一切圆满之前，我们相信，基督屈尊就卑的计画——祂作为人诞生并受苦——并非毫无理由地被先知们藉着先知性的民族、先知性的人民和先知性的王国如此长久地预告。这信仰教导我们，在那比人更有智慧的天主愚妄和比人更有力量的天主软弱中，隐藏着我们成义和受光荣的奥秘途径。在那里隐藏着一切智慧和知识的宝藏，这些宝藏对任何藐视从他母亲乳房——即宗徒和先知的教导——所传递的滋养的人，或任何自以为年长而不适合幼小食物，因鲁莽自认为已准备好而吞食异端毒物而非智慧食粮的人，都是不开启的。要求单纯的信仰与要求相信先知是完全一致的。单纯的信仰正是要人一开始就相信先知，而在心智净化并坚固之后，才能领悟那在先知中发言的那一位。\par

47. 但有人说，如果先知们预言了基督，他们的生活就不符合其职分。你怎么能断定他们是否如此呢？你们对善恶生活的判断力很差，你们的\Quote{正义}在于通过吃掉一个无生命的瓜来解除它的痛苦，而不是给饥饿的乞丐食物。对于在天主教会中尚未认识灵魂的圆满正义、尚未分辨所追求的正义与实际达成的正义之区别的婴孩来说，按照宗徒们有益的教导（义人因信德而生活）来看待这些人就足够了。\BibleQuote{亚巴郎相信了天主，并以此算为他的义德。因为圣经预见天主将由信德使外邦人成义，就预先向亚巴郎传报福音说：\InnerQuote{万民都要因你的后裔蒙受祝福。}} \BibleRef{迦 3:6} 这些是宗徒的话。倘若你们愿意在他清晰熟悉的声音中从你们无益的梦中醒来，就会追随我们的父亲亚巴郎的足迹，与万民一起在他的后裔中蒙受祝福。因为，如宗徒所说：\BibleQuote{他领受了割损的标记，作为他未受割损时因信德而获得的义德的印证，使他成为一切未受割损而信德之人的父亲，也为那些受割损的人作父亲，不仅是为那些受割损的，也是为那些追随我们父亲亚巴郎在未受割损时所走信德之路的人。} \BibleRef{罗 4:11} 既然亚巴郎的信德之正义如此被树立为我们的榜样，使我们也能因信德成义，与天主和平相处，我们就应当理解他的生活方式而不加以指责；免得因过早脱离母亲教会，成为流产的胎儿，而不是在受孕完成后适时出生。\par

48. 这是对福斯图斯的简短答复，为圣祖和先知们的品格辩护。这是我们信德中的婴孩的答复，我自认为其中之一，因为我即使不理解古代圣人们生活的奥秘性质，也不会指责他们的生活。宗徒们在他们的福音中赞许地向我们传报了圣人们的生活，正如圣人们自己在预言中预言了未来的宗徒们，好使两约像撒拉弗一样彼此呼喊：\BibleQuote{圣！圣！圣！万军的上主！} \BibleRef{依 6:3} 当福斯图斯不是像这里提出模糊的普遍指控，而是谴责圣祖和先知们生活中的具体行为时，他们的天主，也是我们的天主，会帮助我适当且恰当地回应各项指控。目前，读者必须选择是相信宗徒保禄的称赞，还是相信摩尼教徒福斯图斯的指控。\par

\SourceRecord{Translated by Richard Stothert.；Nicene and Post-Nicene Fathers, First Series；Vol. 4.；Edited by Philip Schaff.；Buffalo, NY: Christian Literature Publishing Co.,；1887.；<http://www.newadvent.org/fathers/140612.htm>.}

% structure: p0001 source=h1 target=section reason=page-title

\section{\WorkTitle{驳福斯托斯}，第十三卷}

\Emph{福斯托斯断言，即使旧约能显示包含预言，那也只对犹太人有趣，而外邦人的异教文献也同样服务于此目的。奥古斯丁展示预言对外邦人和犹太人的共通价值。}\par

1. \Emph{福斯托斯}说：有人问我们，当我们拒绝那些宣告基督来临应许的先知时，我们如何敬拜基督。经审查，能否证明希伯来先知预言了我们的基督，即天主子，这是可疑的。但即便如此，这与我们何干？倘若你们所说的这些先知证言能使任何人从犹太教皈依基督教，而他日后却忽视这些先知，那他确实错了，且可被控忘恩负义。但我们生来是外邦人，未受割礼；如保禄所说，生在另一种法律之下。外邦人称之为诗人的，是我们最初的宗教导师，我们从他们以后才皈依基督教。我们并非先成为犹太人，再藉着相信他们的先知而达至基督教；我们是唯独被我们的解放者耶稣基督的名声、德行和智慧所吸引。若我仍在我祖先的宗教中，而某位传道人使用先知作为支持基督教的证据，我会认为他疯了，竟想用对一个完全异教的外邦人更可疑的事来支持可疑的事。他首先要说服我相信先知，然后再藉先知相信基督。而要证明先知的真确，又需要别的先知。因为若先知为基督作证，谁又为先知作证？你或许会说，基督与先知彼此支持。但一个与两者都无关的外邦人，既不会相信基督对先知的见证，也不会相信先知对基督的见证。若外邦人成了基督徒，他只需感谢自己的信德，别无其他。为举例说明，让我们假设自己正与一位外邦探询者交谈。我们告诉他相信基督，因为祂是天主。他要求证据。我们指向先知。他问：什么先知？我们回答：希伯来人。他微笑说，他不信他们。我们提醒他基督为它们作证。他笑着回答说，我们必须先让他相信基督。这番对话的结果是，我们哑口无言，而探询者离开，认为我们热心有余而智慧不足。再者，我说，教会既由外邦人多于由犹太人组成，便不会欠希伯来证人任何东西。若如所说，任何关于基督的预言能在西彼拉、或称为赫耳墨斯·特里斯墨吉斯忒斯、或俄耳甫斯、或任何异教诗人那里找到，它们或许能帮助那些像我们一样从异教皈依基督教的人的信德。但希伯来人的见证对我们皈依前是无用的，因为那时我们不能相信它们；皈依后则是多余的，因为我们不信它们也能相信。\par

2. \Emph{奥古斯丁}回答说：上一卷已作长篇答复，此处简短回答就够了。读过那答复的人，必定认为福斯托斯坚持否认希伯来先知预告了基督，是疯狂的，因为希伯来民族是唯一一个在君王和司祭身上\Quote{基督}这名号具有特殊神圣意义的民族；在这个意义上，它一直沿用至那些君王和司祭所预表的祂来临。摩尼教徒从哪里学到基督之名？若来自摩尼，那就非常奇怪了，非洲人（更不用说其他人）竟会相信波斯人摩尼，因为福斯托斯指责罗马人和希腊人以及其他外邦人相信属于另一民族的希伯来先知。照福斯托斯看来，西彼拉、俄耳甫斯或任何异教诗人的预言更适合引导外邦人相信基督。他忘了这些在教会中并未被诵读，而希伯来先知的声音四处响起，吸引成群的人归向基督教。当显而易见人们到处藉希伯来先知被引向基督时，说这些先知不适合外邦人，是极大的荒谬。\par

3. 希伯来先知所预告的基督不令你喜悦；但这正是外邦民族所相信的基督，照你看来，希伯来预言对他们不应有分量。他们接受了福音，如保禄所说：\BibleQuote{这福音是天主先前藉着祂的先知在圣经中所应许的，论及祂的圣子，按血肉是达味后裔所生的。}（\BibleRef{罗 1:2}）我们在\BibleBook{依撒意亚}{书}中读到：\BibleQuote{叶瑟的根将兴起，要统治万民；外邦人将寄望于祂。}（\BibleRef{依 11:10}）又说：\BibleQuote{看，一位贞女将怀孕生子，人要称祂的名字为厄玛奴耳。}（\BibleRef{依 7:14}）这名字译出来就是\Quote{天主与我们同在}。莫让摩尼教徒以为希伯来先知仅将基督预告为人；因为福斯托斯似乎暗示这点，当他说：\Quote{我们的基督是天主子}，仿佛希伯来人的基督不是天主子。我们能够证明希伯来预言中那贞女之子基督是天主。因为主亲自教导属血肉的犹太人，不要因为祂被预告为达味之子，就以为祂仅此而已。祂问：\BibleQuote{你们对基督怎么看？祂是谁的儿子？}他们回答：\BibleQuote{达味的。}然后，为提醒他们厄玛奴耳，即天主与我们同在的名字，祂说：\BibleQuote{那么，达味怎么在圣神里称祂为主，说：主对我主说：你坐在我右边，等我使你的仇敌做你的脚凳？}（\BibleRef{玛 22:42}）这样，基督在希伯来预言中显现为天主。摩尼教徒能展示什么含有基督之名的预言呢？\par

4. \Person{摩尼}确实不是\Person{基督}的先知，却自称为宗徒，这是无耻的谎言；因为众所周知，这个异端不仅始于\Person{戴尔都良}之后，而且是在\Person{西彼廉}之后。\Person{摩尼}在所有书信中这样开头：\Quote{\Person{摩尼}，耶稣基督的宗徒。}你为何相信\Person{摩尼}关于基督所说的话？他有什么证据证明他的宗徒身份？基督这名号，我们只从犹太人得知；犹太人将这名号用于他们的君王和司祭，他们并非个人，而是整个民族作为基督和基督国度的先知。他有什么权利使用这名号？他禁止你们相信希伯来先知，好使你们成为假基督的异端门徒，正如他自己是假的和异端的宗徒一样。如果\Person{福斯图斯}举出一些据他预言基督的先知作为支持自己的证据，他如何满足他那假设的询问者？这询问者既不相信先知，也不相信\Person{福斯图斯}。他会以我们的宗徒为证人吗？除非他能找到一些在世宗徒，否则他必须阅读他们的著作；而这些著作全都反对他。它们教导我们的教义：基督由童贞玛利亚诞生，祂是天主之子，按肉身是达味后裔。他不能声称这些著作被篡改过，因为那就等于攻击他自己证人的可信度。或者，如果他拿出他自己版本的宗徒著作，他还必须为这些著作取得由宗徒亲自建立的教会的权威，证明这些著作经他们认可被保存和流传。一个人要我根据那些其全部权威都来自他本人之言的著作（而我不相信他）来相信他，这是很难的。\par

5. 但或许你们相信关于基督的普遍传闻。\Person{福斯图斯}微弱地提出这种建议作为最后手段，以避免要么提出他毫无价值的权威，要么落入反对他者的控制。好吧，如果传闻是你们的权威，你们就当考虑信赖这种证据的后果。有许多关于你们的坏事被传，你们不希望人们相信。将同样的证据用于基督即真、用于你们自己即假，这合理吗？事实上，你们否认关于基督的普遍传闻。因为最广为传播、每个人反复听到的传闻，明确断言基督按在希伯来圣经中向\Person{亚巴郎}、\Person{依撒格}和\Person{雅各伯}所作的许诺，是达味的后裔：\Quote{地上万民都要因你的后裔蒙受祝福。}你们不接受这希伯来见证，但你们似乎没有其他见证。我们书籍的权威，得到如此众多民族的一致确认，有宗徒、主教和会议的传承支持，反对你们。你们的书籍没有权威，因为只有少数人维护它，而这些人是虚假天主和基督的崇拜者。如果他们不效法他们所崇拜者的榜样，他们的见证就必须反对他们自己的虚假教义。此外，普遍传闻对你们评价很差，并且始终坚持与你们相反，说基督是达味的后裔。你们没有听见来自天上的父的声音。你们没有看见基督为自己作证的事迹。你们声称接受记载这些事的书籍，以维持一种虚假的基督徒外表；但当有人引用反对你们的话时，你们就说书籍被篡改过。你们引用基督的话：\Quote{你们若不信我，当信那些事业}；又说：\Quote{我为自己作证，派遣我的父也为我作证}；但你们不允许我们引用这些经文作为答复：\Quote{你们查考圣经，因你们认为其中有永生，正是这些为我作证}；\Quote{如果你们相信梅瑟，你们必会相信我，因为他论及我而写作}；\Quote{他们有梅瑟和先知，让他们听从他们}；\Quote{如果他们不听梅瑟和先知，纵使有人从死者中复活，他们也不会信服。}你们有什么可说的？你们的权威在哪里？如果你们拒绝这些圣经章节，尽管它们有重大的权威支持，你们能显示什么奇迹？然而，即使你们行了奇迹，我们也当警惕不接受它们在你们身上的证据；因为主预先警告了我们：\BibleQuote{许多假基督和假先知将要兴起，显大征兆和奇事，以致如果可能，连被选的人也要被迷惑；看，我预先告诉了你们。} \BibleRef{玛 24:24} 这表明，圣经的既定权威必须超越一切其他权威；因为事件进展的进程给它带来新的确认，正如圣经所证明的，这些事件应验了在它们发生之前许久所作的预言。\par

6. 那么，你们的教义是否如此显而易见地真实，以至于不需要奇迹或任何证言的支持？如果你们有任何这类东西可展示，就请向我们展示这些不证自明的真理吧。如我们所见，你们的传说又长又愚蠢，是供妇女儿童消遣的老妇之言。开端与其余部分脱节，中间不稳固，结尾则是可悲的失败。如果你们始于不死、不可见、不朽坏的天主，那么祂何必与黑暗种族争战呢？至于你们理论的中段，当祂在水果和蔬菜中的肢体因你们的咀嚼和消化而被净化时，天主的不可朽坏和不变性又成了什么？至于结局，既然天主无法按照祂自己的使命，将可怜的灵魂因外在邪恶而沾染的污秽洗净，那么这灵魂因执行天主的使命而被永久囚禁在黑暗团块中受惩罚，难道公平吗？你们无言以对。看你们自诩数量众多、价值连城的手稿是多么无用！唉，古文字学家们的辛劳！唉，不幸拥有者的财产！唉，受骗跟随者的食粮！你们既无圣经权威，又无奇迹之能，也无道德美德，还无纯正教义，就羞愧地离开吧，怀着悔改之心归回，宣认那位真正的基督，祂是所有信祂之人的救主，祂的名和祂的教会如今彰显出来，正如古时所预言的，不是从地底黑暗里出来的某个存在，而是由一个为此目的而建立的特殊王国中的民族所预言，好使那些关于基督的预表性预言得以在那里写下，如今在现实中应验，先知们笔下预言的，正是宗徒们现今在宣讲中展示的。\par

7. 那么，让我们设想与一位外邦寻道者的对话，\Person{浮士德}描述我们表现得很差劲，尽管他自己的表现更可悲。如果我们对外邦人说：\Quote{相信基督吧，因为祂是天主。}他若要求证据，我们就提出先知们的权威。如果他说他不相信先知，因为他们是希伯来人而他是外邦人，我们就能从先知预言的应验本身证明其真实。他很难不知道早期基督徒所受的这世界君王的迫害；或者如果他不知道，可以从历史和帝国法律记录中得知。但这正是先知久远以前所预言的，他说：\BibleQuote{外邦为什么怒吼，万民为什么妄想？世上的君王起来，首领聚在一起，反对上主，反对祂的基督。}圣咏的其余部分表明这不是指\Person{达味}说的。因为接下来的内容足以说服最顽固的不信者：\BibleQuote{上主对我说：\InnerQuote{你是我的儿子；我今日生了你。你向我请求，我就将外邦人赐你为基业，将地极赐你为产业。}}这事从未发生在犹太人身上，他们曾是君王\Person{达味}的子民；但现在却明显应验在所有民族臣服于基督之名上。这一预言以及许多类似的预言（在此引述会太长）一定会给寻道者留下深刻印象。他会看到这些世上的君王如今都快乐地臣服于基督，万民族都事奉祂；他会听到圣咏中如此久远之前预言的话语：\BibleQuote{地上的众王都要向祂俯伏，万国都要事奉祂。}如果他读那整个圣咏（预表性地应用于\Person{撒罗满}），他会发现基督才是真正的和平之王，因为撒罗满意思就是和平；他会在圣咏中发现许多适用于基督的事情，而这些与字面的撒罗满王毫无关系。还有另一篇圣咏，天主被天主傅了油，傅油一词正指向基督，表明基督是天主，因为天主被描述为被傅油。在阅读这篇圣咏关于基督和教会的内容时，他会发现其中预言的已在当今世界应验。他会看到地上的列国偶像正在消亡，他会发现这正是先知所预言的，如\Person{耶肋米亚}所说：\BibleQuote{那时，你们要对他们说：\InnerQuote{那没有创造天地的神明，必从地上、从天下消灭。}}\BibleRef{耶 10:11} 又说：\BibleQuote{上主，我的力量，我的保障，我在患难之日的避难所，万民必从地极来到你面前，说：\InnerQuote{我们的祖先承受的实在是虚假，是虚无，是无益之物。}人岂能为自己制造神祇？其实那不是神。因此，看哪，那时我要使他们知道，我要使他们知道我的手和我的大能，他们就知道我是上主。}\BibleRef{耶 16:19} 听到这些预言并看到它们实际应验，不用说，他一定深受感动；因为我们凭经验知道，信徒们看到古代预言在现今实现，内心是如何得到坚固的。\par

8. 在同一先知书中，询问者会找到明确的证据，证明\Person{基督}不仅仅是世上出现的伟人之一。因为\Person{耶肋米亚}接着说：\BibleQuote{凡信赖世人、以血肉为臂膀、心怀远离上主的人，是可咒骂的：他必像荒野中的柽柳，即使幸福来到，也看不见；他必住在旷野干旱之处、无人居住的咸地。信赖上主、以上主为依靠的人，是可祝福的：他必像栽在水边的树，在河畔扎根；炎热来到，他不恐惧，叶子常青；干旱之年，他毫无忧虑，且不停止结果。}\BibleRef{耶 17:5} 听到这以预言形象语言发出的咒诅，落在信赖世人的人身上；以及类似方式的祝福，落在信赖天主的人身上；询问者可能对我们的教义产生疑问，我们教导的不仅是\Person{基督}是天主，所以我们信赖的不是人，而且他也是人，因为他取了我们的本性。所以有人错误地否认\Person{基督}的人性，却承认他的神性。另一些人则肯定他的人性，却否认他的神性，因而要么成为不信者，要么犯下信赖人的罪。于是询问者可能会说，先知只说了\Person{基督}是天主，丝毫未提他的人性；而在我们的宗徒教义中，\Person{基督}不仅是我们可以安心信赖的天主，而且是天主与人之间的中保——那人\Person{耶稣}。先知在似乎自我纠正并补充遗漏的话中解释了这一点：\BibleQuote{他的心诡诈至极；且是人，谁能识透？}\BibleRef{耶 17:9} 他是人，为的是以仆人的形象医治心硬的人，使他们承认那位为他们而成为人的天主，好让他们的信赖不在人身上，而在神人身上。他是人，取了仆人的形象。\Quote{谁能识透他？}因为\BibleQuote{他虽具有天主的形体，却没有以自己与天主同等为应当把持不舍的。}\BibleRef{斐 2:6} 他是人，因为\BibleQuote{圣言成了血肉，居住在我们中间。}\Quote{谁能识透他？}因为\BibleQuote{在起初已有圣言，圣言与天主同在，圣言就是天主。}\BibleRef{若 1:1} 他的心确实满忧患。即使对于自己的门徒，他心忧，当他说：\BibleQuote{这么长久的时间，我同你们在一起，你们还不认识我吗？}\BibleRef{若 14:9} \Quote{这么长久的时间}对应\Quote{他是人}，而\Quote{你们还不认识我吗？}对应\Quote{谁能识透他？}而且同一位就是他说：\BibleQuote{谁看见了我，就看见了父。}\BibleRef{若 14:9} 所以我们的信赖不在人身上，以致受先知咒诅，而是在神人身上，即天主子、救主\Person{耶稣基督}，天主与人之间的中保。在仆人的形象中，父比他大；在天主的形象中，他同父平等。\par

9. 在\Person{依撒意亚}中我们读到：\BibleQuote{人的高傲必被贬抑，上主独在那日受尊崇。人们必将他们手所造之物，抛入岩石的裂缝和岩穴及地洞中，为上主的威严和他光荣的威能，当他起来震撼大地时。在那日，人必将他们的金银偶像，即他们制造来朝拜的，抛给鼹鼠和蝙蝠，视为无用和有害。}\BibleRef{依 2:17} 或许询问者本人——\Person{法乌斯图斯}以为他会嘲笑说不相信希伯来先知——已将手造的偶像藏在某条裂缝、或用岩穴或地洞中。或者他认识一位朋友、邻居或同胞，因敬畏上主而做了这事；上主借世上君王严厉的禁令，如今君王都服事并俯伏在他面前，正如先知所预言的，他震撼大地，即粉碎世俗人顽固的心。询问者不太可能不相信希伯来先知，当他发现他们的预言应验了，也许就在他自己身上。\par

10. 人们更可能担心，询问者在如此丰富的证据面前会说：基督徒们是在所述事件已开始发生时编写了这些著作，为的是使这些事看起来不像是出于纯粹人的意图，而仿佛是神圣预言过的。人们可能担心这一点，要不是因为犹太人这个广泛分布且广为人知的民族，这个\Person{加音}——身上带着印记使任何人不得杀他；这个\Person{含}——他兄弟的奴仆，背负着教导他们的书籍作为重担。我们从犹太人的手抄本证明，这些事并不是我们事后编写来迎合事件的，而是早已作为预言在犹太民族中发表和保存的。这些预言现在在应验中得到解释：其中隐晦之处——因为那些事发生在他们身上，作为预象，并且是为了我们这些时代终结的人而写的——如今已变得明白；那些隐于未来阴影中的事，现在在现实经验的光照中清晰可见。\par

11. 询问者可能会提出一个难点：这些预言在其书中被发现的人，并没有在福音上与我们合而为一。但是，当他确信这一点也被预言了时，他就会感到证据是多么有力。关于犹太人不信的预言，没有人能避免看见，没有人能假装对它们视而不见。没有人能怀疑依撒意亚说以下的话时是指犹太人说的：\BibleQuote{牛认识自己的主人，驴也认识自己主人的槽，以色列却不知情，我的百姓却不理会。}\BibleRef{依 1:3} 或者再如，宗徒所引用的话：\BibleQuote{我整天向悖逆和抗拒的人民伸出我的手。} 特别当他说道：\BibleQuote{天主给了他们悔恨的精神，眼不能见，耳不能听，也不能明白。} 以及许多类似的段落。如果询问者反对说，如果天主使犹太人眼瞎，以致他们不认识基督，那不是犹太人的错，我们就当以最简单的方式尽力使他明白：这眼瞎是天主所知道的其它隐秘罪过的公正惩罚。我们应当证明，宗徒论到某些人时承认了这一原则，他说：\BibleQuote{天主任凭他们随从心中的情欲，陷于可弃绝的心思，行不合理的事。}\BibleRef{罗 1:28} 并且先知们自己也说到这一点。因为，回头来看耶肋米亚的话：\BibleQuote{祂是人，谁能认识祂？} 以免犹太人因不认识而有所推诿——因为如果他们认识了，如宗徒所说，\BibleQuote{他们就不会把荣耀之主钉在十字架上了}\BibleRef{格前 2:8}——先知接着指出，他们的无知是隐秘罪行的结果；因为他说：\Quote{我上主洞察人心，考验肺腑，依照各人的行径和作为的果实，赐给每人应得的。}\par

12. 如果询问者心中的下一个难点来自那些称为基督徒的人中的分裂和异端，他就会知道这一点也被先知们注意到了。因为，仿佛很自然，在关于犹太人的盲目得到满足之后，这个来自基督徒分裂的反对意见就会出现，耶肋米亚在预言中注意到这个顺序，立即在所引用的段落中补充说：\BibleQuote{鹧鸪吵闹，聚集它没有孵化出来的，不凭理智聚敛财富。} 因为鹧鸪以好争吵闻名，并且常常因其好斗的热情而被捕。所以异端者讨论不是为了寻找真理，而是抱着固执的决心一定要以某种方式取胜，以便他们能像先知所说的，聚集他们没有产生出来的。因为他们所误导的人已经是因福音而生的基督徒，异端者的基督徒身份误导了他们。所以他们聚敛财富不是凭理智，而是轻率匆忙。因为他们不考虑，他们聚集作为自己财富的跟随者是从真正原始的基督徒团体中取来的，并剥夺了其益处；正如宗徒用以下的话描述这些异端者：\BibleQuote{就如\Person{雅乃斯}和\Person{扬布勒}曾反对\Person{梅瑟}，照样他们也反对真理：心术败坏，在信德上可弃绝。但他们不能再有进展，因为他们的愚昧将暴露在众人面前，正如那两个人的愚昧一样。}\BibleRef{弟后 3:8} 所以先知继续论到聚集它没有孵化出来的鹧鸪说：\BibleQuote{在它年日当中，它们要离开它，最后它成了一个愚人。} 这就是说，那起初以貌似高明智慧误导人的人，将成为愚人，即将被视为愚人。当他的愚昧暴露在众人面前时，他就会被看见，并且对那些起初视他为智者的人，他那时将是一个愚人。\par

13. 仿佛预见到询问者接下来会问，一个年轻的门徒，尚不能在众多错误中分辨真理，凭何种明显的标记可以找到基督的真教会，因为如此多预言的清楚应验迫使他相信基督，先知在接下来的话中回答了这个问题，并教导他所预言的基督教会是显而易见可见的。他的话是：\BibleQuote{荣耀的高位是我们的圣所。}\BibleRef{耶 17:12} 这荣耀的宝座就是教会，宗徒论到它说：\BibleQuote{天主的殿是圣的，你们就是这殿。}\BibleRef{格前 3:17} 主也预见到教会的显著是帮助可能被误导的年轻门徒的，他说：\BibleQuote{建在山上的城，是不能隐藏的。}\BibleRef{玛 5:14} 既然如此，荣耀的高位是我们的圣所，就不应理会那些要带领我们进入宗派主义的人，他们说：\Quote{看，基督在这里}，或\Quote{看，在那里}。看这里，看那里，说的是分裂；但真正的城是在山上，而这山正如我们在\Person{达尼尔}先知书中所读到的，\Quote{从一块小石头长起来，直到充满全地。}\BibleRef{达 2:34} 也不应理会那些宣称某种只有少数人知道的隐秘奥秘的人，说：\Quote{看，祂在内室；看，在旷野}；因为建在山上的城是不能隐藏的，荣耀的高位是我们的圣所。\par

14. 在考虑了这些预言的应验实例——关于君王和人民先是作为迫害者，后来成为信徒，关于偶像的毁灭，关于犹太人的盲目，关于他们对所保存经卷的见证，关于异端者的愚昧，关于真基督徒教会的尊贵——之后，询问者最有理由接受这些先知关于基督的天主性的见证。毫无疑问，如果我们一开始就敦促他相信尚未应验的预言，他可能会公正地回答：\Quote{我与这些先知有何相干？我没有任何证据证明他们的话是真的！} 但是，鉴于如此多显著预言的明显应验，任何诚实的人都不会轻视那些在早期被如此庄严地认为值得预言的事物，也不会轻视那些作出预言的人。无论是关于久远过去的事件，还是关于将来仍要发生的事，我们最值得信赖的，莫过于那些其话语得到如此多著名预言已应验之证据支持的人。\par

15. 若在\Person{西比尔}（或众西比尔），或在\Person{俄耳甫斯}，或在\Person{赫尔墨斯}（若真有此人），或在任何其他外教诗人、神学家、智者或哲学家中，有任何关于天主或天主子的事被教导或预言，这可能有助于驳斥外教谬误，但不能引导我们相信这些作者。因为他们虽然不由自主地谈论我们所敬拜的天主，但他们要么教导同胞崇拜偶像和魔鬼，要么允许他们这样做而不敢抗议。然而，我们的神圣作者，依靠天主的权威和帮助，在他们的人民中建立并维护了一个政体，在该政体下，外教习俗被谴责为亵圣。若这人民中有人陷入崇拜偶像或崇拜魔鬼，他们要么受到法律惩罚，要么遭到先知可怕的谴责。他们崇拜独一天主，天地的创造者。他们拥有礼仪；但这些礼仪是预言的，或是未来事物的预象，并将在所指事物出现时停止。整个国家是一个伟大的先知，其君王和司祭象征性地受傅；这傅油中断了，并非由于犹太人自己的意愿——他们因不信而处于无知中——唯有在祂来临时，祂是天主，以属神的恩宠傅油，超越祂的同伴，至圣者，真正的君王来治理我们，真正的司祭来为我们奉献自己。总之，外教才智关于基督来临的预言与神圣预言之间的差别，如同魔鬼的承认与天使的宣告之间的差别。\par

16. 藉着这些论证——若我们与一个在外教中长大的人讨论，这些论证可以扩展，并可以用更多的证据支持——\Person{福斯图斯}带来我们面前的这个探究者肯定会被引导相信，除非他宁愿选择自己的罪过而不要自己的救恩。作为一名信徒，他会被接纳，安放在天主教会的怀抱中，并在适当的时候得到关于他所应遵守行为的教导。他会看到许多人不履行所要求的义务；但这不会动摇他的信德，即使这些人属于同一教会，并与他共享同一圣事。他会明白，很少有人分享天主的产业，而许多人分享其外在标记；很少有人结合于圣洁的生活，并结合于那藉着赐予我们的圣神倾注在我们心中的爱德之恩——这是一股隐秘的泉水，没有外人能接近；而许多人参与圣事的隆重礼仪，谁若不相称地吃这饼、喝这杯，便是吃喝自己的审判，而谁若忽略不吃，在他内就不会有生命（\BibleRef{若 6:54}），并因此永远不得永生。他也会明白，与恶人的众多相比，善人被称为少数，但散布在世界各处，有许多人在莠子中生长，与糠秕混杂，直到收割和簸扬的日子。正如这在福音中所教导的，也由先知所预言。我们读到：\BibleQuote{如百合花在荆棘中，我的爱人在女子中也是如此。}\BibleRef{歌 2:2}又：\BibleQuote{我曾居住在\Place{刻达尔}的帐幕中；在那憎恨和平的人中，是和平的。}又：\BibleQuote{要在额头上标记那些为了我人民在他们中间所行的恶事而叹息哀哭的人。}\BibleRef{则 9:1}这个探究者会被这样的段落所坚定；现在他已是与圣人同为公民、是天主家里的人，不再是远离以色列的外人，而是真正的以色列人，心中毫无诡诈；他会从一个无诡诈的心学会说出\Person{耶肋米亚}已经引用的段落中的话：\BibleQuote{上主，以色列的忍耐：凡离弃你的，都要惊慌。}在谈到那只喧闹的鹧鸪，聚集它所没有生下的之后；在赞扬那建在山上的城，不能隐藏，以防止异端者将人从天主教会拉走之后；在说了\BibleQuote{我们的圣所是荣耀的崇高宝座}之后，他似乎自问：我们当如何看待那些与教会混杂的恶人，他们随着教会的扩展而越发增多，当万民在基督内合而为一时？随后的话是：\BibleQuote{上主，以色列的忍耐。}忍耐是必需的，以便遵守命令：\BibleQuote{让两者一起生长直到收获。}\BibleRef{玛 13:30}对恶人的不耐烦可能导致离弃善人——严格地说，他们是基督的身体——而离弃他们就是离弃祂。所以先知继续说：\BibleQuote{凡离弃你的，都要惊慌；那些离弃你而转向尘世的人，要蒙羞。}大地是那信赖自己并引诱他人信赖他的人。所以先知补充说：\BibleQuote{愿他们倾覆，因为他们离弃了上主，生命的泉源。}这就是鹧鸪的呼喊：它得到了生命的泉源，并将赐予它；于是人们被聚集到它那里，离开基督，好像基督——他们曾宣称祂的名字——没有实现祂的应许。鹧鸪聚集它所没有生下的人。为了做到这一点，它宣称：基督所应许的救恩在我这里；我将赐予它。与此相反，先知说：\BibleQuote{上主，医治我，我便痊愈；拯救我，我便得救。}所以我们在宗徒处读到：\BibleQuote{不可在人身上夸耀。}\BibleRef{格前 3:21}或先知的话：\BibleQuote{你是我的赞美。}\BibleRef{耶 17:14}这就是宗徒和先知教义的一个教导样本，人可藉此被建造在宗徒和先知的根基上。\par

17. \Person{福斯图斯}没有告诉我们，他如何向外邦人证明基督的天主性；他让外邦人说：\Quote{我不相信先知们为基督作证，也不相信基督为先知们作证。}认为这样的人会相信基督关于祂自己所说的话，而却不信祂关于别人所说的话，那是荒谬的。因为若是他认为基督在一件事上不值得信任，那么他必定认为祂在所有事上都不值得信任，或者至少当祂谈论自己时比谈论别人时更不可信。也许，在此之外，\Person{福斯图斯}会为他读\Person{西彼拉}和\Person{俄耳甫斯}，以及他能找到的任何关于基督的外邦预言。但他怎能做到这一点，当他承认自己一无所知？他说：\Quote{倘若正如人们所说，在\Person{西彼拉}，或在被称为\OriginalTerm{三倍伟大的}{\GreekText{Τρισμέγιστος}}\Person{赫尔墨斯}，或\Person{俄耳甫斯}，或任何外邦诗人那里，能找到关于基督的任何预言。}他怎能阅读自己一无所知、且仅凭传闻认为存在的著作，来给一个既不信先知又不信基督的人听？那么，他会做什么呢？他会拿出\Person{摩尼}作为基督的见证人吗？摩尼教徒所做的恰恰相反。他们利用基督之名广泛传播的芬芳来使\Person{摩尼}被人接受，好让他们毒杯的边缘被这蜂蜜所甜化。他们抓住基督对门徒的应许，即祂要派遣\OriginalTerm{护慰者}{\GreekText{παράκλητος}}，也就是安慰者或辩护者，说这位护慰者就是\Person{摩尼}，或在\Person{摩尼}内，从而悄悄进入那些不知道基督所应许的那位何时真正来临的人心中。那些读过称为\BibleBook{宗徒}{大事录}的正典书卷的人，会找到基督应许的提及及其应验的记载。因此，\Person{福斯图斯}无法给那探求者任何证明。没有人会如此愚昧，以致当拒绝基督的权威时，却接受\Person{摩尼}的权威。难道他不会以嘲笑（若非愤怒）回答说：你为何要我相信波斯书籍，而你却禁止我希伯来书籍？摩尼教徒无法抓住那探求者，除非他在某种程度上已先被基督徒的真理所说服。当他发现那人愿意相信基督时，他就用\Person{摩尼}所呈现的基督形象来迷惑他。所以，鹧鸪聚集它未曾孵出的。你们这些被他聚集的人何时会离开他？你们何时会看出他是个愚人，竟告诉你们希伯来的见证对不信者毫无价值，而对信者则是多余的？\par

18. 若是信者要抛弃所有引导他们相信的书卷，我看不出他们为何还要继续阅读福音本身。福音对那探求者而言也必毫无价值，因为按\Person{福斯图斯}可怜之假设，他嘲笑并拒绝基督的权威。而对信者，它必是多余的，倘若关于基督的真实记载对信者是多余的话。并且，若是信者该读福音，为不忘记他所信的，那么他也该读先知书，为不忘记他为何相信。因为若他忘记这个，他的信德就无法坚定。根据这个原则，你们该抛弃\Person{摩尼}的书卷，你们正是根据其权威已经相信那光明——也就是天主——与黑暗争斗，并且为了捆缚黑暗，光明先被吞没捆缚，被黑暗玷污和撕裂，然后通过你们的饮食被恢复、解放、净化、医治，你们为此得到赏报，免受永远被投入黑暗团体的惩罚，与那部分无法被解脱的光明一同。这虚构之事已充分由你们的实践和言语宣扬出来。你们为何寻求书卷的见证，并通过消耗力气在书写手稿的不必要任务上，来加重你们天主的困窘？烧掉你们所有的羊皮卷，连同它们精美的装订；这样你们就摆脱了一个无用的负担，而那位被囚禁在书卷中的天主也将获得释放。若你们能将书煮了吃掉，那将是何等慈悲之举，以利于你们天主的肢体！然而，动物食品的禁令可能带来困难。那么书写必须分担羊皮的污秽。事实上，这应归咎于你们，因为正如你们所说在光明与黑暗的首次战争中所做的那样，你们将笔中洁净之物带到了羊皮的不洁之上。或者，为了颜色的缘故，我们可以反过来看；那么黑暗就是你们的，即你们带来的墨水，对抗白页的光明。如果这些话惹恼了你们，你们更该生自己的气，因为你们相信了这些必然导致这样后果的教义。至于宗徒和先知的书卷，我们阅读它们作为我们信德的记录，以鼓励我们的望德，激发我们的爱德。这些书卷彼此完全和谐；它们的和谐，犹如天上号角的音乐，将我们从世俗的麻木中唤醒，催促我们向着高天召叫的奖赏奔去。宗徒在引用了先知的话\Quote{辱骂你者的辱骂都落在我身上}之后，接着谈到阅读先知书的益处：\Quote{凡是先前所写的，都是为教训我们而写的，为使我们借着圣经所给予的忍耐和安慰，怀着希望。}（\BibleRef{罗 15:4}）如果\Person{福斯图斯}否认这点，我们只能与\Person{保禄}一同说：\Quote{若有人给你们宣讲的福音，与你们所接受的不同，当受诅咒。}（\BibleRef{迦 1:9}）\par

\SourceRecord{Translated by Richard Stothert.；Nicene and Post-Nicene Fathers, First Series；Vol. 4.；Edited by Philip Schaff.；Buffalo, NY: Christian Literature Publishing Co.,；1887.；<http://www.newadvent.org/fathers/140613.htm>.}

% structure: p0001 source=h1 target=section reason=page-title

\section{驳斥福斯托斯，第十四卷}

\Emph{福斯托斯因梅瑟对基督所发的可怕咒诅而憎恶梅瑟。奥斯定比较旧约与新约章节，阐述关于受苦救主的基督宗教教义。}\par

1. \Emph{\Person{福斯托斯}}说：若你们问我们为何不信梅瑟，那乃是因我们对基督的爱与尊敬。最鲁莽的人不会乐于见到一位咒诅他父亲的人。因此我们憎恶梅瑟，与其说因为他亵渎一切人神之事，不如说因他对为了我们救恩而挂在树上的天主圣子基督所发的可怕咒诅。无论梅瑟是有意还是无意这样做，那是你们的问题。无论哪种情况，他都不能被原谅，也不值得相信。他的话是：\BibleQuote{凡挂在树上者，都当受咒诅。}\BibleRef{申 21:23}你们要我信这人，可是，如果他受了默感，他必定是明知而有意地咒骂基督；如果他出于无知，那他就不可能具有神性。两者任选。梅瑟不是先知，而是用他惯常的方式咒骂，无知地陷入了亵渎天主的罪中。或者他确实是神性之人，并且预见了未来；由于对我们得救的恶意，他把咒骂的毒液指向那要在树上完成救恩的那一位。这般伤害圣子的人肯定没有看见或认识圣父。对圣子最终升天一无所知的人，肯定没有预告祂的来临。此外，还要考虑这个咒诅所造成的伤害程度。因为它谴责了所有以这种方式死去的义人、殉道者和各种受苦者，如\Person{伯多禄}、\Person{安德肋}等。如果梅瑟是先知，他绝不会发出如此残忍的谴责，除非他是这些受苦者的死敌。因为他宣布他们不仅在人面前，而且在天主面前都受咒诅。那么，对基督、对祂的宗徒、或对我们若为基督的缘故被钉十字架，还剩下什么祝福的希望呢？这显示梅瑟非常轻率，缺乏一切神圣默感，因为他忽略了人们被挂在树上的原因千差万别，有些是因罪行，有些则是为了天主和正义的事业而受苦。他以这种轻率的方式不加区分地将所有人一同置于同一咒诅之下；而如果他稍有理智，更不用说默感了，如果他希望将十字架的刑罚与其他刑罚区别开来，作为特别可憎的，他应该说：\Quote{凡有罪和不敬虔之人挂在树上者，都当受咒诅。}这样就能区分有罪和无罪。然而，即使这样也不正确，因为基督将那个罪犯从十字架上与祂一同带入了祂父的乐园。那么，对每一个挂在树上之人的咒诅又有什么用呢？那个臭名昭著的强盗巴拉巴肯定没有被挂在树上，而是应犹太人的请求从监狱中被释放，难道他比那个与基督同从十字架到天堂的罪犯更有福吗？再者，有咒诅落在拜日或拜月的人身上。那么，如果我被外教君主强迫拜日，并且因为害怕这个咒诅而拒绝，我是否会因为遭受钉十字架的刑罚而招致另一个咒诅呢？也许梅瑟惯于咒诅一切好事。我们对待他的谴责就如同对待一个老妇的责骂。所以，我们发现他咒诅所有青年男女，当他说：\BibleQuote{凡在以色列中不留下后裔者，都当受咒诅。}\BibleRef{申 25:5}这是直接针对耶稣，按你们的说法，祂生在犹太人中，却没有留下后裔延续家族。这也指向祂的门徒，其中一些祂从所娶的妻子中带走，另一些未婚的祂禁止娶妻。你们看，我们有充分的理由表达我们对梅瑟以大胆风格将咒骂投向基督、投向光明、投向贞洁、投向一切神圣之事的憎恶。你们不能过于强调挂在树上与被钉十字架的区别，正如你们常试图辩解的那样；因为保禄否认这种区别，他说：\BibleQuote{基督由法律的咒骂中赎回了我们，祂自己成了为我们而受的咒骂，因为经上记载：凡挂在树上者，都当受咒诅。}\BibleRef{迦 3:10}\par

2. \Emph{\Person{奥斯定}}答说：虔诚的福斯托斯因基督被梅瑟咒诅而痛苦。他对基督的爱使他憎恨梅瑟。在解释\BibleQuote{凡挂在树上者，都当受咒诅}这话的神圣含义和虔诚之前，我要问这些虔诚的人，既然梅瑟的咒诅并不影响他们的基督，他们为何生梅瑟的气？如果基督挂在树上，祂必定是被钉子钉在树上，祂复活后曾将钉痕给那怀疑的门徒看。因此，祂必定有一个可受伤、可死的身体，而摩尼教徒却否认这一点。如果他们说那些伤痕和钉痕是假的，那么祂挂在树上也是假的。这个基督不受咒诅影响，也就没有理由对发出咒诅的人发怒。如果他们假装对梅瑟咒骂他们所认为的基督的假死而生气，那么我们对他们本人该作何想？他们虽然不咒骂基督，但更糟糕的是，他们使基督成了说谎者。如果咒骂可死性是错的，那么玷污真理的纯洁就是更严重的罪。然而，让我们以这些异端者的刁难为机会，向信徒解释这个奥迹。\par

3. 死亡临到人，作为罪的惩罚，因此它本身也被称为罪；这并不是说人在死亡时犯罪，而是因为罪是死亡的原因。同样，\Quote{舌头}一词，原意是牙齿与上颚之间的肉质器官，但引申指舌头的活动结果。在这个意义上，我们谈论拉丁语和希腊语。\Quote{手}一词，既指我们工作中使用的身体部位，也指用手完成的书写。在这个意义上，我们说笔迹被证实是某人的手，或认出朋友的手。书写当然不是身体的一部分，但被称为手，因为是手完成的。同样，罪既指应受惩罚的恶行，也指作为罪之结果的死亡。基督在以应受死亡的意义上没有罪，但祂为了我们承担了罪，即因罪而临到人性的死亡。这就是挂在木头上的；这就是梅瑟所诅咒的。因此死亡被定罪，好使其统治终止；被诅咒，好使其被消灭。基督这样承担了我们的罪，其定罪就是我们的解救，而留在罪的管辖之下就是被定罪。\par

4. 福斯图斯对那加于罪、死亡和人的必死性——基督因人的罪而承受了这些，尽管祂自己无罪——的诅咒有什么可惊奇的？基督的身体来自亚当，因为祂的母亲童贞玛利亚是亚当的后裔。但天主在乐园中说：\BibleQuote{在吃的那一天，你必定要死亡。}这就是挂在木头上的诅咒。一个否认基督死亡的人也会否认基督被诅咒。但相信基督死了并承认死亡是罪的果实、本身也被称为罪的人，当他听到宗徒说\BibleQuote{我们的旧人已与祂同钉十字架。}（\BibleRef{罗 6:6}）时，就会明白梅瑟所诅咒的是谁。宗徒大胆地说论基督：\BibleQuote{祂为我们成了诅咒；}因为他也敢说\BibleQuote{祂为众人死了。}\Quote{祂死了}和\Quote{祂被诅咒}是同一回事。死亡是诅咒的后果；而所有的罪——无论是指应受惩罚的行为，还是随之而来的惩罚——都是被诅咒的。基督虽无罪，却担当了我们的惩罚，为要消除我们的罪责，并除掉我们的惩罚。\par

5. 这些不是我的猜测，而是宗徒不断以足够的强调肯定的，足以唤醒疏忽者并封住反对者的口。他说：\BibleQuote{天主派遣了祂的儿子，带着罪恶肉身的形状，为要以罪恶肉身在肉身上定罪。}（\BibleRef{罗 8:3}）基督的肉身不是有罪的，因为它不是由玛利亚以普通的方式所生；但由于死亡是罪的后果，这肉身因有死，便有了罪恶肉身的形状。在以下的话中，这被称为罪：\BibleQuote{为要以罪恶肉身在肉身上定罪。}他又说：\BibleQuote{祂使那不认识罪的，替我们成了罪，好叫我们在祂内成为天主的正义。}（\BibleRef{格后 5:21}）为什么梅瑟不能称呼保禄称为罪的东西为诅咒？在预言中，先知与宗徒一同分享异端者的辱骂。因为凡在先知那里对\Quote{诅咒}一词找岔子的，也必须在宗徒那里对\Quote{罪}一词找岔子；因为诅咒和罪是相连的。\par

6. 如果我们读到\BibleQuote{凡挂在木头上的，都是被天主诅咒的}，加上\Quote{被天主}一词并无困难。因为如果天主不憎恨罪和我们的死亡，祂就不会派遣祂的儿子来承担和消灭它们。天主诅咒祂所憎恨的，这并不奇怪。因为祂愿意赐给我们那将在基督来临时拥有的不朽，是与祂在基督死亡时，当我们的死亡挂在十字架上时所怀的怜悯相称的。如果梅瑟诅咒凡挂在木头上的，那肯定不是因为他没有预见义人将被钉十字架，而是因为他预见异端者会否认主的死亡是真实的，并试图反驳这诅咒应用于基督，以便他们可以反驳祂死亡的现实。因为如果基督的死不是真实的，那么当祂被钉十字架时，就没有任何被诅咒的东西挂在十字架上，因为钉十字架不可能是真实的。梅瑟从遥远的过去向这些异端者呼喊：你们否认基督死亡之现实的推诿是徒劳的。凡挂在木头上的都是被诅咒的；不是这个或那个，而是绝对每一个。什么？天主子？是的，确实。这正是你们反对的，也是你们急于逃避的。你们不允许祂为我们被诅咒，因为你们不允许祂为我们而死。豁免亚当的诅咒意味着豁免亚当的死亡。但是，正如基督作为人并为了人而忍受了死亡，同样，作为天主子，祂常在自己的正义中活着，却为了我们的过犯而死去，祂作为人并为了人，承受了伴随着死亡的诅咒。并且，正如祂在承担我们的惩罚时所取的肉身中死去，同样，祂在自己的正义中永远受祝福，却在为我们承担惩罚而受的死亡中，为我们的过犯而成了被诅咒的。并且\Quote{每一个}这个词是为了制止那些无知的好事者，他们否认这诅咒指向基督，并因此，因为诅咒伴随着死亡，就会导致否认基督真实死亡。\par

7. 真正相信福音教理的人会明白，\Person{梅瑟}说基督是可咒骂的，并非在祂的天主性尊严方面责备祂，而是说祂作为我们的替身挂在木架上，承担我们的惩罚。这正如\OriginalTerm{摩尼教信徒}{Manichaei}否认基督有可朽的身体、从而真正受死，这并不是对基督的赞美。在先知的咒骂中，有对基督谦卑的赞美；而在异端者虚假的敬意中，却有一项虚假的指控。因此，如果你否认基督曾受咒骂，就必须否认祂曾死亡；然后你就要面对宗徒们，而不是\Person{梅瑟}。承认祂死了，你也可以承认祂没有取我们的罪，却取了罪的惩罚。罪的惩罚不可能是可祝福的，否则它就是可渴望之事。咒骂是由天主的正义所宣布的，我们若从其中被救赎，便是好事。那么，承认基督死了，你就可以承认祂为我们承担了咒骂；并且当\Person{梅瑟}说：\BibleQuote{凡挂在树上的，都是可咒骂的}，他其实是说：挂在树上就是有死，或实际死亡。他也许会说：\Quote{凡是有死的，都是可咒骂的}，或\Quote{凡是死亡的，都是可咒骂的}；但先知知道基督将要死在十字架上，并且异端者会说祂挂在树上只是外表，并没有真的死。所以他喊出\Quote{可咒骂的}，意思是祂真的死了。他知道，罪人的死亡——基督虽无罪却承担了——来自那句咒骂：\Quote{你摸它，必定要死}。同样，挂在那杆上的蛇也表明基督并没有假装死亡，而是蛇通过其致命的计谋使人类陷入的真实死亡，被悬挂在基督苦难的十字架上。摩尼教信徒转而不看这真实的死亡，因此他们没有从蛇的毒液中被治愈，正如我们读到在旷野中，凡仰望的都得了痊愈。\par

8. 确实，有些人愚昧地区分\Quote{挂在树上}和\Quote{被钉十字架}。因此有些人解释这段经文是指\Person{犹达斯}。但他们怎么知道他是挂在木头上还是石头上呢？\Person{福斯特}说宗徒迫使我们把这些话指向基督，这是对的。这种愚昧的天主教徒成了\OriginalTerm{摩尼教信徒}{Manichaei}的猎物。他们抓住这些人，用诡辩纠缠他们。当我们陷入这异端并坚持它时，我们就是这样的人。当我们不是凭自己的力量，而是因天主的仁慈被拯救时，我们就是这样的人。\par

9. \Person{福斯特}指控\Person{梅瑟}不放过任何属于人或属于天主的事，他说的这种对神圣事物的攻击是什么？他提出指控却不加证明。相反，我们知道\Person{梅瑟}对一切真正神圣的事物给予了应有的赞美，在人事上，考虑到他的时代和他那时期的恩宠，他是一位公正的统治者。当我们看到\Person{福斯特}指控的任何证据时，再证明这一点也不迟。谨慎地提出指控也许很聪明，但这种聪明却使拥有它的人毁灭，这是极大的不谨慎。站在真理一边聪明是好的，但反对真理聪明则是可悲的。\Person{福斯特}说\Person{梅瑟}不放过任何人事或天主事；不是说他不放过任何神或人。如果他说\Person{梅瑟}不放过天主，很容易回应指出\Person{梅瑟}处处尊崇真天主，宣称祂是天地的创造者。又如果他说\Person{梅瑟}不放过任何神，他就会向基督徒暴露自己是\Person{梅瑟}所谴责的假神的崇拜者；这样，他就无法聚集他所没有产生的东西，因为幼雏会在慈母教会的翅膀下避难。\Person{福斯特}试图诱骗婴孩，说\Person{梅瑟}不放过任何神圣事物，他不想以信仰众神——这明显与基督教对立——来惊吓基督徒，同时又显得站在外邦人一边反对我们；因为他们知道\Person{梅瑟}说了许多明确而尖锐的话反对异教人的偶像和神，那些都是魔鬼。\par

10. 如果\OriginalTerm{摩尼教信徒}{Manichaei}因此不赞同\Person{梅瑟}，就让他们承认自己是偶像和魔鬼的崇拜者。实际上，这可能是事实而他们自己没有意识到。宗徒告诉我们：\BibleQuote{在末世，有些人将背弃信德，听从欺诈的神和魔鬼的教义，以虚伪的言谈讲论。}\BibleRef{弟前 4:1} 若不是从喜爱谎言的魔鬼那里，还能从哪里产生出基督的苦难和死亡是不真实的、祂所显示伤痕也是虚假的念头呢？这些难道不正是说谎魔鬼的教义，教导基督——真理本身——是一个骗子吗？此外，摩尼教信徒公然教导崇拜——若不是崇拜魔鬼，至少是崇拜受造之物，宗徒对此谴责说：\BibleQuote{他们崇拜事奉受造之物，而不事奉造物主。}\BibleRef{罗 1:25}\par

11. 正如摩尼教徒在其虚幻传说中对偶像和魔鬼有一种无意识的敬拜，同样他们也在对太阳和月亮的敬拜中有意地事奉受造物。在他们所谓的对造物主的事奉中，他们实际上是在事奉自己的幻想，而根本不是事奉造物主。因为他们否认天主创造了那些宗徒明确宣布为天主的受造物的事物，当时宗徒论及食物说：\Quote{\BibleQuote{凡天主的受造物都是好的，若以感恩之心领受，没有一样是可弃绝的。}}（\BibleRef{弟前 4:4}）这是纯正的道理，你们不能容忍，便转向荒唐言。宗徒称赞天主的受造物，但禁止敬拜它；同样，\Person{梅瑟}也对太阳和月亮给予了应有的赞美，同时他又陈述了它们由天主所造、并被他安置在轨道上的事实——太阳管理白日，月亮管理黑夜。或许你们以为\Person{梅瑟}没有放过任何神圣之物，只因为他禁止敬拜太阳和月亮，而你们却在敬拜中向它们四面转面。但太阳和月亮并不因你们虚假的赞美而喜悦。是魔鬼、那悖逆者，以虚假的赞美为乐。天上的诸权势，未因罪而堕落的，愿意它们的造物主在它们内受到赞美；它们真正的赞美是不侵犯造物主的。当它们被说成是他的肢体或他本体的部分时，他受到了冒犯。因为他是完美而自立的，非由他出，不分割或散布于空间，而是不变地自存、自足、且在他自己内享福。在他美善的丰盈中，他藉他的话发言，它们便造成了；他命令，它们便被创造了。如果地上的形体是好的——宗徒论及此事时曾说没有食物是不洁的，因为凡天主的受造物都是好的——那么天上的形体更是如此，太阳和月亮是其中的主要者；因为宗徒又说：\Quote{\BibleQuote{地上形体有一种荣耀，天上形体有另一种荣耀。}}（\BibleRef{格前 15:40}）\par

12. \Person{梅瑟}在禁止敬拜太阳和月亮时，并没有责备它们。他称赞它们是天上的形体；同时他也赞美天主是天上和地上万物的造物主，并且不允许天主因有人将应归于他的敬拜给予那些只因其依赖他而被称颂的事物而受到侮辱。\Person{福斯图斯}以他对\Person{梅瑟}咒骂敬拜太阳和月亮之人的反对意见的巧妙为傲。他说：\Quote{如果我在一个外教君王之下被迫敬拜太阳，并且因惧怕这咒骂而拒绝，那么我若因受十字架之刑而遭受这另一咒骂，我岂不陷入困境吗？}没有外教君王强迫你敬拜太阳；就是太阳本身，若它在地上为王，也不会强迫你，正如如今它也不愿受敬拜。正如造物主容忍亵渎者直到审判，同样这些天体也容忍它们被迷惑的敬拜者直到造物主的审判。应当注意，没有一个基督教君王会强制人敬拜太阳。\Person{福斯图斯}举了一个外教君王的例子，因为他知道他们的太阳敬拜是外教习俗。然而，尽管这与基督教相对立，但鹧鸪却取了基督的名，为要聚集它未曾生下的东西。对这反对意见的回答很容易，真理的力量将很快粉碎这困境的两只角。那么，假设一个基督徒受到王权威胁，若他不敬拜太阳就要被挂在树上。你们说，如果我避免法律对敬拜太阳者的咒骂，我就招致同一条法律对挂在树上者的咒骂。这样你们就陷入了困境；只不过你们是敬拜太阳，并没有任何人强迫。但一个真正的基督徒，建立在宗徒和先知的基础上，能够区分这些咒骂及其原因。他看出一个是指挂在树上的可朽身体，另一个是指敬拜太阳的心灵。因为虽然身体在敬拜中鞠躬——这也是严重的罪行——但所敬拜对象的信念或想象是一种心灵的行为。两处咒骂所指的死，一处是身体的死，另一处是灵魂的死。受身体死亡之咒——这将在复活时被除去——比受灵魂死亡之咒更好，后者连同身体一同被判入永火。主用这言语解决了这个难题：\Quote{\BibleQuote{不要害怕那些杀害身体却不能杀害灵魂的；但要害怕那有能力将灵魂和身体投入地狱之火的。}}（\BibleRef{玛 10:28}）换言之，不要害怕身体死亡的咒骂，它到时候会除去；但要害怕属灵死亡的咒骂，它导致灵魂和身体永远的折磨。请确信，\Quote{\BibleQuote{凡挂在树上的，都是可咒骂的}}不是老妇的咒骂，而是预言性的宣告。\Person{基督}藉着咒骂除去咒骂，正如他藉着死亡除去死亡，藉着罪恶除去罪恶。在\Quote{\BibleQuote{凡挂在树上的，都是可咒骂的}}这句话中，并无更多亵渎，正如宗徒的话：\Quote{他死了}，或\BibleQuote{我们的旧人已与他同钉十字架}（\BibleRef{罗 6:6}），或\BibleQuote{藉着罪恶他定罪了罪恶}（\BibleRef{罗 8:3}），或\BibleQuote{他使那不知罪的替我们成为罪}（\BibleRef{格后 5:21}），以及许多类似的话。那么，当你们对基督的咒骂发出惊呼时，你们就是在对他的死亡发出惊呼。如果这不是你们方面的老妇咒骂，那就是魔鬼般的迷惑，它使你们否认基督的死，因为你们自己的灵魂已死。你们教导人基督的死是虚假的，使基督成为你们谎言的领袖，你们用基督徒的名号误导人。\par

13. 如果\Person{福斯图斯}认为\Person{梅瑟}是节欲或童贞的敌人，因为他说：\BibleQuote{凡在以色列中不传种留后的，是可咒骂的}，那就让他们听听\Person{依撒意亚}的话：\BibleQuote{上主如此对一切阉人说：那些遵守我的诫命、选择我所喜悦的事、持守我盟约的人，我要在我殿中，在我墙内，赐给他们比子女更好的名号和地位；我要赐给他们永远不除去的名。}（\BibleRef{依 56:4}）。尽管我们的对手不同意\Person{梅瑟}，但如果他们同意\Person{依撒意亚}，那也算有所收获。我们足以知道：同一位天主既藉\Person{梅瑟}发言，也藉\Person{依撒意亚}发言；并且，凡在\Place{以色列}中不传种留后的，是可咒骂的，既由于那时在婚姻中生育子女（因为延续人民是公民的义务），也由于现在每一个属灵出生的人不应满足于现状，而应各按自己的能力，藉宣讲\Person{基督}来产生基督徒，以寻求属灵的增衍。因此，两约的时代都可用这句话简要描述：\BibleQuote{凡在以色列中不传种留后的，是可咒骂的。}\par

\SourceRecord{Translated by Richard Stothert.；Nicene and Post-Nicene Fathers, First Series；Vol. 4.；Edited by Philip Schaff.；Buffalo, NY: Christian Literature Publishing Co.,；1887.；<http://www.newadvent.org/fathers/140614.htm>.}

% structure: p0001 source=h1 target=section reason=page-title

\section{驳斥福斯特斯，第十五卷}

\Emph{福斯特斯因旧约没有为基督留有余地而拒绝它。基督，独一的净配，为祂的新娘教会已足够。奥斯定尽其所能地回答，并责备摩尼教信徒妄称自己是基督的新娘。}\par

1. \Emph{福斯特斯}说：我们为何不接受旧约？因为器皿满了，浇在其上的东西不被接纳，而任其溢出；饱足的胃拒绝它不能容纳之物。同样，犹太人满足于旧约，拒绝新约；而我们既从基督领受了新约，便拒绝旧约。你们两者都接受，因为你们对两者都仅半满，且互相不能成全，反而败坏。因为半满的器皿不宜用与已有之物性质不同的任何东西来填满。若其中有酒，当以酒填满；有蜜，以蜜；有醋，以醋。因为将胆汁加于蜜，水加于酒，碱加于醋，不是增加，而是掺假。这就是我们不接受旧约的原因。我们的教会，基督的新娘，一位富足新郎的贫穷新娘，满足于拥有她的丈夫，藐视次等情人的财富，轻视旧约及其作者的恩赐，并出于对自己品格的尊重，只接受她丈夫的书信。我们将旧约留给你们的教会，她像一位对配偶不忠的新娘，喜爱另一个人的书信和礼物。这个败坏你们贞洁的情人——希伯来人的天主，在祂的石版上应许你们金银、丰盛的食物和客纳罕地。如此低劣的赏报诱使你们在基督赐予丰厚嫁妆后对祂不忠。希伯来人的天主借此吸引基督的新娘。你们必须知道，你们被骗了，这些应许是虚假的。这位天主贫乏且行乞，不能成就祂所应许的。因为若祂不能将这些赐给听从祂如仆人的合法妻子——会堂，又怎能赐给你们这些外人，且骄傲地将祂的轭从你们颈上挣脱的人？继续吧，正如你们已开始的，将新布补在旧衣上，将新酒装在旧皮囊里，事奉两个主人而不讨任何一方喜悦，使基督教成为半马半人的怪物；但请允许我们只事奉基督，满足于祂不朽的嫁妆，并效法那说\BibleQuote{我们的充足乃出于天主，祂使我们成为新约的合格仆役}的宗徒。\BibleRef{格后 3:5} 我们对希伯来人的天主毫无兴趣；因祂既不能成就应许，我们也不愿祂成就。基督的慷慨使我们漠视这外人的谄媚。这妻子与丈夫关系的比喻得到保禄的认可，他说：\BibleQuote{有丈夫的女人，丈夫活着的时候，她被律法约束；丈夫若死了，她就脱离了丈夫的律法。所以，丈夫活着的时候，她若归于别的男人，便被称为淫妇；丈夫若死了，她就脱离了律法，虽归于别的男人，也不是淫妇。}\BibleRef{罗 7:2} 这里他表明，在与律法的作者结合之前离弃他，视他如死，而归于基督，便是属灵的奸淫。这主要适用于信基督的犹太人，他们应当忘记从前的迷信。我们这些从外邦归向基督的人，视希伯来人的天主不仅如死，且从未存在，无需被告知要忘记他。一个犹太人，当他相信时，应将\OriginalTerm{阿多奈}{Adonai}视如死；一个外邦人应将他的偶像视如死；对一切在归化前曾视为神圣的事物也是如此。一个在放弃偶像崇拜后，既敬拜希伯来人的天主又敬拜基督的人，就像一个淫荡的女人，在一位丈夫死后又嫁给两位丈夫。\par

2. \Emph{奥斯定}答道：凡全心归于基督的人，请说他们能否耐心听这些事，除非基督亲自使他们能够。福斯特斯充满新蜜，拒绝旧醋；而保禄充满旧醋，倒出一半，好让新蜜注入，不是为了保存，而是为了败坏。当宗徒自称\BibleQuote{耶稣基督的仆人，蒙召为宗徒，被选拔为天主的福音}时，这是新蜜。但他接着又说：\BibleQuote{这福音是天主先前藉祂的众先知在圣经中所应许的，论到祂的儿子，按肉身是从达味的后裔生的}\BibleRef{罗 1:1}，这便是旧醋。谁能忍受听到这个，除非宗徒亲自安慰我们说：\BibleQuote{在你们中间不免有异端，好叫那些被认可的人在你们中显明出来？}\BibleRef{格前 11:19} 我们何必重复我们已经说过的话？——即新布与旧衣、新酒与旧皮囊并非指两个约，而是指两种生命和两种希望；主用比喻描述两个约的关系时说：\BibleQuote{因此，凡成为天国门徒的经师，就好像一个家主，从他的宝库里拿出新旧的东西来。}\BibleRef{玛 13:52} 读者或许记得之前说过，或可回查而知。因为若有人试图以两种希望事奉天主，一种属地现世的福乐，另一种属天国的永福，这两种希望不能相容；当后者因某种磨难而动摇时，前者也将失去。因此经上说：\BibleQuote{没有人能事奉两个主人}；基督解释如下：\BibleQuote{你们不能事奉天主而又事奉玛门。}\BibleRef{玛 6:24} 但对于正确理解的人，旧约是新约的预言。即便在古老民族中，圣祖和先知们，他们明白自己所扮演的角色或所参与的事工，已有新约中永生的希望。他们属于新约，因为他们理解并热爱它，尽管只以预像启示。属于旧约的人是那些除了暂时应许之外别无所求，且不懂其作为永恒事物预表的人。但这一切已反复强调得足够了。\par

3. 令人惊异的是，不虔不洁的摩尼教派竟敢夸耀自己是基督贞洁的净配。这对于圣教会中真正贞洁的成员来说，只会使他们记起宗徒警告提防欺骗者说：\BibleQuote{我使你们与一个丈夫结合，将你们如同贞洁的童女献给基督。但我害怕，正如蛇以诡计欺骗了厄娃，你们的心智也会腐化，失去那属于基督的纯洁。}\BibleRef{格后 11:2} 除了我们已领受的福音之外，那些宣讲别的福音的人，他们污蔑天主的法律为旧的，赞扬自己的谎言为新的，仿佛一切新的都是好的，一切旧的都是坏的，他们还想做什么，不就是腐化我们为基督保存的纯洁吗？然而，\Person{若望}宗徒称赞旧诫命，\Person{保禄}宗徒嘱咐我们避开教义上的新奇。作为天主教会——真基督的真净配——一个不配的儿子和仆人，受托分粮给同仆，我也要向她进一句忠言：永远继续躲避摩尼教派亵圣的谬误，这些谬误已经被你自己子女的经验所考验，被他们的康复所谴责。我曾经因那个异端与你的共融隔绝，在陷入本应避免的危险之后，我逃脱了。恢复为你服务，我的经验或许对你有益。若不是你真实而信实的新郎，从祂的肋旁你被造成，借着祂自己的真实宝血获得了罪的赦免，谬误的深渊就会吞没我；我本会成为尘土，被蛇吞吃。不要被真理之名误导。真理在你自己的奶和饼中。他们只有名字，没有实质。你成年的子女固然安全；但我向你的婴儿说话，我的兄弟、儿子和师傅们，他们是你，童贞母亲，纯洁而多产，在你的焦虑的翅膀下抚育至生命，或用婴儿的奶喂养。我恳求这些你的柔弱的子孙，不要被喧嚣的虚浮所诱惑，反而要诅咒任何向他们宣讲别的福音，而非他们在你内所领受的福音的人。我恳求他们不要离开真实而信实的基督，在祂内隐藏着一切智慧和知识的宝藏；不要舍弃祂慈爱的丰盛，祂为敬畏祂的人所积蓄的，为信赖祂的人所成就的。他们怎能期望在宣讲不真实基督的人那里找到真实的话语呢？轻视那些加于你们的指责，因为你们深知你们从新郎所渴望的恩赐是永生，因为祂自己就是永生。\par

4. 这是愚蠢的谎言，说你被引诱到另一个天主那里，那应许丰盛的食物和客纳罕地。因为你能看出，古代的圣人们——他们也是你的子女——是如何被这些预象光照，这些预象是你未来的预言。你不必在意那对石版拙劣的嘲笑，因为那石版过去所象征的石心，并不在你内。因为你们是宗徒的书信，是\BibleQuote{不用墨水，而是用生活天主的圣神写的；不是写在石版上，而是写在心内的血肉版上。}\BibleRef{格后 3:2}我们的对手无知地认为这些话有利于他们，认为宗徒在指责旧约的经纶，然而这些话却是先知的话。宗徒的这段话，是先知先前很久所说的话的应验——这位先知被摩尼教人所拒绝，因为他们相信宗徒却不明白他们。先知说：\BibleQuote{我要从他们身上除掉石心，赐给他们一颗血肉的心。}\BibleRef{厄 11:19}这不是\BibleQuote{不在石版上，而在心内的血肉版上}又是什么呢？因为所谓血肉之心和血肉版，并不是指属血肉的理解力；而是如同血肉有感觉而石头不能，石头的无知觉象征无悟性的心，血肉的知觉象征有悟性的心。所以，不是他们嘲笑你，而是那些说土、木、石头有感觉，且它们的生命比动物的生命更有智慧的人，才该被嘲笑。因此，且不说真理，即使他们自己的虚构也迫使他们承认，那写在石版上的法律，比他们神圣的羊皮卷更纯洁。或者，也许他们更喜欢羊皮而不喜欢石头，因为他们的传说把石头当作王子的骨头。无论如何，旧约的约柜是石版更清洁的遮盖物，胜过他们手稿的山羊皮。嘲笑这些事，同时可怜他们，以显明他们的虚假和荒谬。借着已不再是石心，你能在这些石版中看出适合那心硬的百姓之物；同时你也能在那里找到那石头，你的新郎，由\Person{伯多禄}描述为\BibleQuote{一块活石，虽被人弃绝，却是天主所拣选、所珍贵的。}对于他们，他是\BibleQuote{一块绊脚石，一块跌倒的盘石}；但对于你，\BibleQuote{匠人所弃的石头，已成为屋角的基石。}\BibleRef{伯前 2:4}这一切都由\Person{伯多禄}解释，并且是从先知们引用的，这些异端者与先知们毫无关系。所以，不要害怕读这些石版——它们来自你的丈夫；对别人，石头是无知觉的记号，但对你，是能力和稳固的记号。这些石版是天主的手指所写；你的主以天主的手指驱魔；以天主的手指驱散说谎魔鬼的教义，这些教义烫烙良心。用这些石版，你就能挫败那自称护慰者的诱惑者，免得他以一个神圣的名字欺骗你。因为在逾越节后第五十天，石版被赐下；而在你的新郎——逾越节是其预象——受难后第五十天，天主的手指，圣神，那应许的护慰者，被赐下了。不要害怕这些传达给你古老文字的石版，这些文字如今已变得明白。只是不要处于法律之下，免得恐惧阻止你履行它；而要处于恩宠之下，好使爱，即法律的满全，在你内。因为正是在回顾这些石版时，你的新郎的朋友说：\BibleQuote{不可奸淫，不可杀人，不可贪恋，以及别的任何诫命，都包含在这句话里：你应当爱你的近人如你自己。爱不加害于人，所以爱就是法律的满全。}\BibleRef{罗 13:9}一块石版包含爱天主的诫命，另一块包含爱人的诫命。那最初颁赐这些石版者，亲自来吩咐那些诫命，法律和先知都系于其上。\BibleRef{玛 22:37}在第一诫命中是你们订婚的贞洁；在第二诫命中是你们肢体的合一。在前者中，你们与天主性结合；在后者中，你们聚集一个团体。而这两条诫命等同于十条诫命，其中三条关乎天主，七条关乎近人。这就是那贞洁的石版，你的爱人和你的心爱者从古时就在其中为你预象了十弦琴上的新歌；他自己要为你被钉在十字架上，以致借着罪，他能在肉体中定罪罪过，好使法律的义德在你内得以满全。这就是那配偶的石版，它很可能被不忠的妻子憎恨。\par

5. 我现在转向你们，你们这些受迷惑又迷惑人的摩尼会众——与这么多元素结合，或更确切地说，向这么多魔鬼卖淫，且受孕于亵渎的谎言——你们竟敢诋毁天主教会和你们的主的婚姻为不贞。看哪，你们的情郎，一个平衡受造界，另一个像阿特拉斯一样支撑天地。因为按你们的说法，一个掌管元素的源头，将世界悬于太空；另一个则跪着用肩膀承担重量，以保持平衡。这些神灵在哪里？如果他们如此忙碌，又如何能来拜访你们，用你们温柔舒缓的抚摸来减轻他们的肩膀或手指，消磨闲暇时光？但你们被与你们行淫的邪灵所欺骗，使你们怀上谎言，生出虚妄。你们完全有理由拒绝真天主的讯息，因为这与你们的皮卷相悖，在你们放荡思想的虚妄想象中，你们追随了这么多假神。诗人的虚构比你们的更可敬，至少在这方面：他们不欺骗任何人；而你们书中的寓言，因装出真理的样子，误导了幼稚的人，无论老少，并扭曲了他们的心智。如宗徒所说，\BibleQuote{他们耳朵发痒，离弃真理，偏向无稽之谈。}\BibleRef{弟后 4:4} 你们怎能忍受这些诫命的正道教诲，其中第一条是：\BibleQuote{听，以色列！上主我们的天主是唯一的主。}\BibleRef{申 6:4} 当你们堕落的情感在这么多假神中找到可耻的快乐时，你们岂不记得你们的情歌，在其中你们描述那至高统治者，永世威严，头戴花环，面容如火？拥有这样一个情郎已是可耻；因为贞洁的妻子不会寻求头戴花环的丈夫。你们不能说这个描述或表象有预表意义，因为你们夸耀摩尼无非是向你们讲述简单赤裸的真理，不用比喻的伪装。所以你们歌中的神是一个真正的王，手持权杖，头戴花环。当他头戴花环时，他应该放下权杖；因为柔弱与威严不相称。而且他不是你们唯一的情郎；因为歌中还讲到十二个时令，身着花衣，充满歌声，将鲜花抛向父亲的脸庞。这些是你们十二位大神，各三位在四方围绕第一位神。无人能说这位神如何可能是无限的，当他如此被限定。此外，还有无数的权能、众神和天使的军队，你们说它们并非由神创造，而是从他的本体中产生。\par

6. 你们这样被定罪为敬拜无数的神；因为你们不能忍受那教导独一天主有独一圣子，以及两者有同一圣神的健全道理。而这些，并不是无数，也不是三位神；因为不仅他们的实体是同一的，而且通过这实体的运作也是同一的，尽管他们在物质受造界中有分别的显现。这些事你们不明白，也不能领受。你们是满盈的，如你们所说，因为你们浸没在亵渎的谬论中。难道你们要继续把自己埋藏在这些粗鄙言论之下吗？那么，唱下去吧，如果你们能的话，睁开你们的眼睛，看看自己的羞耻。在这套说谎的魔鬼的道理中，你们被邀请到幸福气候中天使的仙境，以及芬芳的田野，那里有琼浆从树木、山丘、海洋和河流中永远流淌。这些都是你们愚昧心思的虚构，它在这些无聊幻想中自娱。这样的表述有时被用作对神性享受丰盈的比喻性描述；它们引导学者的心灵去探究其隐藏的含义。有时有对肉身感官的物质性表现，如荆棘中的火焰、棍杖变成蛇、蛇变成棍杖、主的外衣不被迫害者分开、虔诚妇人涂抹他的脚或头、群众在他骑驴时前呼后拥的树枝。有时，在睡眠或神魂超拔中，神魂通过取自物质事物的形象得到启示，如\Person{雅各伯}的梯子、在\Person{达尼尔}中非人手凿出的石头长成一座山、\Person{伯多禄}的器皿、以及\Person{若望}所见到的一切。有时，比喻只存在于语言中，如\BibleBook{}{雅歌}、家主为儿子举行婚宴的比喻、荡子的比喻、以及栽种葡萄园租给农夫的人的比喻。你夸耀\Person{摩尼}是最后来的，不是使用比喻，而是解释比喻。他的诠释光照了古代的预象，没有留下任何未解的问题。这种想法得到了以下断言的支撑：古代的预象，无论是神视、行动还是言语，都指向\Person{摩尼}的到来，他将解释一切；而他，知道没有人会跟随他，便采用了一种完全脱离比喻表达的风格。那么，那些田野、荫蔽的山丘、花冠、芬芳的香气，你肉体思想的欲望在其中取乐的，是什么呢？如果它们不是有意义的比喻，那就要么是无聊的幻想，要么是谵妄的梦。如果它们是比喻，那么抛弃那个骗子吧，他以赤裸真理的承诺诱惑你，然后用无聊的故事嘲弄你。他的侍从和那些可怜的受骗追随者惯于用宗徒的那句话作诱饵：\BibleQuote{我们现在是借着镜子，在比喻中观看，但那时就要面对面了。}\BibleRef{格前 13:9}似乎保禄宗徒只是部分知道、部分预言、借着镜子在比喻中观看；而这一切都在\Person{摩尼}来到时被除去，他带来了那圆满的，并面对面启示真理。哦，堕落而不知羞耻！仍然继续说出这样的蠢话，仍然以风为食，仍然拥抱你自己内心的偶像。那么，你曾面对面地看见过那手持权杖的君王、花冠、诸神的军旅、六面光明闪耀的伟大持世者、以及那位被天使队伍围绕的尊高统治者、右手持矛左手持盾的无敌战士、以及那转动火、水、风三轮的著名主宰、以及万有之首的\Person{阿特拉斯}，双肩负着世界，用手臂支撑自己吗？这些，以及成千上万的其他奇观，你是面对面看见的，还是你的诗歌是从说谎的魔鬼那里学来的教义，虽然你不知道？唉！可悲的娼妓，沉迷于这些梦幻，这些就是你所啜饮代替真理的虚妄；被这致命的毒药灌醉，你竟敢以这碑文的戏言，冒犯独生子净配的主妇般的贞洁；因为她不再处于法律的师傅之下，而是在恩宠的管辖之下，既不因活动而骄傲，也不因恐惧而畏缩，她以信德、望德和爱德生活，是那毫无诡诈的以色列人，她听到所写的：\BibleQuote{你的天主上主是惟一的天主。}\par

7. 这些诫命必然与你们相敌对，因为第二条诫命是：\Quote{不可妄呼上主你的天主的名；}\BibleRef{出 20:7}而你们却将虚妄的谎言归给基督自己，祂为了除去肉性心思的虚妄，以真实的身体复活了，这身体是肉眼可见的。同样，第三条关于安息日休息的诫命也与你们相敌对，因为你们被众多不安的幻想所困扰。这三条诫命如何与爱天主相关，你们既没有能力也不愿意明白。你们可耻地固执和狂暴，已达到了愚蠢、虚妄和无价值的顶峰；你们的美丽被毁，你们的秩序已丧失。我了解你们，因为我也曾与你们一样。如今我怎能教导你们，这三条诫命与爱天主相关，万物都出于祂、借着祂并在祂内？你们怎能明白这一点，当你们有害的教义阻止你们理解和遵守那与人相爱相关的七条诫命，而人相爱乃是人类社会之纽带？这些诫命中的第一条是：\Quote{应孝敬你的父亲和你的母亲；}\BibleRef{出 20:12}，\Person{保禄}引用这条作为第一条带许诺的诫命，并亲自重述这训令。但你们被你们的魔鬼教义教导，视自己的父母为仇敌，因为他们的结合使你们进入肉身的束缚，甚至向你们的神施加了不洁的锁链。关于生育子女为恶的教义，直接反对下一条诫命：\Quote{不可奸淫；}\BibleRef{出 20:14}因为相信这种教义的人，为使妻子不致怀孕，即使在婚姻中也会犯奸淫。他们娶妻，如法律所规定，是为了生育子女；但由于对污染神体的错误恐惧，他们与妻子的结合不具有合法性质，并且他们试图避免生育子女——这原是婚姻的正当目的。正如宗徒很久以前所预言的，你们确实禁止结婚，因为你们试图毁掉婚姻的目的。你们的教义将婚姻变为通奸的关系，将新房变为妓院。这种虚假教义同样导致违背\Quote{不可杀人}的诫命。\BibleRef{出 20:13}因为你们出于恐惧将你们神的肢体囚禁在肉身中，就不给饥饿者面包。出于对幻想杀人的恐惧，你们实际上犯了杀人罪。因为如果你遇到一个因缺乏食物而濒死的乞丐，按天主的律法拒绝给他食物就是杀人；而按照\Person{摩尼}的律法，给食物却是杀人。十诫中没有一条诫命是你们遵守的。如果你们不偷盗，你们就会因容许面包或食物（无论是什么）被一个毫无功德的人吃掉而感到痛苦，而不是把它偷走送到你们选民的胃的实验室里；这样，通过偷盗，你们的神就免于受威胁的囚禁，也免于它已经遭受的囚禁。那么，如果你们在偷窃中被抓住，难道你们不会指着这位神发誓说你们无罪吗？因为当你们对它说：\Quote{我指着你发了假誓，但这是为你好；顾全你的名誉本会害了你}，它会怎样对待你们呢？因此，\Quote{不可作假见证}\BibleRef{出 20:16}这条诫命，不仅会在你们的见证中，也会在你们的誓言中被违背，为的是释放你们神的肢体。\Quote{不可贪恋你近人的妻子}\BibleRef{出 20:17}这条诫命，是你们虚假教义不强迫你们违背的唯一一条。但如果贪恋邻人的妻子是非法的，那煽动别人贪恋又该当何论？回想你们那些美丽的男神和女神，他们故意出现在黑暗的首领男女面前，想要激发欲望，以便满足情欲可以释放这位到处被囚禁、需要如此卑劣帮助的神。最后一条诫命：\Quote{不可贪恋你近人的财物}\BibleRef{出 20:17}，你们根本不可能遵守。难道你们的神不是用许诺在属于别人的区域中创造新世界来迷惑你们，作为你们在幻想征服之后幻想胜利的场所吗？在渴望实现这些狂野幻想的同时，你们又相信这黑暗之地与你们自己的本体近在咫尺，你们当然是在贪恋邻人的财物。难怪你们不喜欢这些包含如此美好诫命而与你们虚假教义相敌对的法版。那三条与爱天主相关的诫命，你们完全置之不理。那七条维系人类社会秩序的诫命，你们只是因为顾及人的意见或畏惧人的法律才遵守；或者良好的习俗使你们厌恶某些罪行；或者你们被\Quote{己所不欲，勿施于人}的自然原则所约束。但无论你们是做你们不愿别人对自己做的事，还是不做你们不愿别人对自己做的事，你们都能看到异端与法律的对抗，无论你们实际是否遵守它。\par

8. 基督的真净配，你们竟敢以石版来讥讽她，她知道文字与圣神，换言之，法律与恩宠之间的区别；不再以文字的旧样事奉天主，而以圣神的新样事奉，她不在法律之下，而在恩宠之下。她不为争辩的精神所蒙蔽，而是谦卑地从宗徒学习，我们不应该处于什么样的法律之下；因为他说：\BibleQuote{它原是因了过犯而赐下的，直到那预许的子孙来到。}\BibleRef{迦 3:19} 他又说：\BibleQuote{法律闯入，为使过犯增多；但罪恶在哪里增多，恩宠在哪里也格外丰富。}\BibleRef{罗 5:20} 不是说法律是罪，尽管没有恩宠它不能赐予生命，反而增加罪责；因为\BibleQuote{没有法律，就没有过犯。}\BibleRef{罗 4:15} 没有圣神的文字，没有恩宠的法律，只能定罪。所以宗徒解释他的意思，以免有人不明白：\BibleQuote{那么，我们可说法律是罪吗？绝对不可！但是，如果不是藉着法律，我就不知道什么是罪。如果不是法律说：\InnerQuote{不可贪恋}，我就不知道什么是贪情。但是罪恶藉着诫命，在我内产生各种贪情；因为没有法律，罪恶是死的。从前没有法律时，我是活着的；但诫命一来，罪恶便活了起来，而我就死了。所以诫命本是引人生命的，但为我却成了死亡；因为罪恶藉着诫命，诱骗了我，并藉着诫命杀了我。所以法律是圣的，诫命也是圣的，是正义的，是善的。那么，那善美的事物为我成了死亡吗？绝对不可！但是罪恶为显示自己是罪恶，竟藉着善美给我带来了死亡。}\BibleRef{罗 7:7} 你们所讥讽的她知道这是什么意思；因为她恳切地求，谦卑地寻，温良地叩门。她看见，当说\Quote{文字叫人死，圣神却叫人活}时，没有人挑剔法律；当说\Quote{知识使人傲慢，爱德却能建立人}时，也没有人挑剔知识。这段话是这样：\BibleQuote{我们知道我们都有知识。知识使人傲慢，爱德却能建立人。}\BibleRef{格前 8:1} 宗徒当然不愿意傲慢；但他有知识，因为知识与爱德结合，不但不使人傲慢，反而坚固人。因此，文字与圣神结合，法律与恩宠结合时，就不再是单靠本身因罪恶增多而致人死亡的那种文字和法律。从这个意义上，法律甚至被称为罪的权势，因为它严格的禁令增加了罪致命的快乐。然而即使如此，法律也不是恶的；而是\Quote{罪恶为显示自己是罪恶，藉着善美带来死亡。}同样，不是恶的事物有时会伤害某些人。摩尼教徒眼睛疼痛时，会排除他们的神——太阳。因此，基督的净配在法律上死了，即向罪死了，这罪因法律的禁令而更加增多；因为脱离恩宠的法律只命令，却不给予能力。她向法律死了，为能嫁给那从死者中复活的那一位，她如此区分，并没有对法律有任何指责，因为指责法律就是亵渎法律的作者。这就是你们的罪行；因为虽然宗徒告诉你们法律是圣的，诫命是圣的、正义的、善的，但你们却不承认它是善者的产物。你们将它的作者当作黑暗的首领之一。在此真理面对你们。这是保禄宗徒的话：\BibleQuote{法律是圣的，诫命也是圣的，是正义的，是善的。}这就是那位为了伟大的象征性用途而设立了你们愚昧地嘲笑之石版的那一位所赐的法律。这由梅瑟所赐的法律，透过耶稣基督成为恩宠和真理；因为圣神与文字结合，使法律的正义开始得以实现，而这正义在未实现时只增加了过犯的罪责。那圣的、正义的、善的法律，就是罪藉以带来死亡的法律，我们也必须向这法律死去，才能嫁给那从死者中复活的那一位。请听宗徒又说什么：\Quote{但罪恶为显示自己是罪恶，竟藉着善美给我带来了死亡，为使罪恶藉着诫命成为格外罪恶的。}聋子和瞎子，你们现在难道还听不见、看不见吗？他说：\Quote{罪恶藉着善美在我内带来了死亡。}法律总是好的：无论它伤害缺乏恩宠的人，还是有益于充满恩宠的人，它本身总是好的；正如太阳总是好的，因为天主的每个受造物都是好的，无论它伤害软弱的目光还是使健康的人赏心悦目。恩宠使心灵适合遵守法律，如同健康使眼睛适合看太阳。正如健康的眼睛不是向着观看太阳的快乐而死，而是向着光线照射眼睛增加黑暗的痛苦效果而死；同样，被圣神的爱治愈的心灵，不是向着法律的正义而死，而是向着法律在没有恩宠时所伴随的罪责和过犯而死。因此说\BibleQuote{法律原是好的，只要人合法地使用它}；并紧接着关于同一法律说\BibleQuote{我们知道法律不是为义人定的。}那以正义本身为乐的人，不需要文字的约束。\par

9. 基督的新娘因圆满救恩的望德而喜乐，并切愿你们从寓言转向真理，获得幸福的皈依。她切愿那仿佛陌生情人的\OriginalTerm{阿多奈乌斯}{Adoneus}之恐惧，不会阻止你们逃脱狡猾之蛇的诱惑。\OriginalTerm{阿多奈}{Adonai}是一个希伯来词语，意义是\Quote{主}，只用于天主。同样，希腊词\OriginalTerm{钦崇}{latria}意指\Quote{服务}，是指对天主的服务；\OriginalTerm{阿们}{Amen}意指\Quote{真实的}，具有特殊的圣洁意义。这只能从希伯来圣经或译本中学到。基督的教会理解并珍爱这些名字，而不理会那些因无知而嘲笑者的邪恶。她尚未理解的事，她相信可以像其他相似之事已向她解释那样得到解释。若有人指控她爱慕\Person{厄玛奴耳}，她嘲笑指控者的无知，并紧握这名字的真理。若有人指控她爱慕\OriginalTerm{默西亚}{Messiah}，她藐视无力的对手，并紧紧跟随她受傅的师傅。她为你们祈祷，愿你们也脱离错谬而得治愈，并被建造在宗徒和先知的根基上。你们无知地加于真道的荒谬怪物，实际上存在于你们幻想故事所描述的世界里，这世界部分由你们的神，部分由黑暗世界构成。这个半野蛮半神圣的世界，比怪物更可怕。看到这样的愚妄，应该使你们谦卑痛悔，并引领你们逃避那引诱你们陷入这些错误的蛇。若你们不相信梅瑟关于蛇的诡诈所说的话，你们可以受保禄的警告，当他说起将教会作为童贞新娘献给基督时，他说：\BibleQuote{我怕你们的心意，也许像蛇用它的狡猾诱惑了\Person{厄娃}一样，受了败坏，失去那在基督内的纯朴和纯洁。}\BibleRef{格后 11:2} 尽管有此警告，你们却受了如此误导、如此迷惑于蛇的致命咒惑，以至于他既已说服其他异端者相信各种谎言，又说服你们相信他就是基督。其他人虽陷入纷繁错谬的迷宫，却仍承认宗徒警告的真理。但你们堕落如此之深，颜面尽失，竟将那位宗徒宣告曾诱惑\Person{厄娃}、并以此来警告基督的童贞新娘提防的对象，当作基督。你们的心被那欺骗者弄瞎了，他使你们沉醉于闪烁丛林的梦境。这些应许除了是梦，还能是什么？有什么理由相信它们是真的？啊，醉了，但非因酒！\par

10. 你们竟胆大包天，指控先知的天主未向祂的仆人犹太人履行祂的应许。然而，你并未提及任何未履行的应许；否则，或可表明这应许已应验，只是你不理解；或它尚未应验，只是你不相信。有什么应许应验在你身上，使你有可能获得从黑暗区域夺来的新世界？若有先知预言称赞\OriginalTerm{摩尼教徒}{Manichaeans}，并说该教派的存在是这预言的应验，那么首先必须证明这些预言不是摩尼为了赢得追随者而伪造的。他并不认为说谎有罪。若他赞美基督，说基督在祂的身体上显示了虚假的伤痕，那么他自然不会顾忌在羊皮卷中显示虚假的预言。诚然，在先知书中关于摩尼教徒的预言不那么清晰，但在宗徒著作中却极其明确。例如，他说：\BibleQuote{圣神明明地说：在末后的时期，有人要背弃信德，听从诱惑人的神和魔鬼的道理，这道理是出于伪善者的谎言，他们的良心已被烙焦；他们禁止嫁娶，戒食某些食物，而这些食物是天主所创造，给信者及认识真理的人以感恩的心领受的。因为天主所造的每样受造物都是好的，没有一样是可摈弃的，只要以感恩的心领受。}\BibleRef{弟前 4:1} 这个预言的应验在摩尼教徒身上，对所有认识他们的人来说，如同白昼一样清楚，而且已在我们所允许的时间内得到充分证明。\par

11. 宗徒警告要提防蛇的诡计的那一位，——你已被这诡计所败坏——好使他能把她如同贞洁的童女献给基督，她唯一的夫君，她承认众先知的天主是真天主，也是她自己的天主。祂的许多应许已经在她身上实现了，所以她满怀信心地期待其余应许的完成。没有人能说这些预言是为了迎合现时而伪造的，因为它们在犹太人的书中可以找到。还有什么比所有民族都要因亚巴郎的后裔而蒙受祝福——正如所应许的那样——更不可能的呢？然而这个应许现在多么明显地在实现！最后的应许体现在下面这段简短预言中：\BibleQuote{居住在你殿宇的，真是有福的；他们要不断讚美你。}当考验过去，最后仇敌——死亡——被消灭后，在持续讚美天主的安息中，将不再有进与出。因此先知在别处说：\BibleQuote{耶路撒冷，你要讚颂上主；熙雍，你要讚颂你的天主！因为祂坚固了你的门闩，祝福了你中间的子女。}门被关上，以致无人能进或出。新郎自己在福音中说，祂不会给愚拙的童女开门，即使她们敲门。这个耶路撒冷，即圣教会、基督的新娘，在\BibleBook{若望}{默示录}中有详尽的描述。而这促使这贞洁童女相信未来福乐应许的是：她现在已拥有了同一些先知所预言关于她的事。因为她被这样描述：\BibleQuote{女儿，请听，请看，请侧耳细听：忘记你的民族和你父的家！因为君王恋慕你的美艳；祂是你的天主。提洛的女儿要前来朝贡，民间的富豪也要向你进贡。君王公主的光荣在内；她的衣服是金线绣的。她身穿锦衣，被引到君王面前；陪伴她的同伴，也要被引到你面前。她们要欢欣鼓舞，被引到君王宫阙。你的儿子将继承你的祖先；你要立他们为全地的王子。你的名字要万世流传；所以万民要永远赞美你。}不幸的蛇诡计的牺牲品，君王女儿的内在荣美甚至不是你能想到的。因为这种心灵的纯洁正是你失去的——当你睁开眼睛去爱慕和朝拜太阳与月亮时。因此，按天主公义的审判，你被隔绝于生命树——即永恒的、内在的智慧；在你看来，除了进入邪恶的眼睛的光——它在你不洁的心中扩展并形成虚幻的形象——以外，没有什么是真理或智慧。这些都是你可憎的淫行。然而真理仍召唤你反省并回头。回到我这里，你将得到洁净和恢复——如果你的羞耻引导你悔改。要听真正真理的这些话，祂既没有以虚假的形状与黑暗种族作战，也没有以虚假的血救赎了你。\par

\SourceRecord{Translated by Richard Stothert.；Nicene and Post-Nicene Fathers, First Series；Vol. 4.；Edited by Philip Schaff.；Buffalo, NY: Christian Literature Publishing Co.,；1887.；<http://www.newadvent.org/fathers/140615.htm>.}

% structure: p0001 source=h1 target=section reason=page-title

\section{驳斥福斯图斯，第十六卷}

\Emph{福斯图斯愿意相信，不仅犹太人的先知，而且所有外邦人的先知都曾论及基督，只要能够证明；但他仍然坚持拒绝他们的迷信。奥古斯丁坚持认为，梅瑟所写的全部都是关于基督的，并且他的著作必须要么全部接受，要么全部拒绝。}\par

1. \Emph{福斯图斯}说：你们问我们为什么不信梅瑟，既然基督说过：\Quote{梅瑟论及了我；你们若信梅瑟，也必信我。}我很高兴，不仅梅瑟，而且所有的先知，无论是犹太人还是外邦人，都曾论及基督。这对我们的信仰不会造成障碍，反而会有帮助，如果我们能从各方收集符合我们天主的见证。你们可以从我们所憎恶的迷信中提取关于基督的预言，但我们仍然憎恶那迷信。我很愿意相信，梅瑟，尽管与基督如此对立，似乎也论及了祂。没有人不乐意从每根荆棘中找出一朵花，从每棵植物中找到食物，从每只昆虫中找到蜂蜜，尽管我们不以昆虫或草为食，也不以荆棘为冠冕。没有人不希望在每个深渊中发现珍珠，在每片土地中发现宝石，在每棵树上发现果实。我们可以从海中取鱼，却不喝海水。我们可以取有用的，抛弃有害的。我们为什么不能从一种宗教中取基督的预言，而同时谴责其礼仪为无用呢？这未必会使我们陷入错误的束缚；因为我们并不因为那些不洁之神明明白白地承认耶稣是天主子而减少对它们的憎恨。如果在梅瑟的著作中发现任何类似的见证，我会接受。但我不会因此顺从于他的律法，在我看来那纯粹是异教。完全没有理由认为我会拒绝从任何神灵接受关于基督的预言。\par

2. 既然你们已经证明基督说过梅瑟曾论及祂，我将非常感激你们能指给我看梅瑟写了什么。我查阅了圣经，正如我们被告知要做的，但没有发现关于基督的预言，要么是没有，要么是我无法理解。在这种困惑中，只有两种可能的结论：要么这节经文是伪造的，要么耶稣是撒谎者。由于敬虔不允许认为天主是撒谎者，我宁愿将虚假归于作者，而非真理的源头。况且，祂自己说过，那些在祂以前来的人都是贼和强盗，这首先适用于梅瑟。并且，当祂论及自己的威严，称自己为世界的光时，犹太人愤怒地反驳说：\Quote{你为自己作证，你的见证不真实。}我发现祂并没有像预想的那样引用梅瑟的预言。相反，祂好像与犹太人毫无关联，也不接受他们祖宗的见证，祂回答说：\Quote{你们的律法上记载说：两个人的见证是真的。我为自己作证，还有派遣我的父也为我作证。}祂指的是从天上传来、众人都听见的声音：\Quote{这是我的爱子，你们要听从他。}我认为，如果基督说过梅瑟曾论及祂，那么犹太人出于巧妙的敌意，定会立刻问祂，祂认为梅瑟写了什么。犹太人的沉默证明耶稣从未说过这样的话。\par

3. 然而，我怀疑这节经文真伪的主要原因是，正如我之前所说，在我查阅梅瑟著作时，没有发现任何关于基督的预言。但现在我遇到了一位您这样智慧超群的读者，希望能学到一些东西；我承诺将心怀感激，只要没有恶意阻止您让我受益于您更高的学识，正如您高尚的责备风格使我期待于您的。我请求您教导我，梅瑟著作中关于我们的天主和主的内容，有哪些是我在阅读时遗漏的。我恳求您不要使用那种无知的论点，即基督肯定梅瑟曾论及祂。因为假设您打交道的不是我（就我而言，有义务相信我所跟随的那位），而是与一个犹太人或外邦人打交道，当您说梅瑟曾论及基督时，他们会要求证据。我们该对他们说什么呢？我们不能引用基督的权威，因为他们不信祂。我们只能指出梅瑟所写的。\par

4. 那么，我们该指出什么呢？是否指您经常引用的那段，其中梅瑟的天主对他说：\BibleQuote{我要从他们弟兄中给他们兴起一位像你一样的先知？}\BibleRef{申 18:15} 但犹太人可以看出这不指向基督，而且我们有充分的理由认为并非如此。基督不是先知，也不像梅瑟：因为梅瑟是人，基督是天主；梅瑟是罪人，基督是无罪的；梅瑟是正常生育的，基督按你们所说由童贞女所生，或按我所信根本未有出生；梅瑟因得罪他的天主而在山上被杀；基督是自愿受苦，且父悦纳了祂。如果我们断言基督是像梅瑟一样的先知，犹太人要么会嘲笑我们无知，要么会说我们不诚实。\par

5. 或者我们来看看你们所喜欢的另一段经文：\BibleQuote{他们必看见自己的生命悬而不定，且不相信自己的生命？}\BibleRef{申 28:66} 你们在其中加了\Quote{悬在树上}一语，原文中并没有。要证明这与基督无关，再容易不过了。\Person{梅瑟}在警告百姓若偏离他的法律时将遭受可怕的灾祸时，特别提到他们会被敌人掳去，昼夜期待死亡，对征服者所容许的生命毫无信心，以致他们的生命因恐惧迫近的危险而悬于不定。这段经文并不适用，我们必须另找。我不能承认\Quote{凡悬在树上的，都是被诅咒的}这句话是指着基督说的；同样，当圣经说必须处死那试图引诱百姓离开他们的天主、或违犯任何诫命的首领或先知时，也不适用于基督。\BibleRef{申 13:5} 我不得不承认基督确实行了这事。但若你们断言这些事是论及基督而写的，那么就可以反问：是什么神灵默感这些预言，使\Person{梅瑟}诅咒基督并下令杀死祂？若\Person{梅瑟}有圣神，这些事就不是论及基督而写的；若这些事是论及基督而写的，他就没有圣神。圣神不会诅咒基督，也不会下令杀死祂。要为\Person{梅瑟}辩护，你们就必须承认这些段落也与基督无关。所以，如果你们没有别的经文可引用，那就没有了。如果没有，基督就不可能说过有；如果基督没有说过，那节经文就是伪造的。\par

6. 下一节经文也令人怀疑：\Quote{如果你们相信梅瑟，也必会相信我；} 因为\Person{梅瑟}的宗教与基督的宗教完全不同，以致犹太人若相信一个，就不能相信另一个。\Person{梅瑟}严格禁止在安息日做工，并给出禁令的理由说，天主在六天内创造了世界和其中的万物，在第七天安息了，这一天便是安息日；因此他祝福或圣化这一天，作为他劳作后的安息之所，并命令违犯安息日者应处死。犹太人顺从\Person{梅瑟}，对此坚持不渝，因此甚至不愿听基督说天主一直在工作，没有指定哪一天停止祂纯洁而不知疲倦的活动，因此祂自己也必须在安息日不停地工作。祂说：\BibleQuote{我父一直工作，我也必须工作。}\BibleRef{若 5:17} 再者，\Person{梅瑟}把\OriginalTerm{割损礼}{circumcision}列为悦乐天主的仪式，命令所有男性都要割去包皮，并宣告这是天主与\Person{亚巴郎}所立盟约的必要记号，凡未受割损的男性必从本族中被铲除，也无分于应许给\Person{亚巴郎}及其后裔的产业。\BibleRef{创 17:9} 犹太人在此事上也十分热心，因而不能相信基督，因为基督轻视这些事，并宣告\BibleQuote{一个人受割损后，就成了加倍的地狱之子。}\BibleRef{玛 23:15} 再者，\Person{梅瑟}对动物食物的区别非常讲究，像美食家一样论及鱼、鸟和四足动物的优劣，命令某些可吃，视为洁净，某些不洁则不可触摸。在不洁之物中，他列举了猪和野兔，无鳞的鱼，以及不分蹄或不反刍的四足动物。犹太人在此事上也谨守\Person{梅瑟}的规条，因此不能相信基督，因为基督教导一切食物都是相同的，虽然祂不允许自己的门徒吃任何动物肉，却赐给平信徒充分的自由，随意吃什么，并教导人不是入口的使人污秽，而是从口里出来的恶事。在这些事和许多其他事上，\Person{耶稣}的教导，众所周知，与\Person{梅瑟}的教导相抵触。\par

7. 不必列举所有的分歧，只要提这一事实就足够了：大多数基督徒教派，尤其众所周知的天主教徒，对\Person{梅瑟}著作中所规定的事并不重视。若这并非起源于错误，而是来自基督及其门徒正确传授的道理，那么你们必定承认\Person{耶稣}的教导与\Person{梅瑟}的教导相反，犹太人正是因为依恋\Person{梅瑟}才不信基督。\Person{耶稣}对犹太人说\Quote{如果你们相信梅瑟，也必会相信我；} 当他们的相信\Person{梅瑟}显然阻止了他们相信\Person{耶稣}——如果他们不再相信\Person{梅瑟}，反倒有可能相信\Person{耶稣}——时，这话怎能不假呢？我再请你们指出\Person{梅瑟}所写关于基督的任何事。\par

8. 此外，\Person{Faustus}说：当你们找不到任何经文可引用时，便用这个软弱而不恰当的论据，即基督徒必须相信基督所说\Person{梅瑟}曾论及祂的话，凡不信此事的就不是基督徒。你们还不如立刻承认你们找不到任何经文。这个论据或许对我有用，因为我敬畏基督，不得不相信祂所说的。但仍有疑问：这是基督自己的宣告，要求绝对相信，还是作者的陈述，需要仔细考察？不信谎言不是冒犯基督，而是冒犯骗子。但无论这个论据对基督徒有何用处，在讨论中面对犹太人或外邦人时，它完全不合用。即使对基督徒，这个论据也有问题。当宗徒\Person{多默}怀疑时，基督并没有摒弃他。基督没有说\Quote{你若是个门徒，就相信；谁不信就不是门徒，} 而是通过让\Person{多默}看祂自己身上伤口的痕迹来治愈他心灵的创伤。那么，你怎么能对我说，我不是基督徒，因为我怀疑——不是怀疑基督，而是怀疑一句归给基督的话的真实性？但你说，祂特别祝福那些没有看见而相信的人。如果你们认为这是指不用判断和理性而相信，那么你们尽管享有这种盲目的福气；我却满足于理性的福气。\par

9. \Person{奥斯定}回答说：你说从梅瑟的著作中取任何关于基督的预言，就像从海中取一条鱼，而把取鱼的水扔掉，这个比喻很巧妙。但既然梅瑟所写的一切都关乎基督，或指涉基督，或是用言语行动预言祂，或是彰显祂的恩宠与荣耀，而你，相信摩尼著作中不真实且不说真话的基督，又不信梅瑟，你连鱼也吃不到。此外，虽然你真心敌视梅瑟，你却虚伪地赞美鱼。你怎么能说吃一条从海里捞出的鱼无害呢？而你的教义却说这种食物如此有害，以至于你宁愿饿死也不食用它？如果所有血肉都是不洁的，如你所说，并且你们可怜的神的生命被禁锢在所有水或植物中，通过你们将其用作食物而得以释放，按照你们自己卑劣的迷信，你必须扔掉你赞美的鱼，喝那水，吃你说无用的蓟草。至于你将天主的仆人比作魔鬼，好像他对基督的预言类似魔鬼的承认，这仆人不拒绝承受主人的羞辱。如果一家之主被称为贝尔则布，何况祂的家人呢！\BibleRef{玛 10:25} 你已经从基督的仇敌那里学到了这种羞辱；你比他们更坏。他们不信耶稣是基督，因此认为祂是骗子。而你唯一相信的教义竟敢使基督成为说谎者。\newline{}\par

10. 你有什么理由说梅瑟的法律是纯粹的异教？是因为它谈到圣殿、祭台和司铎吗？但这些名称在新约中也有。基督说：\BibleQuote{拆毁这座圣殿，三天之内我要把它重建起来。}\BibleRef{若 2:19} 又说：\BibleQuote{当你在祭台前献你的礼物时，}\BibleRef{玛 5:24} 又说：\BibleQuote{去，把你自己给司铎查看，并奉上梅瑟所吩咐的祭献，作为对他们的证据。}\BibleRef{玛 8:4} 这些事所预表的是什么，主自己部分地告诉我们，称祂的身体为圣殿；我们也从宗徒那里学到，他说：\BibleQuote{天主的圣殿是圣的，这圣殿就是你们}；\BibleRef{格前 3:17} 又说：\BibleQuote{所以我以天主的仁慈恳求你们，将你们的身体献上，当作生活的、圣洁的、悦乐天主的祭品}；\BibleRef{罗 12:1} 以及类似的经文。正如同一宗徒在值得反复引用的话中所说：这些事是我们的预像，因为它们不是魔鬼的作为，而是那创造天地的唯一真天主的作为，祂虽不需要这些，却按照时代的要求，使古代的仪式成为未来实体的表记。既然你假装厌恶异教，其实只是为了欺骗那些不熟悉或信心不稳固的基督徒，请你在基督教的书籍中给我们指出任何为崇拜和侍奉太阳月亮而有的权威。你的异端比梅瑟的法律更像异教。因为你敬拜的不是基督，只是你称之为基督的某种东西，一个你自己幻想的虚构；你所侍奉的神明要么是天上的可见的物体，要么是你自己捏造的群魔。如果你不为这些毫无价值的幻象偶像建造神龛，你却使你的心成为它们的殿宇。\newline{}\par

11. 你问我指出梅瑟关于基督写了什么。许多段落已经被指出了。但谁能指出所有呢？此外，当提出任何引文时，你总是执拗地试图赋予这些字句另一种含义；或者，如果证据太强而无法抗拒，你就会说你把那段落当作从咸水中取出的甜鱼，因此你不会同意喝下梅瑟书中的所有咸水。那么，只引用福斯图斯从希伯来法律中选择出来批评的段落就足够了，并表明，当正确理解时，它们适用于基督。因为如果我们的对手嘲笑和谴责的事情被证明为他自己被基督真理所谴责，那么要么仅仅引用其他段落，要么仔细检查，就足以表明它们与基督信仰一致。好吧，那么，你这充满诡诈的人，当主在福音中说：\BibleQuote{如果你们相信梅瑟，你们也必相信我，因为他所写的是关于我的，}\BibleRef{若 5:46} 没有必要陷入你假装感到的极大困惑，也不必在两个选择中抉择，要么宣称这节经文是伪造的，要么称耶稣为说谎者。这节经文既真实，其话语也真实。福斯图斯说：我宁愿将虚假归于作者，而不归于真理的源泉。你对基督作为真理的源泉能有什么信心呢？你的教义说祂的肉体、祂的死亡、祂的创伤和它们的印记都是虚构的？如果你敢于对为祂作见证的人——他们的见证经同时代的人证实而传给我们——归于虚假，那么你有什么权威说基督是真理的源泉？你没有见过基督，祂也没有像对宗徒那样与你交谈，也没有从天上像对扫禄那样召唤你。我们对基督能有什么认识或信仰，除了圣经的权威？如果那已经在万民中广泛传播、自基督之名首次被宣讲就为所有教会极其尊崇的福音有虚假，我们哪里能找到可信的关于基督的记录？如果福音在普遍的同意下仍被质疑，那么一个人若不愿相信，他就可以宣称任何著作是伪造的。\newline{}\par

12. 你们继续引用基督的话，说凡在祂以先来的，都是贼和强盗。你们怎么知道这些话是基督说的，除非从福音中得知？你们相信这些话，好像亲耳从主口中听到一样。但若有人宣称这节经文是伪造的，否认基督说过这话，你们就必须努力维护福音的权威。不幸的人！你们所拒绝相信的，与你们所引用为主说的话，写在同一个地方。我们两者都信，因为我们相信同时包含这两者的神圣叙述。我们既相信梅瑟写过基督，也相信凡在基督以先来的都是贼和强盗。祂说他们\Quote{来}，意思是指他们不是被派遣的。那些被派遣的，如梅瑟和圣先知们，并不是在祂以先来，而是与祂同来。他们并不骄傲地想要走在祂前面，而是谦卑地传递祂借他们所说的话。按你们对主的话的解释，显然你们那里不可能有先知。所以你们为自己造了一个基督，这基督要预言未来的基督。如果你们有自己的先知，他们当然没有权威，因为不被其他人承认；但如果有你们敢于引用的先知，他们预言基督将以虚幻的身体来临，受虚幻的死亡，并向疑惑的门徒显示虚幻的伤痕，且不说这类预言的邪恶性质，以及那些在基督里称赞谎言之人的明显虚假，按你们自己的解释，那些先知必定是贼和强盗，因为他们除非在基督以先来，否则不可能以任何方式谈论基督的来临。如果我们将\Quote{在基督以先来的}理解为那些不愿与祂同来（即与天主圣言同来）的人，而是未经天主派遣，将自己的谎言带给人类，那么你们自己，虽然生在基督死而复活之后的这个世界，也是贼和强盗。因为你们不等祂的光照来宣讲祂的真理，就在祂以先来，宣扬自己的欺骗。\par

13. 在经文中，我们读到犹太人向基督说：\Quote{你为自己作证，你的证据不真实}，你们没有看到基督以梅瑟写过祂来回答，只是因为你们没有虔敬的眼光。基督的回答是这样：\BibleQuote{\Emph{你们的律法}上记载着：两个人的见证是真的。我为自己作证，并且那派遣我来的圣父也为我作证。}\BibleRef{若 8:17} 如果正确理解，这意味着什么呢？难道不就是律法所要求的证人数量是藉着预言之神所指定和祝圣的，为预表圣父和圣子的启示，而他们的神就是不可分的天主圣三的圣神？经上写着：\BibleQuote{凭两个或三个证人的口，一切事都可以成立。}\BibleRef{申 19:15} 事实上，一个证人通常说真话，而许多人说谎。世界在归依基督教时，相信了一位宗徒宣讲的福音，胜过逼迫他的错误群众。要求这个数量的证人有特别理由，而在祂的回答中，主暗示梅瑟预言了祂。你们批评祂说\Quote{\Emph{你们的律法}}而非\Quote{天主的律法}？但众所周知，这是圣经中的常见表达。\Quote{你们的律法}意思是赐给你们的律法。同样，宗徒提到他的福音，但同时声明他不是从人领受的，而是藉着耶稣基督的启示。你们也可以说基督否认天主是祂的父亲，当祂使用\Quote{你们的父}而不是\Quote{我们的父}时。此外，你们应该拒绝相信你们所提到的从天上来的声音：\Quote{这是我的爱子，你们要听从他}，因为你们没有听见。但如果你们因为在圣经中找到了而相信它，你们也会在那里找到你们所否认的，即梅瑟写过基督，还有其他许多你们不承认真实的事。你们难道没有看到，你们自己的有害论证可以用来证明这声音从未来自天上？为了你们自己的毁灭，以及人类福祉的损害，你们试图削弱福音的权威，论证说基督说梅瑟写过祂不可能是真的；因为如果祂说过这话，犹太人巧妙的敌意会立刻质问他们认为梅瑟写了关于祂的什么。同样，也可以不敬地论证说，如果那声音真的来自天上，所有听见的犹太人都会相信。为什么你们如此不合理，不考虑正如犹太人听了天上的声音后仍可保持硬心不信，同样，当基督说梅瑟写过祂时，他们也可以不问梅瑟写了什么，因为在他们巧妙的敌意中，他们害怕被证明是错的？\par

14. 除此以外，这个论调是对福音的邪恶攻击，\Person{福斯图斯}本人也意识到它的软弱，因此他更加强调他所谓的首要难题——在他对\Person{梅瑟}著作的一切搜索中，他没有找到任何关于\Person{基督}的预言。显而易见的答复是：他不理解。若有人问为何他不理解，答案是：他怀着敌意、不信的心阅读；他并非为求知而探询，却自以为知，其实无知。这种虚荣的狂妄要么蒙蔽了他理智的眼睛，使他看不到任何东西，要么扭曲他的视线，导致他的赞许或反对的言论都错失目标。他说：\Quote{我要求指点，告诉我\Person{梅瑟}著作中关于我们的天主和主的一切，这些是我阅读时忽略了的。}我立即回答：这一切都从他眼前溜走了，因为全部都是关于\Person{基督}而写的。由于我们不能通篇详述，我将借着天主的助佑，尽我先前所承诺的程度，满足你的要求，表明你所特别批评的段落正是指向\Person{基督}。你告诉我不要使用无知的论据，即\Person{基督}肯定\Person{梅瑟}曾为祂写作。但我若使用这个论据，并非因为我无知，而是因为我是一个信徒。我承认这个论据无法说服外邦人或犹太人。但是，尽管你百般回避，你不得不承认它对你有利，因为你夸耀自己拥有一种基督教。你说：\Quote{假设你并不是与我对峙，而是与犹太人或外邦人对峙，在我的情形中，有义务相信我所宣称追随的那一位。}这等于是说，至少对于目前与我打交道的人来说，你已经满意地认为\Person{梅瑟}写了关于\Person{基督}的事；因为你还不够大胆到完全抛弃福音中记载\Person{基督}自己宣言的可靠权威。即使你间接攻击这个权威，你感到你正在攻击自己的立场。你意识到，如果你拒绝相信那如此广为人知并被接受的福音，那么你在试图用任何可靠的\Person{基督}言行记录来取代它时必然完全失败。你害怕失去基督徒的名号会导致你的荒谬被普遍嘲笑和谴责。因此你试图恢复自己，说你对基督教的宣认迫使你相信福音的这些话语。所以，至少你——这正是我们现在所关心的——被这针对你错误的致命一击所捕获和击杀。你被迫承认\Person{梅瑟}写了关于\Person{基督}的事，因为福音——你的宣认迫使你相信它——陈述了\Person{基督}如此说。至于与外邦人或犹太人的讨论，我已经尽可能展示了我认为应该如何进行。\par

15. 我仍然认为，在你挑选出来反驳的段落中，有一个对基督的指涉，即天主对梅瑟说：\BibleQuote{我要从他们兄弟中给他们兴起一位像你一样的先知。} \BibleRef{申 18:15} 你用一系列华而不实的对比试图装饰你乏味的论述，但这丝毫影响不了我对这真理的相信。你试图通过比较基督与梅瑟来证明他们不相似，因此\Quote{我要兴起一位像你一样的先知}这句话不能理解为指基督。你列举了许多你发现不同之处的细节：一位是人，另一位是天主；一位是罪人，另一位无罪；一位是普通生育所生，另一位按我们信是童贞女所生，而按你所信甚至不是童贞女所生；一位招致天主愤怒并在山上被处死，另一位自愿受苦，且始终获得祂圣父的赞许。但事物虽非在每方面都相似，仍可说相似。除了同性质事物之间的相似，如同两个人之间，或父母与子女之间，或一般人与人之间，或任何动物种类之间，或树木中，一棵橄榄树与另一棵橄榄树之间，或一棵月桂树与另一棵月桂树之间，不同性质的事物之间也常有相似，如一株野生橄榄与一株栽培橄榄之间，或小麦与大麦之间。这些事物在某种程度上是相关的。但天主子——万物藉祂而受造——与一头野兽或一块石头之间有最大的距离。然而在福音中我们读到：\BibleQuote{请看，天主的羔羊，} \BibleRef{若 1:29} 在宗徒那里：\BibleQuote{那磐石就是基督。} \BibleRef{格前 10:4} 除非假设某种相似，这些说法是不可能的。那么，基督屈尊成为像梅瑟一样，又有什么奇怪呢？祂曾被造成像一只羔羊，那是天主藉梅瑟命令祂的子民作为基督的预像而吃的，并命令用它的血作为保护的手段，且称它为逾越节，人人都必须承认这在基督身上应验了。我承认，圣经显示了不同之处；圣经也——我要求你承认——显示了相似之处。两种要点都有，且两者都可以被证明。基督不像人，因为祂是天主；关于祂写道：\BibleQuote{祂是在万有之上，永受赞美的天主。} \BibleRef{罗 9:5} 基督也像人，因为祂是人；同样关于祂写道：\BibleQuote{在人与天主之间的中保，就是基督耶稣，祂也是人。} \BibleRef{弟前 2:5} 基督不像一个罪人，因为祂永远是圣的；祂又像一个罪人，因为\BibleQuote{天主派遣了自己的儿子带着罪恶肉身的形状，为使因着罪恶在肉身中定罪。} \BibleRef{罗 8:3} 基督不像一个普通生育所生的人，因为祂由童贞女所生；然而祂又相似，因为祂也是由女人所生，对那女人曾说：\BibleQuote{那要诞生的圣者，将称为天主子。} \BibleRef{路 1:35} 基督不像一个因自己的罪而死的人，因为祂是无罪而死，且出于自己的自由意志；而祂又相似，因为祂也经历了真实的肉身死亡。\par

16. 你不应该说毁谤梅瑟的话，说他是个罪人，并且因为他的天主向他发怒，就把他杀死在山上。因为梅瑟能因主——他的救主而夸耀，这救主也是那说\BibleQuote{基督耶稣来到世界，是为了拯救罪人，其中我是罪魁}\BibleRef{弟前 1:15}之人的救主。事实上，梅瑟被天主的声音指责，因为当他奉命从磐石中取水时，他的信德显出了软弱的迹象。\BibleRef{户 9:10}在这事上他可能像伯多禄那样犯了罪，伯多禄因信德的软弱，在波浪中害怕起来。\BibleRef{玛 14:30} 但是我们不能因此认为，那位按福音所载，曾与主及圣厄里亚一同在显圣容的山上被算为堪当与主同在的人，竟被隔绝于圣人们的永恒共融之外。神圣历史显示，即使在他犯罪之后，他仍深得天主欢心。既然你可能会问，为什么天主说这罪应当受死亡的惩罚，而我已答应在你所挑剔的那些段落中指出基督的预像，那么我将依靠主的助佑，努力表明你所反对的关于梅瑟之死的事，若正确理解，乃是基督的预像。\par

17. 我们常在天主经书的象征性段落中发现，同一个人在不同场合以不同角色出现。因此，在这场合，梅瑟代表并预像了置于法律之下的犹太人民。就如梅瑟用杖击打磐石时怀疑天主的能力，同样，领受梅瑟所赐律法的百姓，当他们把基督钉在十字架上时，也不相信祂是天主的能力。正如被击打的磐石为口渴的人流出水来，同样，生命因主的苦难来临到信徒。宗徒在这点上的见证是明确而决定性的，他说：\BibleQuote{这磐石就是基督。}\BibleRef{格前 10:4} 在天主命令梅瑟的肉体死亡应发生在山上这事上，我们看到神圣的安排：对基督神性的肉性怀疑应在基督受高举时死去。如同磐石是基督，山也是基督。磐石是祂谦卑的勇毅；山是祂受高举的高峰。因为正如宗徒说\BibleQuote{这磐石就是基督}，同样基督自己说：\BibleQuote{建在山上的城是不能隐藏的。}\BibleRef{玛 5:14}显示祂是山，而信徒是建在祂圣名光荣上的城。肉性的心灵活着时，正如被击打的磐石，基督在十字架上的谦卑被藐视。因为被钉的基督对犹太人来说是绊脚石，对希腊人来说是愚妄。而肉性的心灵死去时，正如山顶，基督在祂的高举中被看见。\BibleQuote{因为对被召的人，无论是犹太人或是希腊人，基督是天主的能力和天主的智慧。}\BibleRef{格前 1:23} 因此梅瑟上了山，为了在肉体的死亡中被活灵所接纳。如果福斯图斯上了山，他就不会从死去的心灵发出肉性的反对。正是肉性的心灵使伯多禄畏惧击打磐石，当主预言祂的苦难时，他说：\BibleQuote{主，愿这事远离祢；求祢怜悯自己。}这罪也受到了严厉的斥责，主回答说：\BibleQuote{撒殚，退到我后面去！你是我的绊脚石，因为你所体会的不是天主的事，而是人的事。}\BibleRef{玛 16:22} 这肉性的不信除了在基督的光荣中，如同在山顶上，又在哪里死去呢？如果当伯多禄胆小地否认基督时它活着，那么当他无畏地宣讲基督时它死了。它在扫禄身上活着，当他因厌恶十字架的绊脚石而摧毁基督信仰时，而除了在这山上，它又在哪里死了，当保禄能够说：\BibleQuote{我如今活着，不再是我，而是基督在我内活着}\BibleRef{迦 2:20}？\par

18. 你异端的愚昧还有什么理由认为在\Quote{我要从他们兄弟中给他们兴起一位像你一样的先知}这句话中没有基督的预言？你指出基督不像梅瑟，这不是理由；因为我们可以表明在其他方面祂是相似的。你怎能反对基督被称为先知，既然祂屈尊成为人，并实际上预言了许多未来事件？先知是什么，不过是一个预言超越人类预见之事的人。基督自己论自己说：\BibleQuote{先知除了在自己的本乡外，是没有不受尊重的。}\BibleRef{玛 13:57} 但我要转身向你，既然你已经承认你信奉基督宗教的义务是相信福音，我向犹太人说话，他享有不受基督轭下束缚的可怜特权，因此认为可以说：你们的基督说了谎言；梅瑟没有写过关于祂的任何事。\par

19. 让犹太人说说，\Person{天主}对\Person{梅瑟}的应许中所指的是哪一位先知：\BibleQuote{我要从他们弟兄中，给他们兴起一位像你一样的先知。} \Person{梅瑟}之后出现了许多先知，但此处特指某一位。犹太人或许自然会想到\Person{梅瑟}的继承者，他带领\Person{梅瑟}带出埃及的百姓进入了预许之地。既然记载了谁在统治和带领百姓上接续了\Person{梅瑟}，他若心中想着这位继承者，或许会嘲笑我问那应许的话指的是哪一位先知。当他嘲笑我的无知时——正如\Person{福斯托斯}设想的那样——我将继续追问，并希望我那嘲笑我的对手认真回答我：为什么\Person{梅瑟}改变了这位继承者的名字？这位继承者比他自己更优先被选为带领百姓进入预许之地的领袖。这是为了说明：由\Person{梅瑟}所给予的法律不能拯救人，只能指证罪人，不能领我们进入天堂，唯有藉着\Person{耶稣基督}的恩宠和真理才能。这位继承者名叫\Person{欧瑟亚}，\Person{梅瑟}给他起名叫\Person{若苏厄}。那么，当他从\Place{帕兰}山谷派他进入那他要带领百姓进入的土地时，为什么给他这个名字？真正的\Person{耶稣}说：\BibleQuote{我去为你们预备地方以后，我必再来接你们到我那里去。}\BibleRef{若 14:3} 我要问那犹太人：当先知说\Quote{\Person{天主}将来自\Place{阿非利加}，圣者来自\Place{帕兰}}时，难道不是揭示了这些事情预言的意义吗？这不正意味着圣洁的\Person{天主}将带着那从\Place{阿非利加}经由\Place{帕兰}而来者的名字来临，也就是带着\Person{若苏厄}的名字来临吗？再者，当\Person{天主}应许为\Person{梅瑟}提供这位继承者时，是圣言亲自说话，称他为天使——在\WorkTitle{圣经}中，这名字通常用于传达信息者。这话是：\BibleQuote{看，我派遣我的使者在你面前，在路上保护你，并领你到我所誓许给你们的地方。你们应谨慎，听从他，不可背叛他，因为他不会不正当地夺取你们任何东西，因为我的名在他内。}\BibleRef{出 23:20} 深思这些话吧。让那犹太人——更不必说摩尼教徒了——在\WorkTitle{圣经}中找出还有其他哪一位天使适用这些话，就只有这位将带领百姓进入预许之地的领袖。然后让他查明是谁接续了\Person{梅瑟}，并带领百姓进入。他将发现那就是\Person{若苏厄}，而且起初这并不是他的名字，而是在他改名之后。由此可知，说\BibleQuote{我的名在他内}的那位，就是真正的\Person{耶稣}，那位带领祂的子民进入永生的基业的领袖，按照\WorkTitle{新约}——\WorkTitle{旧约}是其预象。没有任何事件或行动比这更具有明确的预言特性，因为连名字本身就是一个预言。\par

20. 由此可见，如果这犹太人愿意做一个内心上的、精神上的犹太人，而非仅在文字上，如果他要被看作一个真正的以色列人，心中毫无诡诈，他就会在这死了的\Person{若苏厄}身上——他带领百姓进入了必死的土地——认出那真正永活的\Person{耶稣}的预象，后者是他可以跟随进入生命之地的。如此，他不再以敌对的精神抗拒这样明显的预言，而是因着\WorkTitle{旧约}中那位\Person{若苏厄}的预指所感，愿意温良地听从那\Person{若苏厄}所预表的那一位，祂的名是\Person{若苏厄}所背负的，而祂带领我们进入真正的预许之地；因为祂说：\BibleQuote{温良的人是有福的，因为他们要承受土地。}\BibleRef{玛 5:4} 外邦人也是如此，如果他的心不太坚硬，如果他是那些\Person{天主}从其中给\Person{亚巴郎}兴起子女的石头之一，就必须承认这是奇妙的：在\Person{耶稣}所出自的民族的古老书籍中，竟然记载了如此明显的预言，甚至包括祂的名字；同时必须指出，那预言的并非任何名叫\Person{若苏厄}的人，而是一位神圣的人物，因为\Person{天主}曾说祂的名在那位被指派统治百姓、带领他们进入王国的人里面，那人因改名而被称为\Person{若苏厄}。当祂带着这个新名字被派遣时，祂带来了伟大而神圣的信息，因此被称为天使，正如任何初学希腊文的人所知，天使意为使者。因此，任何外邦人，若不是执拗顽固，都不会因为这些书属于希伯来人、他自己不受希伯来人法律的约束而轻视这些书；反而会重视这些书，无论它们属于谁，因为他在其中发现了如此古老的预言，以及他现在看到正在发生的事。他不会因\Person{耶稣基督}在\WorkTitle{希伯来圣经}中被预言而轻视祂，反而会认为，那位被视为配得在预言中被描述的对象，无论作者是谁，在祂降世之前如此多的世代——有时以明确的宣告，有时以象征性的行为和话语作为预象——必定值得以深深的钦佩和敬畏来看待，并以坚定不移的信赖来跟随。因此，基督徒历史的事实将证明预言的真理，而预言也将证明\Person{基督}的宣称。如果事实上全世界的人过去和现在并非因阅读这些书而信\Person{基督}，那就称之为幻想吧。\par

21. 鉴于来自万民的众多人群已成为这些书的热心信徒，若告诉我们不可能说服外邦人从犹太人的书学习基督信仰，那是荒谬可笑的。事实上，如此重要的见证出自敌人之口，正是我们信德的一大确证。信主的外邦人不能以为这些关于基督的见证是近期的伪造，因为他们发现这些见证存在于那些钉死基督的人所尊为神圣许多世代的书卷中，并且仍被那些每日亵渎基督的人以最高的崇敬对待。若基督的预言是基督的宣讲者所造，我们可能会怀疑其真实性。但现在，宣讲者却解释亵渎者的经文。至高天主就这样为了圣人的益处而安排不敬虔者的盲目，在祂公义的治理中从恶中引出善，使那些因自身选择而生活邪恶的人，在祂公正的审判下，成为祂旨意的工具。因此，为免那些要向世界宣讲基督的人被认为伪造了论及基督的预言——即基督将诞生、行奇迹、无辜受苦、死亡、复活、升天、在万民中传布永生福音——犹太人的不信已为我们带来显著的利益。如此，那些内心不接受这些真理为己益处的人，手中却为我们利益而携带着包含这些真理的著作。犹太人的不信反而增加了这些书的权威，而非减少，因为这盲目本身已被预言。他们因不理解真理而为真理作证。因不理解那些预言他们不会理解的书，他们证明这些书是真实的。\par

22. 在\BibleQuote{你将看见你的生命悬挂着，却不相信你的生命}这段经文中（\BibleRef{申 28:16}），\Person{福斯图斯}被文字的歧义所欺骗。文字可以有不同解释；但\Person{福斯图斯}并没有说它们不能理解为指基督，任何不否认基督是生命，或不否认犹太人看见祂挂在十字架上，或不否认他们不信祂的人，也不能这样说。因为基督亲自说：\BibleQuote{我是生命}（\BibleRef{若 14:6}），而且无疑祂曾被不信的犹太人看见悬挂着，我看不出有什么理由怀疑这是指着基督写的；因为正如基督所说，\Person{梅瑟}曾论及祂。既然我们已经驳斥了\Person{福斯图斯}的论点，他试图藉此说明\Quote{我要从他们兄弟中兴起一位像你一样的先知}这句话不适用于基督，因为基督与\Person{梅瑟}不像，所以我们无需坚持这个别的预言。因为，在第一种情况下，他的论点是基督与\Person{梅瑟}不同，同样，在这里他应当论证基督不是生命，或者祂没有被不信的犹太人看见悬挂着。但他没有这样说，而且现在也没有人敢这样说，那么接受这也是由祂的仆人所说的、关乎我们的主和救主\Person{耶稣基督}的预言，就不应有困难。\Person{福斯图斯}说，这些话出现在一个诅咒的章节中。但为什么因为它出现在预言中间就成为较少的预言呢？或者为什么它不能是关于基督的预言，即使上下文似乎不涉及基督？事实上，在犹太人因罪恶的骄傲而招致的所有诅咒中，没有什么比这更糟糕的了：他们看见自己的生命——即天主之子——悬挂着，却不相信他们的生命。因为预言的诅咒并非敌意的咒骂，而是即将来临审判的宣告。敌意的咒骂是被禁止的，因为经上说：\BibleQuote{祝福，不可咒骂}（\BibleRef{罗 12:14}）。但先知性的宣告在圣人的著作中常见，例如宗徒\Person{保禄}说：\BibleQuote{铜匠\Person{亚历山大}加害了我许多恶事；主必按照他的行为报应他}（\BibleRef{弟后 4:14}）。所以或许会认为宗徒是被愤怒的情绪驱使说出这句咒骂：\BibleQuote{我恨不得那些扰乱你们的人甚至把自己阉割了}（\BibleRef{迦 5:12}）。但如果我们记得作者是谁，我们可以在这种模棱两可的表达中看到一种巧妙的祝福方式。因为有些阉人是为了天国的缘故自阉的（\BibleRef{玛 19:12}）。如果\Person{福斯图斯}对基督徒的食粮有虔诚的胃口，他会在\Person{梅瑟}的话中发现类似的歧义。犹太人可能将\BibleQuote{你将看见你的生命悬挂着，却不相信你的生命}这句话理解为，他们看见自己的生命因敌人的威胁和阴谋而处于危险中，并预期无法存活。但福音之子，听到基督说\Quote{他论及了我}，便在预言的歧义中区分什么是投给猪的，什么是给人说的。他心中立刻浮现基督作为人的生命而悬挂，以及犹太人正因为看见祂悬挂而不信祂的想法。至于反对意见，即\BibleQuote{你将看见你的生命悬挂着，却不相信你的生命}这些话是一个包含不适用于基督的诅咒的段落中唯一指向基督的话，有些人可能会承认这是真的。因为这个预言很可能出现在先知对不敬虔百姓的诅咒中，因为这些诅咒有不同种类。但是，我和那些与我一起更仔细考虑主在福音中所说的话的人——不是\Quote{他也论及了我}，仿佛承认\Person{梅瑟}写了其他不涉及基督的事，而是\Quote{他论及了我}，教导我们在研读圣经时应视其为完全旨在阐明基督的恩宠——认为此段其余部分也指向基督。但在这里解释会花费太多时间。\par

23. 法乌斯图所引的话与其他诅咒一同出现，远不能证明它不指基督；相反，这些诅咒只有作为基督光荣的预言才能正确理解，人的幸福就在于这光荣。至于这些诅咒所说的事，对于这句引文更是如此。即使可以说梅瑟的话有超出他本意的含义，我也宁愿认为他不自觉地预言，而不承认\Quote{你必将看见你的生命悬而不定，且必不信你的生命}这话不适用于基督。同样，盖法的话也有他本意之外的含义；当他怀着敌意说：\Quote{一个人替百姓死，免得整个民族灭亡，是有益的}时，福音作者补充说，他说这话并非出于自己，而是因为他身为大司祭而预言了。\BibleRef{若 11:49} 但梅瑟不同于盖法；因此当梅瑟对希伯来人说\Quote{你必将看见你的生命悬而不定，且必不信你的生命}时，他不仅确实论及基督——即使他说话时可能不理解自己所说的含义——而且他知道自己在论及基督。因为他是预言奥迹（即使人认识基督之名的\OriginalTerm{司铎傅油}{priestly unction}）最忠实的管家；就连邪恶的盖法，也在这奥迹中不自知地预言了。\OriginalTerm{先知傅油}{prophetic unction}使他能预言，但邪恶的生活阻止他明白。那么，谁能说梅瑟中没有基督的预言呢？梅瑟开始了那使我们得以认识基督之名的傅油，甚至基督的迫害者盖法也靠这傅油而不自知地预言了基督。\par

24. 关于挂在树上受诅咒的话，我们已经说了足够多的话。已经充分表明，杀害任何试图使以色列子民离弃他们的天主或违犯任何诫命的先知或首领的命令，并不针对基督。我们越思考我主耶稣基督的言行，这一点就越清楚；因为基督从未试图使任何以色列人离弃他们的天主。梅瑟教导人民要热爱和事奉的天主，就是亚巴郎的天主、依撒格的天主、雅各伯的天主，我主耶稣基督也以此名称呼祂，并用这名称驳斥否认复活的撒杜塞人。祂说：\BibleQuote{关于死人复活，你们没有读过天主在荆棘中对梅瑟所说的话吗？说：\InnerQuote{我是亚巴郎的天主，依撒格的天主，雅各伯的天主。} 天主不是死人的天主，而是活人的天主；因为众人为祂而生活。} 用基督回答撒杜塞人的同样话，我们也可以回答摩尼教徒，因为他们也否认复活，尽管方式不同。此外，当基督赞扬百夫长的信德时说：\BibleQuote{我实在告诉你们：在以色列我从未遇见过这样大的信德。我又告诉你们：将有许多人从东方和西方来，与亚巴郎、依撒格、雅各伯一同坐席于天国；但本国的子民将被驱逐到外面的黑暗中。} \BibleRef{玛 8:10} 那么，如果正如法乌斯图必须承认的，梅瑟所讲的天主是亚巴郎、依撒格、雅各伯的天主，基督也谈到同一天主，正如这些经文所证明的，那么基督并没有试图使人民离弃他们的天主。相反，祂警告他们将进入外面的黑暗，因为祂看到他们离弃了他们的天主，而祂说从四方被召叫的外邦人将与亚巴郎、依撒格、雅各伯一同坐席于天主的国，这意味着他们将信仰亚巴郎、依撒格、雅各伯的天主。同样，宗徒也说：\BibleQuote{圣经预先看见天主要使外邦人因信德成义，就向亚巴郎预报福音说：\InnerQuote{万民都要因你获得祝福。}} \BibleRef{迦 3:8} 这意味着因亚巴郎的后裔而受祝福的人将效法亚巴郎的信德。因此，基督并没有试图使以色列人离弃他们的天主，反而谴责他们离弃了天主。认为基督违犯了梅瑟所颁布的诫命，这并非新说法，因为犹太人曾这样想；但这是个错误，因为犹太人错了。让法乌斯图提出他以为主违犯的诫命，我们会指出他的错误，就像我们之前被要求时已经做过的那样。与此同时，这样说就够了：如果主违犯了任何诫命，祂就不能责备犹太人那样做。因为当犹太人责怪祂的门徒不洗手吃饭时——他们违反的不是天主的诫命，而是长老的传统——基督说：\BibleQuote{你们为什么因你们的传统而违反天主的诫命呢？因为天主曾说：\InnerQuote{应孝敬你的父亲和你的母亲，以及咒骂父母者应处死。但你们却说：谁若对父亲或母亲说：我所能供养你的，已成了献仪，他就不必再孝敬父亲了。这样，你们就因你们的传统，废弃了天主的话。}} \BibleRef{玛 15:3} 从这几方面可以学到：基督并没有使犹太人离弃他们的天主；祂不仅自己没有违反天主的诫命，反而责备那些违反的人；而且正是天主亲自藉着梅瑟赐下了这些诫命。\par

25. 为履行我们的承诺，即证明福斯图斯从梅瑟著作中挑选出来加以批评的那些段落中有关基督的指涉（因为我们无法在此指出我们相信存在于梅瑟所有著作中有关基督的指涉），我们有责任表明，梅瑟这条诫命，即凡试图使百姓离开他们的天主或违背任何诫命的先知或首领都应被处死，指的是保存在基督的教会中所教导的信仰。梅瑟无疑以预言的精神，并从他亲自从天主所听到的，知道会有许多异端者兴起，教导反对基督道理的种种错误，并宣讲另一位基督而非真正的基督。因为真正的基督就是那位由梅瑟本人和该国其他圣人所发出的预言中所预告的基督。因此，梅瑟命令，凡试图教导另一位基督的人都应被处死。服从这一命令，天主教会的呼声，就如用新旧约的属灵双刃剑，处死所有试图使我们离开我们的天主或违背任何诫命的人。其中为首的就是摩尼本人；因为法律和先知书的真理揭露了他的错误，他试图使我们离开我们的天主，即亚巴郎、依撒格和雅各伯的天主，就是基督所承认的那位，并试图违背法律的诫命，而这些诫命即使只是预像，我们也视之为基督的预言。\par

26. 福斯图斯使用了一个要么非常欺骗性要么非常愚蠢的论证。既然福斯图斯并不愚蠢，他很可能有意使用这个论证，旨在误导粗心的读者。他说：如果这些事不是关于基督写的，而且你们不能指出其他任何事，那么结论就是根本没有。这个命题是正确的；但仍有待证明，既这些事不是关于基督写的，也无法指出其他任何事。福斯图斯并未证明这一点；因为我们既指出了如何将这些事理解为关于基督的，也指出了还有许多其他事只有应用于基督才有意义。因此，并不像福斯图斯所说的那样，梅瑟没有写过关于基督的任何事。让我们重复福斯图斯的论证：如果这些事不是关于基督写的，而且你们不能指出其他任何事，那么结论就是根本没有。完全正确。但由于我们已经表明这些事以及许多其他事是关于基督写的，或是与基督有关的，真正的结论是福斯图斯的论证毫无价值。在福斯图斯引用的段落中，他试图（尽管不成功）表明它们不是关于基督写的。但为了得出根本没有的结论，他首先应该证明无法指出其他任何事。相反，他假设他书的读者是盲目的，或听众是聋的，以致于遗漏会被忽视，于是继续说：如果没有，那么基督就不可能断言有任何。如果基督没有这样断言，那么这节经文就是伪造的。这里有一个人如此看重自己所说的话，以至于他根本不考虑另一个人说相反话的可能性。你的才智在哪里？这就是你为一件坏事所能说的一切吗？但如果事情的恶劣使你说了蠢话，那坏事是你自己的选择。要证明你的前提为假，我们只需指出一些关于基督写的其他事。如果存在一些，那么\Quote{没有}就不成立。如果存在一些，基督可能断言过有。如果基督可能这样断言，那么这节福音经文就不是伪造的。回到福斯图斯的命题：如果你们不能指出其他任何事，那么结论就是根本没有。这需要证明我们不能指出其他任何事。我们只需参考我们之前所展示的，足以证明福音经文中基督所说的话的真实性：\BibleQuote{若是你们相信梅瑟，也必相信我，因为他写的是关于我。} 即使因心智迟钝在梅瑟著作中找不到任何关于基督所写的内容，然而福音权威的证据如此有力，以至于我们有责任相信梅瑟所写的不仅有些事，而且一切事都指向基督；因为祂不是说\Quote{他也写了一些关于我}，而是说\Quote{他写的是关于我}。事实是，即使对这节经文的真实性有疑问（但愿没有！），这个疑问也会因我们在梅瑟中找到的众多关于基督的见证而被消除；另一方面，即使我们找不到任何见证，我们仍然应当相信这些见证是存在的，因为对福音中任何一节的经文都不能允许有任何怀疑。\par

27. 至于你的论证，即梅瑟的道理与基督的道理不同，因此若他们相信梅瑟，就不太可能也相信基督；反而会得出相信其中一个必然意味着反对另一个——如果你稍微用你心灵的眼睛看一下全世界的人，当他们不被争辩的精神蒙蔽时，无论有学问的、没有学问的、希腊人、野蛮人、聪明的、愚拙的（宗徒自称是欠他们债的\BibleRef{罗 1:14}），他们都同时相信基督和梅瑟。如果犹太人相信基督和梅瑟是不可能的事，那么全世界都这样做更是不可思议。但既然我们看到所有民族都相信两者，并以共同而坚固的信仰坚持一方的预言与另一方的福音相符合，那么当基督对他们说\BibleQuote{若是你们相信梅瑟，也必相信我}时，呼召这个民族去做的事并非不可能。相反，我们应该对犹太人罪恶的顽固感到惊讶，他们拒绝做我们看见全世界都已经做了的事。\par

28. 关于安息日和割礼以及食物的区别，你们认为梅瑟的教导与基督教导基督徒的不同，我们已表明，如宗徒所说：\BibleQuote{所有这些都是我们的鉴戒}。\BibleRef{格前 10:6} 区别不在于道理，而在于时间。曾有一个时期，这些事应用预像来预言；现在是一个不同的时期，它们应被公开宣布并完全实现。无怪乎犹太人按属肉体的意义理解安息日，就反对基督——祂开始揭示其属神的意义。\Person{福斯图斯}，你若能，就请回答宗徒，他宣告安息日的休息不过是未来事物的阴影。\BibleRef{哥 2:16} 如果犹太人因为不明白真正的安息日是什么而反对基督，那么没有理由让你反对祂，或拒绝学习真正的无罪是什么。因为在那个场合，耶稣似乎特别要废除安息日，当时祂的门徒饿了，就摘下沿途的麦穗吃了，耶稣在回答犹太人时，宣告祂的门徒是无罪的。祂说：\Quote{假如你们了解\BibleQuote{我喜爱仁爱胜过祭献}是什么意思，你们就决不会判断无罪的人了。}\BibleRef{玛 12:7} 他们本应怜悯门徒的急需，因为饥饿迫使他们这样做。但摘麦穗，在基督的教导中是无罪的，在\Person{摩尼}的教导中却是杀人。或者，宗徒摘麦穗是爱德的行为，为的是在吃的时候解放天主的肢体，如你们愚蠢的想法？那么，你们不这样做就是残忍。\Person{福斯图斯}废除安息日的理由是他知道天主的能力永不停止，也不疲倦。这是那些相信所有时间都是天主永恒意志行动之产物的人所说的话。但你会发现这很难与你的教义协调：黑暗种族的反叛打破了你们的神的安息，这安息也因敌人的突然袭击而被打扰；或许神从未有过安息，因为祂从永远就预见了这一切，面对如此可怕的战争、如此大的损失和祂肢体的灾难，祂无法心安。\par

除非基督曾视这安息日——就是你们因无知和冷漠而讥笑的那一个——为关于祂自己的预言之一，否则祂不会如此为它作证。因为当你们赞美基督说祂自愿受苦、因而能选择受苦和复活的时候，祂使自己的身体在安息日于坟墓中安息，不再工作，并使祂第三日的复活——我们称之为主日，即安息日后的第一天，因而也是第八日——证实了第八日的割损礼也预表了祂。因为割损礼意味着什么，岂非消除来自我们血肉之诞生的死亡性吗？所以宗徒说：\BibleQuote{祂脱去自己的肉身，显示了掌权者和权势者的无能，在祂内战胜了他们。}\BibleRef{哥 2:15}这里所说被脱去的肉体，就是那使身体被称为肉身的肉身死亡性。肉身就是死亡性，因为在复活的不死中将没有肉身；正如经上所载：\BibleQuote{肉和血不能承受天主的国。}你们惯常引用这些话来反对我们关于身体复活的信德道理——这复活已发生在主自己身上。你们却忽略了下文，宗徒在其中解释了他的意思。为了表明他这里所说的肉身指什么，他加上说：\BibleQuote{朽坏也不能承受不朽坏。}因为这身体因其死亡性而恰当地被称为肉身，在复活中将发生变化，不再有朽坏和死亡。这是宗徒的陈述，并非我们的假设，他接下来的话证明了这一点：\BibleQuote{我给你们揭示一个奥迹：我们都要复活，但我们并不全被改变。在一刹那，一眨眼之间，在最后的号角吹响时；因为号角将要吹响，死人必要复活成为不朽坏的，我们必要被改变。因为这朽坏的必须穿上不朽坏的，这死亡的必须穿上不死的。}\BibleRef{格前 15:50}为穿上不死的，身体脱去死亡性。这就是割损礼的奥迹，按法律它在第八日举行；而第八日，即主日，安息日后的第一天，由主在其真实意义上实现了。因此经上说：\BibleQuote{脱去自己的肉身，祂显示了掌权者和权势者的无能。}因为借着这死亡性，地狱的敌对权势曾统治我们。据说基督显示了这些，因为在祂自己我们的元首内，祂给出一个榜样，这榜样将在末次复活中，当祂整个身体、教会从魔鬼权势下得到解放时完全实现。这是我们的信德。按保罗所引用的先知预言：\BibleQuote{义人因信德而生活。}这就是我们的成义。连外教人也相信基督死了。但只有基督徒相信基督复活了。宗徒说：\BibleQuote{如果你口里承认耶稣是主，心里相信天主使祂从死人中复活，你必得救。}\BibleRef{罗 10:9}此外，因为我们因信基督的复活而成义，宗徒说：\BibleQuote{祂为我们的过犯被交付，又为我们的成义而复活。}\BibleRef{罗 4:25}因为这种我们赖以成义的复活的信德已由第八日的割损礼预表，宗徒论到亚巴郎——这礼仪从他开始——说：\BibleQuote{他接受了割损礼的标记，作为因信德而成义的印证。}\BibleRef{罗 4:11}那么，割损礼就是基督的预言之一，由梅瑟所写，基督曾说：\BibleQuote{他写了我。}主的话说：\BibleQuote{祸哉，你们经师和法利塞人，假善人！因为你们走遍海洋和陆地，使一个人成为归化者；当他成了，你们却使他比你们更坏两倍，成为地狱之子，}\BibleRef{玛 23:15}这里所指的并非归化者的割损礼，而是他模仿经师和法利塞人的行为，主禁止祂的门徒模仿他们，祂说：\BibleQuote{经师和法利塞人坐在梅瑟的座位上：凡他们所吩咐你们的，你们要遵守；但不要效法他们的行为，因为他们能说不能行。}\BibleRef{玛 23:2}主这些话教导我们，既要尊重梅瑟的教导——即使恶人坐在他的座位上也得教导善事——又要明白归化者成为地狱之子的原因，并非因为他从法利塞人那里听到了法律的言语，而是因为他效法了他们的榜样。这样一个受了割损的归化者，可能会被保罗的话提醒：\BibleQuote{割损礼固然有益，如果你遵守法律。}\BibleRef{罗 2:26}他因效法法利塞人不遵守法律而成为地狱之子。而且他比他们更坏两倍，大概是因为他忽视了履行自愿承担的义务，当他生来并非犹太人，却选择成为犹太人时。\par

你们的嘲弄非常不当，当你们说梅瑟像一个贪食者讨论该吃什么，命令一些东西可自由使用为洁净，另一些为不洁甚至不可触摸。贪食者不加区分，只选择最甜美的食物。也许你们想通过假装不知道或忘记猪肉比羊肉味道更好，来向未入门者夸耀你们节制习惯的无辜。但既然这也由梅瑟以预像预言的方式论及基督，其中动物的肉象征那些将要与基督的身体——即教会——联合，或将要被排除的人，你们正是由不洁动物所预表的；因为你们与天主教信德的分歧表明你们没有反刍智慧之言，也没有分蹄，即正确区分旧约与新约。但你们在采纳你们的\Person{阿狄曼图}的错误意见时表现得更加大胆。\par

31. 你追随\Person{阿迪曼图斯}说，基督在食物上没有分别，只是完全禁止祂的门徒食用动物性食物，而允许平信徒吃任何可吃的东西；并宣称他们不被进入口的东西玷污，只有从口里出来的不雅之物才玷污人。\newline{}\newline{}你这些话确实不雅，因为它们表达了臭名昭著的谎言。\newline{}\newline{}如果基督教导，只有从口里出来的恶事才玷污人，那么为什么这些不是唯独玷污祂门徒的事，从而使任何食物都无需被禁止或不洁呢？难道只有平信徒不被进入口的东西玷污，却被出口的东西玷污？那样的话，他们比圣人们更受保护免于不洁，因为圣人既被进口的东西玷污也被出口的东西玷污。\newline{}\newline{}然而，基督将自己与不吃不喝的若翰相比，说祂自己也吃也喝，我想知道祂吃了什么、喝了什么。当揭露指责两者的悖逆时，祂说：\BibleQuote{若翰来了，也不吃也不喝，你们就说：\InnerQuote{他附了魔。}人子来了，也吃也喝，你们就说：\InnerQuote{看，一个贪食好酒的人，是税吏和罪人的朋友。}}\BibleRef{玛 11:18}\newline{}\newline{}我们知道若翰吃了什么、喝了什么。经上并没有说他什么都不喝，而是说他不喝酒和烈酒；所以他必定是喝水。他并非没有食物，\BibleQuote{他的食物是蝗虫和野蜜。}\BibleRef{玛 3:4}\newline{}\newline{}当基督说若翰不吃不喝时，祂的意思是若翰不使用犹太人惯用的食物。因为主使用了这种食物，所以与若翰相对，称祂为又吃又喝。难道要说主吃的是饼和蔬菜，而若翰不吃这些？那真是奇怪：一个人因为吃蝗虫和蜜就被说成不吃，而另一个人因为吃饼和蔬菜就被说成吃。但无论对吃有什么看法，没有人会被称为酒徒除非他喝酒。那么你为什么称酒为不洁呢？\newline{}\newline{}你禁止这些东西并非为了通过克己来制服身体，而是因为它们是不洁的，因为你称它们是黑暗种族的毒秽；而宗徒却说：\BibleQuote{对于洁净的人，一切都是洁净的。}\BibleRef{铎 1:15}\newline{}\newline{}按照这种学说，基督教导所有食物都一样，但禁止祂的门徒使用摩尼教徒称为不洁的东西。你在哪里找到这项禁令？你不怕用谎言欺骗人；但在天主公义的眷顾下，你如此盲目，以至于为我们提供了驳斥你的方法。\newline{}\newline{}因为我不能不引用\Person{福斯图斯}用来反对梅瑟的那段福音全文来查考；好让我们从中看出先是\Person{阿迪曼图斯}、后是\Person{福斯图斯}所说的谎言，即主耶稣禁止祂的门徒食用动物性食物，而允许平信徒食用。\newline{}\newline{}在基督回应门徒用不洁的手吃饭的指控之后，我们在福音中读到如下内容：\BibleQuote{祂叫了群众来，对他们说：你们听，也要明白。不是入口的玷污人，而是出口的，这才玷污人。然后门徒前来对祂说：祢知道法利塞人听了这话就生气了吗？}\newline{}\newline{}在此，当门徒对祂说话时，按摩尼教徒的说法，祂理当给他们特别的指示，要他们戒绝动物性食物，并表明祂的话\Quote{不是入口的玷污人，而是出口的玷污人}只适用于群众。\newline{}\newline{}那么，让我们听听福音作者所记，主如何回答——不是回答群众，而是回答祂的门徒：\BibleQuote{祂回答说：凡不是我天父所栽种的植物，必要连根拔起。任凭他们吧！他们是盲目的向导。若是瞎子领瞎子，两人都要掉在坑里。}\newline{}\newline{}这样说的原因是，他们渴望遵守自己的传统，却不明白天主的诫命。\newline{}\newline{}那时门徒尚未问老师该如何理解祂对群众所说的话。但现在他们这样做了；因为福音作者接着说：\BibleQuote{于是伯多禄回答说：请给我们解释这个比喻。}\newline{}\newline{}这表明伯多禄认为，当主说\Quote{不是入口的玷污人，而是出口的玷污人}时，祂并不是直白而按字面说，而是像往常一样，愿意以比喻的形式传达一些教导。\newline{}\newline{}那么，当祂的门徒私下提出这个问题时，祂是否像摩尼教徒所说的那样，告诉他们所有动物性食物都是不洁的，他们绝不可触碰？相反，祂责备他们不理解祂明白的语言，并以为那是比喻而其实不是。我们读到：\BibleQuote{耶稣说：你们还不明白吗？岂不知凡入口的，是进入肚腹，又排泄到厕所里吗？但从口里出来的，是从心里发出来的，这才玷污人。因为从心里发出来的有恶念、凶杀、奸淫、邪淫、偷盗、妄证、亵渎。这些都是玷污人的事；至于用不洁的手吃饭，那并不玷污人。}\BibleRef{玛 15:16}\par

32. 此处已完全揭露摩尼教徒的虚谎：显然，主在此事上并未对群众教导一套，私下又对门徒教导另一套。这里有充分证据表明，错误与欺骗在于摩尼教徒，而不在于梅瑟，不在于基督，也不在于一约中以预象教导、另一约中明示的教义——一约中预言，另一约中应验。摩尼教徒怎能说天主教徒不遵守梅瑟所写的任何事物？事实上，他们如今遵守这一切，不是以预象的方式，而是以预象所预指的事物。没有人会说，一个人在圣经写成之后阅读它，却因不构成他所读的文字而不遵守它。文字是他所发声音的预象；尽管他不构成文字，但他若不审视文字便无法阅读。犹太人之所以不信基督，是因为他们连梅瑟律法上明确的字面规条也不遵守。因此基督对他们说：\Quote{你们捐献十分之一的薄荷、茴香，却忽略了法律上更重要的事：公义与怜悯。你们滤出蚊虫，却吞下骆驼。这些是你们该做的，而那些也不可放弃。}（\BibleRef{玛 23:23}）祂也告诉他们，因他们的遗传，使天主的诫命——孝敬父母——失效。由于他们因骄傲与执拗而忽视他们所理解的，他们被公正地蒙蔽了心眼，以致无法理解其他事物。\par

33. 你看，我的论点并非：你若是个基督徒，就必须相信基督所说梅瑟论及祂的话；若你不信，便不是基督徒。你自称愿被当作犹太人或外邦人对待，这是你的事。我的努力是为你关闭一切错误的通道。我也将你从最后的悬崖——当你说福音中的这些段落是伪造的——隔离开来；好使你摆脱这有害见解的影响，被迫相信基督。你说你愿像基督徒多默那样受教，基督并未因多默怀疑祂而拒绝他，反而为了医治他心灵的创伤，向他显示了自己身体上伤痕的印记。这是你自己的话。你愿像多默一样受教，这很好。我原怕你也会把这段落说成是伪造的。那么，请相信基督伤痕的印记吧。因为若印记是真的，则伤痕必然是真的。伤痕不可能是真的，除非祂的身体能承受真实的伤痕；这立刻推翻了摩尼教的全部错误。你若说基督向那怀疑的门徒显示的印记是假的，那么祂必是个欺骗人的导师，而你愿受祂教导便是愿受欺骗。但既然无人愿受欺骗，而许多人愿欺骗人，那么你更可能效法你归给基督的那种教导，而非你归给多默的那种学习。因此，你若相信基督用虚假的伤痕印记欺骗了一个怀疑的求问者，那么你自己必被视为危险的骗子，而非安全的导师。另一方面，若多默触到了基督伤痕的真实印记，你就必须承认基督有一个真实的身体。所以，你若像多默那样相信，你便不再是摩尼教徒。你若不连多默那样也不信，你就只能留在你的不信中。\par

\SourceRecord{Translated by Richard Stothert.；Nicene and Post-Nicene Fathers, First Series；Vol. 4.；Edited by Philip Schaff.；Buffalo, NY: Christian Literature Publishing Co.,；1887.；<http://www.newadvent.org/fathers/140616.htm>.}

% structure: p0001 source=h1 target=section reason=page-title

\section{驳福斯图斯，第十七卷}

\Person{福斯图斯}否认\Person{基督}宣称祂不是来废除法律和先知，而是来成全它们的说法，理由是该说法只见于\BibleBook{玛窦}{福音}，而\Person{玛窦}在所谓说出这些话时并不在场。\Person{奥古斯丁}斥责这种拒绝相信\Person{玛窦}却相信\Person{摩尼}的愚蠢，并阐明了这段经文真正的含义。\par

1. \Person{福斯图斯}说：你们问，为什么我们不接受法律和先知，既然\Person{基督}说祂来不是要废除它们，而是要成全它们。我们从何得知\Person{耶稣}说了这话？是从\Person{玛窦}，他宣称\Person{耶稣}在山上传了这话。这话是在谁面前说的？是在\Person{伯多禄}、\Person{安德肋}、\Person{雅各伯}和\Person{若望}面前——只有这四位；因为其余的人，包括\Person{玛窦}本人，当时还没有被拣选。难道这四位中的一位——即\Person{若望}——没有写福音吗？他写了。他提到\Person{耶稣}的这句话吗？没有。那么，怎么会出现这样的情况：\Person{若望}当时在山上，却没有记下这句话，而\Person{玛窦}是在\Person{耶稣}下山后很久才成为\Person{基督}门徒的，却记下了这句话？所以，首先，我们必须怀疑\Person{耶稣}是否说过这些话，因为恰当的见证人对此保持沉默，而我们只有不太可信的见证人的权威。但除此之外，我们还会发现，欺骗我们的并不是\Person{玛窦}本人，而是某个假借他之名的人，这从叙述的间接风格可以明显看出。我们读到：\BibleQuote{耶稣从那里前行，看见一个人在税关那里坐着，名叫玛窦，就对他说：\InnerQuote{跟随我！}他就起来跟随了耶稣。}（\BibleRef{玛 9:9}）没有人在写自己时会说：\Quote{他看见一个人，并叫他；他就起来跟随了他}，而会说：\Quote{他看见了我，并叫了我，我就跟随了他}。显然，这段并不是\Person{玛窦}本人所写，而是某个假借他之名的人写的。那么，即使这段话出自\Person{玛窦}之手，也不会是真的，因为\Person{耶稣}在山上传道时他并不在场；更何况，作者并非\Person{玛窦}本人，而是某个既借用耶稣之名又借用玛窦之名的人，这就更明显地证明了这段话的虚假性。\par

2. \Person{基督}告诉犹太人不要以为祂来是废除法律的这段经文本身，反倒更显示出祂确实废除了法律。因为，如果祂没有做过类似的事，犹太人就不会怀疑祂。祂的话是：\BibleQuote{你们不要以为我来是废除法律。}假设犹太人回答说：\Quote{您的哪些行为会使我们产生这样的怀疑？是因为您揭露割礼、违反安息日、弃绝祭献、不区分食物吗？}这会是\Quote{不要以为}这句话的自然回答。犹太人完全有理由认为耶稣废除了法律。如果这不是废除法律，那什么是呢？但事实上，法律和先知认为自己已经完美无瑕，不想被成全。它们的作者和父亲禁止对它们进行增减，正如我们在《申命纪》中读到的：\BibleQuote{凡我所吩咐你们的话，你们要谨守遵行，不可加添，也不可删减，好使你们能遵守我\Person{上主}你们天主的诫命。}（\BibleRef{申 12:32}）因此，无论耶稣是向右转，通过加添来成全法律和先知，还是向左转，通过删除来废除它们，他都冒犯了法律的作者。所以，这句话要么有其他含义，要么就是伪造的。\par

3. \Person{奥古斯丁}回答说：何等惊人的愚蠢！不信\Person{玛窦}所记关于\Person{基督}的事，却相信\Person{摩尼}！如果因为\Person{玛窦}不在场就不可信他所记\Person{基督}说的\BibleQuote{我来不是废除法律和先知，而是来成全}，那么\Person{摩尼}在场吗？\Person{基督}显现于人时，他出生了吗？所以，按照你的规则，你也不该相信\Person{摩尼}说的任何关于\Person{基督}的话。反之，我们拒绝相信\Person{摩尼}关于\Person{基督}的话，不是因为他没有作为见证人亲临\Person{基督}的言行，而是因为他反对\Person{基督}的门徒和基于他们权威的福音。宗徒藉圣神说话，告诉我们这样的导师会出现。关于他们，他对信徒说：\BibleQuote{若有人给你们宣讲福音，与你们所接受的不同，他就该受诅咒。}（\BibleRef{迦 1:9}）如果没有人亲眼见过、亲耳听过\Person{基督}，就不能说关于\Person{基督}的真话，那么现在没有人可以相信了。但如果信徒现在能说出关于\Person{基督}的真话，因为真理已通过那些亲眼见过、亲耳听过的人以口传或书写的方式传递下来，那么\Person{玛窦}为什么不能从同门\Person{若望}那里听到真理呢？即使\Person{若望}在场而他自己不在，正如我们这些后来出生的人和以后出生的人都可以从\Person{若望}的著作中了解关于\Person{基督}的真理。这样，\Person{路加}和\Person{马尔谷}的福音（他们是宗徒的同伴），以及\Person{玛窦}的福音，都与\Person{若望}的福音有同样的权威。此外，\Person{主}自己也可能将那些先蒙召的人已经见证的事告诉\Person{玛窦}。你的想法是，\Person{若望}应该记下\Person{主}的这句话，因为他当时在场。但也许因为不可能写下所有从\Person{主}那里听到的话，他选择写一些，而省略了这句和其他话。他不是在福音结束时说：\BibleQuote{耶稣所行的还有许多别的事；假使要一一写出来，我想所要写的书，连这世界也容不下}（\BibleRef{若 21:25}）吗？这表明他有意省略了许多事。但如果你选择\Person{若望}作为关于法律和先知的权威，我只要求你相信他关于它们的见证。正是\Person{若望}写道依撒意亚看见了\Person{基督}的荣耀。（\BibleRef{若 12:41}）在他的福音中，我们还找到了已经讨论过的经文：\BibleQuote{若是你们相信梅瑟，必会相信我，因为他是指着我而写的。}（\BibleRef{若 5:46}）因此，你的各种回避都被堵住了。你应当坦白地说，你不相信基督的福音。因为信你所愿信的，不信你所不愿信的，就是信自己，而不是信福音。\par

4. \Person{福斯图斯}认为自己极为聪明，因为他证明\Person{玛窦}不是这部福音的作者，因为当谈及自己的蒙召时，他不是说：\Quote{他看见了我，并对我说：跟随我；}而是说：\Quote{他看见了他，并对他说：跟随我。} 这必然是无知或出于误导的意图才说的。\Person{福斯图斯}不太可能无知到没有读过或听说过，叙述者在谈论自己时，常常使用一种仿佛在谈论别人的表达方式。更可能的是，\Person{福斯图斯}意图迷惑那些比他更无知的人，希望抓住不少对这些事情不熟悉的人。无需引用其他著作中世俗作家的例子来教导我们的朋友并反驳\Person{福斯图斯}。我们在\Person{福斯图斯}本人上面引用的\Person{梅瑟}片段中就能找到例子，他并不否认，甚至断言那些是\Person{梅瑟}写的，只是不是关于\Person{基督}的。那么，当\Person{梅瑟}写他自己时，他是否说：\Quote{我说了这个} \Quote{我做了那个}？难道不是\Quote{\Person{梅瑟}说了} \Quote{\Person{梅瑟}做了}？或者说\Quote{主召叫了我} \Quote{主对我说了}，而不是\Quote{主召叫了\Person{梅瑟}} \Quote{主对\Person{梅瑟}说了}等等？同样，\Person{玛窦}也用第三人称谈论自己。\Person{若望}也是如此；因为在他的书卷末尾，他说：\BibleQuote{\Person{伯多禄}转身，看见\Person{耶稣}所爱的那个门徒，就是晚餐时靠在他胸膛上，问主说：\InnerQuote{主，出卖你的是谁？}的那个门徒。} 难道他说\Quote{\Person{伯多禄}转身，看见了我}吗？或者你会由此论证\Person{若望}没有写这部福音？但他稍后补充说：\BibleQuote{这个门徒就是为\Person{耶稣}作证，并写下这些事的；我们知道他的作证是真实的。}\BibleRef{若 21:20} 难道他说\Quote{我是为\Person{耶稣}作证并写下这些事的门徒，我们知道我的作证是真实的}吗？显然，这种风格在叙事作者中很常见。主亲口使用的例子不胜枚举。他说：\BibleQuote{当人子来临时，他要在世上找到信德吗？}\BibleRef{路 18:8} 而不是\Quote{当我来临时，我要找到吗？}又说：\BibleQuote{人子来了，也吃也喝}\BibleRef{玛 11:19}，而不是\Quote{我来了}。又说：\BibleQuote{时候要到，且现在就是，死人要听见天主子的声音，凡听见的就要活}\BibleRef{若 5:25}，而不是\Quote{我的声音}。还有许多其他地方。这些足以满足探求者并反驳嘲弄者。\par

5. 任何人都能看出这个论点的薄弱：\Person{基督}不可能说\Quote{你们不要以为我来是废除法律或先知；我来不是为废除，而是为成全}，除非他做了什么引起这种怀疑的事。当然，我们承认，未启蒙的犹太人可能将\Person{基督}视为法律和先知的废除者；但他们的怀疑恰恰确证了那位真实而诚信者，在说他来不是废除法律和先知时，所指的正是犹太人的法律，而不是别的法律。这一点由随后的话得到证明：\BibleQuote{我实在告诉你们，天地过去，法律的一撇一划也决不会过去，直到全部实现。所以，谁若废除这些诫命中最小的一条，并这样教训人，在天国里他将称为最小的；但谁若实行并教训这些诫命，在天国里他将称为大的。} 这话是针对法利塞人说的，他们口头上教导法律，行动上却违反法律。\Person{基督}在另一处论到法利塞人说：\BibleQuote{他们对你们说的话，你们要遵行；但不要效法他们的行为，因为他们只说而不做。}\BibleRef{玛 23:3} 所以，他在这里也加上：\BibleQuote{我告诉你们：除非你们的义德超过经师和法利塞人的义德，你们决进不了天国。}\BibleRef{玛 5:17} 意思是，除非你们既实行又教导他们所教导而不实行的事，你们决进不了天国。因此，法利塞人只教导而不遵守的法律，\Person{基督}说他来不是废除，而是成全；因为这就是与\Person{梅瑟}的座位相连的法律，法利塞人坐在其上，他们因为只说而不做，所以应听他们的话，但不应效法他们。\par

6. \Person{福斯图斯}不明白，或假装不明白，什么是成全法律。他以为这个说法是指在法律上添加话语，关于这点，经上记载不可在天主圣经上添加或删减。\Person{福斯图斯}由此论证，既然法律被说得如此完美，不可添加或删减，那么它就不能被成全。必须告诉\Person{福斯图斯}，法律是通过按其所吩咐的生活来成全的。正如宗徒所说：\BibleQuote{爱德是法律的满全。}\BibleRef{罗 13:10} 主曾屈尊显示并传授了这种爱，将圣神赐予祂的信徒。同一位宗徒又说：\BibleQuote{天主的爱，借着赐给我们的圣神，已倾注在我们心中。}\BibleRef{罗 5:5} 主亲自说：\BibleQuote{如果你们彼此相爱，众人因此就可认出你们是我的门徒。}\BibleRef{若 13:35} 因此，法律既通过遵守其规诫，也通过实现其预言而得以成全。因为\BibleQuote{法律是藉梅瑟传授的，恩宠和真理却是由耶稣基督而来的。}\BibleRef{若 1:7} 法律本身通过被成全，成为恩宠和真理。恩宠是爱德的满全，真理是预言的实现。既然恩宠和真理都是借基督而来，那么，他来不是废除法律，而是成全法律；不是填补法律中的任何缺陷，而是服从法律所记载的。基督自己的话证明了这一点。因为他不是说，法律的一撇一划决不会过去，直到它的缺陷被填补，而是说\BibleQuote{直到全部实现}。\par

\SourceRecord{Translated by Richard Stothert.；Nicene and Post-Nicene Fathers, First Series；Vol. 4.；Edited by Philip Schaff.；Buffalo, NY: Christian Literature Publishing Co.,；1887.；<http://www.newadvent.org/fathers/140617.htm>.}

% structure: p0001 source=h1 target=section reason=page-title

\section{驳斥福斯图斯，第十八卷}

\Emph{基督与预言之关系（续）}\par

1. \Person{福斯图斯}说：\Quote{\InnerQuote{我来不是要废除法律，而是要成全。}如果这是基督的话，除非别有意义的，它们对你们和对我都一样不利。你们的基督教和我的基督教都是基于这样的信念：基督来是要废除法律和先知。你们的行为证明了这一点，即使你们口头上否认。正是因此，你们无视法律和先知的规诫。正是因此，我们都承认耶稣是新约的创立者，这就意味着承认旧约已被废除。那么，我们怎能相信基督说了这些话，而不先承认我们迄今完全错误，并通过走上服从法律和先知的道路，仔细遵守它们的一切要求来表明我们的悔改呢？这样做了之后，我们才能诚实地相信耶稣说祂来不是要废除法律，而是要成全。而事实上，你们指责我不相信你们自己也不相信的东西，因此那是假的。}\par

2. 但是，假设我们迄今都错了。现在该怎么办？我们是否要服从法律，因为基督没有废除，反而成全了？我们是否要藉着割礼雪上加霜，并且相信天主喜欢这样的圣事？我们是否要守安息日的休息，把自己捆绑在萨图尔努斯的枷锁下？我们是否要用公牛、公羊和山羊的屠杀——更不用说人了——来餍足犹太人的魔鬼（因为他不是神），并且服从法律和先知，采用那些我们因此而放弃偶像崇拜的习俗，只是更加残忍？最后，我们是否要把一些动物的肉称为洁净的，另一些称为不洁净的，其中按照法律和先知，猪肉有特别的污秽？当然你们会承认，作为基督徒，我们绝不能做这些事，因为你们记得基督说过：\BibleQuote{一个人受割礼后就成了加倍的地狱之子。}\BibleRef{玛 23:15}同样明显的是，基督自己也没有守安息日，也没有命令人守安息日。关于食物，他明确地说：\BibleQuote{人不是被入口的东西污秽，而是被出口的所污秽。}\BibleRef{玛 15:11}关于祭献，他也常说：\BibleQuote{天主喜欢仁爱，而不是祭献。}\BibleRef{玛 9:13}那么，他说\Quote{\InnerQuote{我来不是要废除法律，而是要成全}}这句话又是什么意思呢？如果基督说了这句话，他必定有其他意思，或者——这不可想像——他说了谎，或者他根本没有说过。没有一个基督徒会允许耶稣说假话；因此他要么没有说过这句话，要么是用另一种意思说的。\par

3. 拿我来说，作为一个\OriginalTerm{摩尼教徒}{Manichæan}，这句话对我没有什么困难，因为从一开始我就被教导要相信，许多在圣经中归于救主名下的东西都是伪造的，因此必须检验它们是否真实、正确和纯正；因为那在夜间来的敌人几乎在每一段中撒了莠子。所以这些文字尽管冠以神圣的名号，我却不感到害怕；我仍然有权检验这是出自白天播种的好撒种人，还是夜间播种的恶者之手。但是你们有什么办法逃避这个困难呢？你们不加检验地接受一切，谴责运用理性——这是人性的特权，并且认为区分真理与谬误是不敬虔，就像小孩子怕鬼一样害怕区分好与坏。假如一个犹太人或者其他了解这些话的人问你们，既然基督说他来不是要废除法律，而是要成全，你们为什么不遵守法律和先知的规诫呢？你们要么必须加入犹太人的迷信愚妄，要么宣布这节经文为假，要么否认你们是基督的追随者。\par

4. \Person{奥斯定}回答说：既然你重复那些已屡次暴露和驳斥的东西，我们只好满足于重复反驳。法律和先知中基督徒不遵守的那些东西，只是他们所遵守之事的预像。这些预像是未来事物的预象，当基督完全启示出这些事物本身时，它们必然被除去，以便在具体的除去中成全法律和先知。正如先知书上记载，天主将订立新盟约，\BibleQuote{不像我赐给他们祖先的那样。}\BibleRef{耶 31:32}旧约下的百姓心硬，因此赐给他们许多诫命，不是因为它们好，而是因为它们适合那百姓。尽管如此，在这一切中，未来已被预告和预象，虽然百姓不明白他们自己礼仪的意义。当这些被预表的明显事物出现后，我们不必遵守预像，但阅读它们以了解其意义。同样，先知书中也预言：\BibleQuote{我要拿走他们的石心，赐给他们血肉之心。}\BibleRef{则 11:19}——即一颗有感觉的心，代替无感觉的心。宗徒在以下话语中暗指这一点：\BibleQuote{不是在石版上，而是在血肉的心版上。}\BibleRef{格后 2:3}血肉的心版与血肉之心相同。既然这些仪式的废除已被预告，法律和先知若不通过这种废除就不能被成全。而现在，预言已经实现，法律和先知的成全正体现于表面看来完全相反的东西中。\par

5. 我们不怕面对你们对安息日的嘲弄，你们称它为\Quote{萨图尔努斯的枷锁}。这是一个愚蠢且无意义的说法，之所以出现在你们脑中，只因为你们习惯在你们所谓的\Quote{太阳日}崇拜太阳。你们称为\Quote{太阳日}，我们称为\Quote{主的日子}，在这一天我们不崇拜太阳，而是崇拜主的复活。同样，先祖们遵守安息日的休息，并非因为他们崇拜\Person{萨图尔努斯}，而是因为当时那是必须的，因为它是未来事物的影子，正如宗徒所证实的。\BibleRef{哥 2:17}外邦人——宗徒论及他们说\BibleQuote{他们崇拜事奉受造物而不事奉造物主}\BibleRef{罗 1:25}——将他们的神祇之名赋予一周的日子。到目前为止，你们也做了同样的事，只是你们只崇拜两个最明亮的天体，不像外邦人那样崇拜其余的星辰。此外，外邦人将他们的神祇之名给了月份。为了纪念\Person{罗慕路斯}——他们相信他是\Person{玛尔斯}的儿子——他们将第一个月奉献给\Person{玛尔斯}，称为三月。第二个月，四月，不是以任何神祇命名，而是来自于\Quote{开放}一词，因为花蕾通常在这个月开放。第三个月称为五月，以纪念\Person{墨丘利}的母亲\Person{迈亚}。第四个月称为六月，来自\Person{朱诺}。其余月份直到十二月，过去是按数字命名的。然而，第五和第六个月从被尊为神明的人那里得到了七月和八月的名称；而其他月份，从九月到十二月，继续以数字命名。一月，则是从\Person{雅努斯}而来，二月从称为\Quote{\OriginalTerm{费布鲁埃}{Februæ}}的路佩耳祭仪而来。难道我们必须说你们在三月崇拜\Person{玛尔斯}神吗？但那是你们大张旗鼓举行你们称为\Quote{\OriginalTerm{贝玛}{Bema}}的节日的月份。但如果你们认为可以庆祝三月而不想到\Person{玛尔斯}，那为什么仅仅因为外邦人将那日称为\Quote{萨图尔努斯日}，你们就试图将\Person{萨图尔努斯}之名与圣经所规定的第七天的休息联系起来呢？圣经中那日的名称是安息日，意为休息。你们的嘲弄既无理性又亵渎。\par

6. 关于动物祭献，每个基督徒都知道，那些祭献是为适合一个悖逆的百姓而规定的，并非因为天主喜悦它们。然而，即使在这些祭献中，也有我们所享受之事的预像；因为我们非藉血不能获得净化或天主的平息。这些预像的应验是在基督内，藉着祂的血，我们得到净化并蒙救赎。在这些神圣谕示的图像中，公牛代表基督，因为祂藉着十字架的角驱散恶人；羔羊代表祂无与伦比的无辜；山羊代表祂被造成罪恶肉躯的样式，为能藉着罪而定了罪案。\BibleRef{罗 8:3}无论你们选择哪一种祭献，我都将在其中向你们指出基督的预言。这样，我们论及割礼、安息日、食物的分别和动物祭献，已表明这一切都是我们的预像和预言，基督来不是为废除，而是为应验，藉此应验了所预告的一切。你们的敌对者是宗徒，我用他自己的话陈述他的意见：\BibleQuote{这一切都是我们的预像}\BibleRef{格前 10:6}。\par

7. 如果你们从\Person{摩尼}那里学会了任凭不敬，只接受不与你们错误相矛盾的福音部分，而拒绝其余，那么我们从宗徒那里学会了虔诚的谨慎，视任何向我们宣讲不同于我们所领受的福音的人为可咒骂的。因此，天主教基督徒将你们视为莠子；因为主在解释莠子的意义时，它们不是指圣经中掺入的谎言，而是恶者的子女——即，那些效法魔鬼诡诈的人。天主教基督徒并非什么都信；因为他们不信\Person{摩尼}或任何异端者。他们也不谴责使用人的理性；但你们所谓的推理，他们证明是谬误的。他们也不认为分辨真假是亵渎；因为他们分辨天主教信仰的真理和你们教义的谎言。他们也不害怕分别善恶；相反，他们主张恶不是自然本性，而是反本性的。他们对你们的黑暗种族一无所知，你们说那是从自身的一个原则产生的，并反对天主的国，而你们的神似乎比孩子们怕鬼更害怕它；因为，按照你们所说，他把自己蒙起来，以免看见自己的肢体被仇敌的攻击夺取和掠夺。总之，天主教基督徒在基督的话语上毫无困难，尽管从某种意义可以说他们不守法律和先知；因为藉着基督的恩宠，他们借着爱天主和爱人而遵守法律；因为全部法律和先知都系于这两条诫命。\BibleRef{玛 22:40}此外，他们在基督和教会身上看到旧约所有预言的应验，无论是行为形式、象征礼仪还是比喻语言。因此，我们既不参与迷信的愚行，也不宣告这节经文为虚假；也不否认我们是基督的跟随者；因为根据我尽己所能阐述的原则，基督来不是为废除而是为应验的法律和先知，正是教会所承认的同一法律和先知。\par

\SourceRecord{Translated by Richard Stothert.；Nicene and Post-Nicene Fathers, First Series；Vol. 4.；Edited by Philip Schaff.；Buffalo, NY: Christian Literature Publishing Co.,；1887.；<http://www.newadvent.org/fathers/140618.htm>.}

% structure: p0001 source=h1 target=section reason=page-title

\section{驳斥福斯托斯，第十九卷}

\Emph{\Person{Faustus}愿意承认基督可能说过祂来不是为废除法律和先知，而是为成全；但如果说，祂是为了安抚犹太人，且是在一种修正的意义上。奥古斯丁答复，并进一步阐述天主教关于预言及其应验的观点。}\par

1. \Emph{\Person{Faustus}}说：\Quote{我愿承认基督说过祂来不是为废除法律和先知，而是为成全。但耶稣为何说这话？是为了平息那些因看见自己神圣制度被基督践踏而愤怒、并将祂视为狂妄亵渎者，不愿听祂更不愿跟从祂的犹太人吗？还是为了教导我们这些外邦信徒，好叫我们温顺忍耐地承受犹太人法律和先知所加在我们颈上的轭？您自己几乎不能认为基督的话旨在使我们服从希伯来人的法律和先知。所以，我对此话的另一种解释必定是真的。人人都知道犹太人总以言语和实际行动攻击基督。自然，他们一听到基督在废除他们的法律和先知，便愤怒不已；而为了安抚他们，基督很可能会告诉他们不要以为祂来是废除法律，而是来成全法律。这其中并无谎言或欺骗，因为祂是在一般意义上使用\Quote{法律}一词，并非特指某种法律。}\par

2. 有三种法律。一种是希伯来人的法律，宗徒称之为罪与死的法律。\BibleRef{罗 8:2} 第二种是外邦人的法律，他称之为自然法律。他说：\BibleQuote{外邦人若顺着本性行法律上的事，他们虽然没有法律，自己就是自己的法律；如此显出法律的功用刻在他们心里。}\BibleRef{罗 2:14} 第三种法律是真理，宗徒说：\BibleQuote{在基督耶稣内生命之神的法律，使我脱离了罪与死的法律。}\BibleRef{罗 8:2} 既然有三种法律，我们须仔细查考，基督说祂来不是为废除法律而是为成全时，指的是哪一种法律。同样，有犹太人的先知，有外邦人的先知，也有真理的先知。犹太人的先知自然是人所共知的。若有人对外邦人的先知有疑问，请听保禄在致弟铎书关于克里特人的话：\BibleQuote{他们本族中有一位先知说：\InnerQuote{克里特人常说谎，是恶兽，是好吃懒做的。}}\BibleRef{铎 1:12} 这证明外邦人也有他们的先知。真理也有其先知，我们从耶稣和保禄得知。耶稣说：\BibleQuote{看，我派遣先知和贤士到你们这里来，其中有些你们要各处杀害。}\BibleRef{玛 23:34} 保禄说：\BibleQuote{主自己设立了宗徒，其次是先知。}\BibleRef{弗 4:11}\par

3. 既然\Quote{法律和先知}可能有三种不同含义，耶稣所用的词语意思也就不能确定，尽管我们可以从下文推测。因为如果耶稣继续谈论割损、安息日、祭献以及希伯来人的礼仪，并添加一些作为成全，那么毫无疑问，祂所说\Quote{我来不是为废除，而是为成全}指的是犹太人的法律和先知。但基督并未提及这些，反而只谈到自古以来的诫命：\Quote{不可杀人；不可奸淫；不可作假见证。}可以证明，这些诫命在远古时代就由厄诺客和舍特以及其他义人颁布于世，他们从崇高的天使那里领受了这些规诫，以驯服人类的野蛮本性。由此看来，耶稣说的是真理的法律和先知。所以，我们看到祂成全了上述诫命。祂说：\BibleQuote{你们听见有对古人说：\InnerQuote{不可杀人}；但我对你们说：连生气也不可。}这就是成全。再者：\BibleQuote{你们听见有话说：\InnerQuote{不可奸淫}；但我对你们说：连贪恋也不可。}这也是成全。还有：\BibleQuote{有话说：\InnerQuote{不可作假见证}；但我对你们说：什么誓都不可起。}这也是成全。祂如此既确认了古训，又补足了它们的缺欠。至于祂似乎提到的某些犹太人的规条，祂非但没有成全，反而代之以相反的规定。祂接着说：\BibleQuote{你们听见有话说：\InnerQuote{以眼还眼，以牙还牙}；但我对你们说：不要抵抗恶人；有人打你的右脸，连另一面也转给他。}这不是成全，而是废除。再者：\BibleQuote{有话说：\InnerQuote{当爱你的朋友，恨你的仇敌}；但我对你们说：爱你们的仇敌，为迫害你们的人祈祷。}这也是废除。还有：\BibleQuote{有话说：\InnerQuote{凡休妻的，应当给她休书}；但我对你们说：凡休妻的，若不是为淫乱，就是使她犯奸淫，并且若再娶别的女人，他本身也是犯奸淫。}\BibleRef{玛 5:21} 这些规条显然被废除，因为它们是梅瑟的规条；而其他的则被成全，因为它们是古代义人的规条。如果您同意这解释，我们可以承认耶稣说过祂来不是为废除法律，而是为成全。若您不赞同这解释，请提出您自己的。只是要当心，不要使耶稣成为说谎者，也不要使自己成为犹太人，因为基督没有废除法律，您就被束缚去遵守法律。\par

4. 如果有一个纳匝肋人，或有时被称为西马基人的人，用耶稣的这句话——\Quote{祂来不是为废除法律，而是为成全法律}——来与我辩论，我会觉得难以回答他。因为无可否认，耶稣在来临时，身心都受法律和先知的影响。此外，我所指的那些人，虽然自称基督徒，却遵行割损礼、守安息日、禁食猪肉及诸如此类的事物，按法律而行。他们显然和你一样，被基督所说的这节经文误导，即祂来不是为废除法律，而是为成全法律。若不先摆脱这节棘手的经文，要与这样的反对者辩论并不容易。但与你辩论，我却毫无困难，因为你没有任何依据；你似乎并非使用论证，而只是出于恶意，试图让我相信基督说过连你自己显然也不相信祂说过的话。凭这节经文，你指责我迟钝和逃避，而你自己却没有丝毫遵守法律而非废除法律的迹象。难道你也像犹太人或是纳匝肋人一样，以受割损这令人羞耻的标记而夸耀？你以守安息日而自豪？你能因没有吃猪肉而自我庆幸？或者你能因以祭品的血和犹太祭献的馨香满足了天主的食欲而夸口？如果不能，那你为什么还坚持说基督来不是为废除法律，而是为成全法律？\par

5. 我不断感谢我的老师，他阻止我陷入这个错误，使我至今仍是基督徒。因为我和你们一样，读这节经文时没有充分考虑，几乎决定成为一个犹太人。而且有理由如此；因为如果基督来不是为废除法律，而是为成全法律，而一个器皿要充满，就不能是空的，而应已经部分装满，于是我得出结论：除了以色列人之外，没有人能成为基督徒，以色列人已经被法律和先知几乎装满，他们来到基督那里，以便按自己的容量被完全充满。我还得出结论，他来到时不应毁掉他已拥有的东西；否则就不是成全，而是倒空。那么，似乎我这个外邦人，来到基督那里什么也得不到，因为我没有什么可以让他增加而充满。经查考，这预备好的供应包括安息日、割损礼、祭献、新月、洗礼、无酵节、食物和饮料及衣服的区分，以及其他多不胜数的事物。那么，似乎这就是基督来不是为废除，而是为成全的。自然必须如此：因为法律若没有诫命，或先知若没有预言，还算什么呢？此外，还有那可怕的诅咒，是针对那些不按法律书上所写的一切去行的人。\BibleRef{申 27:15} 一方面有来自天主的对这诅咒的恐惧，另一方面有基督作为天主圣子，似乎说祂来不是为废除这些事，而是为成全它们，那么还有什么能阻止我成为一个犹太人呢？\Person{摩尼}的智慧教导使我免于这个危险。\par

6. 但你怎敢用这节经文来反对我？或者为什么只针对我，而它也针对你自己呢？如果基督不废除法律和先知，那么基督徒也不应废除它们。那么你为什么废除它们？你是否开始意识到你不是基督徒？你怎么敢用各种工作亵渎那在法律和一切先知中被宣告为神圣的日子，他们说是创造世界的天主亲自安息的日子，而你不怕对犯安息日者所宣告的死刑，或对违法者的诅咒？你怎么敢拒绝在你身上接受那不雅的割损标记，而法律和所有先知都宣告它是可敬的，尤其是在\Person{亚巴郎}身上，在他被认为有信德之后；因为犹太人的天主不是宣告说，凡没有这耻辱标记的，必从祂的民中剪除吗？你怎么敢忽视那指定的祭献，这些祭献在\Person{梅瑟}和律法下的先知，以及\Person{亚巴郎}在他的信德中，都非常重视？如果你相信基督来不是为废除这些事，而是为成全它们，你怎么敢因不区分食物而玷污你们的灵魂？你为什么废弃每年的无酵节和指定的羔羊祭献，而按照法律和先知，这应永远遵守？总之，如果基督没有废除它们，你为什么如此轻视新月、洗礼、帐棚节，以及法律和先知中所有其他肉身的规条？所以我很有理由说，要为自己忽视这些事辩护，你必须要么放弃你自称为基督门徒的宣告，要么承认基督自己已经废除了它们；从这承认必然推出，要么这节经文是伪造的，其中说基督来不是为废除法律，而是为成全它，要么这话有与你们所想的完全不同的含义。\par

7. \Person{奥斯定}答说：如果你们因福音的权威而承认基督说过，祂来不是为废除法律和先知，而是为成全它们，那么你们也应因宗徒的权威而承认同样的事，当他说：\BibleQuote{这一切事都是我们的鉴戒；}以及论到基督：\BibleQuote{祂并不是\InnerQuote{是}而又\InnerQuote{非}的，在祂只有\InnerQuote{是}；因为天主的一切恩许，在祂内部都成了\InnerQuote{是}。}\BibleRef{格后 1:19}这就是说，在祂内得以实现并成全。这样，你们就会最清楚地看到基督成全了什么法律，以及如何成全。你们试图用你们的三类法律和三类先知来逃避，这是徒劳的。十分明显，而且新约对此毫不含糊：基督来不是为废除，而是为成全的，是哪一种法律和哪一种先知。由梅瑟所赐的法律，正是藉耶稣基督而成为恩宠和真理的。\BibleRef{若 1:17}由梅瑟所赐的法律，正是基督所说的：\BibleQuote{他为我写了。}\BibleRef{若 5:46}因为无疑地，这就是那法律，它介入是为了使过犯增多；\BibleRef{罗 5:20}你们往往无知地引用这些话来诋毁法律。请读一下那里关于这法律所说的话：\BibleQuote{法律是圣的，诫命也是圣的、公义的、善的。那么，善的竟使我死吗？绝对不可！但罪恶为显出是罪恶，藉着善的使我死。}\BibleRef{罗 7:12}法律的介入使过犯增多，不是因为法律要求了错事，而是因为骄傲和自以为义的人，在认识了法律圣、公义、善的诫命之后，作为违法者增添了额外的罪责；如此，他们因谦卑而学会，只有藉着恩宠，因着信德，才能从作为违法者受法律约束中得到释放，并作为义人与法律和好。所以同一宗徒说：\Quote{但在信德来到以前，我们被看守在法律之下，被包围起来，直到那以后要显示的信德。因此法律成了我们在基督耶稣内的启蒙导师，使我们可以因信德而成义；但信德来到以后，我们就不再在启蒙导师之下了。}这就是说，我们不再受法律的刑罚，因为我们藉恩宠获得了自由。在我们谦卑地领受圣神的恩宠之前，文字对我们只有死亡，因为它要求我们所不能作出的服从。因此保禄也说：\BibleQuote{文字叫人死，圣神却叫人活。}\BibleRef{格后 3:6}他又说：\BibleQuote{如果曾有一种法律能赐予生命，那么义德诚然就是出于法律了；但圣经把众人都禁锢在罪恶之下，为的叫恩许因对耶稣基督的信德，赐给相信的人。}\BibleRef{迦 3:21}再一次：\BibleQuote{法律因肉性而软弱，有所不能；天主却派遣了自己的儿子，带着罪恶肉身的形状，为罪恶而定了罪恶的罪案，使法律所要求的义德，在我们这些不随从肉性，而随从圣神的人身上得以成全。}\BibleRef{罗 8:3}这里我们看到基督来不是为废除法律，而是为成全。正如法律使骄傲的人在过犯的罪责之下，借他们所不能遵守的诫命增加了他们的罪，同样，这同一法律的义德，藉圣神的恩宠，在那些从基督学习温良和心谦的人身上得以成全；因为基督来不是为废除法律，而是为成全。此外，即使对于在恩宠之下的人，在这必死的生命中，完全遵守法律所写的\BibleQuote{不可贪恋}也是困难的；基督，藉着祂肉身的祭献，作为我们的司祭，为我们获得宽恕。在这方面，祂也成全了法律；因为我们因软弱而亏缺的，由祂的圆满所补足，祂是元首，我们是祂的肢体。因此若望说：\BibleQuote{我的孩子们，我写这些事给你们，是要你们不犯罪；但若有人犯了罪，我们在父那里有一位辩护者，就是正义的耶稣基督；祂自己就是为我们罪过的赎罪祭。}\BibleRef{若一 2:1}\par

8. 基督也成全了先知们的预言，因为天主的恩许在祂内得以实现。正如宗徒在上文所引的经句所说：\BibleQuote{天主的恩许在祂内都成了\InnerQuote{是}。}他又说：\BibleQuote{我说，耶稣基督为了天主的真理，成了受割损者的仆役，为证实向先祖们所作的恩许。}\BibleRef{罗 15:8}因此，凡在先知书中应许的，无论是明言还是预象，无论藉言语还是藉行动，都在祂内得以成全；祂来不是为废除法律和先知，而是为成全。你们不明白，如果基督徒继续使用那些预表未来之事的行动和规条，那么唯一的含义就是所预表的事尚未到来。要么预表之事尚未到来，要么如果已经到来，预象就成为多余或误导的。因此，如果基督徒不实行先知在希伯来入中所吩咐的一些事，这非但不是像你们所认为的表明基督没有成全先知，反而表明祂成全了。基督如此完全地成全了这些预象所预表的事，以致不再需要预表了。所以主亲自说：\BibleQuote{法律和先知到若翰为止。}\BibleRef{路 16:16}因为那将违法者禁锢在加增的罪责之中，并指向那以后要显示的信德的法律，藉耶稣基督成了恩宠，藉着祂恩宠更加丰富。如此，那在文字要求上没有成全的法律，在恩宠的自由中得了成全。同样，法律中一切预言救主来临的事，无论藉言语还是藉预表性的行动，都在耶稣基督内成了真理。因为\BibleQuote{法律是藉梅瑟赐予的，恩宠和真理却是由耶稣基督而来的。}\BibleRef{若 1:17}在基督来临时，天主的国开始被宣讲；因为法律和先知到若翰为止：法律，为使违法者渴望救恩；先知，为预告救主。无疑，自基督升天以来，教会中仍有先知。关于这些先知，保禄说：\BibleQuote{天主在教会中设立的，第一是宗徒，第二是先知，第三是教师，}等等。\BibleRef{格前 12:28}\Quote{法律和先知到若翰为止}这句话，不是指这些先知说的，而是指那些预言基督首次来临的先知；既然基督已降临，这些预言显然不能再预言了。\par

9. 因此，当你问为什么基督徒不受割礼，如果基督来不是为废除法律，而是为成全法律时，我的回答是：基督徒不受割礼，恰恰因为割礼所预像的事已在基督内得以成全。割礼是除去我们肉性之预像，这已在基督的复活中得以成全，而圣洗圣事教导我们期待在自己的复活中实现。新生命的圣事并未完全废止，因为我们的肉身复活尚未来到；但此圣事已藉圣洗代替割礼而得以改善，因为如今基督的复活为我们提供了来世永生的模式，而从前却没有这样的事。同样，当你问为什么基督徒不守安息日，如果基督来不是为废除法律，而是为成全法律时，我的回答是：基督徒不守安息日，恰恰因为安息日所预像的事已在基督内得以成全。因为我们是在祂内得安息，祂说：\BibleQuote{凡劳苦和负重担的，你们都到我跟前来，我要使你们安息。你们背起我的轭，跟我学吧！因为我是良善心谦的；这样你们必要找到你们灵魂的安息。}\BibleRef{玛 11:28}\par

10. 当你问为什么基督徒不遵守法律所规定的食物区别，如果基督来不是为废除法律，而是为成全法律时，我回答说：基督徒不遵守这种区别，恰恰因为这样预像的事如今已在基督内得以成全。祂接纳进入祂的身体（即祂在圣人中预定了永生）的，在人类行为中没有与禁食动物的特征相对应之处。当你再次问为什么基督徒不献上宰杀动物的血肉作为祭品给天主，如果基督来不是为废除法律，而是为成全法律时，我回答说：如今这样预像的事已在基督奉献自己的体血中得以成全，基督徒再献那样的祭品就很不合适。当你问为什么基督徒不像犹太人那样守无酵节，如果基督来不是为废除法律，而是为成全法律时，我回答说：基督徒不守这个节，恰恰因为这样预像的事已在基督内得以成全，祂藉除去旧酵引导我们度新生活。\BibleRef{格前 5:7}当你问为什么基督徒不守逾越节羔羊的节日，如果基督来不是为废除法律，而是为成全法律时，我的回答是：他不守这个节，恰恰因为这样预像的事已在无玷羔羊基督的苦难中得以成全。当你问为什么基督徒不守法律所规定的新月节，如果基督来不是为废除法律，而是为成全法律时，我回答说：他不守这些节，恰恰因为这样预像的事已在基督内得以成全。因为新月节预像新受造物，宗徒说：\BibleQuote{所以谁若在基督耶稣内是新的受造物，旧的已成过去；看，一切都成了新的。}\BibleRef{格后 5:17}当你问为什么基督徒不遵守法律中各种不洁的洗礼，如果基督来不是为废除法律，而是为成全法律时，我回答说：他不遵守这些，恰恰因为它们是将来之事的预像，基督已予以成全。因为祂来是藉圣洗使我们与祂同埋葬于死亡，好使基督如何从死者中复活，我们也如何在新生活中行走。\BibleRef{罗 6:4}当你问为什么基督徒不守帐棚节，如果法律不是被废除而是被基督所成全时，我回答说：信徒是天主的帐幕，在他们因爱被联合并建筑在一起时，天主屈尊居住其中，所以基督徒不守这个节，恰恰因为这样预像的事如今已在基督内、在祂的教会中得以成全。\par

11. 我只是极其简短地提及这些事，以免全然忽略。这些主题分别而论，已写满了许多巨卷，证明这些礼仪都是基督的预像。由此可见，旧约中所有那些你认为不被基督徒遵守是因为基督废除了法律的事，实际上不被遵守正是因为基督成全了法律。这些礼仪本身的目的就是预像基督。如今基督既已来到，用来预像祂来临的事不再举行，非但不奇怪或荒谬，反而是完全正确合理的。那些预像基督来临的礼仪，如果未被基督的来临所成全，本应继续遵守；因此，我们现在不遵守这些礼仪，远非证明它们未被基督的来临所成全。任何宗教团体，无论其宗教真实或虚假，都不能没有某种圣事或可见的标记作为结合纽带。这些圣事的重要性不可夸大，只有嘲弄者才会轻视它们。因为如果虔诚要求它们，那么忽视它们必是不虔诚。\par

12. 诚然，不虔诚的人也可以领受虔诚的可见圣事，正如我们读到西蒙·玛古斯领受了圣洗。宗徒论到这样的人说：\BibleQuote{他们有虔诚的外貌，却否认了它的能力。}\BibleRef{弟后 3:5}虔诚的能力就是诫命的终极目标，即出自纯洁的心、善良的良心和真诚无伪的信德的爱。\BibleRef{弟前 1:5}所以伯多禄宗徒论到方舟的圣事（诺厄一家藉此从洪水中得救）时，说：\BibleQuote{与此预像相应的圣洗，现在也拯救你们。}唯恐他们满足于可见的圣事（藉此有了虔诚的外貌），并以放荡的生活否认其生活中的能力，他立刻加上：\BibleQuote{不是除去肉体的污秽，而是向天主许下纯洁的良心。}\BibleRef{伯前 3:21}\par

13. 如此，旧约的圣事，那些按法律举行的，是预表将要来临的基督；当基督藉祂的来临满全了它们，它们就被废除了，并且它们被废除是因为被满全了。因为基督来不是为废除，而是为满全。如今，信德的正道已彰显，天主的子女被召叫进入自由，那曾为属血肉和硬颈的人民所必需的奴役之轭已被除去，其他的圣事被建立，效能更大，使用更有益，施行更容易，数量更少。\par

14. 如果古时的义人，在他们时代的圣事中看到未来信德启示的应许，这应许甚至在当时他们的虔诚使他们在预言的朦胧光中辨别，并赖以生活，因为义人只能因信德而生活；\BibleRef{罗 1:17} 如果这些古时的义人准备忍受，正如许多人确实忍受了，为了那些预表未来事物的典型圣事而受尽一切考验和折磨；如果我们称赞三位青年和达尼尔，因为他们拒绝因国王餐桌上的食物而被玷污，出于对他们当时圣事的尊重；如果我们对玛加伯感到最强烈的钦佩，他们拒绝触碰基督徒合法使用的食物；\EditorialNote{玛加伯下第七章} 那么，我们当代的基督徒岂不更应该准备为基督的圣洗、为基督的圣体、为基督的神圣标记而忍受一切，因为这些是那先前圣事仅指向未来的事已实现的证明！因为仍应许给基督的身体——教会的事，既已清楚显明，而且在救主本身，身体的元首、天主与人之间的中保、那人基督耶稣内，已经实现了。永生不是藉着从死者中复活的应许吗？我们看见这在祂的肉身上实现了，经上说：圣言成了血肉，住在我们中间。\BibleRef{若 1:14} 过去的日子，信德是朦胧的，因为那些时代的圣人和义人都相信并盼望同样的事，所有这些圣事和仪式都指向未来；但现在我们有了信德的启示，这信德是人民被关闭在法律之下的；\BibleRef{迦 3:23} 如今在审判中应许给信徒的，已在祂的榜样中实现了，祂来不是为废除法律和先知，而是为满全它们。\par

15. 在圣经学者中有一个问题：基督受难和复活之前的信德，即古时义人藉启示得知或从预言中收集的信德，是否与现在基督已受苦并复活后的信德具有同样的效能；或者，天主羔羊之血的真正流出，正如祂自己所说，是为许多人而流，为得罪之赦，\BibleRef{玛 26:28} 是否以净化或增加纯洁的方式，对那些怀着信德期待基督死亡、但在它发生前已离世的人带来任何益处；事实上，基督的死亡是否延伸到死者，以达成他们的解放。在此讨论这个问题，或证明已确定的主题，将花费太长时间，且与我们当前的目的无关。\par

16. 与此同时，足以证明，针对\Person{福斯特}的无知吹毛求疵，他们大错特错，那些从标记和圣事的改变中推断，在预言的时期之仪式中所预表的事物与在福音仪式中宣告已实现的事物必定有分别的人；或者另一方面，那些认为既然事物相同，宣告其实现的圣事就不应与预告那实现的圣事有所不同的人。因为在语言中，动词的形式根据时态在字母和音节的数量上改变，正如\Emph{已完成}表示过去，\Emph{将完成}表示未来，那么为什么宣告基督死亡和复活已实现的象征，不应与预言其实现的象征有所不同呢？正如我们看见词语的形式和声音的差异：过去和未来，受苦和将受苦，复活和将复活。因为物质的象征不过是可见的言语，虽然是神圣的，却是可变和短暂的。因为天主是永恒的，圣洗的水，以及圣事中一切物质的东西，都是短暂的：甚至在祝圣时必须念出的\Quote{天主}这个词，也是一种瞬间消逝的声音。行为和声音消逝，但其效能保持不变，这样传递的神恩是永恒的。因此，说如果基督没有废除法律和先知，法律和先知的圣事就会继续在基督徒教会的会众中遵守，就等于说如果基督没有废除法律和先知，祂就仍会被预言为将要出生、受苦和复活；而事实上，已证明祂没有废除它们，而是满全了它们，因为祂的出生、受难和复活的预言，曾在这些古代圣事中预表，已经终止，而基督徒现在遵守的圣事包含宣告：祂已出生、已受苦、已复活。祂来不是为废除法律和先知，而是为满全它们，藉此满全，祂除掉了那些预告那已被证明现已实现之事的实现的事物。正如同样地，祂可以用\Quote{祂已出生、已受苦、已复活}这些表达来代替\Quote{祂将出生、将受苦、将复活}这些当时合适的表达，现在其他事已实现，所以被废除了。\par

17. 与词语上的这种改变相对应，在取代旧约圣事而设立新圣事的过程中，自然也发生了改变。对于最初作为犹太人来信的第一批基督徒，他们是逐渐被引导改变习俗并清晰认识真理的；宗徒准许他们保留他们出生并成长的世袭敬拜和信仰，并要求那些与他们打交道的人对这种不愿接受新习俗的态度予以体谅。因此，当宗徒和弟茂德（其母为犹太人，其父为希腊人）一起前往这类人群中间时，宗徒给弟茂德行了割损礼；宗徒自己也将他的做法适应于他们，并非出于虚伪，而是出于明智的目的。因为对于那些生来就在这些习俗中长大的人，这些做法是无害的，尽管它们不再需要预象将来的事。如果把它们当作有害的加以谴责，对于那些它们本应继续存留到其时代的人，反而会造成更大的伤害。基督来应验所有这些预言，祂发现这些人已在自己的宗教中受过训练。但对于那些未受这种训练、却从未受割损的对墙被带到基督这基石面前的人，他们没有义务采纳犹太习俗。的确，如果他们像弟茂德那样，选择使自己适应那些仍依恋旧圣事的受割损者的看法，他们可以自由这样做。但如果他们认为自己的希望和救恩取决于这些法律行为，就要警告他们这是致命的危险。因此宗徒说：\BibleQuote{请看，我保禄告诉你们：如果你们受割损，基督对你们毫无益处。}\BibleRef{迦 5:2}这就是说，如果他们受割损，正如他们打算做的那样，是顺从某些腐化的教师——这些教师告诉他们没有这些法律行为就不能得救。因为，当外邦人主要通过宗徒保禄的宣讲而来到基督的信仰中时，他们应适当地前来，而不必背负犹太教规的重担；因为那些成年人因害怕不习惯的礼仪，尤其是割损礼，而被信仰所吓退；而且，如果这些从小未受这些教规训练的人以通常的方式成为归依者，那就意味着基督的来临仍需要被预言为未来的事件。因此，当外邦人没有这些礼仪就被接纳时，那些信了的受割损的人不明白为什么外邦人不必采用他们的习俗，也不明白为什么他们自己仍然被允许保持这些习俗，便开始用属世的争论扰乱教会，因为外邦人被接纳为天主的子民，而没有以通常的方式通过割损和其他法律教规成为归依者。还有些归化的外邦人因害怕同住的犹太人而执意遵守这些礼仪。宗徒保禄常写信反对这些外邦人，当伯多禄被他们的虚伪所裹挟时，保禄用弟兄般的责备纠正了他。\BibleRef{迦 2:14}后来，宗徒们在会议中议定，这些法律行为对外邦人没有约束力，\BibleRef{宗 15:6}一些受割损的基督徒感到不快，因为他们不明白，这些规条只对在信仰启示之前已受这些训练的人允许，以便使那些在预言应验之前从事先知生活的人结束那种生活，以免通过强制放弃，似乎是被谴责而非终结；而将这些加于外邦人，要么意味着它们不是为预象基督而设立，要么意味着基督仍需要被预象。在古代，天主的子民在基督来成全法律和先知之前，必须遵守所有预象基督的这一切。对于明白遵守意义的人是自由，对于不明白的人则是束缚。但在后来的时期，那些相信基督已经来临、受苦、复活的人，若这信仰发现他们已受这些圣事的训练，则既不要求他们遵守，也不禁止他们遵守；而对于那些不受习俗束缚或任何必要去适应他人做法的人，则有禁令，以便显明这些事是为预象基督而设立的，并且在祂来临之后应该停止，因为恩许已经应验。一些受割损的信徒不明白这一点，对圣神通过宗徒所实行的宽容安排感到不满，并顽固地坚持要求外邦人成为犹太人。这就是\Person{福斯图斯}以\Quote{西马基派}或\Quote{纳匝肋派}之名所提及的人。他们的人数现在很少，但这个派别仍然存续。\par

18. 因此，摩尼教徒没有理由说基督来是为了毁坏律法和先知，而非成全它们，因为他们说基督徒不遵守其中所吩咐的；其实，基督徒不遵守的只是那些预表基督的事，这些事之所以不被遵守，是因为它们的成全在基督内，而已经成全的事不再被预表；那些典型性的圣事在那些受过此类训练、并已相信基督为它们的成全的人的时代，已经恰当地结束了。难道基督徒不遵守圣经的诫命\BibleQuote{以色列，你要听：上主你的天主是唯一的天主}，\BibleQuote{你不可为你自己制作雕像}等等吗？难道基督徒不遵守\BibleQuote{你不可妄呼上主你天主的名}的诫命吗？难道基督徒不按照诫命守安息日，即使是在真正安息的意义上？难道基督徒不按照诫命孝敬父母吗？难道基督徒不戒绝奸淫、杀人、偷盗、作假见证、贪恋邻人的妻子和财产吗？——这一切都写在律法中。这些道德诫命有别于典型性的圣事：前者藉天主恩宠的帮助得以成全，后者藉其所应许之事的实现而得以成全。两者都在基督内得以成全，祂从未停止赐予这恩宠，这恩宠如今也在祂内启示出来，祂如今显明了祂昔日所应许之事的实现；因为\BibleQuote{法律是藉梅瑟传授的，恩宠和真理却是由耶稣基督而来的}\BibleRef{若 1:17}。再者，这些关乎保持良好良心的事，是在那藉爱德工作的信德中得以成全\BibleRef{迦 5:6}；而未来之事的预表在实现时就过去了。然而，即使是预表，也不是被毁坏，而是被成全；因为基督揭示预表所象征的事，并没有证明它们是虚妄或虚假的。\par

19. 因此，福斯图斯错误地认为主耶稣成全了某些在梅瑟律法之前义人的诫命，例如\BibleQuote{不可杀人}，基督非但没有反对，反而通过禁止发怒和辱骂来确认了它；而祂毁坏了一些似乎是希伯来律法特有的东西，例如\BibleQuote{以眼还眼，以牙还牙}，基督似乎否定了而非确认了它，当祂说：\BibleQuote{我却对你们说：不要抵抗恶人；而且，若有人掌击你的右颊，你把另一面也转给他}\BibleRef{玛 5:38}等等。\par

20. 首先，请让我问我们的对手：这些古代的义人，如福斯图斯特别提到的哈诺客和舍特，以及任何其他在梅瑟之前，或者甚至，如果你愿意，在亚巴郎之前的人，他们是否曾无缘无故地向自己的弟兄发怒，或对自己的弟兄说\Quote{傻子}？如果没有，他们为何可能没有教导这些事，如同他们宣讲这些事一样？如果他们教导了这些事，那么基督又怎能说是成全了他们的义德或他们的教导，超过梅瑟的教导，通过加上\Quote{我却对你们说：凡向自己弟兄发怒的，或说\InnerQuote{拉卡}的，或说\InnerQuote{傻子}的，就要受裁判，或受公议会，或受地狱之火的惩罚}呢？既然这些人自己就做了这些事，并吩咐别人也这样做。是否可以说他们不知道义人有责任抑制自己的怒气，不因愤怒的辱骂而激怒弟兄？或者，他们知道这点，却无法照此行事？那样的话，他们就该受地狱的惩罚，就不可能成为义人。但没有人敢说，在他们的义德中，对义务的无知和自制力的缺乏到了使他们该受地狱惩罚的地步。那么，基督怎能说是成全了这些人赖以生活的律法，通过加上那些没有它们他们就毫无义德可言的东西呢？是否可以说，急躁的脾气和恶语只在基督时代以后才是罪，而以前，这种内心和言语的素质是允许的——正如我们发现某些制度因时代而异，以至于在一个时代合适的事在另一个时代不合适，反之亦然？你不会愚蠢到作出这样的断言。但即使你这样做，回答将是：按照这种想法，基督来不是要成全旧律法中的缺陷，而是要制定一个先前不存在的律法；如果古代的义人向自己的弟兄说\Quote{傻子}不是罪，而基督却宣告这样做是极其有罪的，以至于如此行的人要受地狱的惩罚。所以，你未能找到任何可以称之为基督通过这些添加弥补了其缺陷的律法。\par

21. 是否可以说，在早期时代，律法在\BibleQuote{不可奸淫}方面是不完全的，直到主添加了\BibleQuote{不可贪恋妇女}而使之完成？这正是你引用经文时所暗示的：\BibleQuote{你们听见有话说：\InnerQuote{不可奸淫}，只是我告诉你们：连贪恋也不可。}你说：\Quote{这就是成全。}但让我们照福音的原话，不加任何你的修改，看看你给那些古代义人的描述是什么。原话是：\BibleQuote{你们听见有话说：\InnerQuote{不可奸淫。}我却对你们说：凡注视妇女，有意贪恋她的，他已在心里奸淫了她。}\BibleRef{玛 5:27}那么，按你的意见，哈诺客和舍特等人心里犯了奸淫；他们的心要么不是天主的宫殿，要么他们在天主的宫殿里犯了奸淫。但如果你不敢这样说，你怎能说基督来临时成全了律法，而这律法在那些人的时代就已经是完全的呢？\par

22. 关于不发誓，你们还称基督为此完善了赐给古时这些义人的法律。我不能确信他们从未发誓，因为我们发现宗徒\Person{保禄}曾发誓。在你们这里，发誓仍是常事，因为你们凭着光发誓，你们爱光如同苍蝇爱光；至于那照亮每个人的心灵之光，就是来到世上的真光，你们却一无所知。你们也凭着你们的师傅摩尼发誓，他本名玛内斯。因为玛内斯这名字似乎与希腊语中\Quote{疯狂}一词相关，你们通过添加后缀改变了它，这只会使事情更糟，因为它有了\Quote{倾泻疯狂}的新含义。你们教派中一人告诉我，摩尼这个名字本意是从希腊词\Quote{灌注吗哪}而来；因为\OriginalTerm{\Quote{灌注}}{\GreekText{χέειν}}的意思是灌注。但事实上，你们只是更加强调了疯狂的概念。因为添加了两个音节，而忘了在词首插入另一个字母，你们得到的是摩尼而非玛尼，这一定意味着他在他那冗长无益的讲论中倾泻疯狂。再者，你们常凭着护慰者发誓——不是基督应许并派遣给祂门徒的护慰者，而是这位疯狂倾泻者本人。既然你们经常发誓，我想知道你们是以什么意义认为基督履行了律法的这一部分，就是你们提及属于最早时期的那一部分。你们对宗徒的誓言又作何解释？因为你们的权威对你们自己而言分量不大，更不用说对我或任何其他人了。因此很明显，基督的话：\Quote{我来不是要废除法律，而是要成全它}并不具有你们所赋予的意义。基督在这些话中根本没有提及祂对所引用的古语的评论，祂的演讲是对这些古语的解释，而非成全。\par

23. 因此，关于杀人，原来仅指夺去肉体生命，使人丧命；主解释说，对弟兄的任何不义意图都是一种杀人。\Person{若望}也说：\BibleQuote{凡恨自己弟兄的，便是杀人的。}\BibleRef{若一 3:15}。而且，因为以前以为奸淫仅指与妇女的非法交合，师傅表明祂所描述的贪欲也是奸淫。再者，因为发虚誓是重罪，而不发誓或说实话却绝不是罪，主想通过完全不发誓来保护我们不偏离真理，而不是靠诚实发誓的习惯使我们陷入发虚誓的危险。因为从不发誓的人比习惯于诚实发誓的人更容易犯虚誓。因此，在宗徒的记载讲论中，他从未使用过誓言，以免因习惯发誓而无意中犯虚誓。另一方面，在他的著作中，他有更多闲暇和谨慎的机会，我们发现他在几处使用了誓言，以教导我们说实话发誓没有罪，但鉴于人性的软弱，完全不发誓言最能保护我们免于虚誓。这些考虑也将会说明，\Person{浮斯图斯}认为属于\Person{梅瑟}的独特事物，并未被基督废除，如他所说那样。\par

24. 试以古人的这一句话为例：\Quote{你应爱你的近人，恨你的仇人}，\Person{福斯图斯}怎能证明这是\Person{梅瑟}独有的呢？\Person{宗徒保禄}不是也曾说到有些人是天主所憎恨的吗？\BibleRef{罗 1:30} 事实上，主也在这句话上命令我们效法天主，祂说：\BibleQuote{好使你们成为你们在天之父的子女，因为祂使太阳上升，光照恶人，也光照善人；降雨给义人，也给不义的人。}\BibleRef{玛 5:45} 在某种意义上，我们必须恨我们的仇人，效法天主——保禄说有些人被天主所恨；同时，我们也必须爱我们的仇人，效法天主——祂使太阳上升光照恶人和善人，降雨给义人和不义的人。如果我们明白这一点，就会发现主在向那些没有正确理解\Quote{你应恨你的仇人}这句话的人解释时，利用这句话来教导他们应当爱他们的仇人，这对他们而言是一个新的概念。要在此说明这两者的一致性会花费太多篇幅。但当摩尼教徒毫无例外地谴责\Quote{你应恨你的仇人}这条诫命时，很容易用一个问题来反驳他们：他们的神是否爱黑暗种族？或者，如果我们现在应当爱我们的仇人，因为他们有一部分善，难道我们不应也恨他们，因为他们有一部分恶吗？即使这样，看起来古时的话\Quote{你应恨你的仇人}与福音的\BibleQuote{爱你们的仇人}之间并不矛盾。因为每个恶人，就其恶而言，应当被恨；同时，作为人，他应当被爱。我们在他身上合理憎恨的恶行应当被谴责，以便通过消除它，我们在他身上合理喜爱的人性能够得到修正。这正是我们所坚持的原则：我们应当恨我们的仇人，因为在他身上的恶，即他的邪恶；同时我们也爱我们的仇人，因为在他身上的善，即他作为社会性和理性存有的本性。我们与摩尼教徒的区别在于，我们证明人是恶的，不是由于本性（无论是他自己的还是其他的），而是由于他自己的意志；而他们认为人是恶的，是由于黑暗种族的本性，按照他们的说法，这位天主在黑暗种族完整时对它感到恐惧，并且还被它部分地征服了，以至于他不能完全解脱。那么，主的意图是要纠正那些人，他们由于不理解而只知道古时的话\Quote{你应恨你的仇人}，就去恨他们的同胞，而不是只恨他们的邪恶；为此，祂说：\BibleQuote{爱你们的仇人。}祂并没有废除法律中关于恨仇敌的记载——祂说过：\BibleQuote{我来不是为废除法律，而是为成全}——祂要我们通过爱仇人的责任，学会对同一个人如何既因他的罪而恨他，又因他的本性而爱他。指望我们顽固的对手理解这一点未免太难了。但我们可以让他们哑口无言，表明他们由于不理性的反对而谴责了他们自己的神，因为不能说他们的神爱黑暗种族；因此，在命令每个人爱仇人时，他们不能引用祂的榜样。黑暗种族似乎比摩尼教徒的神更有爱仇人之心。故事是这样的：黑暗种族觊觎与他们领土接壤的光明领域，并出于占有它的欲望，制定了入侵的计划。渴望真善和真福并不是罪。因为主说：\BibleQuote{天国是以猛力夺取的，以猛力夺取的人就攫取了它。}\BibleRef{玛 11:12} 那么这个虚构的黑暗种族，由于它的美丽和诱人的外貌，希望用强力夺取他们所渴望的善。但天主没有回报那些渴望拥有祂的人的愛，反而恨他们，试图消灭他们。因此，如果恶人爱善人，是出于占有它的欲望，而善人恨恶人，是出于害怕被玷污，我问摩尼教徒，这两者中哪一个遵守了主的诫命\BibleQuote{爱你们的仇人}？如果你们坚持认为这两条诫命彼此对立，那么就会得出一个结论：你们的神服从了梅瑟法律中所记载的\Quote{你应恨你的仇人}；而黑暗种族却服从了福音中所记载的\BibleQuote{爱你们的仇人。}然而，你们从未成功地解释白天飞行的苍蝇和夜间飞行的蛾子之间的区别；因为按照你们的说法，两者都属于黑暗种族。为什么一种（苍蝇）违反本性喜爱光明，而另一种（蛾子）却逃避光明，偏爱它们所从出的黑暗？奇怪，肮脏的下水道竟然能繁殖出比黑暗橱柜更洁净的物种！\par

25. 再者，古时所说的\BibleQuote{以眼还眼，以牙还牙}与主所说的\BibleQuote{我却对你们说：不要抵抗恶者；若有人打你的右颊，你把另一面也转给他}等等并无对立。旧诫命与新诫命都是为抑制仇恨的猛烈，约束怒气的冲动。因为谁会自愿满足于与伤害相等的报复？我们岂不是看见人受了一点轻伤，就渴求杀戮、渴望流血，仿佛永远不能让敌人受够痛苦？若有人挨了一记耳光，他岂不是召唤攻击他的人，好让他在法庭上被判罪？或者，他若宁愿还手，岂不是用手脚，或者若能得到武器，就用武器扑向那人？为了约束这种因过度而不义的报复，法律规定了赔偿原则，使刑罚与所受的伤害相称。因此，\BibleQuote{以眼还眼，以牙还牙}这条诫命，与其说是点燃已熄之火的火把，不如说是防止已燃之火蔓延的盖子。因为攻击者对受伤害者负有正义的报复之责；所以当我们宽恕时，我们放弃了本可正义要求的东西。因此，在主祷文中，我们被教导要宽恕别人的债，好让天主宽恕我们的债。要求偿还债务并非不义，但宽恕债务则是仁慈。然而，正如在发誓时，一个发誓的人，即使是真实的，也有发虚誓的危险，而从不发誓的人则没有这种危险；虽然真实发誓不是罪，但不发誓则更远离罪；因此，不发誓的命令是防止发虚誓的保障：同样，既然渴望以不正义的过度报复是罪，而在正义范围内渴望报复则不是罪，那么根本不想报复的人就更远离不正义报复的罪。要求超过应得的东西是罪，而要求债务则不是罪。防止不正义要求之罪的最好保障是什么都不要求，尤其考虑到必须向那无所负欠者偿还债务的危险。因此，我这样解释这段经文：古时曾说：你们不可行不正义的报复；但我说：你们根本不要报复：这就是成全。福斯特斯在引用了\BibleQuote{你们不可虚誓}之后，加上\BibleQuote{我却对你们说：什么誓都不可起}，并说：这就是成全。我也许可以用同样的说法，如果认为基督加上这些话是补充法律的缺陷，而不是不如说法律防止不正义报复的意图，通过根本不报复得到最好的保障，正如防止发虚誓的意图通过根本不发誓得到最好的保障一样。因为若\BibleQuote{以眼还眼}与\BibleQuote{若有人打你的脸，把另一面也转给他}相对立，那么\BibleQuote{你要向主履行你的誓愿}与\BibleQuote{什么誓都不可起}之间岂不也有同样的对立？\BibleRef{玛 5:33} 如果福斯特斯认为在一种情况中不是毁灭，而是成全，那么他也应认为另一种情况也是如此。因为若\BibleQuote{不可起誓}是\BibleQuote{真实起誓}的成全，那么为何\BibleQuote{不可报复}不是\BibleQuote{正义报复}的成全呢？\par

所以，按照我的解释，两种情况都是防止罪的保障，或防止发虚誓，或防止不正义的报复；不过，关于放弃报复的权利，还有一个额外的考虑：借着宽恕这些债务，我们将获得我们债务的宽恕。旧诫命对于一个固执的百姓是必要的，教导他们不要过度要求。这样，当那渴求无节制报复的怒气被抑制时，任何有此心志的人便有闲暇思考借着主取消自己债务的可取性，并因此受引导宽恕他同仆的债务。\par

26. 再者，我们查考时就会发现，主关于不可休妻的诫命，与古人所说的\Quote{谁若休妻，就该给她休书}，两者之间并无矛盾。主阐明了法律的意图，即凡休妻者均须出具休书。不可休妻的诫命，与说人可以随意休妻，是相反的；法律并未如此说。相反，为了防止妻子被休，法律规定了这个中间步骤，以期藉由书写休书来抑制休妻的急切心情，并使丈夫有时间思考休妻之恶；尤其因为，据说在希伯来人中，除经师以外，任何人都不得书写希伯来文：因为经师声称拥有卓越智慧；若他们是正直虔诚的人，他们的身份足以使他们有理由提出这一要求。因此，要求男人在休妻时给她休书，其用意是使他不得不去求助那些可以期望从中获得法律谨慎解释及适当劝诫以免分离的人。由于没有其他方法获得休书，男人便须服从他们的指导，并允许他们努力在他与妻子之间恢复和平与和谐。在仇恨无法克服或抑制的情况下，休书当然会被写出。当明智的劝告未能恢复丈夫心中应有的感情和爱慕时，妻子就有理由被休。若妻子不被爱，她就应被休。而为了她不被休，丈夫有责任爱她。如今，虽然不能强迫人违背意愿去爱，但人可以通过劝告和说服受到影响。这是经师作为明智正直之人的职责；法律规定凡发生争吵时，丈夫都必须去找他，以便让他书写休书，这便给了经师这个机会。除非是顽固厌恶以致无法和解的情况，否则没有一个善良明智的人会书写休书。但是，根据你们不敬神的观念，休妻毫无意义；因为在你们看来，婚姻是一种犯罪纵欲。\Quote{婚姻}一词表明男子娶妻是为了使她成为母亲，这在你们看来是一种邪恶。按你们的说法，这意味着你们的神的一部分被黑暗族类击败、俘虏，并束缚在血肉的桎梏中。\par

27. 但是，为了解释当下之点：如果基督在引用古语时附加了\Quote{但我告诉你们}这句话，祂的补充既没有成全远古的法律，也没有用相反诫命废除给梅瑟的法律，反而在祂所作的一切引用中对希伯来法律表现出如此尊重，以致祂自己的话主要是对法律所陈述不清之处的解释，或是确保法律所意图之设计得以实现的手段，那么，从\Quote{我来不是为废除法律，而是为成全}这句话，我们不应理解为基督以祂的诫命填补了法律所欠缺的；而是指，由于人的骄傲和不服从，字面命令所未能成就的，藉着恩宠在那些悔改和谦卑的人身上得以完成。成全不在于添加话语，而在于服从的行动。所以宗徒说：\Quote{信德以爱德行事}\BibleRef{迦 5:6}；又说：\Quote{爱别人的，就满全了法律}\BibleRef{罗 13:8}。这爱德，藉之也能满全法律之义，由基督在祂来临时，通过祂按应许派遣的圣神，赋予其意义；因此祂说：\Quote{我来不是为废除法律，而是为成全}。这是新约，在这新约中，天国的应许是赐予这爱德的；这爱德在旧约中已被预表，适合那时期的经济。所以基督又说：\Quote{我给你们一条新诫命，你们要彼此相爱}\BibleRef{若 13:34}。\par

28. 因此，我们在旧约中发现了基督用\Quote{但我告诉你们}这句话所引入的全部或几乎全部劝言与规诫。论愤怒，经上记载：\BibleQuote{我的眼睛因愤怒而困扰}；又说：\BibleQuote{征服自己愤怒的人，比攻取一座城的人更优秀}。\BibleRef{箴 16:32} 论粗暴言语：\BibleQuote{鞭子的抽打造成伤口，舌头的抽打却折断骨头。} 论心中的奸淫：\BibleQuote{不可贪恋你邻人的妻子。}\BibleRef{出 20:17} 经上不是\BibleQuote{不可行奸淫}，而是\BibleQuote{不可贪恋}。宗徒引用此句时说：\Quote{若不是律法说\InnerQuote{不可贪恋}，我就不知何为贪欲。}\BibleRef{罗 7:7} 关于忍耐、不抵抗，一个人受到称赞，他\BibleQuote{把自己的面颊给打他的人，受尽羞辱。}\BibleRef{哀 3:30} 关于爱仇敌，经上说：\BibleQuote{如果你的仇敌饿了，就给他吃；渴了，就给他喝。}\BibleRef{箴 25:21} 宗徒也引用了这句。\BibleRef{罗 12:20} 在圣咏中也有话说：\BibleQuote{我在那仇恨和平的人中间是一个缔造和平的人。} 此外，关于我们效法天主，不寻报复，甚至爱恶人，有一段经文完全描述了天主在这方面的属性；因为经上写道：\BibleQuote{唯独祢永远拥有大能，谁能抵抗祢手臂的力量？整个世界在祢面前，如同天秤上的一粒微尘，又如同清晨降在地上的朝露。然而，祢怜悯众生，因为祢无所不能，祢对人所犯的罪过视而不见，是为使他们悔改。因为祢爱万有，不憎恨祢所造的；因为如果祢憎恨什么，绝不会造它。怎能存留，若不是祢愿意？怎能保存，除非由祢所召？但祢宽恕一切，因为万物都属于祢，主，祢是爱灵魂的主。因为祢的至善圣神充满万有，故此祢逐渐责罚那些犯罪的人，并以提醒他们何处犯了罪来警告他们，使他们认识自己的邪恶，从而信靠祢，主。} 基督劝勉我们要效法天主这种长久忍耐的良善，祂使太阳升起，光照恶人和善人，降雨给义人和不义的人；这样我们就不必急于报复，反而要善待恨我们的人，从而得以成全，如同我们在天之父是成全的。\BibleRef{玛 5:44} 从这些古书的另一段经文我们得知：若不要求应得的报复，我们便获得自己罪过的赦免；若不宽恕别人的债，当我们祈求赦免自己的债时，就有被拒绝的危险：\BibleQuote{报复的人，必遭主的报复，主必牢牢记住他的罪过。宽恕你邻人对你造成的伤害，这样当你祈祷时，你的罪过也会得到赦免。一人怀恨另一人，他能向主求得宽恕吗？他对与自己相似的人毫无怜悯，而能祈求自己的罪过获得宽恕吗？如果这个血肉之躯还怀恨在心，并向主求恩惠，谁能为他的罪过求情呢？}\BibleRef{德 28:1}\par

29. 关于不可休妻，无需引用旧约的其他段落，只需引用主在回答犹太人就此问题的询问时最恰当的引文。当他们问：男人是否可以因任何理由休妻，主回答说：\BibleQuote{你们没有念过：那起初造人的，是造男造女，并且说：\InnerQuote{为此，人要离开父母，与妻子结合，两人成为一体}吗？这样，他们不再是两个，而是一体了。所以，天主所结合的，人不可拆散。}\BibleRef{玛 19:4} 在这里，那些犹太人以为他们按照\Person{梅瑟}律法的意思休妻，却从\Person{梅瑟}的书卷中看到妻子不应被休。而且，我们在此从基督自己的宣告得知：天主造了男人和女人并使他们结合；因此，若摩尼教徒否认这一点，他们就既反对基督的福音，也反对\Person{梅瑟}的著作。假设他们的教义是真的，即魔鬼造了男人和女人并使他们结合，那么我们就看到\Person{福斯特}的魔鬼般的狡诈：他指责\Person{梅瑟}允许休妻书来解散婚姻，却称赞基督以福音的诫命加强了结合。相反，\Person{福斯特}应与其愚蠢不敬的见解一致，称赞\Person{梅瑟}分离了魔鬼所造所结合的事物，而责备基督批准了魔鬼的作品的结合。言归正传，让我们听一听良善的导师解释：\Person{梅瑟}既然在最初男女的结合中写到婚姻的贞洁是如此神圣不可侵犯，为何后来允许人民休妻。因为当犹太人回答说：\BibleQuote{那么\Person{梅瑟}为什么吩咐给休书就可以休她呢？} 基督对他们说：\BibleQuote{\Person{梅瑟}因为你们的心硬，才容许你们休妻。}\BibleRef{玛 19:7} 这段话我们已经解释过。这心硬确实非常严重，以至于即使通过休书提供了让明智正直的人劝告恢复夫妻之爱的机会，也无法促使他们接受。然后主引用了同一部律法，以显示对善者所命令的和对心硬者所容许的；因为从关于男女结合所写的，祂证明妻子不可休，并指出结合的神圣权威；并从同一圣经指出，休书之所以要给予，是因为心硬，这心硬可能被克服，也可能不被克服。\par

30. 既然\Person{法乌斯都}试图证明与希伯来人的古书相悖的这些主的卓越训诫，实际上都出现在这些书中，那么主来不是毁灭法律，而是成全法律的唯一含义就是：除了预表性的实现（这些预表因实际成就而被搁置）之外，那法律中圣洁、公义、良善的训诫也在我们身上得以成全，不是借着命令并借着违命的附加罪责增加骄傲者过犯的陈旧文字，而是借着帮助我们且谦卑顺服的圣神的新性，借着使我们获得自由的救恩恩宠。因为，虽然所有这些崇高的训诫都出现在古书中，但它们所指的目的并未在那里启示出来；虽然预见到这启示的圣人们按照它生活，或按照时代将它隐藏在预言中，或亲自发现这样隐藏的真理。\par

31. 经过仔细考察，我倾向于怀疑主常使用的\Quote{天国}一词是否能在这类书中找到。的确，经上记着说：\BibleQuote{爱慕智慧，好使你永远为王。}\BibleRef{智 6:22} 如果永生没有在旧约中清楚显明，主就不会甚至对不信的犹太人说：\BibleQuote{你们查考圣经，以为其中有永生；正是给我作证的。}\BibleRef{若 5:39} 与此相同的是圣咏作者的话：\BibleQuote{我必不死，而活下去，并要宣扬主的作为。} 又说：\BibleQuote{求你照亮我的眼睛，免得我沉睡至死。} 再者，我们读到：\BibleQuote{义人的灵魂在主手中，痛苦不能触及他们。} 紧接着：\BibleQuote{他们处于平安中；即使他们受了人的折磨，他们的希望也充满了不朽；经过少许患难，他们要享受许多赏报。}\BibleRef{智 3:1} 又在另一处：\BibleQuote{义人永远活着，他们的赏报在主那里，他们的关切在至高者那里；因此他们要从主手中领受荣耀的王国和华美的冠冕。}\BibleRef{智 5:16} 这些以及许多类似关于永生的、或显或隐的宣告，都在这些书卷中。连身体的复活也被先知们提及。因此，法利塞人激烈反对不信复活的撒杜塞人。这不仅从正典的\WorkTitle{宗徒大事录}得知（摩尼教徒因它讲述主所应许的护慰者的来临而拒绝它），也从福音得知，当撒杜塞人问及主关于那个嫁过七个兄弟（一个接一个死去）的妇人，她在复活时是谁的妻子。\BibleRef{玛 22:23} 那么，关于永生和死人复活，在这部圣经中可以找到许多见证。但我没有在那里找到\Quote{天国}一词。这个说法确切属于新约的启示，因为在复活时，我们尘世的身体将由\Person{保禄}充分描述的改变成为属神的身体，如此属天，以致我们可以拥有天国。并且这个表述是为那一位保留的，祂作为君王治理、作为司铎圣化祂的信众，祂的来临被旧约的一切象征——其谱系、预表性的行为和言语、祭献、礼仪和庆节，以及一切预言性的话语、事件和形象——所宣告。祂满溢恩宠和真理，以祂的恩宠帮助我们遵守训诫，以祂的真理确保应许的实现。祂来不是毁灭法律，而是成全法律。\par

\SourceRecord{Translated by Richard Stothert.；Nicene and Post-Nicene Fathers, First Series；Vol. 4.；Edited by Philip Schaff.；Buffalo, NY: Christian Literature Publishing Co.,；1887.；<http://www.newadvent.org/fathers/140619.htm>.}

% structure: p0001 source=h1 target=section reason=page-title

\section{\WorkTitle{驳福斯特，卷二十}}

\Emph{福斯特反驳了太阳崇拜的指控，并坚持认为，虽然摩尼教徒相信天主的权能居于太阳、智慧居于月亮，但他们依然敬拜唯一的神，即圣父、圣子、圣神。他们不是外邦人的分裂，也不是一个宗派。奥古斯丁强调多神论的指控，并对摩尼教与异教神话进行了细致的比较。}\par

1. \Emph{福斯特}说：你们问为什么我们崇拜太阳，如果我们是一个宗派或独立的宗教，而非异教徒，或只是外邦人的一个分裂？因此，最好探究此事，看看异教徒之名究竟更适合你们还是我们。或许，当我以友好的方式简明陈述我的信仰时，看上去就像是在为此道歉，好像我羞于礼敬神圣的光体，愿天主禁止。你们可以随意理解；但如果我能至少让一些人明白，我们的宗教与异教徒毫无共同之处，我将不后悔我所做的。\par

2. 所以，我们敬拜唯一的神，具有三重称谓：全能天主圣父、祂的子基督，以及圣神。虽然三者是同一的，但我们相信圣父适当地居于最高或首要的光中，保禄称之为\Quote{不可接近的光}\BibleRef{弟前 6:16}，圣子则居于祂的第二或可见之光。既然圣子本身是双重的，按照宗徒所说，基督是天主的德能和天主的智慧\BibleRef{格前 1:24}，我们相信祂的德能居于太阳，祂的智慧居于月亮。我们也相信圣神，第三位尊威，在全部大气圈中拥有祂的宝座和家园。藉着祂的影响和灵性倾注，大地孕育并生出可朽的耶稣，祂悬挂在每一棵树上，是人的生命和救恩。虽然你们如此激烈地反对这些教义，但你们的宗教在赋予饼和酒以神圣性方面与我们相似，而我们则赋予一切事物神圣性。这就是我们的信仰，如果你们愿意，将有机会听到更多。同时，有一点颇有力量：你们或任何被问及他的神住在哪里的人，都会说祂住在光中；所以，支持我敬拜的见证几乎是普遍的。\par

3. 至于你们称我们为外邦人的分裂，而非宗派，我想\Quote{分裂}一词适用于那些与其他人有相同教义和敬拜，只是选择单独聚会的人。而\Quote{宗派}一词则适用于那些教义与他人完全不同，并且建立了独特敬拜方式的人。如果这些词是此意，那么首先，我们的教义和敬拜与异教徒毫无相似之处。我们稍后看看你们是否相似。异教徒的教义是：一切善与恶、卑劣与荣耀、衰败与不朽、可变与不变、物质与神圣的事物，只有一个本原。与此相反，我的信仰是：天主是一切善的事物的本原，而\OriginalTerm{原质}{Hyle} [质料]则是相反事物的本原。\OriginalTerm{原质}{Hyle}是我们的神学大师给邪恶的本原或本性起的名字。因此，异教徒认为应当用祭坛、神龛、雕像、祭品和香料来敬拜神。而我的做法与他们完全不同：因为我视自己为天主的理性殿堂，如果我有幸如此；我认为基督祂的儿子是祂活生生的尊威的活形象；我持守一个良好修养的心灵是真正的祭坛，纯洁而简单的祈祷是敬拜神和祭献的正确方式。这算是异教徒的一个分裂吗？\par

4. 至于对全能天主的敬拜，你们可能会称我们为犹太人的一个分裂，因为所有犹太人都大胆地声称这种敬拜，若不是我们在敬拜方式上的差异的话，尽管有疑问犹太人是否真的敬拜全能者。但我所提到的教义在异教徒对太阳的崇拜和犹太人对全能者的崇拜中都是共同的。甚至与你们相比，我们也不恰当地称为一个分裂，虽然我们承认基督并敬拜祂；因为我们的敬拜和教义与你们不同。在一个分裂中，与原始相比很少或没有变化；例如，你们在从外邦人的分裂中，带来了单一本原的教义，因为你们相信万物出自天主。你们将祭献变为爱筵，将偶像变为殉道者，你们向他们祈祷如同他们向偶像祈祷。你们用酒食安抚亡灵的阴魂。你们遵守与异教徒相同的节日，例如朔日和至日。在生活方式上，你们没有改变。显然你们只是一个分裂；因为与原始的唯一区别是你们单独聚会。在此你们效法了犹太人，他们脱离了外邦人，但仅在无偶像上不同。因为他们使用圣殿、祭献、祭坛、司铎和整套仪式，与异教徒相同，只是更加迷信。像异教徒一样，他们相信单一本原；因此你们和犹太人都是外邦人的分裂，因为你们有同样的信仰，几乎同样的敬拜，而你们自称宗派只是因为你们单独聚会。事实上，只有两个宗派：异教徒和我们。我们和异教徒在信仰上如同真理与谬误、白昼与黑夜、贫穷与财富、健康与疾病那样对立。而你们，无论与真理还是与谬误相比，都不是一个宗派。你们只是一个分裂，并且不是真理的分裂，而是谬误的分裂。\par

5. \Emph{\Person{奥斯定}}回答说：可憎的无知与狡诈的混合！为什么你们要把无人会用的论点放在你们对手的口中？我们并不称你们为异教徒，也不是异教徒的分裂派；而是说你们在崇拜多神方面与他们相似。但你们比异教徒更糟糕，因为他们崇拜存在之物，尽管它们不该被崇拜：因为偶像存在，尽管对于救恩它们是无用的。所以，用祈祷来崇拜一棵树，而不以耕作来改善它，并不是崇拜虚无，而是以错误的方式崇拜。当宗徒说\BibleQuote{外邦人所祭献的，是祭献给魔鬼，而不是给天主}\BibleRef{格前 10:20}时，他表示这些魔鬼是存在的，祭献是献给他们的，他不愿我们与他们有份。同样，天和地、海洋和空气、太阳和月亮以及其他天体，都是具有感性存在的对象。当异教徒将它们作为神祇，或作为一位伟大之神的部分来崇拜时（因为他们中有些人将宇宙与至尊神等同），他们崇拜的是存在之物。在与异教徒辩论时，我们并不否认这些事物的存在，而是说它们不该被崇拜；我们推荐崇拜所有这些事物的无形造物主，只有在他内，人才能找到所有人都承认他所渴望的幸福。相反，对于崇拜那无形无质但仍是受造之物（如人的灵魂或心智）的人，我们说幸福甚至在这种形式下也找不到，我们必须崇拜真天主，他不仅无形，而且不变；因为只有他配受崇拜，在享受他时，朝拜者获得幸福，没有他，灵魂无论拥有什么，必定是悲惨的。而你们呢，你们崇拜的东西除了在你们虚构的传说中存在外，根本不存在；如果你们是异教徒，或者崇拜存在之物（尽管不是适当的崇拜对象），你们会更接近真虔诚和宗教。事实上，你们并没有真正崇拜太阳，尽管太阳在绕天运行时携带着你们的祈祷。\par

6. 你们关于太阳本身的言论如此虚假荒谬，以至如果他报复你们对他的伤害，他会将你们烧死。首先，你们称太阳为一艘船，所以你们不仅如俗语所说偏离世界，而且漂流。其次，尽管每个人都知道太阳是圆的，这是因其完美性而与他在天体中所处位置相称的形状，你们却主张他是三角形的，也就是说，他的光是通过天上的一个三角形窗户照到地上的。因此你们向太阳弯腰低头，而你们所崇拜的不是这可见的太阳，而是某种你们假想通过三角形开口发光的虚构之船。的确，如果为异端虚构所需的词语必须像建造船梁所需的木材一样付费，这船就永远不会被听到。这一切相对来说是无害的，无论多么可笑或可怜。截然不同的是你们邪恶的幻想：有一群男女少年从这船中出来，他们的美貌激起黑暗王子公主的强烈欲望；因此，你们的神的成员借着罪恶的情欲和感官的贪欲，从黑暗种族成员身上的这一耻辱束缚中被释放出来。而你们竟然将天主圣三的奥迹与这些污秽的破布联结起来；因为你们说圣父住在隐秘之光中，圣子的能力在太阳里，圣子的智慧在月亮里，圣神在空气里。\par

7. 至于这三重或更确切地说四重的虚构，关于圣父的隐秘光明，我还有什么可说的呢？只不过你们除了所见的光明之外，无法想象任何光明。从你们对可见光明的认识，即野兽、昆虫以及人类都熟悉的这种光明，你们在脑海中形成某种模糊的概念，称之为天主圣父与祂的臣民所居住的光明。你们如何区分我们所见的光明与我们所理解的光明？因为按照你们的想法，理解真理无非就是形成物质形态的概念，这些形态或是有限的，或有时是无限的；而你们竟然相信这些狂野的幻想！很明显，我心中想到你们那并不存在的光明之域，与我想到虽然未曾眼见却实际存在的\Place{亚历山大}，是全然不同的。再者，形成对未曾见过的\Place{亚历山大}的概念，与想到我所认识的\Place{迦太基}，又大不相同。但这种差异与以下差异相比微不足道：即想到我凭眼见而认识的物质事物，与理解正义、贞洁、信德、真理、爱德、善良等本性之事。你们能描述这种理智之光吗？它使我们清晰地认识它自身与其他事物的区别，以及那些事物彼此间的区别。然而即便如此，也不能说天主就是这种光，因为这种光是被造的，而天主是造物主；光是被造的，而祂是创造者；光是可以变化的。因为理智从厌恶变为渴望，从无知变为知识，从遗忘变为回忆；而天主在意志、真理和永恒中保持不变。我们由天主获得存在的开端、知识的原则、情感的法度。所有理性的和非理性的动物，从天主获得其生命的本性、感觉的能力、情感的功能。所有的物体，从天主获得其在广延中的存在、在数目中的美、在重量中的秩序。这光是一个神圣的存在，在不可分的天主圣三中；然而，你们却为那非物质、属灵、不变的本体中的各部分指派了不同的位置，仿佛它是可分的，尽管它并没有采取任何形体。你们没有为天主圣三安排三个位置，而是四个：一是不可接近的光明（你们对此一无所知）给圣父；二是太阳和月亮给圣子；又一个是大气圈给圣神。关于圣父的不可接近的光明，我在此暂不多言，因为正统信徒并不将圣子和圣神与圣父在这光明中分离。\par

8. 很难理解你们如何被这种荒谬的想法所吸引，将圣子的能力置于太阳，将祂的智慧置于月亮。因为，圣子不可分割地存在于圣父内，祂的智慧和能力不能彼此分离，以致一个在太阳，另一个在月亮。只有物质事物才能如此被指派到不同的地方。你们若明白这一点，就不会将病态幻想的产物当作如此众多虚构的素材了。但你们的观念中不仅有虚假，还有矛盾和不可信。因为，按照你们的说法，智慧的座位在亮度上低于能力的座位。然而，活力与生产力是能力的属性，而光则教导和显明；因此，如果太阳有更大的热，月亮有更大的光，这些荒谬之处对于只凭物质概念认识事物、属血肉的人而言，或许看起来有几分可能性。从高热与运动的联系，他们可能将能力等同于热；而光因其明亮，且使事物可辨别，他们可能将光代表智慧。但将能力置于在光方面远为卓越的太阳，而将智慧置于在亮度上远为逊色的月亮，这是何等愚蠢和亵渎！此外，当你们将基督与祂自己分离时，你们没有区分基督与圣神；而基督是独一的，是天主的能力，也是天主的智慧，圣神则是另一位。但按照你们的说法，你们所指定的圣神之座——大气，充满并渗透宇宙。因此，太阳和月亮在运行中始终与大气联合。但月亮时而靠近太阳，时而远离。因此，如果我们相信你们，或者更确切地说，如果我们允许自己被你们欺骗，那么智慧就远离能力半个圆周，然后又靠近另外半个圆周。而当智慧圆满时，它却远离能力。因为当月亮满月时，两个天体之间的距离如此之大，以致月亮在东方升起时，太阳正在西方落下。但由于能力的缺失产生软弱，月亮越圆，智慧必定越弱。如果——这确实是真实的——天主的智慧在能力中不可改变，天主的能力在智慧中不可改变，那么你们如何能将它们分离，指派给不同的地方呢？当本体相同时，地方怎能不同？这难道不是受制于物质幻想的迷狂吗？显示出能力和智慧的如此匮乏，以致你们的智慧软弱如你们的愚笨能力？这可憎的荒谬之处将基督分割在太阳和月亮之间——祂的能力在一个之中，祂的智慧在另一个之中；这样祂在两者中都不完全，在太阳中缺乏智慧，在月亮中缺乏能力，而在两者中祂都提供青年男女，以激发黑暗王子和公主的情感。这就是你们学习并承认的信条。这就是指导你们行为的信仰。你们怎能奇怪自己被视为可憎恶的呢？\par

9. 然而，除了你们对这些显眼而熟悉的发光体所犯的错误之外——你们崇拜它们，不是因其本来面目，而是因你们狂野的幻想所赋予的样子——你们的其他荒谬之处比这更糟糕。你们那显赫的\Person{世界承载者}和帮助支撑他的\Person{阿特拉斯}都是虚幻的存在。如同你们幻想中无数其他造物一样，它们并不存在，而你们却崇拜它们。因此我们说你们比异教徒更恶劣，尽管你们在崇拜多神方面与他们相似。你们更恶劣，因为他们崇拜的是存在之物（尽管不是神），而你们崇拜的既非神也非任何其他东西，因为它们根本不存在。异教徒也有神话，但他们知道那是神话；要么视之为有趣的诗歌幻想，要么试图将其解释为事物的本性或人类生活的象征。因此他们说\Person{伏尔甘}是跛脚的，因为普通火焰中的火苗运动不规则；\Person{福尔图娜}是盲目的，因为所谓偶然事件具有不确定性；\OriginalTerm{命运三女神}{Parcae}手持纺线杆、纺锤和手指纺羊毛成线，那是因为有\Emph{三个时间}——过去，已纺好绕在纺锤上；现在，正从纺者手中穿过；未来，仍缠在纺线杆上，即将经过手指到纺锤，即从现在进入未来；\Person{维纳斯}是\Person{伏尔甘}的妻子，因为快乐与热量有天然联系；她是\Person{玛尔斯}的情妇，因为快乐并不恰当地伴随武士；\Person{丘比特}是带翅膀持弓的男孩，因为轻率无常的情欲在不幸福者心中造成创伤；其他许多神话也是如此。在做出这些解释之后，他们仍继续崇拜这些存在，这本身就是极大的荒谬；因为不加解释的崇拜虽然也有罪，但罪行较轻。恰恰是这些解释证明，他们并不崇拜那惟独能赐予幸福的天主，而是崇拜祂所创造的事物。而且在受造物中，他们不仅崇拜德行，如从\Person{朱庇特}头中诞生的\Person{弥涅尔瓦}，她代表理智，根据\Person{柏拉图}之说，理智位于头部；他们也崇拜恶习，如\Person{丘比特}。因此他们的一位戏剧诗人写道：\Quote{罪恶的情欲，为支持恶习，使爱成了神。}甚至身体的病患在罗马也有神庙，如苍白和发热。且不说这些偶像崇拜者的罪过，他们在某种程度上受形体形状的影响，以至于当看到这些偶像被置于高处、倍受尊崇时，便将它们当作神祇敬拜；更大的罪恰恰在于那些试图为这些哑、聋、瞎、无生命的对象辩护的解释之中。尽管如此，如我所说，这些事物在救恩或实用方面毫无价值，但它们以及它们所象征的事物毕竟是真实存在的。但你们的\OriginalTerm{原人}{Primus Homo}与五大元素争战；你们的\OriginalTerm{大能之灵}{Magna Virtus}用俘虏的黑暗种族身体建造世界，更确切地说，是用你们的神受制受缚的肢体；你们的\OriginalTerm{持世者}{Mundum Tenens}手中握着这些肢体的残余，哀叹其余部分被俘、被缚、被玷污；你们的\Person{阿特拉斯}巨人将\OriginalTerm{持世者}{Mundum Tenens}扛在肩上，免得他因疲惫而丢弃重负，从而阻止最终模仿黑暗种族之团块的完成——这将是你们戏剧的最后一场；——这些以及其他无数荒谬之处，并未以绘画、雕刻或任何解释来表现；然而你们却相信并崇拜那些并不存在的事物，同时嘲笑基督徒轻信，因为他们以能抚慰心灵的信心相信真实的事物。你们崇拜的对象可以通过许多证据证明其不存在，这里我不再赘述，因为虽然我可以毫无困难地从哲学上论述世界的构造，但在此详述将耗费太多时间。有一个证据足以说明：如果这些事物是真实的，那么天主必定会变化、腐败、受污染；这种假设既亵渎神明又违背理性。因此，这一切都是虚妄、虚假、不真实的。因此你们比那些异教徒更恶劣得多，而那些异教徒的行为人人皆知，他们至今仍保留着连他们自己都引以为耻的旧习俗痕迹；因为他们崇拜的不是神的事物，而你们崇拜的却是根本不存之物。\par

10. 如果你认为你们的教义是真实的，因为它们不同于异教徒的错误，而我们之所以是错的，是因为我们与你们的差异或许比与他们的差异更大，那么你不妨说：一个死人之所以健康，是因为他没有生病；或者，健康之所以不可取，是因为它与疾病之间的差异小于与死亡之间的差异。又或者，如果异教徒在许多情况下被视为死人而非病人，那么你不妨赞美坟墓中的骨灰，因为它们不再具有人形，相比之下，活着的身体与尸体的差异并不如与骨灰的差异那样大。因此，我们被指责为更接近异教徒的死尸，而非摩尼教的骨灰。但在划分中，常常发生这样的情况：同一事物被归入不同类别，这取决于划分所依据的相似点。例如，若将动物分为会飞的和不会飞的，那么在此划分中，人和兽类就被归为一类，以区别于鸟类，因为二者都不会飞。但若将它们分为有理性的和无理性的，那么兽类和鸟类就被归为一类，以区别于人类，因为二者都缺乏理性。福斯图斯没有想到这一点，他说：\Quote{实际上只有两个派别，即外邦人和我们，因为我们在信仰上与他们直接对立。}他所指的对立是：外邦人相信一个本原，而摩尼教徒还相信黑暗种族的本原。当然，根据这一划分，我们总体上与异教徒是一致的。但若将所有有宗教的人分为崇拜一位天主者和崇拜多位神者，那么摩尼教徒必须与异教徒归为一类，而我们则与犹太人归为一类。这是另一种划分，也可以说只有两个派别。也许你会说，你们认为你们的所有神明都是同一实体，而异教徒则不然。但你们至少在与异教徒相似这一点上，为你们的神明分配了不同的能力、职能和职司。一位与黑暗种族作战；另一位用俘获的部分建造世界；又一位，站在上方，将世界握在手中；还有一位从下方托住他；另一位转动下方火、风、水的轮子；又一位在天际巡行，用他的光芒从粪池中收集你们神明的肢体。事实上，根据你们的神话描述，你们的神明有无数职司，这些你们既未解释，也未以可见的形式表现出来。但再一次，若将人们分为相信天主关心人事者和不相信者，那么异教徒、犹太人、你们以及所有带有一些基督信仰成分的异端，都将被归为一类，而与伊壁鸠鲁派及其他持相同观点者相对。由于这是一项重要原则，这里我们再次可以说只有两个派别，而你们与我们属于同一派别。你们几乎不敢与我们相异，认为天主不关心人事，因此在这件事上，你们因反对伊壁鸠鲁派而站在我们一边。因此，根据划分的性质，同一事物此时在一类，彼时在另一类；此处相连者，彼处分离；在某些事上我们与他人归为一类，他人与我们归为一类；在另一些事上我们被分开，独树一帜。如果福斯图斯想到这一点，他就不会说出如此雄辩的胡言乱语。\par

11. 但我们对\Person{福斯图斯}的话应作何理解：\Quote{圣神借其影响和灵性灌注，使大地受孕并生出有死的耶稣，他从每棵树上悬挂下来，乃是人的生命与救恩？}暂时放过此言的荒谬，我们注意到，相信这有死的耶稣可借圣神的能力由大地受孕而生，而非由童贞\Person{玛利亚}所生，这是何等愚蠢。你们胆敢将那位贞洁童女的子宫与任何长有树木和植物的土地相提并论吗？你们假装憎恶一位纯洁的童女，却毫不畏惧地相信耶稣是在由城市污秽的排水灌溉的花园中产生？因为各种植物都在这样的潮湿中发芽生长。你们要让耶稣如此诞生，同时却反对他由童贞女所生的观念。你们认为肉体比其本性所排出的排泄物更不洁吗？排泄物比排出它的肉体更洁净吗？你们不知道田地是如何施肥以使其多产吗？你们的愚昧到了这样的地步：按照你们的说法，圣神既蔑视\Person{玛利亚}的子宫，就使大地结果更丰，因为大地被动物粪便精心施肥。你们回答：圣神处处保存其不朽的纯洁。我再问：为何不在童贞女的子宫中呢？撇开受孕不谈，至于有死的耶稣——你们说他由大地而生，因圣神的能力而受孕——你们说他以果实的形状悬挂在每棵树上：因此，除了这种污染外，他还从无数吃果子的动物的肉身上受到额外的玷污；除了被你们吃掉而净化的那一小部分以外。我们相信并宣认基督是天主之子、天主圣言，变成肉体而未受玷污，因为神性实体不被肉体玷污，正如它不被任何事物玷污一样；而你们的幻想将使耶稣甚至在树上悬挂时就已经被玷污，然后才进入任何动物的肉身；因为如果他未被玷污，就不需要被你们吃他来净化。如果所有的树都是基督的十字架，正如\Person{福斯图斯}所暗示的那样，当他说耶稣挂在每棵树上时，你们为何不摘果子，将耶稣从树上取下，埋葬在你们的胃里，这岂不相当于\Place{阿黎玛特雅}人\Person{若瑟}的善行，他从十字架上取下真实的耶稣并埋葬了他？\BibleRef{若 19:38}为什么从树上取下基督是亵渎，而将他放在坟墓中却是虔诚？也许你们要将\Person{保禄}从先知引用的那句话应用于自己：\BibleQuote{他们的喉咙是敞开的坟墓}\BibleRef{罗 3:13}，因此你们张着嘴等待，好使某人的喉咙成为基督最好的坟墓。再者，你们造了多少个基督？有一个你们称为有死的基督，由大地借着圣神的能力受孕而生；另有一个在\Person{般雀比拉多}手下被犹太人钉在十字架上；第三个你们分给太阳和月亮？还是同一位，其一部分被囚在树上，要靠未囚的另一部分来释放？若是如此，并且你们承认基督在\Person{般雀比拉多}手下受苦，虽然很难看出他如何能没有肉体而受苦，因为你们说他确实没有肉体，那么重大问题就是：他与谁留下了你们所说的那些船只，好让他下来承受这些苦难，而这些苦难他显然不能没有某种身体而承受。仅仅是灵性的临在不会使他遭受这些苦难；而他在身体的临在中不可能同时在太阳、月亮和十字架上。因此，若他没有身体，他就没有被钉十字架；若他有身体，问题是他从哪里获得身体：因为按照你们的说法，所有身体都属于黑暗种族，尽管你们除了物质形态之外不能思考神性实体。因此，你们必须说：或者基督没有身体而被钉十字架，这是完全荒谬的；或者他钉十字架是外表而非现实，这是亵渎；或者并非所有身体都属于黑暗种族，神性实体也有身体，且不是不朽的身体，而是可被钉十字架且可死的，这又是完全错误的；或者基督从黑暗种族那里得到一个可死的身体，因此，你们既不允许基督的身体来自童贞\Person{玛利亚}，却将它来自魔鬼种族。最后，正如在\Person{福斯图斯}的陈述中，他以尽可能简短的方式暗示了摩尼教编造的长篇故事：大地借圣神的能力受孕并生出有死的耶稣，他从每棵树上悬挂下来，乃是人的生命与救恩。为什么这位救主以悬挂之物为代表，因为他曾被挂在树上，而不以诞生之物为代表，因为他曾诞生？但如果你们的意思是，在树上的耶稣、在\Person{般雀比拉多}手下被钉死的耶稣、分给太阳和月亮的耶稣，都是同一个实体，你们为何不给你们整群的神祇都加上耶稣的名号？为何你们的持世者不可以也是耶稣，\Person{阿特拉斯}、尊荣之王、大力神、原人，以及其余所有以不同名号和职能存在的呢？\par

12. 既然如此，论到圣神，当你们提到的\Quote{位}不可胜数时，你们怎能说祂是第三位呢？又为何祂不是耶稣自己？\Person{福斯图斯}企图在他自己与真正的基督徒之间制造一种共识（其实他与他们相去甚远），说什么\Quote{我们在全能天主圣父、祂的圣子基督和圣神这三重称谓下崇拜独一天主}，他为何要误导人？为何称谓只是三重，而不是多重？如果名称有多少，位格就有多少，那么为什么区别只在名称上，而不在实质上？因为你们并非给同一事物三个名称，如同同一武器可称为短剑、匕首或匕首；或者你们给同一事物取名为月亮、较小的船、夜晚的光体等等。因为你们不能说\OriginalTerm{原人}{First Man}与\OriginalTerm{大能之神}{Mighty Spirit}或\OriginalTerm{持世者}{World-Holder}或巨人\Person{阿特拉斯}相同。它们都是不同的位，而你们并不称它们中任何一个为基督。一个神性如何能有相反的功能？或者，如果基督以同一实体悬挂在树上，被犹太人迫害，又存在于太阳和月亮中，为何基督自己不是唯一的位？事实上，你们的幻想全都偏离正道，与疯狂的梦境无异。\par

13. 当\Person{福斯图斯}认为我们与\OriginalTerm{摩尼派信徒}{Manichæans}相似，将尊崇归于饼和酒时，摩尼派信徒认为尝酒是亵渎，他怎能这样想？他们在葡萄中承认自己的神，却在杯中不承认；或许他们对神被践踏和被装瓶感到震惊。我们并非把任何饼和酒当作自然产物来尊崇，仿佛基督被限制在谷物或葡萄藤中，如同摩尼派信徒所幻想的那样；而是尊崇那真正被祝圣为标记的东西。未被祝圣的饼和酒，尽管是饼和酒，只是滋养或提神之物，毫无神圣性；虽然我们为每一种恩赐，不论是身体的还是精神的，都祝福并感谢天主。按你们的观念，基督被限制在你们所吃的一切中，并通过消化从你们肠道的附加限制中释放出来。所以，当你们吃时，你们的神受苦；当你们消化时，你们因祂的恢复而受苦。当祂充满你们时，你们的所得就是祂的所失。这或许可被视为祂的仁慈，因为祂在你们内为你们的益处受苦，要不是祂逃离你们而获得自由，使你们空虚。我们对饼和酒的敬意与你们的教义毫无相似之处，你们的教义毫无真理。将两者相比，甚至比有些人说我们在饼和酒中崇拜\Person{色列斯}和\Person{巴克斯}更为愚蠢。我现在提及此事，是为了指出你们从哪里得来那种愚蠢的想法，即我们的祖先守安息日是为了敬拜\Person{萨图尔努斯}。因为正如在我们的饼和酒圣事的遵守中，与对这些外邦神灵色列斯和巴克斯的崇拜毫无瓜葛，你们对此事如此赞许，以至希望在这方面与我们相似，同样，我们的祖先守安息日的休息，以适合先知时代的方式，也并非服从萨图尔努斯。\par

14. 你们或许已经在外邦人的宗教中找到了与你们相似的关于\Emph{\OriginalTerm{质料}{Hyle}}的东西，外邦人常提到它。相反，你们却主张，在对你们的神学导师称为\Emph{\OriginalTerm{质料}{Hyle}}的那个恶原则的信仰上，你们直接反对他们。但这里你们只表现出自己的无知，并且卖弄学问，使用这个词却不知其含义。希腊人在谈论自然时，将\Emph{\OriginalTerm{质料}{Hyle}}这个名称赋予事物的基底材质，它本身没有形式，但接受一切有形体的形式，只有通过这些可变的现象才能被认识，本身并非感官或知觉的对象。确实，有些外邦人错误地使这种质料与天主同永恒，认为并非出自祂，尽管有形体的形式是出自祂。在这个明显的错误上，你们与这些外邦人相似，因为你们认为\Emph{\OriginalTerm{质料}{Hyle}}有它自己的本原，并非来自天主。只是无知使你们否认这种相似。当外邦人说\Emph{\OriginalTerm{质料}{Hyle}}没有自己的形式，只能从天主获得形式时，他们接近了真理，这真理正是我们相信的，与你们的错误截然相反。你们不知道什么是\OriginalTerm{质料}{Hyle}或事物的基底材质，就把它当作\OriginalTerm{黑暗种族}{race of darkness}，在其中不仅放置了五种不同类别的无数有形体的形式，还放置了一个\OriginalTerm{赋形的心智}{formative mind}。事实上，你们的无知或疯狂如此之甚，以至于你们称这个心智为\OriginalTerm{质料}{Hyle}，使之赋予形式而非接受形式。如果确实有像你们所说的这样一个赋形的心智，以及能够赋形的有形体的元素，那么\OriginalTerm{质料}{Hyle}一词本应适用于有形体的元素，即被心智塑造的质料，而你们则将心智当作恶的原则。即便如此，这个词的使用也不太准确，因为\OriginalTerm{质料}{Hyle}没有任何形式；而这些元素虽然能够接受新形式，却已经具有元素的形式，并属于不同的种类。尽管如此，尽管你们无知，这个词的此种用法也不至于太错；因为它适用于接受形式的东西，而非赋予形式的东西。然而，即便如此，你们的愚蠢和不敬仍然表现在：你们将那么多美好的东西追溯至恶的原则，因为你们不知道所有各种类型的自然、比例恰当的一切形式、秩序有序的一切重量，只能来自圣父、圣子和圣神。事实是，你们既不知道什么是\OriginalTerm{质料}{Hyle}，也不知道什么是恶。但愿我能说服你们不再误导那些比你们更无知的人！\par

15. 人人皆可见尔等自夸优于外教人之愚妄，盖因外教人于敬拜天主时使用祭台、圣殿、偶像、祭献与馨香，而尔等则否。仿佛为一块有某种存在的石头建造祭台并献祭，胜于以狂热的想象敬拜全然不存在之物。尔等自称是\Quote{天主的理性圣殿}，此何意？那部分由魔鬼建造的，能是天主的圣殿么？尔等岂非如此？因尔等声称，尔等所有肢体与整个身体皆由称为\Quote{原质}（\OriginalTerm{原质}{Hyle}）的恶原则所造，且此成形心神的一部分与尔等之神的一部分一同居于身体内。既然尔等之神的部分被捆绑束缚，尔等应被称为天主的监狱，而非祂的圣殿。或许尔等的灵魂才是天主的圣殿，因它来自光明之境。但尔等通常不称灵魂为圣殿，而称其为天主的一部分或肢体。故此，当尔等自称是\Quote{天主的圣殿}时，必指尔等的身体，而尔等声称这身体是由魔鬼所造的。这样，尔等亵渎天主的圣殿，不仅称其为撒殚的制造品，且称为天主的牢狱。反之，宗徒说：\Quote{天主的圣殿是圣的，这圣殿就是你们。}（\BibleQuote{天主的圣殿是圣的，这圣殿就是你们。}）为表明这不单指灵魂，他明确说：\Quote{岂不知你们的身体是住在你们内之圣神的宫殿么？这圣神是天主赐给你们的？}（\BibleQuote{岂不知你们的身体是住在你们内之圣神的宫殿么？这圣神是天主赐给你们的？}）尔等竟将魔鬼的制造品称为天主的圣殿，并在其中——用福斯特的话说——安置了基督、天主之子、永生威严的活像。尔等的不敬虔足以为一个神话般的基督虚构一个神话般的殿宇。尔等所说的像，必得如此称呼，因为它不过是尔等想象的产物。\par

16. 倘若尔等的心智是一座祭台，那么尔等当知它是谁的祭台。从尔等声称受训的教义与职分中，即可看出。尔等受教不得给乞丐食物；因此，尔等的祭台冒出残忍的祭烟。这样的祭台，主要予以摧毁；因为主引用法律的话告诉我们，什么祭献悦乐天主：\Quote{我喜欢仁慈，不喜欢祭献。}（\BibleQuote{我喜欢仁慈，不喜欢祭献。}）请观察主在何种场合说这话。那是当祂经过田地时，门徒因饥饿而摘了麦穗。尔等的教义会叫你们称此为谋杀。尔等的心智是一座祭台，并非天主的祭台，而是说谎魔鬼的祭台，藉其教义，邪恶的良心如同被烙铁烙过一般，\BibleRef{弟前 4:2}将真理称为无罪的事称作谋杀。因为基督在祂对犹太人的话中，预先给了你们致命的一击：\Quote{假如你们了解\InnerQuote{我喜欢仁慈，不喜欢祭献}是什么意思，你们就决不定罪无罪的人了。}（\BibleQuote{假如你们了解\InnerQuote{我喜欢仁慈，不喜欢祭献}是什么意思，你们就决不定罪无罪的人了。}）\BibleRef{玛 12:7}\par

17. 尔等也不能声称，你们以纯洁简朴的祈祷作为祭献来尊崇天主；因为尔等对神性本质与实体的卑下、羞辱的观念，使你们的神成为外教人祭献中的牺牲；如此，尔等远非以你们的祭献悦乐真天主。因为尔等主张，天主不仅被囚禁在树木、植物或人体中，也存在于动物的肉中，那肉以其不洁玷污了祂。再者，当尔等称自己的灵魂是黑暗种族所俘获的天主本体微粒时，岂不是在责备天主，仿佛天主若不藉其肢体的败坏和这可耻的俘获，就无法维持这场争战？尔等非但在祈祷中不尊崇天主，反而侮辱了祂。因为当你们属于祂时，你们犯了什么罪，致使你们向你们呼号的神这样被惩罚？不是因你们有罪的自由选择而离开了祂，而是祂自己把你们给了祂的仇敌，好为祂的王国取得和平？你们甚至不是作为被体面守卫的人质交出的。也不像牧人设陷阱捕野兽：他不将自己的肢体放入陷阱，而是从他的羊群中取一只动物；通常，野兽在被捕获之前动物不会受伤。而你们，既是天主的肢体，却被交给仇敌，你们阻止仇敌的凶暴，只有靠沾染他们的不洁、感染他们的败坏，而你们自己并无罪过。你们不能在祈祷中使用这些话：\Quote{主，求祢为祢圣名的荣耀释放我们；并为祢圣名的缘故赦免我们的罪。}（\Quote{主，求祢为祢圣名的荣耀释放我们；并为祢圣名的缘故赦免我们的罪。}）你们的祈祷是：\Quote{藉祢的技巧释放我们，因为我们在这里受压迫、折磨、玷污，只是为了祢能在祢的王国中不受侵扰地哀悼。}（\Quote{藉祢的技巧释放我们，因为我们在这里受压迫、折磨、玷污，只是为了祢能在祢的王国中不受侵扰地哀悼。}）这是责备的话，而非恳求。你们也不能用真理导师教导我们的话：\Quote{宽免我们的罪债，如同我们也宽免得罪我们的人。}（\BibleQuote{宽免我们的罪债，如同我们也宽免得罪我们的人。}）\BibleRef{玛 6:12} 因为谁是得罪你们的债主？若是黑暗种族，你们并不宽免他们的债，反而叫他们被彻底逐出并囚禁于永久的牢狱。而天主又如何能宽免你们的债呢？祂把你们送到这种境地，反而是祂得罪了你们，而不是你们得罪了祂——你们服从祂去了。若这不是祂的罪，因为祂被迫这样做，那么这借口更适用于你们（现在你们已在争战中被打败）而非祂（在争战开始之前）。你们现在因恶的混杂而受苦，而祂却没有，但祂仍然被迫差遣你们。所以，要么祂该要求你们宽免祂的债；要么，若祂并不欠你们，你们就更不欠祂了。由此可见，尔等的祭献与纯洁简朴的祈祷，实为虚假而卑鄙的亵渎。\par

18. 顺便问一下，你们怎么使用\Quote{圣殿}、\Quote{祭台}、\Quote{祭献}这些词来称赞你们的行为？假如这些东西可以被说成是正当属于真宗教，那么它们必定构成对真天主的真崇拜。而如果存在对真天主的真祭献，这包含在\Quote{尊神之礼}一词中，那么必定有一些真正的祭献，其余的都是模仿。一方面，我们在假而说谎的神祇——也就是魔鬼——那里有虚假的模仿，他们傲慢地要求受骗的信徒给予尊神之礼，正如外邦人的圣殿和偶像中曾经或现在所行的。另一方面，我们有对那最真实的祭献的预表，这祭献要为所有信徒的罪而献上，正如天主命令我们先祖献上的祭献那样；与之相伴的还有象征性的傅油，预象\OriginalTerm{基督}{Christ}，因为\Quote{基督}这个名字本身就有\Quote{受傅者}的意思。因此，动物祭献被魔鬼僭妄地据为己有，它们是对那唯独属于唯一真天主、且唯独基督在其祭台上献上的真祭献的模仿。因此宗徒说：\Quote{外邦人所献的祭，是祭献给魔鬼，而不是献给天主。} \BibleRef{格前 10:30} 他并非指摘祭献本身，而是指摘祭献给魔鬼。再者，希伯来人在他们多种多样的动物祭献中，按照这礼仪的意义，预表了基督所献的祭献。基督徒也在神圣的奉献和领受基督体血中纪念这祭献。摩尼派人士既不理解外邦人祭献的罪恶性，也不理解希伯来人祭献的重要性，更不理解基督徒祭献礼仪的用处。他们自己的谬误就是他们献给那欺骗他们的魔鬼的供物。因此他们离弃了信德，听从了引诱人的邪神和魔鬼的教义，用伪善说谎言。\par

19. 或许福斯特，或至少那些被福斯特著作所吸引的人，应当知道单一原则的教义不是我们从外邦人那里得来的；因为信仰一位真天主——万有的本性都由祂而来——是外邦人中所保留的原初真理的一部分，尽管他们已经堕落信奉许多假神。因为外邦哲学家认识天主，正如宗徒所说：\Quote{自从天主创世以来，祂那看不见的事，祂永恒的大能和神性，都可通过所造之物被人明了，因此他们无可推诿。} \BibleRef{罗 1:20} 但宗徒又接着说：\Quote{他们既然认识了天主，却没有当作天主光荣祂，也不感谢祂，反而在思想中成为虚妄，他们愚昧的心就昏暗了。他们自称为智者，反而成了愚拙，将不朽天主的荣耀换成了必朽的人、飞鸟、四足走兽和爬行昆虫的偶像。} \BibleRef{罗 1:25} 这就是外邦人的偶像，他们除了用天主所造的受造物来解释之外，无法解释它们；因此，他们对偶像崇拜的这种解释——更开明的外邦人常以此自夸，视为高人一等的证明——恰好证明了宗徒下面的话是真实的：\Quote{他们敬拜事奉受造物，而不敬拜事奉造物主，祂是永受赞美的。} 在你们与外邦人不同的地方，你们是错误的；在你们与他们相似的地方，你们比他们更坏。你们不像他们那样相信单一原则；所以你们陷入不虔敬，相信唯一真天主的本体可能被征服和败坏。至于对众神的敬拜，说谎魔鬼的教义引导外邦人敬拜许多偶像，而引导你们敬拜许多幻象。\par

20. 我们并没有像福斯特所说的那样，把外邦人的祭献改成\OriginalTerm{爱宴}{love-feasts}。我们的爱宴反而是对主所说的那祭献的替代，如已引用的话：\Quote{我要的是仁慈，不是祭献。} 在我们的爱宴中，穷人得到蔬食或肉食；如此，天主的受造物按其适宜的方式被用来滋养人——他也是天主的受造物。你们被说谎的魔鬼引导，不是出于克己，而是出于亵渎的谬误，\Quote{要禁戒那些天主所造、为使信者并认识真理的人以感恩之心领受的食物。因为天主的每样受造物都是好的，没有一样是可弃绝的，只要以感恩之心领受。} \BibleRef{弟前 4:3} 对于造物主的恩赐，你们却忘恩负义地用不敬来侮辱祂；并且因为在我们的爱宴中常常把肉给穷人，你们就把基督徒的爱德比作外邦人的祭献。这确实是你们在某些方面与外邦人相似的另一点。你们认为杀害动物是一种罪行，因为你们以为人的灵魂会进入它们里面；这种思想在一些外邦哲学家的著作中可以找到，尽管他们的后继者似乎有不同的看法。但在这方面你们又是最大的错误：因为他们害怕在动物中屠杀亲属；而你们却害怕屠杀你们的神，因为你们甚至认为动物的灵魂是他的肢体。\par

21. 关于我们尊敬殉道者的纪念，以及\Person{Faustus}指控我们以敬拜他们来代替偶像，我本不愿回答这样的指控，若不是为了显示\Person{Faustus}在渴望责备我们时，已经超越了摩尼教的虚构，而轻率地陷入了外邦诗歌中流行的观念，尽管他如此渴望与外邦人区别。因为他说我们把偶像变成了殉道者，他声称我们用类似的礼仪敬拜他们，并用酒食安抚逝者的幽灵。那么，你相信幽灵吗？我们从未听你说过这类事，也未在你的书中读过。事实上，你们通常反对这种观念：因为你们告诉我们，死者的灵魂，如果他们是恶的，或未净化，就得经历各种变化，或遭受更严厉的惩罚；而善良的灵魂则被置于船中，航行过天，到达他们为之战斗至死的那虚构的光明领域。按照你们的说法，没有灵魂留在埋葬身体的墓地附近；那么逝者的幽灵怎么可能存在？它们是什么，在哪里？\Person{Faustus}好谗言，使他忘记了自己的信条；或者他是在梦中谈论鬼魂，甚至看到自己的话被写在纸上也没有醒来。确实，基督徒以宗教敬礼尊敬殉道者的纪念，既为激励我们效法他们，也为分享他们的功德及他们祈祷的帮助。但我们不向任何殉道者建筑祭台，而是向殉道者的天主建筑祭台，尽管这是为了纪念殉道者。在圣人墓地主持祭礼的人从不说：\Quote{我们向您，\Person{伯多禄}！或\Person{保禄}！或\Person{西彼廉}！呈献供品。}而是向赐予殉道荣冠的天主呈献，同时纪念那些获得荣冠的人。场所的联想增加了情感，爱心被激发，既朝向那些作我们榜样的人，也朝向那藉其帮助我们能跟随这样榜样的人。我们以同样的亲密感情看待殉道者，如同我们在今生看待天主的圣人们，当我们知道他们的心已准备好为福音的真理忍受同样的苦难。我们对殉道者的感情有更多的虔敬，因为我们知道他们的战斗已经结束；我们可以更自信地赞美那些已在天堂得胜的人，而不是那些仍在世上战斗的人。真正的神圣敬拜，希腊人称之为\OriginalTerm{钦崇}{latria}，拉丁文中没有对应的词，无论在教义还是实践上，我们都只归于天主。属于这种敬拜的包括献祭；正如我们在\OriginalTerm{偶像崇拜}{idolatry}一词所见，意指将这种敬拜给予偶像。因此，我们从不向殉道者、或圣洁的灵魂、或任何天使献祭，或要求任何人这样做。任何陷入这种错误的人，都会通过教义得到教导，或是纠正或是警告。因为神圣的存有本身，无论是圣人还是天使，都拒绝接受他们所知唯独归于天主的。我们在\Person{保禄}和\Person{巴尔纳伯}身上看到这一点，当\Place{吕考尼雅}人因他们行的奇迹想要向他们献祭如神时，他们撕裂衣服，阻止百姓，向他们呼喊，说服他们不是神。我们也看到天使，如我们在\BibleBook{}{默示录}中所读，天使不允许自己被敬拜，对敬拜者说：\BibleQuote{我是你的同僕，和你的弟兄们。}\BibleRef{默 19:10}那些要求这种敬拜的是骄傲的邪灵，魔鬼及其使者，正如我们在外邦人的所有庙宇和礼仪中所见。有些骄傲的人也模仿他们的榜样；如同巴比伦的几个国王的记载。圣\Person{达尼尔}因此被控告和迫害，因为当国王颁布法令，不得向任何神，只可向国王祈祷时，他被发现敬拜和祈祷自己的天主，即唯一真天主。\EditorialNote{达尼尔第六章}至于那些在殉道者宴会上酗酒的人，我们当然谴责他们的行为；因为即使在他们自己家里这样做，也违背纯正的教义。但我们必须既改正坏的，也规定好的，并且必然暂时容忍一些不符合我们教导的事。基督徒行为的准则不应从放纵无度的人或软弱者的放任中取得。即使如此，放纵的罪过远比不敬的罪过小。向殉道者献祭，即使空腹，也比从他们的宴会醉酒回家更坏：我说向殉道者献祭，这与向天主献祭以纪念殉道者不同，后者是我们经常所做的，按照新约启示以来所要求的方式，因为这属于唯独归于天主的敬拜或\OriginalTerm{钦崇}{latria}。但想使这些异端者完全理解圣咏作者这些话的意义是徒劳的：\BibleQuote{那以赞颂为祭献的人光荣我，并且我必以此向他显示我的救恩。}在基督来临之前，这祭献的血肉在宰杀的牲畜中预表；在基督的苦难中，预表由真实的祭献实现了；在基督升天之后，这祭献在圣事中纪念。外邦人的祭献与希伯来人的祭献之间的区别，正如虚假的模仿与典型的预兆之间的区别。我们不因有维斯塔贞女就鄙视或谴责圣洁妇女的童贞。同样，外邦人有祭献并不构成对我们先祖祭献的指责。基督徒的童贞与维斯塔贞女之间的区别很大，却完全在于向谁许愿和还愿；同样，外邦人和希伯来人祭献的对象之间的区别，使它们之间有巨大的差别。前者是献给魔鬼，它们僭妄地要求这祭献以便被当作神，因为祭献是神圣的尊荣。后者是献给唯一真天主，作为真正祭献的预象，这祭献也将在基督身体和血的苦难中献给祂。\par

22. 福斯图斯错误地说，我们的犹太祖先在与外邦人分离时，保留了圣殿、祭献、祭台和司祭职，而只放弃了雕像或偶像；因为他们本可以在没有雕像的情况下，像某些人那样向树木和山丘，甚至向太阳、月亮和星辰献祭。如果他们向这些对象献上称为\OriginalTerm{崇拜}{latria}的敬礼，他们就会事奉受造物而非造物主，从而陷入异教迷信的严重错误；即使没有偶像，他们也必会发现魔鬼乐于利用他们的错误，并接受他们的供品。因为这些骄傲而邪恶的灵体并不像某些愚昧者所想像的那样，以祭献的气味和烟雾为食，而是以人的错误为食。他们享受的不是肉体的满足，而是一种恶意的快感，当他们以任何方式欺骗人，或者大胆地借取借来的威严，自夸接受神圣的尊荣时。因此，我们的犹太祖先所放弃的不仅仅是外邦人的偶像。他们既不向大地也不向任何地上的事物献祭，也不向海洋、苍天或天上的万军献祭，而是将祭品放在唯一的天主、万有的创造者的祭台上，祂要求这些供品作为预表真正牺牲的方式，那牺牲就是我们的主耶稣基督，藉着祂，天主在赦免罪恶中使我们与自己和好。因此，保禄对信徒——他们成为以基督为元首的身体——说：\Quote{所以我劝你们，弟兄们，因天主的仁慈，献上你们的身体作为生活、圣洁、中悦天主的祭品。}\BibleRef{罗 12:1} 相反，摩尼教徒说人的身体是黑暗种族的创造物，是囚禁被俘神明的牢狱。因此福斯图斯的教义与保禄的相去甚远。但既然谁若给你们宣讲一个与你们所领受的不同的福音，就该受咒诅，那么基督在保禄内所说的是真理，而在福斯图斯内的摩尼则受咒诅。\par

23. 福斯图斯还说，我们保留了外邦人的生活方式，但他并不知道自己在说什么。既然义人因信德而生活，命令的终极是出于纯洁的心、良好的良心和真诚的信德的爱德，而这三样——信德、望德、爱德——常存以形成信徒的生命，那么在这些事上不同的人，其生活方式就不可能相同。那些信德不同、望德不同、爱德不同的人，生活也必然不同。如果我们在使用饮食、房屋、衣服、沐浴等事物上，以及我们中结婚的人在娶妻、生子、养育子女作为继承人的事上，与外邦人相似，那么，为自己某种目的而使用这些事物的人，与在使用时感谢天主、对天主没有不配或错误观念的人之间，仍存在巨大差异。因为按照你们自己的异端，尽管你们与其他人吃同样的面包，以同样的植物产品和同一泉源的水为生，穿着同样的羊毛和亚麻衣物，但你们过着不同的生活，不是因为你们吃喝或穿着不同，而是因为你们在观念和信德上与他人不同，并且在这一切事上着眼于你们自己的目的——即你们虚假教义所提出的目的；同样，我们虽然在这些事物和其他事物的使用上与外邦人相似，但在生活中却不与他们相似；因为事物相同，目的却不同：我们着眼的目的，按照天主的公义命令，是出于纯洁的心、良好的良心和真诚的信德的爱德；有些人偏离了这些，转向了\OriginalTerm{虚妄的争辩}{vain jangling}。在这虚妄的争辩中，你们独占鳌头，因为你们没有注意到，即便拥有的和使用的事物相同，不同的信德所产生的生命差异是如此之大，以致尽管你们的信徒有妻子，并不得已地生育儿女，为他们积蓄财富；他们吃肉、饮酒、沐浴、收割庄稼、收获葡萄、从事贸易、担任高级官职，你们却仍将他们算作属于你们，而不属于外邦人，尽管他们的行为更接近外邦人而非你们。即使某些外邦人——例如那些因迷信的虔诚而戒绝肉食、酒和婚姻的人——在某些事上比你们自己的追随者更与你们相似，你们仍将自己的追随者——即使他们使用这一切事物，因而与你们不同——算作属于摩尼的羊群，而不是那些在实践上相似于你们的人。你们认为一位相信摩尼的女人——尽管她是一位母亲——属于你们，而不是一位西比尔——尽管她从未结婚。但你们会说，许多称为天主教基督徒的人是奸夫、强盗、吝啬鬼、醉汉，以及其他一切与健全教义相悖的人。我问，在你们那几乎小到不能称为团体的团体中，是否也找不到这样的人？而因为在外邦人中有一些不是这种品性的人，你们是否认为他们比你们更好？然而事实上，你们的异端是如此亵渎，以致即使你们的追随者中并非这种品性的人，也比那些是这种品性的外邦人更坏。因此，对于唯一的、大公的健全教义，不能因为许多人希望取它的名字却不愿服从它的有益影响而加以指责。我们必须牢记主在小团体与散布全世界的群众之间所作对比的真正含义；因为圣人和信徒的团体是小的，正如谷粒与糠秕堆相比是小的；然而好谷粒的数量足以远远超过你们——好人和坏人加在一起，因为好人和坏人都远离真理。简言之，我们不是外邦人的分裂，因为我们与他们有极大的不同，且是更好的；你们也不是，因为你们与他们有极大的不同，且是更坏的。\par

\SourceRecord{Translated by Richard Stothert.；Nicene and Post-Nicene Fathers, First Series；Vol. 4.；Edited by Philip Schaff.；Buffalo, NY: Christian Literature Publishing Co.,；1887.；<http://www.newadvent.org/fathers/140620.htm>.}

% structure: p0001 source=h1 target=section reason=page-title

\section{驳福斯图斯，第二十一卷}

\Emph{\Person{福斯图斯}否认摩尼教徒信仰两位神祇。\OriginalTerm{祂力}{Hyle}并非天主。\Person{奥古斯丁}详细讨论天主与祂力的道理，并将二元论的指控加于摩尼教徒。}\par

1. \Person{福斯图斯}说：我们信仰一位天主还是两位？当然是一位。如果我们被指控制造了两位神祇，我回答说，从未有人能证明我们说过类似的话。你们为何怀疑我们？因为，你们说，我们相信善与恶两个本原。不错，我们相信两个本原；但我们称其中一个为天主，另一个为\OriginalTerm{祂力}{Hyle}，或者用通常通俗的说法，魔鬼。如果你们认为这意味着两位神祇，你们也可以认为医生所说的健康与疾病是两种健康，或者善与恶是两种善，或者财富与贫乏是两种财富。如果我描述两样东西，一白一黑，或一热一冷，或一甜一苦，若你们说我在描述两个白色、两个热、或两个甜的东西，那就会显得像白痴或疯癫。同样，当我断言有两个本原，天主与\OriginalTerm{祂力}{Hyle}时，你们没有理由说我相信两位神祇。你们是否认为我们必须两者都称为神，因为我们将一切恶的力量归于\OriginalTerm{祂力}{Hyle}，一切善的力量归于天主？如果是这样，你们也可以说毒药和解药都必须称为解药，因为两者各有其能力，且两者的作用产生一定效果。同样，你们可以说医生和投毒者都是医生；或者公正与不公正的人都是公正的，因为两者都做了某事。如果这是荒谬的，那么说天主与\OriginalTerm{祂力}{Hyle}必须都是神，因为它们都产生一定效果，就更加荒谬。这是一种非常幼稚和无力论证的方式，当你们无法反驳我的陈述时，就在名目上争吵。我承认，我们有时也称那敌对的本性为神；并非我们相信它是神，而是因为这一名称已被此本性的崇拜者采用，他们在错误中以为它是神。因此宗徒说：\BibleQuote{这世界的神已蒙蔽了不信者的心思。}\BibleRef{格后 4:4}他称它为神，因为它的崇拜者称它为神；他同时说它蒙蔽了他们的心思，以表明它不是真天主。\par

2. \Person{奥古斯丁}回答说：你们常在言论中谈到两位神祇，正如你们承认的，虽然起初你们否认了。你们引宗徒的话作为如此说的理由：\BibleQuote{这世界的神已蒙蔽了不信者的心思。}我们大多数人对此句标点不同，解释为真天主蒙蔽了不信者的心思。他们在\Quote{神}字后加句点，然后把后面的词连读。或者不改变标点，为了解释起见，可以调整词序，读作：\Quote{在这世界，天主蒙蔽了不信者的心思。}这给出同样的含义。蒙蔽不信者心思的行为在某种意义上可归于天主，这不是出于恶意，而是出于正义。因此保禄自己在别处说：\BibleQuote{天主不义，竟要降怒吗？}\BibleRef{罗 3:5}又说：\BibleQuote{那么我们说什么呢？难道天主有不义吗？决不！因为梅瑟说：我要恩待谁，就恩待谁；要怜悯谁，就怜悯谁。}请注意他在断言与天主没有不义这一不可否认的真理之后，又补充说：\BibleQuote{但如果天主愿意显示祂的忿怒，彰显祂的权能，就大量容忍了那为毁灭所预备的怒器，并为要彰显祂对那已经预备得光荣的怜悯之器的恩宠的丰富，那又怎么样呢？}等等。这里显然不能说，显示祂的忿怒并在为毁灭所预备的怒器上彰显祂权能的是一位天主，而在怜悯之器上显示祂恩宠丰富的是另一位天主。按照宗徒的道理，是同一位天主做了这两件事。因此他再次说：\BibleQuote{因此，天主随他们陷入心中的情欲，陷入不洁，以致彼此侮辱自己的身体。}紧接着：\BibleQuote{因此，天主随他们陷入可耻的情欲。}又说：\BibleQuote{他们既然不肯认识天主，天主就让他们陷入可鄙的心思。}在这里我们看到真正而公义的天主如何蒙蔽不信者的心思。因为以上所引的宗徒的话中，所理解的天主，无非就是祂所派遣的圣子，那位来的时候说：\BibleQuote{我是为了审判来到这世上，叫看不见的看得见，叫看得见的反而盲目。}\BibleRef{若 9:39}这里，对信者而言，再次清楚地看到天主如何蒙蔽不信者的心思。因为在那些包含天主公义审判原则的奥秘事中，有一个奥秘决定了一些人的心思要被蒙蔽，另一些人的心思要被光照。关于这一点，论到天主时说得很好：\BibleQuote{你的判断如同深渊。}宗徒在感叹这无法测度的深渊时，高呼：\BibleQuote{啊，天主的智慧和知识的财富何其深奥！祂的判断何其难测，祂的道路何其难寻！}\BibleRef{罗 11:33}\par

3. 你不能分辨天主在怜悯中所行的事与祂在审判中所行的事，因为你既不能理解也不能使用我们\WorkTitle{圣咏集}的话：\BibleQuote{我要向祢歌颂怜悯和审判，上主。}因此，凡在你脆弱人性的软弱中似乎不妥的事，你都完全与天主的旨意和判断分开：因为你为自己准备了另一个邪恶的神，不是藉着发现真理，而是藉着愚昧的发明；你将你行事不义的以及你受苦难当的，都归给这个神。这样，你将赐福的恩赐归于天主，却将降罚的审判从祂夺去，仿佛基督说祂为恶人预备了永火的那一位，与祂叫日头上升照善人也照恶人、降雨给义人也给不义人的那一位是不同的存在。为什么你不明白这伟大的良善和伟大的严厉属于同一位天主呢？岂不是因为你还没有学会歌颂怜悯和审判吗？那叫日头上升照恶人也照善人、降雨给义人也给不义人的，难道不是那折断天然的枝子、逆性接上野橄榄树的同一位吗？宗徒论到这一点时，论及这同一位天主说：\BibleQuote{可见天主又慈善又严厉：对那跌倒的人是严厉，对你却是慈善，只要你存留在祂的慈善里。}\BibleRef{罗 11:17}在此要注意，宗徒既没有使天主失去审判的严厉，也没有使人失去自由意志。这是一个深邃的奥迹，人的思想无法参透：天主如何定恶人的罪，又使恶人成义；因为这两件事都在\WorkTitle{圣经}的真理中论及祂。那么，天主旨意的奥秘性，就成为我们喜欢挑剔反对的理由吗？多么更适合我们有限的能力，怀着宗徒所感受到的敬畏而感叹说：\BibleQuote{啊，天主的智慧和知识的富饶，何等深奥！祂的判断，何其难测，祂的道路，何其难寻！}这样更好：对你无法解释的事感到惊叹，而不是因为不能理解唯一的善神，就试图在真神之外再制造一个恶神。因为这不仅仅是名号的问题，而是行动的问题。\par

4. \Person{福斯图斯}巧妙地为自己辩护说：\Quote{我们不说有两个神，而是说神和\OriginalTerm{物质}{Hyle}。}但当你追问\OriginalTerm{物质}{Hyle}的意义时，你会发现它实际上是另一个神。如果摩尼教徒像古人一样，给可接受物体形式的未成形质料起名为\OriginalTerm{物质}{Hyle}，我们就不会指责他们制造了两个神。但赋予物质形成物体的能力，或者否认拥有这种能力的是神，这纯粹是愚昧和疯狂。当你将属于真神的、使物体、元素和动物按其各自模式存在的性质和形式的能力，给予另一个存在，无论你给这存在取什么名字，你都难免制造了另一个神。确实，在这亵渎的教义中有两个错误。首先，你把天主的作为归给一个你羞于称之为神的存在；然而，只要你让他做只有神才能做的事，你就必须称他为神。其次，善神所行的善事，你却称之为恶，并归给一个恶神，因为你对于那打击堕落人性之脆弱的一切有幼稚的恐惧，而对相反的事物有幼稚的喜悦。所以你以为蛇是恶者造的；而你却认为太阳是如此之伟大，以致你以为它不是天主的受造物，而是祂本体的一种流出。你必须知道，真神——可惜你尚未相信祂——既造了蛇和低等受造物，也造了太阳和其他高贵的受造物。此外，在更高贵的受造物中，不是天体，而是属神的存在，祂造了远超过太阳之光的东西，这是任何属血肉的人都不能感知的，更不用说你了，你在谴责肉体时，恰恰谴责了你用以判断善恶的原则。因为你对于恶的唯一观念，来自于某些事物对肉性感觉的不适；而你对善的唯一观念，则来自于感官的满足。\par

5. 当我思虑自然秩序中最低微、为我们所常见、虽属尘世、脆弱、会死、却仍是天主所造的事物时，我不禁惊叹造物主：祂在伟大的工程中如此伟大，在微小的工程中也同样伟大。因为在形成天地间一切受造物时所显现的神圣技艺始终如一，即使在不同事物之间也是如此；它处处完美，按各自种类中的完美赋予一切。我们看到每个受造物并非单独造作，而是与受造的其余部分相关；因此，在形成每一个时都展现了全部神圣技艺，将各物安排在其应有的位置和秩序中，为所有事物单独并联合提供合适之物。请看，这里在等级中最低微的：飞、游、行、爬的动物。这些都是会死的受造物，其生命，正如经上所载，\BibleQuote{如蒸气，出现片刻} \BibleRef{雅 4:15}。其中每一个，按其种类的容量，都在造物主的良善中为整体的圆满贡献所定之度量，使最低微者也分有最高者以更大程度拥有的善。你们若可能，请指出任何一个受造动物，无论多么卑贱，其灵魂会恨自己的肉身，而不是养育并珍爱它、以活力运动促进其成长、引导其活动、并对其自身的\Quote{小宇宙}行使某种管理、使之服从于自身的保存。即使在一个理性受造者对自己身体的管束中，他制服身体，免得世俗情欲阻碍他对智慧的知觉，其中也存在着对自己肉身的爱，他随后使肉身服从，这正是其适当状态。实际上，你们自己，虽然你们的异端教给你们一种对肉身的肉身性厌恶，却不能不爱自己的肉身，并关心其安全与舒适，既避免一切来自打击、跌倒、恶劣天气的伤害，又寻求保养健康的手段。所以，自然律比你们的虚假教义强大得多。\par

6. 再看肉身本身，难道我们在其生命器官的构造、形态的对称、行动肢体和感觉器官的位置与安排（这一切都和谐运作）中，在尺度的调整、数字的比例、重量的秩序中，看不到真天主的工程吗？关于祂，确实记载着：\BibleQuote{祢以度量、数目和重量安排了一切} \BibleRef{智 11:21}。如果你们的心没有被谎言硬化并败坏，你们就能从祂所造之物中，甚至在这些软弱的肉身受造物中，理解天主无形的事。因为谁是我所提及之事的作者，岂不是那独一者是所有尺度的标准、其智慧是一切美的模型、其法律是一切秩序的法则吗？如果你们对这些事视而不见，至少听听宗徒的话。\par

7. 因为宗徒，在论及丈夫对妻子的爱时，举了灵魂对肉身的爱为例。他说：\BibleQuote{爱自己妻子的，就是爱自己，因为从来没有人恨过自己的肉身，而是培养抚育它，如同基督对教会一样} \BibleRef{弗 5:28}。请看一切受造动物，在每一动物的本能自我保存中，你都能发现这一对自己肉身的爱的自然原则。不仅在人——他们若正当地生活，既为肉身的安全提供保障，又将肉欲置于理性运用之下——如此，就是在牲畜中，它们也躲避痛苦、畏惧死亡，并尽快逃离任何可能破坏其身体结构或解除精神与肉身联系的事物；因为牲畜也养育并珍爱自己的肉身。如宗徒所说：\BibleQuote{因为从来没有人恨过自己的肉身，而是培养抚育它，如同基督对教会一样}。请看宗徒从何处开始，又上升到何处。请思量，若可能，受造物从其造物主那里获得的伟大：它涵盖了从天上万军直到血肉的一切，以多样形态的美和依次递进的秩序。\par

8. 同一位宗徒再次论及神恩虽各有不同，却趋向和谐行动时，为了阐明如此伟大、神圣且奥秘的事，便以人体作比喻——从而清楚暗示，这肉身是天主的化工。整段经文见于\WorkTitle{致格林多人书}，与主题如此贴切，虽然很长，我认为全部引用也无妨：\BibleQuote{弟兄们，论到神恩，我不愿意你们不明白。你们知道，你们当外邦人的时候，被牵引迷惑，去拜那哑巴偶像。所以我告诉你们，凡受天主圣神感动的，没有说\InnerQuote{耶稣是可咒可诅的}；除非受圣神感动，也没有人能说\InnerQuote{耶稣是主}。神恩虽有区别，却是同一圣神；职分虽有区别，却是同一主；功效虽有区别，却是同一天主，在众人身上行一切事。圣神显现在每个人身上，是为了众人的好处。这人从圣神领受智慧的话，那人从同一圣神领受知识的话；有人从同一圣神领受信德，有人从同一圣神领受治病的恩赐；有人能行异能，有人能说先知话，有人能辨别神恩，有人能说各种语言，有人能解释语言。这一切都是同一圣神所运行，随自己的意愿分给各人。就如身体是一个，却有许多肢体；身体上的肢体虽多，仍是一个身体：基督也是这样。因为我们众人，不拘是犹太人或外邦人，或为奴的或自主的，都因一个圣神受洗成了一个身体，并都被饮于一个圣神。身体原不是一个肢体，而是许多。如果脚说：\InnerQuote{我不是手，所以不属于身体}，它因此就不属于身体吗？如果耳说：\InnerQuote{我不是眼，所以不属于身体}，它因此就不属于身体吗？若全身是眼，听觉在哪里？若全身是听觉，嗅觉在哪里？但天主照自己的意思，把肢体一一安置在身体上。若全是同一个肢体，身体在哪里？但如今肢体虽多，身体却只是一个。眼不能对手说：\InnerQuote{我不需要你}；头也不能对脚说：\InnerQuote{我不需要你}。相反，身体上那些看似软弱的肢体，更是必需的；那些我们认为不体面的肢体，我们越发给它加上体面；不雅观的肢体，越发显得雅观。我们体面的肢体，本不需要；但天主将身体配合在一起，把加倍的体面给了那有缺欠的肢体，免得身体有分裂，反而使各肢体彼此同样关怀。若一个肢体受苦，所有肢体都一同受苦；若一个肢体得荣耀，所有肢体都一同欢乐。}见：\BibleRef{格前 12:1} 完全撇开基督信仰不谈——它会使你们相信宗徒——如果你们有常识能觉察自明之理，就让每个人自己察看并看见宗徒所说之事的明显真理：最小的有何等伟大，最低的有何等良善；因为宗徒所赞扬的正是这些事，为的是借这些平常可见的身体物件，来阐明最高层次、不可见的属神实在。\par

9. 凡是否定我们的身体及其肢体（宗徒如此赞许和称扬）是天主手笔的人，你看他是在反驳谁，宣讲与你们所领受的相悖的道理。因此，与其驳斥他的意见，不如让他被所有基督徒诅咒。宗徒说：天主调和了身体。福斯特说：不是天主，而是\OriginalTerm{原质}{Hyle}。对于这样的矛盾，诅咒比辩论更合适。你不能说这里称天主为今世的神。如果有人理解出现这辞句的经文意指魔鬼蒙蔽不信者的心智，我们承认他确实借着恶毒的暗示如此行，人因屈从这些暗示，在天主公义的报应中失去正义的光明。这一切都符合圣经。宗徒亲自讲到外来的诱惑：\BibleQuote{我担心，正如蛇以狡猾欺骗了厄娃，你们的心也这样败坏，失去了在基督里的单纯和纯洁。}\BibleRef{格后 11:3}同样的话也有：\BibleQuote{败坏的风俗，败坏善行。}\BibleRef{格前 15:33}当他论到自欺的人时说：\BibleQuote{人若以为自己有什么，其实什么也没有，就是欺骗自己。}\BibleRef{迦 6:3}或又在之前已经引用的天主审判的话：\BibleQuote{天主就任由他们存邪僻的心思，去行那些不合理的事。}\BibleRef{罗 1:28}同样，在旧约中，\BibleQuote{天主并没有制造死亡，也不喜欢生灵丧亡}之后，我们读到：\BibleQuote{因魔鬼的嫉妒，死亡才进入了世界。}关于死亡，为使人们不将罪责推卸自己：\BibleQuote{恶人用手和声音招请死亡，将其视为朋友而陷入其中。}\BibleRef{智 1:16}然而，在别处却说：\BibleQuote{善与恶，生与死，贫穷与财富，都来自上主天主。}\BibleRef{德 11:14}这对于那些不明白的人似乎困惑，他们不知道，除了将来对每一恶行明显的审判之外，还有当下的审判；所以在一个行为中，除了欺骗者的诡计和自愿者的邪恶，还有审判者的公正刑罚：当魔鬼诱惑，人同意，天主就放弃。所以，如果你将\Quote{今世的神}连在一起，理解为魔鬼以恶毒的欺骗蒙蔽不信者，含义并不坏。因为\Quote{神}字并非单独使用，而是附有条件\Quote{今世的}，即指那些只求在今世亨通的恶人。在这个意义上，世界也被称为邪恶，经上记载：\BibleQuote{为救我们脱离这邪恶的现今世代。}\BibleRef{迦 1:4}同样，在\Quote{他们的肚腹是神}这句话中，只有与\Quote{他们的}一词相连时，肚腹才被称作神。同样，在圣咏中，若非加上\Quote{外邦的}，魔鬼不会被称为神。但在我们现在讨论的经文中，并没有说\Quote{今世的神}，或\Quote{他们的肚腹是神}，或\Quote{外邦的神是魔鬼}，而只是说\Quote{天主调和了身体}，这只能理解为真天主，万有的创造者。这里没有像其他情况那样贬低性的附加语。但也许福斯特会说，天主调和了身体，并非作为身体的制造者安排肢体，而是将自己的光混入其中。这样，福斯特将会把身体的构造和肢体的安排归功于其他而非天主，而天主则以祂的善性调和了身体的邪恶。摩尼教徒正是用这样的发明塞满软弱的心灵。但天主为了帮助软弱的人，借着圣经作者的口责备这种意见。因为我们在前几节读到：\BibleQuote{天主将每一个肢体都安置在身体上，如祂所愿意的。}显然，说天主调和了身体，是因为祂用许多肢体构造了身体，这些肢体在联合中保持着各自功能的多样性。\par

10. 难道摩尼教徒以为，按照他们狂野的想法，由\OriginalTerm{希勒}{Hyle}在黑暗族系中所构造的动物，在天主将祂的光与它们混合之前，没有宗徒所称赞的这种肢体和谐运作，以至于那时头对脚说，或眼对手说：\Quote{我不需要你}？这并非也不可能是摩尼教徒的教义，因为他们描述动物使用所有这些肢体，并说它们各自按其种类爬行、行走、游泳、飞翔。它们也能看见、听见，并使用其他感官，以适当的手段和工具滋养和照顾自己的身体。此外，它们还有繁殖能力，因为据说它们有后代。所有这一切，福斯图斯轻蔑地称之为\OriginalTerm{希勒}{Hyle}的作为，若没有宗徒所称赞并归于天主的和谐安排，是不可能做到的。现在难道还不清楚谁该被跟随，谁该被宣布为可咒骂的吗？事实上，摩尼教徒告诉我们，有些动物能说话；它们的言论被所有受造物，无论爬虫、四足兽、飞鸟还是鱼类，所听见、理解并赞同。惊人而超自然的雄辩！尤其因为它们没有文法家或演说家教导它们，也没有经历过藤条和桦木条的痛苦经验。哎呀，福斯图斯自己晚年才开始学习修辞学，以便能雄辩地谈论这些荒谬之事；尽管他聪明过人，但因学习而损害了健康，他的宣讲只赢得了寥寥无几的追随者。可惜他生在光明中，而非那黑暗区域！如果他曾在那里宣讲反对光明，整个动物界，从两足动物到百足虫，从龙到贝类，都会热切倾听并立即服从；然而，当他在这里宣讲反对黑暗族系时，他更常被称为有口才而非有学问，更常被称为最恶劣的假教师。而且在少数尊他为伟大教师的摩尼教徒中，他没有一个低等动物作为门徒；甚至连他的马也没有因主人的教导而变得更聪明，以至于部分神性的混合似乎只让动物更加愚蠢。这是何等的荒谬！这些受迷惑的众生何时才能有智慧，将摩尼教虚构中描述动物过去在自己区域的情形，与现在在这个世界的情形相比较？那时它们的身体强壮，现在却虚弱；那时它们的视力如此之好，以致因所见之美而被诱导入侵天主的区域，现在却弱得连太阳光线都抵挡不住；那时它们有足够的智慧理解讲给它们的道理，现在却没有这种能力；那时这种惊人而有效的雄辩是天生的，现在最贫乏的雄辩也需要勤奋的学习和准备。黑暗族系因善的混合而失去了多少善事啊！\par

11. \Person{福斯图斯}在我现在回应的评论中展示了他的才智，他为自己列出了一长串对立面——健康与疾病、富有与贫穷、白与黑、冷与热、甜与苦。关于黑与白，我们无需多言。或者，如果颜色中有善或恶的特性，以至于白色必须归于天主，黑色归于\OriginalTerm{希勒}{Hyle}；那么，如果天主在摩尼教所说的\OriginalTerm{希勒}{Hyle}创造鸟类时，将白色投在鸟的翅膀上，那么当天鹅变白时，乌鸦到哪里去了？冷与热也无需讨论，因为两者适度则善，过度则危险。至于其余，\Person{福斯图斯}可能意图将善与恶（他本可将其放在首位）理解为包含其余，因此健康、富有、白、热、甜应属于善；而疾病、贫穷、黑、冷、苦属于恶。这种看法之无知和愚蠢是显而易见的。如果我分别讨论白与黑、热与冷、甜与苦、健康与疾病，那可能会显得是辱骂。因为如果白和甜都是善，黑和苦是恶，那么为什么大多数葡萄和所有橄榄在变甜时变成黑色，从而通过获得恶而得到善？又如果热和健康都是善，冷和疾病是恶，为什么身体在受热时会生病？发烧难道是健康的吗？但我放过这些，因为它们可能出于仓促，或者只是作为对立的例子给出，并非真的善与恶，尤其因为从未说过黑暗族系中的火是冷的，所以此处的热无疑必须是恶。\par

12. 那么，我们接着谈健康、富有和甜美，\Person{福斯图斯}在他的对立中显然视之为善。难道在黑暗族系中就没有身体的健康吗？那里动物出生、成长并繁殖，有如此活力，以至于有些怀孕的动物被抓走（如传说所述）并被捆绑在天上时，甚至早产的流产后代从天上坠落到地上，却仍然活着、成长，并产生了现在存在的无数种动物？难道那里没有财富吗？那里树木不仅能在水和风中生长，还能在烟和火中生长，并能结出如此丰富的果实，以至于动物按各自种类从果实中生出，并从那些肥沃的树木中获得生计，并通过众多后代证明它们营养良好？而所有这一切没有耕种的劳苦，没有从夏到冬的严酷变化，因为没有太阳以年度运行给季节带来变化。必定有持久的丰产，那里树木不仅在自己元素中出生，而且有适当的养分供应使它们持续肥沃；正如我们所见，橙树如果浇水充足，全年结果。财富必定充足，且免受伤害；因为当没有你们寓言中驱动雷电的聚光者时，就不必担心冰雹。\par

13. 黑暗种族中的生灵也不会寻求食物，除非食物是甘甜愉悦的，否则他们会因缺乏而死亡。因为我们发现所有身体都有其特定的需要，据此食物或是合宜的或是令人生厌的。如果合宜，就被称为甘甜或愉悦；如果令人生厌，就被称为苦或酸，或以某种方式令人不快。我们发现在人类中，一个人渴望的食物另一个人却厌恶，这是由于体质、习惯或健康状态的差异。更有甚者，构造完全不同的动物可以在我们觉得不适口的食物中找到乐趣。否则山羊何以如此急切地啃食野橄榄？这种食物对它们是甘甜的，就像在某些疾病中蜂蜜对我们来说是苦的一样。对于深思的探究者，这些事暗示了安排的优美，其中每个事物都找到适合它的东西，以及从最低到最高、从物质到精神的善的伟大。至于黑暗种族，如果从某元素生出的动物以该元素所产之物为食，毫无疑问，食物必然因其适口而甘甜。再者，如果这动物找到另一元素的食物，不适口就会在味觉上表现为令人不快。这种不快被称为酸、苦、不适口或类似物；或者如果其有害性质足以破坏身体构造的和谐，从而夺去生命或削弱力量，则被称为毒药，仅仅因为这种不适口，而它可能滋养与其相宜的生命。因此，如果一只鹰吃了我们日常食用的面包，它会死；而我们如果吃了藜芦（牛常吃的草，且其某种形式可用作药物）也会死。如果福斯图斯知道或想到这一点，他就不会以毒药和解毒剂作为善恶两种本性的例子，好像天主是解毒剂，而\OriginalTerm{原始物质}{Hyle}是毒药。因为同一事物，同一种本性，按适当或不适当的使用会致死或治愈。在摩尼教的传说中，他们的神可以说是黑暗种族的毒药；因为他如此损害他们的身体，以致他们从强壮变得极其虚弱。但另一方面，由于光明本身被夺取并遭受损失和伤害，它也可以说是自身的毒药。\par

14. 你们似乎不是建立了一个善的本原和一个恶的本原，而是使两者皆善或皆恶，或者更确切地说，两个善和两个恶；因为它们在自身内是善的，而彼此相对是恶的。我们稍后可以看出哪一个更好或更坏；但与此同时，我们可以认为它们在自身内都是善的。因此天主在一区域统治，而\OriginalTerm{原始物质}{Hyle}在另一区域统治。两个王国都有健康，都有丰富的出产；两者都有众多的后裔，都品尝了与其本性相宜的快乐的甘甜。但摩尼教徒说，黑暗种族除了与毗邻的光明相对的邪恶部分外，对它自身也是邪恶的。然而，既然我已经指出了其中许多善的事物，如果你能指出它的恶，仍然有两个善的王国，只不过没有恶的那个是两者中较好的。那么，你称之为它的恶的是什么呢？据福斯图斯说，\Quote{他们彼此掠夺、杀害和吞噬。} 但如果他们不做别的，这么多的大军如何能出生并长大成熟？他们也必然享受过和平与安宁。但是，即使承认没有纷争的王国是两者中较好的，它们仍然都应被称为善，而不是一个善一个恶。因此较好的王国是既不杀害自身也不彼此杀害的；较坏或较不善的王国是其中虽然彼此争斗，但每个单独的动物保持其本性的健康与安全。但我们不能在你的神和黑暗之君之间做出太大区别，没有人反对他，他统治被所有人承认，他的提议被一致同意。所有这些都意味着伟大的和平与和谐。那些所有的人都衷心服从国王的王国是幸福的。此外，这位君主的统治不仅延伸到他的同类，或延伸到你所创造的人类父母的双足生物，而且延伸到各种动物，它们等候在他的面前，服从他的命令，相信他的宣言。你以为人们会如此愚蠢，以至于认不出你对这位君主的描述中的神性属性，或者认为你有可能有另一个神吗？如果这位君主的权威建立在他的资源上，他一定非常强大；如果建立在他的名声上，他一定声名显赫；如果建立在他的爱上，他一定受到普遍尊重；如果建立在他的恐惧上，他一定维持了最严格的秩序。那么，如果一些邪恶与这么多善的事物混合在一起，谁懂得词语的意义会称之为邪恶的本性呢？此外，如果你称之为邪恶的本性，因为它不仅对另一种本性是恶的，而且在自身内也是恶的，你认为在你的神与对立本性混合之前，在他的可怕必然性中就没有恶吗？他被强迫与之战斗，并如此无情地派遣自己的成员被吞没，以致于完全康复无望。这是那种本性在与被你允许为恶的唯一事物混合之前的巨大恶。你的神要么有能力不受黑暗种族的伤害和玷污，在这种情况下他自己的愚蠢必定使他陷入困境；要么如果他的实体易于腐败，你所崇拜的对象就不是宗徒所说的不可腐朽的天主。\BibleRef{弟前 1:17}那么，这种腐败的可能性，即使不考虑实际经验，对你来说不就是你的神中的一种恶吗？\par

15. 显然，要么他缺乏预见——这在天主而言确是重大缺陷，不知将来之事；要么若他有预见，则他绝不可能感到安全，而必常怀恐惧，此乃严重恶事，你必承认。时刻都可能有这样的恐惧：那场冲突的时刻已到，他的肢体遭受如此损失和玷污，以致要解放并净化它们需耗费无限劳力，且最终只能部分完成。若将此惊恐状态归于天主自身过于过分，那么至少他的肢体必曾畏惧遭受这一切恶事。再者，若他们对将要发生之事一无所知，则你们的天主的本质必在某种程度上缺乏预见。你在你的至善中计有多少恶事？或许你会说，他们无所畏惧，因为他们既预见受苦，也预见他们的解放和胜利。但他们仍必为其同伴担忧，若他们知道他们将被逐出王国，永远束缚于黑暗之团中。\par

16. 难道他们没有爱德去同情那些注定要承受永罚、却没有犯任何罪的人吗？这些将要被捆绑在团块中的灵魂，难道不也是你们神的一部分吗？它们不是同一来源、同一实体吗？至少它们对自己将来永远的束缚应当感到悲伤或恐惧。若说它们不知道将要发生的事，而其他灵魂知道，那就是说同一实体部分预知未来，部分无知。如果在这实体与恶原则混合之前就已经有如此之恶，你们怎能称它为纯洁、完美、至高的善呢？你们必须承认你们的两原则要么都是善的，要么都是恶的。如果你们设定两个恶，你们可以随意称其中哪个更恶。但如果你们设定两个善，我们就需要探究哪个更善。与此同时，你们关于一善一恶两原则的教义——实际上就是两个神，一善一恶——就终结了。然而，若伤害他人是恶，那么两者互相伤害。也许首先发起攻击的原则是更大的恶。但若一方先造成伤害，另一方则以更严厉的方式报复，而非你们愚昧指责的\Quote{以眼还眼}的补偿法则。你们必须选择称哪一个更恶——是首先伤害的一方，还是有意愿和能力造成更大伤害的一方。一方试图分享光明之乐；另一方则彻底推翻其对手。若一方得到了它想要的，它当然不会伤害自己。而另一方在挫败其对手时，却对自身的一部分造成了巨大伤害；这使人想起那句众所周知的激愤之辞，据记载确实有人说过：\Quote{愿我们的朋友灭亡，如果这能让我们摆脱敌人。}因为你们神的一部分被派去遭受绝望的污染，以便为永远埋葬敌人的团块提供遮盖。即使在被征服和捆绑之后，他仍将如此被畏惧，以至于神性一部分的安全——尽管有限——必须以其他部分的永苦为代价。这就是善原则的无害性！你们的神看来也犯了你们指控黑暗种族所犯的罪——同时伤害朋友和敌人。从最终囚禁其敌人而自己的子民也卷入其中的团块来看，对你们神的指控是成立的。事实上，你们称为神的原则对朋友和敌人更具伤害性。就\OriginalTerm{质料}{Hyle}而言，它并没有毁灭对立王国的欲望，只想占有它；尽管其一些子民被其他子民的暴力杀死，但他们会以其他形式再次出现，因此在生死交替中，他们历史上有享受的间歇。而你们的神，以你们归于他的全能和完善卓越，却判定敌人永死，朋友永罚。而疯狂之极在于相信内在争斗伤害了\OriginalTerm{质料}{Hyle}的成员，而胜利却给神的成员带来惩罚。这种愚蠢意味着什么？用福斯图斯对神和\OriginalTerm{质料}{Hyle}的比喻——解药和毒药——来看，解药似乎比毒药更有害。我们没有听说\OriginalTerm{质料}{Hyle}把神永远关在黑暗团块中，或把自己的成员赶入其中；更没有听说它诽谤这不幸的残余，作为不净化它的借口。因为摩尼在其\WorkTitle{基础书信}中说，这些灵魂理应受此惩罚，因为它们允许自己被诱离原初的光明，成为圣光的仇敌；而实际上是神自己派它们迷失于黑暗区域，以便以光对抗光：如果神强迫它们违背意愿，那是不义的；如果它们愿意去，神惩罚它们就是忘恩负义。这些灵魂绝不可能幸福，若它们在冲突前因知道将成为原初原则的仇敌而恐惧，冲突中绝望地被污染，之后又被永远定罪。另一方面，它们绝不可能曾是神圣的，若它们在冲突前因缺乏预知而不知将要发生什么，冲突中表现软弱，之后遭受痛苦。对它们如此，对神也是如此，因为它们同实体。你们是否还有希望看到这些亵渎的愚蠢？你们确实试图通过断言\OriginalTerm{质料}{Hyle}被囚禁后不再能伤害自己来证明神的良善。看来\OriginalTerm{质料}{Hyle}将得到一些善，即当它不再有任何善混合时。也许，正如冲突前神有恶的必然性（当善未与恶混合时），冲突后\OriginalTerm{质料}{Hyle}将有休息的善（当恶未与善混合时）。这样，你们的两原则要么是两个恶，其中一个更恶；要么是两个善，都不完美，但一个比另一个更善。然而，较善者更可怜；因为若这场大冲突的结果是敌人因\OriginalTerm{质料}{Hyle}中相互伤害停止而得到一些善，而神自己的子民却遭受被赶入黑暗团块的严重恶，那么我们要问谁得了胜利。我们要理解，毒药是\OriginalTerm{质料}{Hyle}，其中动物生命却找到了丰富的生长和繁殖手段；而解药是神，他能定罪自己的成员，却不能恢复他们。事实上，称一个为\OriginalTerm{质料}{Hyle}和称另一个为神同样荒谬。这些都是不能忍受纯正道理而转向无稽之谈的人们的愚蠢。\par

\SourceRecord{Translated by Richard Stothert.；Nicene and Post-Nicene Fathers, First Series；Vol. 4.；Edited by Philip Schaff.；Buffalo, NY: Christian Literature Publishing Co.,；1887.；<http://www.newadvent.org/fathers/140621.htm>.}

% structure: p0001 source=h1 target=section reason=page-title

\section{反福斯图斯，第二十二卷}

\Emph{福斯图斯陈述他对法律和先知之道德的反对意见，而奥斯定试图通过预象和寓意法的应用来消解旧约的道德困难。}\par

1. \Emph{\Person{福斯图斯}}说：我们远非宣称或感到对法律和先知的任何敌意，以至于我们准备，若你允许我们，宣布所有使法律和先知看起来可憎的著作都是虚假的。但你拒绝承认这一点，通过维持你的作者们的权威，你给先知带来了或许是不应有的责备；你诽谤圣祖，并侮辱法律。你如此不合理，以致否认你的作者们是虚假的，同时你却维护那些被这些著作描述为犯有最恶劣罪行、过着邪恶生活者的虔诚与圣洁。这些观点是不一致的；因为要么这些人是坏品性，要么作者们是不诚实的。\par

2. 那么，假设我们同意谴责这些作者，我们也许就能成功地辩护法律和先知。通过法律必须被理解为不是割礼、安息日、祭献或其它犹太礼仪，而是真法律，即：\Quote{你不可杀人}，\Quote{你不可奸淫}，\Quote{你不可作假见证}，等等。对于这部法律——它在全世界颁布，即是在世界现存秩序的起初——希伯来作者们施加了暴力，用他们关于割礼和祭献的可憎规诫所携带的污秽感染了它。作为法律的朋友，你应当与我一同谴责犹太人将不适合的规条混合进去从而损害了法律。显然，你必定意识到这些规诫不是法律，也不是法律的任何部分，因为你声称自己是义人，虽然你并不试图遵守这些规诫。在寻求过义人的生活时，你非常重视禁止罪恶行为的诫命，却对犹太礼仪毫不留意；如果它们同属一部法律，这便是不可原谅的。你将被称为疏忽\Quote{你不可杀人}或\Quote{你不可奸淫}的规诫视为一种可耻的责备。如果你同样将被称为未受割礼或疏忽安息日视为可耻的责备，那么显然你认为两者都是天主的法律和命令。然而事实上，你认为遵守其中一项的尊荣与光荣绝不因忽视另一项而受损。正如我所说的，显然这些礼仪不是法律，而是对法律的毁损。如果我们谴责它们，并非视之为真实的，而是视之为虚假的。在这谴责中，既无对法律或其作者天主的责备，而只有对那些以这些名义发表其骇人听闻迷信者的责备。如果我们有时在攻击犹太规条时滥用了法律的可敬名称，那是你们的过错，因为你们拒绝区分希伯来礼仪与法律。只要你们除去这些肮脏的以色列污点，恢复法律应有的尊严；承认这些作者犯有毁损法律之罪，你们就会立刻看到，我们不是法律的敌人，而是犹太主义的敌人。你们被\Quote{法律}这个词误导了；因为你们不知道这个名字原本属于什么。\par

3. 就我而言，我看不出你有任何理由认为我们亵渎了你的先知和圣祖。如果我们直接或间接地是记载他们行为叙述的作者，那么这项指控确实有根据。但由于这叙述要么是他们自己写的——出于一种因恶行而成名的罪恶欲望——要么是他们的同伴和同时代人所写，你为何要责怪我们？你因憎恶那些他们自愿承认有罪的邪恶行为而谴责他们，尽管并无必要作这样的承认。或者，如果这叙述只是恶意的虚构，那就让它的作者受惩罚，让这些书被谴责，让先知之名从这可耻的责备中澄清，让圣祖恢复因其作为管理者的纯朴与纯洁而应得的尊重。\par

4. 此外，这些书含有对天主本身的骇人诽谤。我们被告知，他从永远存在于黑暗之中，并在他看见光时赞美光；他对未来如此无知，以至于给\Person{亚当}一条命令，却没有预见它会被人违反；他的感知如此有限，以至当\Person{亚当}因认识自己的赤身而躲在\Place{乐园}一角时，他看不见\Person{亚当}；嫉妒使他害怕，唯恐他的受造物人尝了生命树上的果子，而永远活着；之后他贪求血和各样祭献的脂油，并且如果它们献给除他以外的任何人，他便嫉妒；他有时对自己的敌人发怒，有时对自己的朋友发怒；他因轻微的过犯或无端毁灭成千上万的人；他威胁要带着剑来，不饶恕任何人，无论义人还是恶人。这些对天主进行如此大胆诽谤的作者，很可能也诽谤了天主的人。如果你想为先知辩护，你必须与我们一同把责备放在这些作者身上。\par

5. 再者，对于以下关于亚巴郎的事，我们不承担责任：他因不理性的渴望生子，且不相信天主应许他的妻子撒辣将生子，便在妻子知情的情况下与妾通奸，这使事情更糟；\BibleRef{创 16:2} 或者，他因贪婪和贪欲，屡次亵渎婚姻，将妻子撒辣出卖给阿彼默肋客王和法郎王，供他们淫乐，谎称她是妹妹，因为她十分美貌。这记述不是我们的，其中记载亚巴郎的兄弟罗特从索多玛逃出后，在山中与两个女儿同寝\BibleRef{创 19:33}（他宁可死于索多玛的火灾，也不该被乱伦的情欲焚烧）；或者依撒格如何效法父亲的行为，称妻子黎贝加为妹妹，以便靠她获得可耻的生计；\BibleRef{创 26:7} 或者他的儿子雅各伯，有四个妻子——两亲姐妹辣黑耳和肋阿，以及她们的婢女——像公山羊一样在她们中间生活，以致女人们每日争抢，谁该在田间归来时首先抓住他，有时甚至互相雇用他过夜；或者，他的儿子犹大如何与儿媳塔玛尔同寝，塔玛尔先后嫁给了他的两个儿子，据说犹大被她以妓女装扮所骗，因为塔玛尔知道公公常与这类女子交往；\EditorialNote{创世纪第三十八章} 或者，达味在已有众多妻子的情况下，勾引士兵乌黎雅的妻子，并致使乌黎雅阵亡；\BibleRef{撒下 11:4} 或者，他的儿子撒罗满有三百个妻子、七百个嫔妃，以及无数公主；\BibleRef{列上 11:1} 或者，首位先知欧瑟亚从妓女得子，更糟的是，据说这可耻的行为是天主所命；\BibleRef{欧 1:2} 或者，梅瑟杀人\BibleRef{出 2:12}、掠夺埃及\BibleRef{出 12:35}、发动战争，并命令或亲自施行许多残酷行为。\BibleRef{出 17:9} 他也不满足于一个妻子。对于这些及类似记载于族长和先知书中的叙述，我们既非直接也非间接的作者。要么是你们的作者伪造了这些事，要么是先祖们确实有罪。请任选其一；无论哪种，罪行都可憎，因为恶行和谎言同样可憎。\par

6. \Person{奥古斯丁} 回答道：你们既不明白法律的预象，也不明白先知的行为，因为你们不知道什么是圣洁或正义。我们已多次详尽指出，旧约的规诫和预象既包含通过新约所赐恩宠而因服从得以成全的内容，也包含因在现已显明的真理中得以成全而应予废除的内容。因为爱天主和爱近人确保了法律规诫的实行，而割损及其他昔日实行的预表性礼仪的废除则显示了法律预许的应验。规诫使人因罪责感而渴望救恩；预许则使人通过预表性礼仪确信救主将要来临。所渴望的救恩要通过新约显现时所赐的恩宠获得；期待的应验使预象不再必要。由梅瑟所传的法律，在耶稣基督内成了恩宠和真理。借着赦罪的恩宠，规诫在那些蒙受神助的人身上继续有效。借着真理，象征性礼仪被废除，以使预许在那些信赖天主信实的人身上得以成就。\par

7. 因此，那些因不理解而指责法律之预表性制度为畸形与赘疣的人，如同不晓其用而厌弃事物之人。好比聋子见别人开口说话，便指责嘴唇的动作多余难看；又好比盲人听人称赞房屋，便用手触摸墙壁表面来验证，当摸到窗户时，便指责那是墙面的缺陷，或者以为墙壁塌陷了。\par

8. 我如何能使那些满心虚妄的人明白，先知的行为也是奥秘而预言的？他们心中的虚妄表现为：他们认为我们相信天主曾一度处于黑暗之中，因为经上记载：\BibleQuote{黑暗笼罩在深渊之上。}\BibleRef{创 1:2} 仿佛我们将那曾有黑暗的深渊称为天主，因为天主用祂的圣言创造光明之前，那里并无光明。由于他们不能区分那本身就是光明的天主与天主所创造的光明，他们便认为天主在创造光明之前必在黑暗中，因为黑暗笼罩在深渊之上，直到天主说：\Quote{有光！就有了光。}在新约中，这两者都归于天主。因为我们读到：\BibleQuote{天主是光，在祂内毫无黑暗；}\BibleRef{若一 1:5} 又：\BibleQuote{那吩咐\InnerQuote{光从黑暗中照耀}的天主，曾照耀在我们心中。}\BibleRef{格后 4:6} 同样在旧约中，天主的智慧被称为\BibleQuote{永光的辉耀}\BibleRef{智 7:26}，这智慧当然不是受造的，因为万物都是藉着它受造的；至于那仅作为这智慧之产物而存在的光明，经上说：\Quote{上主，你必点亮我的灯；我的天主，你必照明我的黑暗。}同样，起初当黑暗笼罩在深渊之上时，天主说：\Quote{有光！就有了光。}这光是那赐予光明的光明——即天主本身——所造的。\par

9. 因为天主既是祂自己的永恒幸福，又是幸福的赐予者；同样，祂既是祂自己的永恒光明，又是光明的赐予者。祂不嫉妒任何善，因为祂本身就是一切善的幸福之源；祂不惧怕任何恶，因为一切恶者的丧失都在于被祂抛弃。祂既不能从祂所赐予幸福的人那里获益，也不惧怕那些祂自己的审判所判定的悲惨结局。啊，摩尼，你所崇拜的对象大不相同。你已离开天主，追求你自己幻想出来的东西，这些东西在你愚蠢而游荡的心中不断增多繁殖，通过视觉吸收天体的光。这光虽也是天主所造，却不能与虔诚者心灵中所造的光相比——天主使他们出黑暗入光明，正如祂使他们出罪恶入正义。它更不能与那不可接近的光相比，一切光都源于那光。这光并非对所有人不可接近；因为\BibleQuote{心里洁净的人是有福的，因为他们要看见天主。}\BibleRef{玛 5:8} \BibleQuote{天主是光，在祂内毫无黑暗；}但恶人不得见光，如依撒意亚所说。\BibleRef{依 8:20} 对他们而言，这赋予生命的光是不可接近的。从这光不仅产生虔诚者心灵中的属灵之光，也产生物质之光，它不拒绝恶人，而是上升照耀恶人和善人。\par

10. 所以，当黑暗笼罩深渊时，那光本身说：\BibleQuote{有光。}这光来自什么光，是清楚的；因为经文说：\BibleQuote{天主说。}那被造的光是什么，并不那么清楚。因为圣经学者之间曾有过友好的讨论，究竟天主当时是在天使的心灵中造了光——换句话说，就是这些理性精神本身——还是在宇宙较高区域中某种我们无法知晓的物质光。因为在第四天，祂造了天上的可见光体。还有一个问题：这些天体是否与它们的光同时被造，还是从已造的光中点燃而成？但任何人以理解圣经所需的虔诚精神阅读圣经，必须确信：当黑暗笼罩深渊时，天主说\BibleQuote{有光}，无论那被造的光是什么，都是受造之光，而创造之光则是它的制造者。\par

11. 也不能因此断定，天主在造光之前曾居于黑暗中，因为经文说黑暗笼罩深渊，然后天主的圣神在水面上运行。深渊是深不可测的水渊。属肉体的心可能会以为圣神曾居于覆盖深渊的黑暗中，因为经文说祂在水面上运行。这是因为不明白光如何在黑暗中照耀，而黑暗却不接受光，直到藉着天主的话，那些曾是黑暗的人成为光，并对他们说：\BibleQuote{你们从前曾是黑暗，但现在在主内是光明。}\BibleRef{厄 5:8} 但如果因罪恶意志而在黑暗中的理性心灵，尽管光无处不在，却不能理解天主智慧之光，因为它们在心性上而非空间上与之分离：那么天主的圣神在水面上运行时，虽与水的黑暗有不可测量的距离（不是空间上的，而是本性上的），为什么不能在水面上运行呢？\par

12. 我知道，这一切我是对聋子唱歌；但主，我们所讲论的真理由祂而来，可以开启一些耳朵来听这曲调。但我们将如何评说那些批评圣经的人，他们反对天主喜悦自己的作品，并指责\BibleQuote{天主看见光是好的}，好像这意味着天主像欣赏什么新奇之物一样赞赏光？天主看见自己的工作为好的，意味着造物主认可自己的作品，认为它们使自己喜悦。因为天主不能被强迫做任何违背祂意愿的事，以致祂不喜悦自己的工作；祂也不能出错，以致后悔做了某事。为什么摩尼人反对我们的天主看见祂的工为好的，而他们的神在将自己肢体与黑暗混合时，却在自己面前放置了一个遮盖物？因为他不看自己的工作为好的，反而因它是邪恶而拒绝看它。\par

13. 福斯图斯说我们的天主感到惊奇，但圣经并没有这样说；也不意味着一个人看见某物为好就必须惊奇。有许多善事我们看见却不惊奇，仿佛它们比我们预期的更好；我们只是认可它们符合它们应有的样子。然而，我们可以给出一个天主惊奇的例子，不是来自旧约——摩尼人以不当的指责攻击旧约——而是来自新约，他们声称相信新约以诱捕粗心的人。因为他们承认基督是天主，并以此作为诱饵引诱基督的跟随者陷入他们的圈套。因此，当基督惊奇时，天主也惊奇。因为我们在福音中读到，当基督听到某位百夫长的信德时，祂感到惊奇，并对祂的门徒说：\BibleQuote{我实在告诉你们，在以色列我也没有见过这样大的信德。}\BibleRef{玛 8:10} 我们已经解释过\BibleQuote{天主看见那是好的}这句话。更好的人可能会有更好的解释。与此同时，让摩尼人来解释基督为何对祂预见会发生、在听到之前就已经知道的事感到惊奇。因为虽然看见一件事是好的与对之惊奇大为不同，但在这种情况下有一些相似之处，因为耶稣是对祂自己在百夫长心中所创造的信德之光感到惊奇；因为耶稣是真光，照亮每一个来到世上的人。\par

14. 因此，一位不虔敬的外邦人可能会用同样的责难来反对福音中的基督，正如福斯图斯用旧约中的天主来反对一样。他可能会说基督缺乏远见，不仅因为祂对百夫长的信心感到惊奇，而且因为祂选择了犹大作为门徒，而此人不听从祂的命令；正如福斯图斯反对在乐园中给出的诫命，结果并未被遵守一样。他还可能对基督不知道谁摸了祂吹毛求疵，当患血漏的妇人摸了祂的衣穗时；正如福斯图斯责怪天主不知道亚当躲在哪里。如果这种无知隐含在天主所说的\BibleQuote{亚当，你在哪里？}\BibleRef{创 3:9}中，那么基督问\BibleQuote{谁摸了我？}\BibleRef{路 8:44}也可同样说。外邦人还可能称基督胆怯和嫉妒，因为祂不愿让十个童女中的五个通过进入祂的王国获得永生，并把她们关在门外，以至于她们徒劳地敲门请求开门，仿佛忘记了祂自己的应许\BibleQuote{敲门，就给你们开门}\BibleRef{玛 7:7}；正如福斯图斯指责天主恐惧和嫉妒，不允许人在犯罪后获得永生。此外，他可能称基督贪恋人的血，而非动物的血，因为祂说\BibleQuote{凡为我而丧失生命的，必得着生命直到永生}\BibleRef{玛 10:39}；正如福斯图斯在关于那些预表我们得赎的流血祭献的动物牺牲上指责天主。他还可能指责基督嫉妒，因为在叙述祂将买卖人赶出圣殿时，福音作者引用了适用于祂的话\BibleQuote{我对祢殿的热心把我耗尽}\BibleRef{若 2:17}；正如福斯图斯指责天主嫉妒，禁止向其他神祇献祭。他可以说基督对祂的朋友和敌人都发怒：对朋友，因为祂说\Quote{那知道主人旨意而不去做的仆人，必多受鞭打}；对敌人，因为祂说\BibleQuote{凡不接待你们的，你们出去的时候，跺掉脚上的尘土；我实在告诉你们，在审判的日子，所多玛所受的比那城还容易忍受}\BibleRef{玛 10:14}；正如福斯图斯指责天主有时对朋友发怒，有时对敌人发怒；宗徒论及这两类人说：\BibleQuote{凡没有律法而犯罪的，也必不带律法而灭亡；凡在律法之下犯罪的，必因律法受审判}\BibleRef{罗 2:12}。或者他可以说基督毫无怜悯地流了许多人的血，为了一点小过犯或无缘无故。因为对外邦人来说，在婚宴上没有穿礼服似乎没什么大碍或根本无害，而我们的君王在福音中因这事命人把那人手脚捆起来，丢到外面的黑暗中；或者不愿基督作王，这是基督所说的罪：\BibleQuote{那些不要我作他们君王的人，把他们带到这里来，在我面前杀掉}\BibleRef{路 19:27}；正如福斯图斯在旧约中指责天主因轻微过犯（如福斯图斯所称）或无缘无故屠杀了成千上万的人。再者，如果福斯图斯指责天主威胁要带刀来，不饶恕义人或恶人，难道外邦人不会同样指责宗徒保禄的话吗？当他论及我们的天主时说：\BibleQuote{祂连自己的儿子都没有怜惜，反而为我们众人把祂交出}\BibleRef{罗 8:32}；或者伯多禄的话，当他劝勉圣人们在逼迫和屠杀中忍耐时说：\BibleQuote{时候到了，审判要从天主的家开始；如果先从我们开始，那些不信主福音的人的结局将如何？如果义人仅仅得救，那不信主和罪人将出现在哪里？}\BibleRef{伯前 4:17}。有什么比独生子更义的？然而圣父并没有怜惜祂。还有什么比这更明白的：义人也不被怜惜，反而用许多患难来鞭策，正如\BibleQuote{义人仅仅得救}的话清楚暗示的？正如在旧约中所说：\BibleQuote{主所爱的，祂必管教；凡收纳的祂都鞭打}\BibleRef{箴 3:12}；\BibleQuote{我们从主手里得福，不也受祸吗？}\BibleRef{约 2:10}。同样，我们在新约中也读到：\BibleQuote{凡我所爱的，我就责备管教}\BibleRef{默 3:19}；\BibleQuote{我们若是先审判自己，就不至于受主的审判；但受审判时，乃是受主的管教，免得我们和世人一同被定罪}\BibleRef{格前 11:31}。如果一位外邦人对新约提出这样的异议，摩尼教徒岂不会试图回答它们？尽管他们自己对旧约提出类似的异议。但假如他们能够回答外邦人，那么为了一部圣经辩护却对另一部圣经吹毛求疵是多么荒谬！但如果他们不能回答外邦人的异议，他们为什么不允许在两部圣经中（而不是仅一部）承认：凡因无知而显得错谬的事，虔敬的读者即使也不明白，却应相信是正确的？\par

15. 也许我们的对手会坚持认为，这些从新约引用的平行经文本身既无权威也不真实：因为他们拥有不敬的自由，主张和教导凡他们认为有利于自己异端的话都由基督和宗徒所说；同时他们无耻地宣称，同一著作中凡不利于他们的话都是伪造的插入。我已经在本着当前作品意图所需的篇幅内，相当详细地揭露了这种对整部圣经权威的攻击的不合理性。\par

16. 目前我要提请注意这个事实：当摩尼教信徒虽然以基督教之名掩盖其亵渎的荒谬，却对基督宗教的圣经提出这类的反对意见时，我们不得不既针对外教人也针对摩尼教信徒，捍卫两部圣约中天主之记载的权威。外教人可能像福斯图斯指责旧约中所谓对天主不体面的描绘那样，挑剔新约中的段落；而引用其本族作者中的类似段落，如保禄在雅典的演讲那样，即可答复外教人。\BibleRef{宗 17:28} 即使在外教人的著作中，我们也能找到这样的教义：天主创造并建构了世界，祂是光的赐予者，这并不意味着在光被造之前祂居于黑暗；当祂完成工程后，祂充满了喜乐，这远胜于说祂看到那是好的；祂制定法律，顺从有赏，违逆有罚，但他们并不因此说天主对未来无知，因为祂把法律赐给了那些将要违反它的人。他们也不能把发问当作缺乏预见力的证据，即便对人类也是如此；因为在他们书籍中，许多问题提问的目的只是为了用对方的回答来说服受话的人：提问者不仅知道他想要什么答案，而且知道实际上会得到什么回答。再者，如果外教人试图证明天主嫉妒任何人，因为祂不愿把幸福赐给恶人，他会在其本族作者的著作中找到许多支持这一神圣治理原则的段落。\par

17. 外教人关于祭献可能提出的唯一反对意见，是指责我们既然批评外教人的祭献，为何在旧约中天主却要求人向祂奉献祭献。如果我要详细答复这个问题，我可以向他证明：祭献只应归于唯一真神；而这祭献是由唯一真司铎、天主与人之间的中保所奉献的；并且这祭献应以动物祭献为预像，以预示那将为赦免因血肉所犯的罪而献上的唯一祭献的血肉——血肉不能承受天主的国：因为属血气的身体将获得属天的属性，正如祭献中的火象征了死亡在胜利中被吞灭。那些礼仪本属于那个民族，其王国和司祭职预示着那位将来的君王和司铎，祂要治理并圣化万民中的信徒，带领他们进入天国、天使的圣团和永生。正如这真正的祭献给虔敬地呈现在希伯来人的礼仪中，同样又被外教人亵渎地歪曲，因为正如宗徒所说：\BibleQuote{他们所奉献的是献给魔鬼，而不是献给天主}。\BibleRef{格前 10:20} 祭献中流血仪式的典型从最早的时代就开始了，从人类历史的起点就指向中保的苦难。因为圣经提到亚伯尔是第一个奉献这种祭献的人。\BibleRef{创 4:4} 因此，我们不必惊讶那些占据空中的堕落天使，其主要罪恶是骄傲和说谎，会向那些他们希望被视为神明的崇拜者要求他们知道只应归于天主的东西。人心的愚昧助长了这种欺骗，尤其是因对死者的追念而制造了肖像，进而使用了偶像。\BibleRef{智 14:15} 因此，祭献只应归于天主的本质不仅在天主正当地要求时显明，也在假神骄傲地僭取时显现。如果外教人对这些事难以相信，我应当从预言论证，并指出尽管预言是许久以前发出的，如今已应验。如果他仍然不相信，这与其说是奇怪，不如说是意料之中；因为预言本身就暗示了并非所有人都相信。\par

18. 倘若外教人继而指责两部圣约书将嫉妒归给天主和基督，那他只不过暴露了自己对文学的愚昧无知，或是健忘。因为尽管他们的哲人区分欲望与激情、喜悦与满足、谨慎与恐惧、温和与柔肠、明智与狡猾、勇敢与鲁莽等等，将每对中的第一个名称赋予善，第二个赋予恶，但他们的书籍仍然充满实例，因滥用这些词语，德行被冠以原本属于恶习的名称，例如激情用于欲望，满足用于喜悦，恐惧用于谨慎，柔肠用于温和，狡猾用于明智，鲁莽用于勇敢。言语中类似的不精确之处不可胜数。此外，每一种语言都有自己的习语。因为在宗教著作中，我不记得有将柔肠一词用于坏意义的实例。而日常用法中也有类似的词语使用特性。在希腊语中，一个词代表两件不同的事：劳作和痛苦；而我们则各有不同的名称。同样，我们用一个词表示两种意思，例如我们称未死之物有生命；又例如我们说某人生活端正，而在希腊语中，每个意义各有其词。所以，撇开所有语言中普遍存在的词语滥用不谈，用嫉妒一词表示两种意义可能是希伯来习语：一种称丈夫嫉妒，是因妻子不忠而痛苦所导致的心态病态，此义不能用于天主；另一种意义是表现出警觉以守护婚姻的贞洁，在此意义上，我们不只可以毫不犹豫地承认，更应感恩地宣告，天主对自己的子民怀有嫉妒，因为祂称他们为自己的配偶，并警告他们不可向众多假神行奸淫。同样也可以说及天主的义怒。因为当天主在愤怒中巡视人时，祂并不经受扰动；而是藉着词语的滥用或习语的特殊性，愤怒被用于表示惩罚之意。\par

19. 屠杀众多群众对外教人来说似乎并不奇怪，除非他否认天主的审判——而外教人并不否认；因为他们承认宇宙中从最高到最低的一切都由天主眷顾治理。但若他不承认这一点，便会被外教作家的权威或更冗长的论证所说服；若他仍固执乖僻，便会被交给他所否认的审判。那时，倘若他举出无罪或仅有轻微过失而遭毁灭的例子，我们将会表明那些是过失，并且并非轻微。例如，就前述婚宴礼服一事，我们将证明，一个人参加圣筵，寻求的不是新郎的荣耀，而是自己的荣耀（或无论那礼服在更佳解释下可能象征什么），这是大罪。而在那不愿那人作王统治之人在王面前被屠杀的事例中，我们或许容易证明，拒绝顺服同为人者，虽未必是罪，但拒绝那唯独在其统治中才有正义、幸福与永存者的统治，则既是过失，且是重大的。\par

20. 最后，关于福斯图斯狡猾的暗示，即旧约圣经错误地描述天主威胁要带着剑来，这把剑既不会放过义人也不会放过恶人。如果向外邦人解释这些话，他或许既不会不同意旧约也不会不同意新约；他或许会看到福音中比喻的美妙之处，而那些假装是基督徒的人或因盲目而误解，或因悖逆而拒绝。伟大的葡萄园主人对结果实的枝条和不结果实的枝条使用不同的修剪钩；然而他既不放过好的也不放过坏的，修剪一个而砍掉另一个。\BibleRef{若 15:1}没有一个人如此正义，以至于不需要通过苦难来试炼，以增进、巩固或证明他的德行。摩尼教徒们难道不认为保禄是义人吗？他谦卑而诚实地忏悔自己过去的罪，但仍因因信耶稣基督而成义而感恩。那么，保禄难道被他（愚昧之人误解的那位）赦免了吗？当祂说\Quote{我必不放过义人也不放过罪人}时，听听宗徒自己怎么说：\Quote{因了启示的浩大，恐怕我过于高抬自己，就有一根刺加在我的肉体上，就是撒殚的使者来击打我。为这事我三次求过主，叫这刺离开我。祂对我说：\InnerQuote{我的恩宠够你用的，因为我的德能在软弱中才全显出来。}}\BibleRef{格后 12:7}这里，一个义人没有被放过，为的是使他的德能因软弱而全显出来，因为那赐给他一位撒殚的使者来击打他的那位。如果你说是魔鬼给了这个使者，那么魔鬼就是试图阻止保禄因启示浩大而过于高抬自己，并成全他的德能。这不可能。因此，那将这个义人交给撒殚的使者击打的，与那通过保禄将恶人交付给撒殚本身的是同一位，保禄说：\Quote{我已经把他们交给撒殚，叫他们受管教不再亵渎。}\BibleRef{弟前 1:20}你现在看到至高者既不放过义人也不放过恶人了吗？或者是剑让你害怕？因为被击打不如被处死那么糟糕。但成千上万的殉道者不是以各种方式遭受死亡吗？难道迫害者能有这样的权柄加害于他们，除非是从天主那里得来的？天主是这样既不放过义人也不放过恶人的。因为主自己，这位首要的殉道者，明确对比拉多说：\Quote{若不是从上头赐给你，你根本无权柄管教我。}\BibleRef{若 19:11}保禄除了记录他自己的经历外，还说义人的苦难和迫害显明了天主的审判。\BibleRef{得后 1:5}这个真理在伯多禄宗徒已经引用的那段话中详细说明了，他说：\Quote{时候已经到了，审判必须从天主的家开始；若是先从我们起，那些不信从天主福音的，将有什么结局呢？若是义人仅仅得救，那不虔敬和犯罪的人，将有何地可站呢？}\BibleRef{伯前 4:17}伯多禄也解释了恶人如何不被放过，因为他们是折断下来要被焚烧的枝条；而义人也不被放过，因为他们的净化要达至完全。他将这些事归于那在旧约中说\Quote{我必不放过义人也不放过罪人}的那一位的旨意；因为他说：\Quote{如果圣神的旨意要这样，更好是因行善受苦，而不是因行恶受苦。}\BibleRef{伯前 3:17}所以，当因圣神的旨意，人为行善而受苦时，义人没有被放过；当人为行恶而受苦时，恶人没有被放过。在这两种情况下，都是按照那说\Quote{我必不放过义人也不放过罪人}的那一位的旨意；他管教义人如同儿子，惩罚恶人如同悖逆者。\par

21. 我已尽力表明，我们所敬拜的天主并非从永远住在黑暗中，而是祂自己就是光，在祂内毫无黑暗；祂自居于不可接近的光中；这光的光辉就是祂的永恒智慧。从我们所说的看来，天主并非因光的突然出现而措手不及，而是光的存在归功于祂是创造者，光的持续存在也归功于祂的认可。天主也并非不知未来，而是祂既颁布诫命，又惩罚悖逆；祂通过对过犯显明义怒，为当下提供约束，为将来提供警戒。祂发问也不是出于无知，而是通过祂的询问显明祂的审判。祂也不是出于好奇或胆怯，而是将悖逆者排除在永生之外，这是对顺服的公正赏报。祂也不是贪图血和脂油；而是向一个肉性的民族要求适合他们性格的祭献，从而以某些预像预示真正的祭献。祂的嫉妒并非苍白焦虑的情绪，而是安静慈爱的表现，渴望保持灵魂——它只应向唯一真天主保持贞洁——不被侍奉许多假神所玷污和出卖。祂的愤怒也不是类似人怒的情感，而是愤怒，不是出于报复的欲望，而是以独特的方式完全执行公义报应的判决。祂也不是为了微不足道的过犯或无缘无故地消灭数千人，而是通过已死之人的暂时死亡，向世界彰显敬畏祂的益处。祂也并非不加区分地惩罚义人和罪人，而是为了义人的益处管教他们，以成全他们，并给予罪人应得的惩罚。这样，你们摩尼教徒，你们的猜疑误导了你们，因为误解我们的圣经，或听了不良的解释者，你们对公教徒形成了错误的判断。因此你们离开了纯正的道理，转向不虔敬的寓言；因着你们的悖逆和远离圣徒的团体，你们拒绝新约的教导，而正如我们已经表明的，新约含有与你们在旧约中谴责的相似的说法。所以我们不得不为两约辩护，既对你们，也对外邦人。\par

22. 但是假设有某个人如此被肉性迷惑，以至于不去崇拜我们所崇拜的、唯一而真实的天主，却去崇拜你们猜疑或诽谤的虚构之物——你们说我们所崇拜的就是它，难道这个神不比你们的神更好吗？我恳求你们，请注意那些连最微弱的理解力也能明白的事；因为这里不需要很大的洞察力。我向所有智慧和不智慧的人说话。我向所有人的常识和判断力呼吁。请听、请思考、请判断。难道你们的神从永恒中留在黑暗里，不是比将与他同永恒且同性质的光投入黑暗中更好吗？难道更好的是对出现的新光表示惊讶和赞赏——这新光来驱散黑暗，而不是无法抵挡黑暗的进攻，除非让出他自己的光？如果他是在惊恐中这样做，那是不幸的；如果不需要这样做，那是残忍的。确实，如果他看见由自己创造的光，并赞赏它为善，这比使他自己所生的光变恶更好；更好的是他自己的光在与黑暗势力对抗时不应该成为他的仇敌。因为那些将被永远定罪到黑暗群中的人所受的控诉将是：他们容许自己失去原有的光明，并成了圣光的仇敌。如果他们从永恒中不知道他们将被这样定罪，那么他们必定遭受了永恒的无知之黑暗；或者如果他们知道，则遭受了永恒的恐惧之黑暗。因此，你们的神的一部分实体确实从永恒中留在自己的黑暗里；而不是在它出现时赞赏新光，它只是遇到了另一种敌对的黑暗，这是它一直恐惧的。事实上，天主自己必定在恐惧的黑暗中，为了他自身的这一部分，如果他正在害怕降临其上的邪恶。如果他没有预见这邪恶，他必定在无知之黑暗里。如果他预见了，却不恐惧，那么这种残忍的黑暗比无知或恐惧的黑暗更糟。你们的神似乎缺少宗徒在身体上所称赞的品质——这身体你们疯狂地相信不是由天主而是由\OriginalTerm{质料}{\Emph{Hyle}}所造的：\Quote{若一个肢体受苦，所有的肢体都一同受苦。}\BibleRef{格前 12:26} 但假设他受苦了；他预见了、恐惧了、受苦了，但他无法自救。这样，他从永恒中留在自己悲惨的黑暗里；然后，不是赞赏即将驱散黑暗的新光，他反而遇到了另一种他始终惧怕的黑暗，伤害了他自己的光。再者，我说，与其像天主一样颁布诫命，不如像\Person{亚当}一样接受诫命——遵守就得赏报、违反就受惩罚，两者都凭自己的自由意志——这岂不是好得多，而不是被不可避免的必然性强逼着违背自己的意愿将黑暗放入他的光中？确实，与其让必然而驱使他违背自己的神圣本性去犯罪，不如向人的本性颁布一条诫命，而不知道它会成为有罪的。请想一下，说一个被必然性征服的人如何能征服黑暗？他已经为这个更大的敌人所征服，却在他的征服者的命令下对抗一个较不强大的对手。与其不知道亚当藏在哪里，不如他自己没有任何逃脱的方法——首先摆脱一个严酷且可恨的必然性，然后摆脱一个不同且敌对的种族，这岂不是更好？与其拒绝人的本性得永生，不如将神圣的本性交付悲惨；与其贪求祭品的血和脂油，不如自己被以多种形式屠杀，因为与每一个牺牲品的血和脂油混合；与其因这些祭品也献给其他神而嫉妒不安，不如自己在所有祭台上被献给所有魔鬼——不仅与所有果实混合，而且与所有动物混合？岂不更好吗？即使受人的怒气影响，以致对自己的朋友和敌人因他们的罪而发怒，也比自己受恐惧和怒气的影响（无论这些情感在哪里存在），或者分享所有犯下的罪和所受的惩罚，岂不是好得多？因为这是你们的神那一部分的命运，它到处被禁锢，由他自己判处此刑——不是因为有罪，而是为了征服他所惧怕的敌人。他自己注定遭遇如此致命的必然性，他自己那部分已交付定罪的部分可能会原谅他——如果他像他悲惨那样谦卑的话。但是，当你们的神对犯罪的朋友和敌人发怒时，你们怎能假装指责天主呢？而你们幻想的神先是强迫自己的肢体去被罪吞噬，然后判定他们留在黑暗中。虽然他是这样做的，但你们说这不会是在发怒。但他难道不会羞于惩罚——或显得惩罚——那些他本应以这样的话向他们请求宽恕的人：\Quote{请宽恕我，我恳求你们。你们是我的肢体，我还能这样对待你们吗？除非出于必然。你们自己知道，你们被派到这里是因为一个可怕的敌人起来了；现在你们必须留在这里以防止他再次起来}？再次，与其为了微小或没有的过错杀害成千上万的人，不如将天主的自身肢体、他的实体——实际上就是天主自己——投入罪的深渊，并判定受永远监禁的惩罚？关于天主的真实实体，不能说它有选择犯罪或不犯罪的能力，因为天主的实体是完全不变的。天主不能犯罪，正如他不能否认自己一样。相反，人能犯罪和否认天主，也可以选择不这样做。但是，如果你们的神的肢体像理性的人的灵魂一样，有选择犯罪或不犯罪的能力，那么他们也许可以因重大罪行而被公正地惩罚，被禁锢在黑暗群中。但你们不能将这些部分归因于一种你们否认天主本身所有的自由。因为如果天主没有把他们交出去犯罪，他自己就会被黑暗种族的优势所迫而犯罪。但如果没有被逼迫的危险，那么将这些部分送到一个他们面临这种危险的地方就是一项罪行。确实，凭自由选择这样做是一种罪行，应受你们的神反常施加于自己肢体的折磨——甚至比这些肢体因他的命令而去到一个失去正直生活能力的地方的行为更甚。但如果天主自己因入侵和捕获而有被强制犯罪危险，除非他首先通过自己肢体的不当行为然后是惩罚来保护自己，那么无论是你们的神还是他的肢体，都没有自由意志。让他不要自命为审判者，而应承认自己为罪犯。因为尽管他是违背自己意愿被强迫的，他仍然声称在判处那些他知道是遭受了恶而非行恶的人时，是正义的；他这样做是为了不让人以为他被征服了——好像一个乞丐被称为繁荣幸福就能给他带来什么好处似的。确实，与其毫无区别地惩罚义人和不义之人（这是\Person{福斯图斯}对我们天主的最后指控），不如对待自己的肢体如此残忍——先是让他们遭受无法治愈的污染，然后，好像还不够似的，又诬告他们行为不当。\Person{福斯图斯}宣称，他们公正地遭受这种严厉和永久的惩罚，因为他们容许自己偏离原来的光明，并成了圣光的仇敌。但根据\Person{福斯图斯}，这是因为他们在黑暗王子的第一次攻击中被贪婪地吞噬，以致无法恢复自己，或无法从敌对原则中分离出来。因此，这些灵魂自己并没有作恶，在这一切中他们是无辜的受害者。真正的行为者是那个把他们从自己身边送到这种悲惨中来的那个。他们从父亲那里比从敌人那里受的苦更多。他们的父亲把他们送到这一切悲惨中；而他们的敌人渴望他们，认为他们是好东西，不愿伤害他们，只想享受他们。一个是有意伤害他们，另一个是无知。这位神如此软弱无助，否则他无法保护自己——首先不受威胁进攻的敌人，然后不受囚禁中的同一个敌人。那么，让他不要判罪那些因顺从而保卫他、因死亡而保障他安全的肢体。如果他不能避免冲突，为什么要诽谤他的保卫者？当这些肢体容许自己偏离原来的光明，并成了圣光的仇敌时，这必定是由于敌人的力量；如果他们是被强迫的，违背他们的意愿，那么他们是无罪的；而如果他们本可以抵抗却选择不抵抗，那么就不需要在一个想象的恶本性中寻找恶的起源，因为恶的起源在于自由意志。他们本可以抵抗却没有抵抗，这显然是他们自己的过错，而不是由于任何外来的强迫。因为，假设他们有能力做一件事，做这件事是对的，而不做是重大和可憎的罪，那么他们不做就是他们自己的选择。所以，如果他们选择不做，那么过错在于他们的意志，而不是在于必然性。罪的起源在于意志；因此，恶的起源也在于意志——既指违背公正诫命的行动，也指在公正判决下遭受痛苦。因此，在你们寻找恶的起源时，没有理由陷入如此大的恶，以至于将如此丰富的善的本性称为恶的本性，并将必然性的可怕之恶归于尚未与恶混合的完美善的本性。这种错误信仰的原因是你们的骄傲，你们不必有骄傲，除非你们选择；但你们为了不惜一切代价维护你们已陷入的错误，就从自由意志中取走了恶的起源，并把它放在一个虚构的恶的本性中。因此你们最终说，那些注定要被永远囚禁在黑暗群中的灵魂成为圣光的仇敌不是由于选择，而是由于必然；并使你们的神成为一个审判者，在他的法庭上为你们的当事人证明他们是被迫的毫无用处，也是一个君王，不会宽恕你们的兄弟——他自己的儿子和肢体——你们将他们与你们和他自己敌对归因于必然，而非选择。何等骇人的残忍！除非你们接下来要为你们的神辩护，说他也是出于必然，而非选择而行动。所以，如果还能找到另一个不受必然性束缚的审判者，能以公正原则裁决这个问题，他会将你们的神判决捆绑到这个黑暗群上——不是从外面绑住，而是与可怕的敌人一起关在里面。在必然性的罪中为首者，应在定罪判决中居首。那么，与这样一个神相比，选择我们确实不崇拜、但你们认为或假装认为我们所崇拜的神，岂不是好得多？虽然他毫不区别地不宽恕义人或有罪的人，不区分惩罚和管教，他岂不是比那个神更好吗？那个神不宽恕自己无辜的肢体，如果必然性不是罪过，或者如果他们因顺从他而有罪（如果必然性本身是犯罪），那么他们被他永远定罪，而如果他们因胜利恢复了任何自由，他们本应与他一起被释放；而他本应与他们一起被定罪，如果胜利即使只将必然性的力量减少到给正义如此微小力量的程度。所以，你们所描绘的我们崇拜的神——虽然他不是我们真正崇拜的唯一真神——远比你们的神好。两者都不存在；但两者都是你们想象的产物。但是，根据你们自己的描述，你们称为我们的并指责的那个神，比你们称为你们自己的并崇拜的那个神更好。\par

23. 同样，你们所抨击的族长和先知，并非我们所尊敬的人，而是你们因恶意而凭幻想描绘出来的人物。然而，即使按你们所描绘的那样，我也不仅要证明他们优于你们那些遵守摩尼一切诫命的选民，更要证明他们优于你们的神明本身。但在证明这点之前，我必须藉天主的帮助，以真理的明辨来反驳你们心里的肉性，为我们的圣祖族长和先知辩护。至于你们这些摩尼教徒，只需指出你们归咎于我们祖先的过错，比你们自己所称赞的还要可取，并且加上一句：你们的神明甚至远不如你们所描述的我们的祖先，这就足以使你们蒙羞了。这对你们而言已是充分的答复。然而，即使撇开你们的乖僻不谈，有些人本身在比较旧约先知的生活与新约宗徒的生活时，也会感到困扰——因为他们分辨不出预许尚在帷幕之下的时代与预许已显明的时代之间的差异——所以我首先必须答复那些要么胆敢自夸在节制上胜过先知，要么引用先知来为自己恶行辩护的人。\par

24. 首先，这些人的言语不仅是预言性的，他们的生活也是预言性的；整个希伯来王国就像一个伟大的先知，与所预言的伟大人物相称。因此，关于那些藉天主训导而内心明智的希伯来人，我们可以在他们的言行中发现基督和教会的预言；同样，就天主对整个民族的整体处置而言也是如此。因为正如宗徒所说：\Quote{这一切事都是我们的鉴戒。}\par

25. 那些挑剔先知，例如以他们不能理解的行为指控先知犯奸淫的人，就像那些外邦人亵渎地指责基督愚昧或疯狂，因为他在不合时令的树上找果子；\BibleRef{玛 21:19} 或者指责他幼稚，因为他弯下腰在地上写字，在回答完问话的人后又继续写字。\BibleRef{若 8:6} 这样的批评者无法理解：大人物身上的某些美德，在形貌上与小人物的恶行极为相似——不是实质上，而是外表上。这样对大人物批评，就像学童在学校里，他们的学问包括一条重要规则：如果主格是单数，动词也必须是单数；于是他们挑剔最好的拉丁文作家，因为他写了\Quote{\Emph{Pars in frusta secant}}，他们说本应写\Quote{\Emph{secat}}。又如，他们知道\Quote{religio}拼写时只有一个\Quote{l}，却责怪他写了\Quote{\Emph{relligio}}，当他说\Quote{\Emph{Relligione patrum}}的时候。因此，有理由说：正如诗人用词的习惯与无知者的文法错误和生造词语不同，同样，先知们的象征性行为也与恶人的污秽行为不同。因此，正如一个犯了语法错误的孩子如果以维吉尔的用法来辩护就会挨鞭子，同样，任何人若引用亚巴郎从哈加尔生子来为自己对妻子的婢女怀有不洁的激情辩护，就不仅该受杖责，还应受严厉的鞭笞，以免他在永恒的惩罚中遭受奸淫者的命运。这确实是将大事与琐事相比，并非将言语的特殊用法与圣事相提并论，也非将语法错误与奸淫等同。不过，考虑到主题性质上的差异，在言语的恰当与不当方面所谓的学识与无知，在对美德与恶行这一重要道德区别上的智慧或缺乏，是类似的。\par

26. 我们不必逐一区分可称赞与可指责、有罪与无辜、危险与无害、有咎与无咎、可欲与不可欲——这些全都是罪与义之区别的例证——而应首先思考什么是罪，然后查考圣书中所记载的圣人的行为，以便发现：若我们读到这些圣人犯了罪，则查明将这些罪记录在案并传诸后世的真正原因。此外，若我们读到一些事，它们并非罪，但在愚昧和心怀恶意的人看来却是罪，实际上也未显出任何德行，那么我们也须明白为何将这些事载入圣经——我们相信圣经包含有益于现世生活的教义并使人有资格继承来生。至于圣人的行为中那些正义的榜样，记录它们的恰当性即使愚昧人也应明白。问题是那些既非义行又非罪行的行为，提及它们似乎无用；或者那些实在有罪的行为，则可能危险，因为它们在圣经中未受谴责，故人们可能认为它们不是罪而加以模仿；或者即使它们被谴责，人们也可能因圣人做过而认为它们是小罪，从而照样去做。\par

27. 罪，就是任何在行为、言语或欲望上违犯\OriginalTerm{永恒律法}{eternal law}的事。永恒律法是天主的秩序或旨意，它要求保持\OriginalTerm{自然秩序}{natural order}，并禁止破坏它。但在人身上，这自然秩序是什么呢？我们知道，人由灵魂和身体组成；但野兽也是如此。再者，显然在自然的秩序中，灵魂高于身体。此外，在人的灵魂中有理性，这是野兽所没有的。因此，正如灵魂高于身体，同样在灵魂本身，理性按自然规律高于那些在野兽中也存在的其它部分；而在理性本身中，它部分属于\OriginalTerm{默观}{contemplation}，部分属于\OriginalTerm{行动}{action}，默观无疑是更高的部分。默观的对象是天主的肖像，我们藉着\OriginalTerm{信德}{faith}被更新，得以看见。理性的行动因此应当服从默观的引导，这默观在我们远离主的时候藉着信德实行，正如将来要藉着看见，那时我们要像祂，因为我们必要看见祂实在怎样。\BibleRef{若一 3:2} 那时，在一个属神的身体里，我们要藉祂的恩宠被造成与天使同等，当我们穿上不朽和不坏的衣裳，这有死和可朽坏的将被这样穿上，使死亡被胜利吞灭，当义德藉恩宠达于完美时。因为圣洁而崇高的天使也有他们的默观和行动。他们自愿执行他们所默观的那位的命令，甘甜地服从祂永恒的管理。我们呢，我们的身体因罪而死，直到天主藉着居住在我们内的圣神，也使我们有死的身体活过来，我们按着那保存自然规律的永恒律法，在我们微弱的尺度内度义德的生活，当我们凭着那藉爱德运作的无伪信德而生活时，在良心中存有对天上不朽和不坏的望德，以及对义德完美达到无法言喻的满足的望德，为此我们在旅途中必须饥渴，因为我们凭信德行走，而非凭目见。\par

28. 因此，一个服从那听命于天主的信德而行动的人，抑制一切可朽的私欲偏情，并将它们保持在自然的限度内，调节自己的欲望，以便将更高的置于更低的之前。如果违禁之事中没有愉悦，就没有人会犯罪。犯罪就是放纵这种愉悦，而非抑制它。所谓违禁，就是被那保存\OriginalTerm{自然秩序}{natural order}的律法所禁止的。这是一个重大的问题：是否有任何\OriginalTerm{理性受造物}{rational creature}对于违禁之事没有任何愉悦？如果有这样一类受造物，它不包括人，也不包括那没有持守真理的天使性体。这些理性受造物被造时，具有抑制自己远离违禁之事的潜能；而没有这样做时，他们犯了罪。那么，受造物人真是伟大，因为他藉着这潜能被恢复，如果他当初选择的话，他本不会堕落。创造人的主是伟大的，应受大赞美。因为祂也创造了不能犯罪的较低本性，以及不会犯罪的较高本性。\OriginalTerm{野兽}{beasts}不犯罪，因为它们的本性与永恒律法相符，是由于服从它，而不拥有它。同样，天使不犯罪，因为他们的天上本性如此拥有永恒律法，以致天主是他们欲望的唯一对象，他们服从祂的旨意而不经历任何诱惑。但人，其在这地上的生命由于罪而是一场试探，他将自己与野兽共有的事物降服于自己，并将自己与天使共有的事物降服于天主；直到义德达于完美、不朽获得之时，他将从野兽中被提升，与天使并列。\par

29. 身体欲求的运用或放纵，旨在确保个体或种族的持续存在与活力。如果欲求超出此范围，将人——不再自制——带出节制的界限，它们就成为不法和可耻的情欲，必须由严厉的纪律加以驯服。但如果这种无节制的结果是使人堕入如此深重的恶习，以至他以为自己的罪恶情欲不会受惩罚，因而拒绝那能拯救他的\OriginalTerm{告解}{confession}与\OriginalTerm{悔改}{repentance}的健康纪律；或者，由于更糟糕的麻木，他公然反对天主的永恒律法，为自己的放纵辩护；如果他死在这种状态中，那无误的律法如今判他不再是纠正，而是\OriginalTerm{永罚}{damnation}。\par

30. 那么，让我们参照那命令保持\OriginalTerm{自然秩序}{natural order}、禁止破坏它的\OriginalTerm{永恒律法}{eternal law}，看看我们的祖先\Person{亚巴郎}在\Person{福斯托}指控为极其罪恶的事中，是如何犯罪的，即如何违反这律法的。\Person{福斯托}说，\Person{亚巴郎}因有非理性的得子渴望，且不相信曾应许他的妻子\Person{撒辣}要生一个儿子的上主，就用一个妾玷污了自己。但此处的\Person{福斯托}，因他有非理性的找茬欲望，既暴露了其异端的不敬，又在错误与无知中赞扬了\Person{亚巴郎}与使女的同房。因为正如永恒律法——即创造万有的上主的旨意——为保存自然秩序，容许在理性指引下，在性交中运用身体欲求，不是为了满足情欲，而是为了藉生育子女延续种族；相反地，摩尼人的不义律法，为了阻止他们哀叹为禁锢在一切种子中的神，在子宫中遭受更紧密的禁锢，要求已婚者无论如何不可生育子女，他们的最大愿望是解放他们的神。因此，我们发现在\Person{亚巴郎}身上不是非理性的得子渴望，在\Person{摩尼}身上却是非理性的反对生育的幻想。所以，一个藉在婚姻中只寻求生育子女而保存了自然秩序；另一个，受其异端观念影响，认为没有什么邪恶比禁锢他的神更大了。\par

31. 同样，当福斯特说妻子知悉丈夫的行为使事态更糟时，他只是出于不仁爱之心想要责备亚巴郎和他的妻子，但他无意中却赞扬了双方。因为撒辣并没有纵容丈夫的任何罪恶行为以满足他非法的情欲；而是出于与他相同的自然生育欲望，并且知道自己不孕，她合理地声称拥有婢女的生育能力；并非同意丈夫的罪恶欲望，而是请求他给予他本该允许的事。这也不是骄傲自负的请求；因为人人都知道妻子的本分是服从丈夫。但关于身体，宗徒告诉我们，妻子对丈夫的身体有权力，如同丈夫对妻子的身体有权力一样；\BibleRef{格前 7:4}因此，在其他社会事务中妻子应当服从丈夫，但在男女结为夫妻的身体结合这一件事上，他们彼此拥有权力——男人对女人，女人对男人。所以，当撒辣自己不能生育时，她希望通过自己的婢女获得子女，并且来自同一血脉，如果可能的话，她本会自己生育。任何女人若不是仅仅出于动物性的情欲爱丈夫，就不会这样做；她更会嫉妒一个主母而非使她成为母亲。因此，这里对生育子女的虔诚愿望表明没有罪恶的放纵。\par

32. 如果像福斯特所说，亚巴郎因对天主应许他通过撒辣得子没有信德，而想通过哈加尔得子，那么亚巴郎确实无法辩护。但这完全是错误的：这个应许尚未作出。任何人若阅读前面的章节便会发现，亚巴郎已经得到了土地的应许，有无数居民，\BibleRef{创 12:3}但尚未向他显明所提到的后裔如何产生，是通过他自己的身体生育，还是通过他选择收养一个儿子，或是如果来自他自己的身体，是通过撒辣还是通过别人。凡考察此事者都会发现福斯特犯了轻率的错误或厚颜无耻的歪曲事实。因此，亚巴郎看到自己没有儿子，尽管应许是给他的后裔，他首先想到的是收养。这一点从他向天主谈论他的仆人说：\BibleQuote{这是我的继承人}中可以看出；意即：既然你没有给我自己的后裔，就在这个人身上实现你的应许。因为\Emph{后裔}一词可以用于不是从人自己身体所出的，否则宗徒就不能称我们为亚巴郎的后裔；因为我们确实不是他肉身的后裔；但我们是因追随他的信德而成为他的后裔，因相信那肉身出自亚巴郎肉身的基督。那时主对亚巴郎说：\BibleQuote{这人不会是你的继承人；你亲生的才会是你的继承人。}\BibleRef{创 15:3}收养的想法因此被排除；但仍不确定来自他自身的后裔是通过撒辣还是通过别人。天主乐于保留这个秘密，直到旧约的一个预表在婢女身上应验。这样我们就容易理解，亚巴郎看到妻子不孕，并且她渴望从丈夫和婢女那里得到她自己无法产生的后代，他行动不是出于肉欲，而是出于对婚姻权威的服从，相信撒辣的愿望有天主的许可；因为天主已经应许他一个来自他自身的继承人，但没有预言谁将是母亲。因此，当福斯特以自己不忠信的心指控亚巴郎不信从而显明自己的不信时，他无根据的指控只证明了攻击者的疯狂。在其他情况下，福斯特的不忠使他无法理解；但在这里，出于诽谤的爱好，他甚至没有花时间阅读。\par

33. 再者，当福斯特出于贪婪和吝啬，指控一位正义忠信的人无耻亵渎婚姻，先后两次把妻子撒辣卖给阿彼默肋客和法郎两位君王，告诉他们她是他的妹妹，因为她十分美丽，他并未公正地区分善恶，反而不公地谴责了整件事。那些认为亚巴郎出卖妻子的人，无法以永恒之法的眼光辨明罪与义之间的区别；因此他们称恒心为顽固，称信赖为狂妄，正如在这些和类似情况下，判断错误的人常会指责他们认为错误的行为。亚巴郎并没有通过把妻子卖给他人而成为她的共犯：而是如同她将婢女给丈夫，并非为满足他的情欲，而是为了子嗣，并在与自然秩序相符的权威中，要求履行义务，而非顺从罪恶的欲望；同样在这一事例中，丈夫完全确信妻子对他的贞洁依恋，并知道她的心灵是谦逊和贞洁情感的居所，他称她为妹妹，而没有说她是他的妻子，以免自己被杀，妻子落入外人及恶人手中；因为他确信他的天主不会允许她遭受暴力或耻辱。他对信德和望德没有失望；因为法郎被奇异的事情惊吓，并因她而遭受许多灾祸之后，当天主告诉他撒辣是亚巴郎的妻子时，法郎就完好无损地把她归还给他，并予以尊敬。阿彼默肋客也在梦中得知真相后这样做了。\par

有些人，并非如法乌斯图斯那样的嘲弄者和诽谤者，而是那些对圣经给予应有尊重的人——法乌斯图斯之所以指责圣经，是因为他不理解它们，或者他因指责而未能理解——在评论亚巴郎这一行为时，认为他因信德软弱而失足，并因怕死而否认他的妻子，正如伯多禄否认了主。如果这个看法是正确的，我们必须承认亚巴郎犯了罪；但这罪不应取消或抹掉他的一切功绩，正如宗徒的情况一样。此外，否认他的妻子与否认救主并非同一回事。但既然有另一种解释，为什么不坚持它，反而无缘无故地指责，因为并无证据表明亚巴郎因恐惧而说谎？他并没有在回答任何关于这事的提问时否认撒辣是他的妻子；但当被问到她是谁时，他说她是他的妹妹，并没有否认她是他的妻子：他隐瞒了部分真相，但没有说任何虚假的话。\par

注意法乌斯图斯的评语——亚巴郎虚假地称撒辣为他的妹妹——是浪费时间；好像法乌斯图斯发现了撒辣的家族，虽然圣经并未提及。在一件亚巴郎知道而我们不知道的事情上，相信这位圣祖所说的他所知道的事，当然比相信摩尼所指责的他毫不知情的事更好。那么，既然亚巴郎生活在人类历史的那一时期，当时虽然同一父母所生的子女之间的婚姻已被视为非法，但兄弟的子女或任何更近亲等之间的婚姻习俗并未受到法律或权威的干预，他为什么不能娶他的妹妹——即一位从他父亲血统而来的女子——为妻呢？因为他自己在归还撒辣时告诉国王，她是他的妹妹，出于他的父亲而非母亲。而在这次情况下，他不可能因恐惧而被引导说谎，因为国王知道她是他的妻子，并且因受到天主的警告，正以尊敬的态度归还她。我们从圣经中得知，在古人中，习惯上将表亲称为兄弟姊妹。因此，多俾亚在与他的妻子同房前，在天主面前祈祷说：\BibleQuote{主啊，现在祢知道，我并非出于情欲而娶我的妹妹为妻；}\BibleRef{多 8:9}尽管她并非直接来自同一父亲或同一母亲，而只属于同一家族。而罗特被称为亚巴郎的兄弟，尽管亚巴郎是他的叔父。并且，出于这同一措辞用法，在福音中被称为主之兄弟的那些人，肯定不是童贞玛利亚的子女，而是主的所有血亲。\BibleRef{玛 12:46}\par

有人可能会说：为什么亚巴郎对天主的信赖没有阻止他害怕承认他的妻子？天主本可以为他挡开他所惧怕的死亡，而在外邦人中保护他和他的妻子，使撒辣虽然十分美丽，也不致被任何人垂涎，亚巴郎也不致因她而被杀。当然，天主本可以这样做；否认这一点是荒谬的。但是，如果亚巴郎在回答众人时告诉他们撒辣是他的妻子，他对天主的信赖本应包含自己的生命和撒辣的贞洁。如今，合理的教导之一是：当人拥有任何能力手段时，他不应试探主他的天主。所以，救主告诉祂的门徒们：\BibleQuote{如果你们在一座城里受迫害，就逃到另一座城去。}\BibleRef{玛 10:23}并非因为祂不能保护他们。而且祂自己树立了榜样。因为虽然祂有权舍弃自己的生命，且直到祂选择这样做才舍弃，但祂婴孩时仍由父母抱着逃往埃及；\BibleRef{玛 2:14}而当祂上到庆节时，祂不公开上去，而是秘密地，尽管在其他时候祂公开向犹太人讲话，他们尽管愤怒和敌视却无法抓住祂，因为祂的时辰还未到——不是祂必须死的时辰，而是祂认为适宜被处死的时辰。如此，那位通过公开教导和谴责而显示天主性权力、不容许敌人的愤怒伤害祂的，也通过逃跑和隐藏自己而展现了适合人类软弱的行为，即当他们有办法逃避威胁的危险时，不应试探天主。同样，在宗徒的事例中，他被人用篮子从城墙上缒下，以逃避被敌人捉拿，并非出于对天主援助和保护的绝望，或信德丧失：\BibleRef{宗 9:25}他如此逃脱并非出于信德缺乏，而是因为当逃脱是可能时，不逃脱将是对天主的试探。因此，亚巴郎在外邦人中时，由于撒辣的非凡美丽，他的生命和她的贞洁都处于危险之中，因为他能够保护的并非这两者，而只有一个——即他的生命——为了避免试探天主，他做了他能做的；而对于他不能做的，他信托了天主。他无法隐瞒自己是男人，便隐瞒了他是丈夫，以免被处死；信托天主保存他妻子的纯洁。\par

关于一个微妙之处也可能有意见分歧：即使有人与撒辣交合，她的贞洁是否会受到侵犯？因为她为了救丈夫的性命而服从此事，既得到他的知情，也出于他的权柄。这其中并没有背弃婚姻的忠贞或反抗丈夫的权柄；正如亚巴郎不是奸夫，当他出于对他妻子合法权柄的服从，同意借妻子的使女成为父亲时一样。但是，从关系的本质而言，妻子有两个都在世的丈夫与男子有两个妻子是不同的：因此我们认为已经给出的关于亚巴郎行为的解释是最正确且无可指摘的：即我们的圣祖亚巴郎通过采取他能力范围内的一切措施来保全自己的生命，避免了试探天主；而他将妻子的贞洁托付给天主，显示了他对天主的望德。\par

38. 然而，从这一忠实记载于圣经的叙述中，所有的人都必然获得一种喜乐，当我们审视这一行动的先知性特征，并以虔敬的信德和勤勉叩响奥秘之门时，愿主开启，向我们指明在这位古人的身上预表了谁，以及这位妇人是谁——她在异邦和陌生人中间，不容受到玷污或污损，好使她能被带到自己的丈夫面前，毫无瑕疵或皱纹。因此我们发现，教会正义的生活是为基督的光荣，好使她的美貌为丈夫带来荣耀，正如亚巴郎因撒辣的美貌而受到那异邦居民尊敬一样。对教会，即在\WorkTitle{雅歌}中被称作\BibleQuote{女子中最美丽者}\BibleRef{歌 1:7}的那一位，君王们因她的美貌而献上礼物，正如阿彼默肋客王因欣赏撒辣容颜的恩宠而向她献礼；尤其当他在爱慕时，却不允许亵渎她。同样，圣教会也是主耶稣基督秘密的配偶。因为人的灵魂与天主的圣言在秘密和圣神的隐藏深处结合，以致\BibleQuote{二人成为一体}；关于这一点，宗徒在论及婚姻时称之为一个伟大的奥秘，是指基督和教会说的\BibleRef{弗 5:31}。再者，这个世界的尘世国度——由那些不被允许亵渎撒辣的君王所预表——对教会作为基督的配偶并无认识或经验，即不知她如何忠信地保持与丈夫的关系，直到它试图侵犯她，并被迫屈服于由殉道者的信德所见证的神圣见证，并在后来的君王身上，谦卑地以礼物尊崇那位新娘，而他们的前任曾试图通过征服她来羞辱她，却未能成功。在预像中，发生在同一位君王统治时期的事，也在这一时代的早期君王及其后继者身上应验。\par

39. 再者，当说到教会是基督的姊妹，不是由于母亲而是由于父亲时，我们便认识到这一关系的卓越，它并非尘世血统的暂时性质，而是天主的恩宠，是永远的。藉此恩宠，我们将不再是会死亡的一族，当我们获得能力被称为并成为天主的子女时。这恩宠我们不是从会堂——即基督按血肉的母亲——得到的，而是从圣父得来的。当基督召叫我们进入另一个生命，那里没有死亡时，祂教导我们不要承认，而要否认那尘世的关系，因为在那里死亡紧随出生而来；祂对门徒说：\BibleQuote{不要称地上的人为你们的父，因为你们只有一位父，就是那在天上的。}\BibleRef{玛 23:9}祂也以身作则，当祂说：\BibleQuote{谁是我的母亲？谁是我的兄弟？}并伸手向门徒说：\BibleQuote{这些人就是我的兄弟。}为免有人以为祂指的是尘世关系，祂又加了一句：\BibleQuote{谁遵行我父的旨意，他就是我的兄弟、姊妹和母亲。}\BibleRef{玛 12:48}这等于说：我从我的父天主那里获得这一关系，而非从我的母亲会堂；我召叫你们进入永生，在那里我有不朽的出生，并非进入尘世生命，因为为召叫你们离开这生命，我取了必死的人性。\par

40. 至于为什么，尽管在陌生人中隐藏着教会是谁的配偶，却未隐藏她是谁的姊妹，原因很明显：人的灵魂与天主的圣言如何结合、相融，或用其他任何词句来表达天主与受造物之间的这种联系，这是隐晦难懂的。正是从这种联系中，基督与教会被称为新郎与新妇，或丈夫与妻子。另一种关系——基督与所有圣人因天主的恩宠而非因尘世血统、因父而非因母成为弟兄——则更容易用言语表达，也更容易理解。因为同样的恩宠也使所有圣人彼此成为弟兄；而在他们的团体中，没有一个人是所有其他人的新郎。同样，尽管基督有超越的正义和智慧，祂的人性却被陌生人更明显、更容易地认出，他们确实没有错误地相信祂是人，但他们并不理解祂同时也是神。因此耶肋米亚说：\BibleQuote{祂既是人，谁能认识祂呢？}\BibleRef{耶 17:9}祂是人，因为显明祂是兄弟。而谁能认识祂呢？因为隐藏了祂是丈夫。这足以作为对我们的父亲亚巴郎的辩护，反驳福斯图斯的无耻、无知和恶意。\par

41. 亚巴郎的兄弟罗特，在索多玛为人正义、好客，得以逃脱那预表未来审判的火灾；因为他完全没有参与索多玛人的腐化。他是基督身体的预像，这身体在诸圣身上，如今在不虔诚和邪恶的人中间呻吟，并不赞同他们的恶行，并在世界终结时，当这些人被判定受永火的刑罚时，它将从他们的团体中被解救出来。罗特的妻子是另一类人的预像——即那些蒙天主恩宠召唤时，却回头观看，不像保禄那样忘记背后，努力面前的人。\BibleRef{斐 3:13} 主亲自说：\BibleQuote{手扶着犁向后看的，不适合进天主的国。} \BibleRef{路 9:62} 他也没有忽略提及罗特妻子的事；因为她为了警戒我们，变成了盐柱，好让我们这样被调味后，不会轻率地玩忽这个危险，而是警惕它。因此，当主劝诫每个人要竭力摆脱背后的事，以追求前面的事时，他说：\BibleQuote{你们要记得罗特的妻子。} \BibleRef{路 17:32} 此外，还有第三个预像在罗特身上，就是当他的女儿与他同寝时。因为这里的罗特似乎预表了将来的法律；那些出自法律、并置于法律之下的人，由于误解法律，好似使它麻木，并通过非法使用法律而结出不信的果实。宗徒说：\BibleQuote{法律原是好的，只要人用的合法。} \BibleRef{弟前 1:8}\par

42. 罗特和他女儿的这个行为，并不能因为它预表了日后在某些情况下将显示的邪恶而成为借口。罗特女儿的目的是一回事，天主容许此事发生以显明某些真理的目的是另一回事；因为天主既判决了当时人们的行为，也以祂的眷顾安排了未来的预表。作为圣经的一部分，这个行为是一种预言；作为相关人物的历史，它是罪行。\par

43. 同时，在这个行为中，没有理由承受福斯图斯盲目敌意所倾泻的如潮辱骂。根据永恒法律，它要求维护自然秩序并谴责违反该秩序的行为，此案的判决，与罗特若是因犯罪情欲的驱使而与女儿乱伦，或者女儿们被非自然欲望所燃起的情况不同。公正地说，我们必须询问不仅做了什么，而且出于什么动机，以便公正地看待作为该动机结果的行为。罗特女儿决定与父亲同寝，是出于保存种族的自然传宗接代的欲望；因为她们以为找不到其他男人，认为整个世界已被那场大火烧尽，据她们所知，除了她们自己无人幸存。她们宁可从未成为母亲，也比成为自己父亲的母亲要好。但是，这种愿望的实现与受诅咒的情欲满足截然不同。\par

44. 知道父亲会谴责她们的计谋，罗特的女儿们认为有必要在父亲不知情的情况下实施。我们被告知她们使他喝醉，以至于他不知道发生了什么。因此，他的罪过不是乱伦，而是醉酒。这同样受到永恒法律的谴责，法律只允许为了维护健康而按自然需要食用饮食。确实，醉酒者和习惯性醉汉之间有很大区别；因为醉汉并不总是醉，一个人可能偶尔醉一次而不成为醉汉。然而，对于一个义人，我们即使一次醉酒也需要解释。是什么让罗特同意接收女儿们不断给他倒的或可能给他纯饮的那么多杯酒？她们是否假装过度的悲伤，而他在她们的孤独和失去母亲的情况下，以此为慰藉，以为她们也在喝酒，而她们只是假装喝？但这似乎不是一个义人在朋友遭遇困苦时安慰他们的适当方法。女儿们是否在索多玛学到某种卑劣的技巧，使她们能用几杯酒就使父亲喝醉，使他因无知而犯罪，或者说被犯罪？但圣经不太可能完全忽略这一点，或者天主会允许祂的仆人因自己的过错而无辜地遭此滥用。\par

45. 但是，我们是在为圣经辩护，而不是为人的罪辩护。我们也不是要证明这个行为是正当的，好像我们的天主曾命令或赞同它；或者好像当圣经称人为义人时，意味着他们不可能选择犯罪。而且，在那些批评者所挑剔的书中，天主没有在任何地方表示赞同这个行为，那么从这个叙述中对这些著作提出指控，而在其他地方，这类行为被明确禁令所谴责，这岂不是无知的愚蠢！在罗特女儿的叙述中，行为是被叙述的，而不是被称赞的。并且，天主的判决在某些情况下宣告，而在其他情况下隐藏，这是合适的，这样通过其显明，我们的无知得以启蒙；通过其隐藏，我们的心智可以通过回忆我们已经知道的事而得到锻炼，或者我们的懒惰被激励去寻求解释。因此，天主，祂能从恶中引出善，使各民族从这一起源中产生，正如祂看为好的，但并没有将人的罪过归咎于祂自己的圣经。这是天主的著作，而不是祂的行为；祂不把这些事提出作为我们的仿效，而是作为我们的警戒。\par

四六、\Person{福斯图斯}的厚颜无耻尤其显露于他指责\Person{依撒格}，\Person{亚巴郎}的儿子，假装他的妻子\Person{黎贝加}是他的姊妹。\BibleRef{创 26:7}因为关于\Person{黎贝加}的家族，圣经并非缄默，且显明她在众所周知的词义上是他的姊妹。他隐瞒她是他的妻子并不令人惊讶，也不是无关紧要的，如果他仿效他的父亲而行，那么他能够基于同样理由而被证明为正当。我们只需参考对\Person{福斯图斯}指控\Person{亚巴郎}的答复，因为它同样适用于\Person{依撒格}。或许，有些探究者会问，外邦王看见\Person{依撒格}与\Person{黎贝加}嬉戏而发现她是\Person{依撒格}的妻子，这有何预表意义？因为若不是他看见\Person{依撒格}与\Person{黎贝加}嬉戏，仿佛与不是他妻子的女人行不适当的事，他就不会知道。当圣人们如此以丈夫的方式行事时，他们并非愚昧地，而是有意地这样做：因为他们以温和嬉戏的言行来迁就较弱性别的本性；并非出于柔弱，而是出于克制的刚毅。但这样的行为对妻子以外的任何女人都是可耻的。这是一个有关良好习俗的问题，这里提及它，只因为也许有些严厉的、主张不动心的人会指责圣人甚至与妻子嬉戏。如果这些缺乏人性的人看到一个稳重的人为了以和善的同理心迁就儿童的幼稚理解而与儿童嬉戏交谈，他们就认为他癫狂了；忘记了自己也曾是儿童，或是对自己的成熟忘恩负义。关于基督和祂的教会，在这个伟大的族长与妻子嬉戏并由此发现婚姻关系中所能找到的预表意义，每一个为了避免以谬误教义冒犯教会而仔细研究圣经中教会新郎的奥秘的人都会看到。他将发现，教会的丈夫在仆人的形体中暂时隐藏了祂与圣父同等的尊威，因祂具有天主的形体，为的是使软弱的人性能与祂结合，并这样祂能迁就祂的新妇。这非但不荒谬，反而具有象征性的合适性：天主的先知使用一种属于肉体的嬉戏来迎合妻子的情感，正如天主圣言自己成了血肉，以便居住在我们中间。\par

四七、再者，\Person{依撒格}的儿子\Person{雅各伯}被指控犯了大罪，因为他有四个妻子。但这里没有指控犯罪的根据：一夫多妻在有此习俗时并非犯罪；而现在成为犯罪，是因为不再有此习俗。有违反本性的罪，有违反习俗的罪，有违反法律的罪。那么，\Person{雅各伯}在拥有多个妻子方面，犯了哪一种罪呢？就本性而言，他利用女人不是为了感官满足，而是为了生育子女。就习俗而言，这在当时那些地方是普遍做法。就法律而言，不存在禁止。现在这样做成为犯罪的唯一原因，是因为习俗和法律禁止它。凡藐视这些约束的人，即使他只用妻子来生育子女，仍然犯罪，并且对人类社会本身造成伤害，而生育子女正是为了人类社会。在现今习俗和法律改变的情况下，人们不可能从多个妻子中获得快乐，除非是出于过度的情欲；因此产生了一种错误看法，即认为除了由于私欲和罪恶欲望的强烈之外，没有人能拥有多个妻子。他们无法构想那些心智力量超出其理解的人，便如宗徒所说，\BibleRef{格后 10:12}以自己衡量自己，因而犯错。他们意识到，在与甚至只有一个妻子的交往中，他们常常为纯粹的动物情欲而非理智动机所驱使，便以为一个明显的推论是：如果在一个妻子时都不遵守节制的界限，那么在超过一个时，软弱必然加剧。\par

48. 但那些没有节制之德的人，不可允许他们判断圣人们的行为，正如发烧的人不能判断食物的甘甜与有益健康一样。营养的提供不应根据病态的味觉，而应根据健康的判断与指引，以治愈疾病。因此，如果我们的批评者希望获得并非虚假和做作的，而是真实和健全的道德健康，就让他们在相信圣经记载中找到治疗，即高尚的圣人称号并非无缘无故地赐给那些有多个妻子的男人；原因在于，心灵能如此控制肉体，以致不允许天主眷顾植入我们本性的欲望超出理性意图的界限。由于类似的误解，这种批评（与其说是诚实的判断，不如说是出于不诚实的诽谤）也可能指控圣宗徒们向如此多的人宣讲福音，并非出于为永生生育子女的愿望，而是出于对人间赞誉的爱。对这些我们在福音中的父辈来说，并不缺乏名声，因为他们的赞美通过多种言语在基督的教会中传播。事实上，人给予同类的尊荣与荣耀竟无过于此。正是对教会中这种荣耀的罪恶欲望，导致那弃者西满在盲目中希望用金钱购买那由神圣恩宠白白赐予宗徒们的东西。\BibleRef{宗 8:18} 在主用福音制止那希望跟随祂的人时，必有这种对荣耀的欲望，说：\Quote{狐狸有穴，天空的飞鸟有巢，人子却没有枕头的地方。}\BibleRef{玛 8:20} 主看见他的心思被虚假表象蒙蔽，被突然的情绪激动，且没有信德的基础来为谦卑的导师提供住所；因为那人在追随基督时，寻求的不是基督的恩宠，而是自己的荣耀。那些被宗徒保禄描述为不真诚地宣讲基督，而是出于竞争和嫉妒的人，正是被这种对荣耀的爱所引导；然而宗徒却因他们的宣讲而欢喜，知道可能发生这样的事：当宣讲者满足他们对人间赞誉的欲望时，他们的听众中可能有信徒诞生——并非因使他们渴望与宗徒名声竞争或超越的嫉妒情绪，而是通过他们所宣讲的福音，尽管不真诚；这样天主可以从他们的邪恶中带来善。因此，一个人可能被感官欲望引诱而结婚，而非为了生育子女；然而一个孩子可能出生，是天主的美好造化，由于自然的力量，而非父母的过失。因此，正如圣宗徒们在他们教义被听众接纳时感到欣慰，不是因为他们贪图赞美，而是因为他们渴望传播真理；同样，圣祖们在夫妻行房时，并非受感官欲望的驱动，而是出于对家族延续的理性愿望。因此，听众的数量并未使宗徒们变得野心勃勃；妻子的数量也未使圣祖们变得放荡。但为何要为丈夫辩护？他们的品行得到天主圣言的最崇高见证，事实上妻子们自己也认为与丈夫的结合仅仅是为了得到儿子？所以，当她们发现不生育时，她们把自己的婢女给了丈夫；这样，婢女拥有肉体的母职，而妻子则是意图上的母亲。\par

49. 福斯图斯做了一个毫无根据的陈述，当他指控那四个女人像放荡的人物一样为占有她们的丈夫而争吵时。我不知道福斯图斯从哪里读到这一点，除非是在他自己的心里，如同在一本不敬虔的谬误之书中，在那里福斯图斯自己被那条蛇所诱惑，关于这蛇宗徒曾为教会担忧，他愿将教会作为贞洁的童女献给基督；免得蛇用诡计欺骗了厄娃，也败坏你们的心智，使你们偏离对基督的纯朴。\BibleRef{格后 11:2} 摩尼教徒如此喜爱这条蛇，以致他们断言蛇带来的益处多于害处。福斯图斯必定从这条蛇那里使自己的心智受到注入谎言的败坏，他现在用这些无耻的诽谤重新制造出来，甚至大胆地写下来。并非一个婢女从另一个那里夺走了雅各伯，也不是她们为占有他而争吵。有安排，因为没有放荡的情欲；夫妻权柄的法则更加强大，因为没有肉欲的无律。他被其中一个妻子雇佣（租赁）的事实证明了这里所说的，与摩尼教徒的诽谤明显相反。为什么一个人应该雇佣他，除非按照安排他本应进入另一个那里？不能得出结论说，除非她被雇佣，他永远不会去肋阿那里。他必须总是在轮到她的时刻去她那里，因为他由她有了许多孩子；并且出于对她的服从，他由她的婢女得了孩子，之后，没有任何雇佣，由她自己得了孩子。这一次轮到辣黑耳，所以她有那在新约中宗徒明确提到的权柄：\Quote{丈夫对自己的身体没有权柄，而是妻子有。}\BibleRef{格前 7:4} 辣黑耳与她的姐姐做了一笔交易，既然欠她姐姐的债，她就把她的债务人雅各伯引荐给她。因为宗徒在他说\Quote{丈夫对妻子要尽本分}时使用了这个比喻。\BibleRef{格前 7:3} 辣黑耳把在她权柄下的丈夫所欠的当作应归还的给出，以换取她选择从姐姐那里得到的东西。\par

50. 如果雅各伯像福斯图斯在其不可救药的盲目中所设想的那样，是一个侍奉私欲偏情而非正义的人，那么他难道不会整天急切地盼望夜晚与更美丽的妻子共度吗？他确实更爱这位妻子，并为她付出了无偿服役十四年的代价。然而，在一天结束时，前往他所爱之人那里的路上，他怎能甘心被劝阻，如果他真是无知摩尼派人所描绘的那样？难道他不会不顾妇女们的意愿，坚持去找美丽的辣黑耳吗？那晚她不仅是他的合法妻子，而且也是按正常顺序轮到的人。这样，他就可以行使丈夫的权力——因为妻子对自己的身体也没有主权，而是丈夫——并在这次服从的安排中，利用对美丽爱情的满足来更成功地强制他的权威。那样的话，妇女们就值得称赞了，因为当他考虑自己的快乐时，她们却在为生子而争辩。事实上，这位有德之人，以男子气概控制着感官欲望，更多地考虑自己应尽的义务而非应得的权益，他没有为自己的快乐使用权势，反而只愿在这相互的义务中成为债务人。因此，他同意将债务偿还给那位应当被偿还的人，即她所希望他偿还的对象。当雅各伯因妻子的私下交易而突然意外地从美丽的妻子转向平凡的妻子时，他没有表现出愤怒或失望，也没有试图说服妻子让他按自己的意愿行事；相反，作为一位正义的丈夫和明智的父亲，他看到妻子们关心生儿育女——这也是他自己在婚姻中所渴望的一切——他认为最好服从她们的权威，希望每个妻子都有一个孩子：因为所有的孩子都是他的，他自己的权威并未受损。仿佛他对她们说：你们自己安排谁做母亲；这对我来说无关紧要，因为无论如何我都是父亲。福斯图斯本来足够聪明，能够看到并赞同这种对欲望的控制和单纯的生育愿望，要不是他的心智被其教派骇人听闻的教义所败坏，这些教义导致他对圣经中的一切都吹毛求疵，并且还教导他将生儿育女——这是婚姻的真正目的——谴责为最大的罪行。\par

51. 现在，在辩护了圣祖的品行并驳斥了这些可憎错误所引发的指控之后，让我们借此机会探求其预象意义，并以信德之虔诚叩门，愿主向我们开启雅各伯四位妻子的预象意义——其中两位是自由的，两位是奴隶。我们看到，在亚巴郎的妻子和婢女身上，宗徒领悟了两种盟约。\BibleRef{迦 4:22}但在那里，每位代表一个；而在这里，因为有两个和两个，所以这种应用并不那么合适。在那里，婢女的儿子被剥夺继承权；但在这里，婢女的儿子与自由妇女的儿子一同获得应许之地：因此这个预象必定有不同的含义。\par

52. 假设两个自由之妻象征新约，藉以我们被召叫至自由，那么有两个是什么意思呢？或许因为，正如细心的读者会发现的，在圣经中，我们被告知在基督的身体内有两种生命：一种是在时间中的，在其中我们受苦；另一种是永恒的，在其中我们将要瞻仰天主的福乐。我们看到一个在主受难中，另一个在祂的复活中。这些妇女的名字指向这一含义：据说\Person{肋阿}意为\Quote{受苦}，而\Person{辣黑耳}意为\Quote{可见的\OriginalTerm{第一原理}{First Principle}}或\Quote{使第一原理可见的\OriginalTerm{圣言}{Word}}。那么，我们可朽坏的人类生命的行动——在其中我们凭信德生活，做许多痛苦的工作，却不知道对我们所关心的人可能产生什么益处——就是肋阿，\Person{雅各伯}的第一个妻子。因此，据说她有一双弱眼。因为可朽坏者的意图是胆怯的，我们的计划是不确定的。再者，对永恒默观天主的盼望伴随着对真理确实而愉快的感知，就是辣黑耳。因此，她被描述为美丽而匀称。这是每位虔诚求知者的挚爱，为此他服事天主的恩宠，藉此我们的罪，虽如朱红，却变白如雪。\BibleRef{依 1:18} 因为\Person{拉班}意为\Quote{使之变白}；我们读到雅各伯为辣黑耳服事拉班。\BibleRef{创 29:17} 没有人转向服事正义，服从于赦免的恩宠，除非为能在使第一原理或天主可见的圣言中生活平安；也就是说，他服事是为辣黑耳，而不是为肋阿。因为人在正义的行为中所喜爱的，不是劳苦和受苦的辛劳。没有人为了其自身欲求这种生命；正如雅各伯不曾欲求肋阿，然而她却嫁给了他，成了他的妻子和众子的母亲。虽然她不能因自身被爱，但主使她作为通向辣黑耳的阶梯而被忍受；之后她因她的孩子而被认可。因此，每个有用的天主的仆人，被带入祂的恩宠中——藉此他的罪被变白——当他如此转向天主时，在他心中、心里和情感中，除了智慧的知识外一无所有。我们常常期望通过实践律法中关于爱近人的七条诫命来获得智慧作为奖赏，这七条诫命是：要孝敬父母；不可奸淫；不可杀人；不可偷盗；不可作假见证；不可贪恋近人的妻子；不可贪恋近人的财产。当一个人尽己所能遵守了这些，却未得到他所渴望和盼望的光明真理的喜乐，反而在世界各种考验的黑暗中发现自己必须承受痛苦的忍耐，或者说他拥抱了肋阿而不是辣黑耳时，如果他的爱中存有坚忍，他就忍受一个以求得另一个；就好像有人对他说：\Quote{另服事七年为得辣黑耳}，他听见了七条新诫命——神贫的人，温良的人，哀恸的人，饥渴慕义的人，怜悯的人，心里洁净的人，缔造和平的人。\BibleRef{玛 5:3} 人若可能，会愿意立即获得可爱而完美智慧的喜乐，而无需忍受行动和受苦的辛劳；但在可朽坏的生命中这是不可能的。这似乎是当对雅各伯说：\BibleQuote{在我们地方没有先嫁少女的规矩} \BibleRef{创 29:26} 时所意指的。长女很可能是指时间上在先的。同样，在人的训练中，践行正义之事的辛劳先于理解真理的喜乐。\par

53. 为此写着：\BibleQuote{你渴慕智慧，就应遵守诫命，上主必赐给你。} \BibleRef{德 1:33} 这些诫命是关于正义的，而这正义是因信德而获得的，充满试探的不确定性；因此，理解是对尚未理解之事的虔诚信仰的赏报。我对这些话\Quote{你渴慕智慧，就应遵守诫命，上主必赐给你}的理解，也在\BibleQuote{除非你们相信，你们将不能理解}这句经文中得到印证；这表明正如义德由信德而来，理解由智慧而来。因此，对于那些急切要求明显真理的人，我们不应谴责这种愿望，而应加以引导，使他们从信德开始，通过善行达到所渴望的目标。德行的生活是劳苦的；所渴望的目标是无云的智慧。有人说：为什么我要相信没有明显证明的事？让我听听一些能揭示万物\OriginalTerm{第一原则}{first principle}的话。这是理性灵魂在追求真理时的一大渴望。回答是：你所渴望的是卓越的，值得你爱；但必须先娶\Person{肋阿}，然后才娶\Person{辣黑耳}。你热切的愿望适当地引导你服从正确的方法，而不是反抗它；因为若不通过这种方法，你就不能达到你如此热切渴望的目标。而一旦达到，对知识的优美形式的拥有，在今世将伴随着正义的辛劳。因为在今生，无论我们对不变之善的感知多么清晰和真实，可朽坏的身体仍然是心灵的负担，地上的帐幕压着多思多虑的理智。所以，目标是一个，但为了它必须经历许多事。\par

54. 这样，\newline{}\Person{雅各伯}有两位自由的妻子；因为二人都是罪过赦免或漂白，即\newline{}\Person{拉班}的女儿。一位被爱，另一位是被承受的。但被承受的那位最早且最多产，好使她能被爱，若不是为她自己，至少为她的儿女。因为义人的劳苦特别在那些他们为天国所生的人身上结果实，他们在许多考验和诱惑中宣讲福音；他们称这些人为自己的喜乐和冠冕（\BibleRef{斐 4:1}），为了这些人，他们\BibleQuote{多受劳苦、多受鞭打、多冒死亡}（\BibleRef{格后 11:23}）——为了这些人，他们\BibleQuote{外有争战，内有恐惧}（\BibleRef{格后 7:5}）。这样的生育最轻易且丰盛地来自信德之言，即被钉十字架的基督的宣讲，这也论及祂的人性，只要它能轻易被理解，以免伤害\newline{}\Person{肋阿}的弱眼。反之，\newline{}\Person{辣黑耳}眼光清朗，为天主疯狂（\BibleRef{格后 5:13}），看见在起初\newline{}\OriginalTerm{圣言}{Word}与天主同在，并愿生育，却不能；因为谁能述说祂的世代？所以献身于默观的生活，为了以无软弱的心灵之眼看见血肉所不能见的事物，却藉着受造之物被理解，并辨认天主永恒的大能和神性那不可言喻的彰显，便寻求从一切事务中闲暇，因此不孕。在这种退隐的习惯中，默想的火焰明亮燃烧，却缺少对人性软弱的同情，以及对人在灾难中需要我们的帮助的同情。这种生活也燃起对子女的渴望（因为愿意教导它所知道的事，而不随从嫉妒的败坏（\BibleRef{智 6:23}）），并看见它姊妹的生活完全忙于工作和生育；它悲伤人们追求那关心他们需要和软弱的德行，而非那拥有神圣不朽教训的德行。这就是当经上说\newline{}\BibleQuote{辣黑耳嫉妒她的姊姊}（\BibleRef{创 30:1}）时的含义。此外，正如对那非物质、故非身体感官对象的纯粹理智感知，不能用出自血肉的言语表达，智慧的道理宁愿借着任何出现的物质形象和说明，在天主真理的心灵中寻得某种寓所，也不愿放弃教导这些事物；因此\newline{}\Person{辣黑耳}宁可让她的丈夫由她的婢女得子，也不愿自己无子。彼耳哈是她婢女的名字，据说意思是年老；因此，即使当我们谈论天主的属神不变的本性时，所提示的观念也关乎身体感官的旧生命。\par

55.\newline{}\Person{肋阿}也由她的婢女得了孩子，出于渴望多子多孙。她的婢女齐耳帕，解释为\newline{}\Quote{开口}。所以\newline{}\Person{肋阿}的婢女代表那些被圣经说成以开口宣讲福音，却未以敞开心怀的人。论及一些人，经上记载：\newline{}\BibleQuote{这百姓用嘴唇尊敬我，他们的心却远离我。}（\BibleRef{依 29:13}）对这样的人，宗徒说：\newline{}\BibleQuote{你教训人不可偷窃，你自己却偷窃吗？你说人不可奸淫，你自己却奸淫吗？}（\BibleRef{罗 2:21}）但即使借此安排，\newline{}\Person{雅各伯}那位自由的妻子——劳苦或坚忍的预像——也能得到孩子作为天国的继承者，主说：\newline{}\BibleQuote{凡他们所吩咐你们的，你们都要遵行；但不要效法他们的行为。}（\BibleRef{玛 23:3}）再者，宗徒的生活在忍受监禁时说：\newline{}\BibleQuote{或是出于假意，或是出于诚心，无论怎样，基督究竟被传开了，为此我就喜欢，并且还要喜欢。}（\BibleRef{斐 1:18}）这是母亲因她众多的子女而有的喜乐，虽是由婢女所生。\par

56. 有一次，肋阿之所以成为母亲，是得益于辣黑耳：辣黑耳用一些曼德拉草作为交换，允许丈夫在她姐姐那里过夜。我知道有人以为，吃这种果子能使不孕妇女受孕，而辣黑耳因渴望子女，才执意从姐姐那里得到此果。但即使辣黑耳当时怀了孕，我也不赞同这说法。既然肋阿当时怀了孕，而且在天主开启辣黑耳子宫之前又生了两个孩子，就没有理由在没有经验证明的情况下，假定曼德拉草有这种功效。我将给出我的解释；比我更有能力的人可以给出更好的解释。虽然这种果子不常见，但我曾有一次，由于它与这段圣经的关联，非常满意地看到了它。我尽可能仔细地检查了这果子，不是借助任何对根的性质或植物功效的高深知识，而只是通过视觉、嗅觉和味觉来了解它或任何人所能了解的东西。我认为它外观好看，气味香甜，但淡而无味；我承认很难说为什么辣黑耳如此渴望它，除非是因为它的稀少和香气。至于为什么这段情节被记载在圣经中——其中妇女的幻想若无意让我们从中学到重要的教训，就不会被视为重要——我唯一能想到的是非常简单的观念：这果子代表好名声；不是少数正义和智慧之人给予人的称赞，而是公众的声誉，它赋予人伟大和名望，其本身并不值得追求，但对于善人努力造福同胞是必不可少的。所以宗徒说，\BibleQuote{应当在那些外人有好名声}；\BibleRef{弟前 3:7} 因为他们虽然并非无误，但他们称赞的光辉和善意评价的芬芳，对那些努力有益于他们的人大有帮助。而这公众声誉，除非教会中地位最高的人投身于活跃生活的辛劳和危险，否则无法获得。因此，肋阿的儿子出去到田里时找到了曼德拉草，也就是说，在外人面前正直行事。另一方面，追求智慧的人远离忙碌的人群，沉浸在宁静的默想中，除非通过那些管理公共事务的人——他们不是为了作领袖，而是为了有用——他永远无法获得丝毫这种公众认可。这些行动和事务之人为了公共利益而操劳，通过善用影响力而获得人民的认可，甚至也为安静的学生和真理探索者的生活赢得了认可；这样，通过肋阿，曼德拉草到了辣黑耳手中。肋阿自己是从她的长子那里得到的，也就是说，以她的生育能力为荣，这代表了劳苦生活的一切有益成果，这种生活暴露于常见的变迁之中；许多人因其烦人的事务而避开这种生活，因为虽然他们可能能够领导，但他们一心向学，将所有精力投入到安静的知识追求中，爱慕辣黑耳的美貌。\par

57. 但是，这好学的生活理应通过让自己被人知晓来获得公众认可，而如果它使追随者隐居，不运用其力量管理教会事务，从而妨碍他普遍有益，它就不能正当地获得这种认可。为此肋阿对妹妹说：\BibleQuote{你夺去了我的丈夫，还算小事吗？你又要夺走我儿子的曼德拉草吗？}\BibleRef{创 30:15} 丈夫代表那些虽然适合活跃生活并能够治理教会，向信徒施行信仰的奥迹，却因热爱学问和追求智慧，而想要放弃一切烦扰的事务，把自己埋在书斋里的人。因此，\BibleQuote{你夺去了我的丈夫，还算小事吗？你又要夺走我儿子的曼德拉草吗？}这句话的意思是：\Quote{学问的生活把公共生活所需要的男人留在隐居之中，这还算小事吗？它还要寻求公众声誉吗？}\par

58. 要公正地获得这种声誉，\Person{辣黑耳}就把她的丈夫让给她姐姐过夜；也就是说，那些因经商才华而适合治理的人，必须为了公共利益同意承担公务生活的重担和艰辛；以免他们所倾心从事的智慧追求受到非议，并且不能在民众中获得那以果实为代表的好评，而这种好评对于鼓励他们的弟子是必要的。但经商的生活必须强加于他们。这由\Person{肋阿}在\Person{雅各伯}从田间回来时遇见他，并拉住他，说：\BibleQuote{你要进到我这里来，因为我用我儿子的曼德拉草雇了你。}\BibleRef{创 30:16} 来清楚表明。似乎她在说：你愿意你所爱的知识被人好评吗？就不要逃避工作的辛劳。同样的事在教会中经常发生。我们所读的由我们经验中所遇到的来解释。我们岂不处处看到一些人从世俗工作中出来，寻求闲暇来研究和默想真理，即他们所爱的\Person{辣黑耳}，却在半路上被教会事务截住，这些事务要求他们投入工作，好像\Person{肋阿}对他们说：你必须进到我这里来吗？当这样的人真诚地施行天主的奥秘，以致在这世界的黑夜中生出了在信德中的子女时，民众的认可也获得了那种生活，正是为了热爱这种生活，他们才被引导放弃世俗的追求，并且从接受这种生活之后，他们被召唤去承担仁慈的治理工作。在他们的所有劳苦中，他们主要追求这一点，即他们所选择的生活方式可能获得更大更广的声誉，因为它为民众提供了这样的领袖；正如\Person{雅各伯}同意与\Person{肋阿}同房，好让\Person{辣黑耳}获得那芬芳而悦目的果实。\Person{辣黑耳}也因天主的仁慈，经过一段时间后自己生了一个孩子，但那是过了相当时间以后；因为很少有这样的事：一种健全的、虽只是部分的领悟，而没有肉性的观念，来领会像这样的神圣智慧教训：\BibleQuote{在起初已有圣言，圣言与天主同在，圣言就是天主。}\BibleRef{若 1:1}\par

59. 以上这些必须足以作为对\Person{福斯图斯}针对三位圣祖\Person{亚巴郎}、\Person{依撒格}和\Person{雅各伯}所提出的虚假指控的答复；这些圣祖的名号，正是天主教会所敬礼的天主乐意用来称呼自己的。这里不是讨论这三人的功劳和虔诚的地方，也不是讨论他们先知性尊位的地方，因为这种尊位是属血肉的人心无法理解的。在这篇论文中，就足够为他们辩护，反驳恶毒谎言的诽谤性攻击，以防那些以吹毛求疵和敌对精神阅读圣经的人，因为他们企图抹黑那些在圣经中备受推崇的人物的品格，就以为他们已经证明这些书的圣洁和有益之处中有什么漏洞。\par

60. 还应当补充说，\Person{罗特}，\Person{亚巴郎}的兄弟，即血亲，不应被列为与那些人同等，就是天主所说的：\BibleQuote{我是\Person{亚巴郎}的天主、\Person{依撒格}的天主、\Person{雅各伯}的天主。}他也不属于圣经中见证为始终正义的人，虽然他在\Place{索多玛}度过了虔诚和有德的生活，并显示了值得称赞的好客，因此他从火中被救出，并且天主为了他叔父\Person{亚巴郎}的缘故，赐给他后裔一块土地居住。由于这些原因，他在圣经中得到称扬——并非因为他的放纵或乱伦。但当我们发现同一个人有善行和恶行记录时，我们必须从一个中得到警戒，从另一个中得到榜样。那么，正如\Person{罗特}的罪——我们被告知他在此罪之前是正义的——不但没有给天主的品德或圣经的真实性带来污点，反而要求我们赞许和钦佩其记录如同忠实的镜子，它不仅反映靠近它的人的美善和完美，也反映他们的缺点和畸形；更甚者，在\Person{犹大}与他儿媳同寝的事上，我们可以看到对叙述的责备是多么没有根据。神圣的记录有一种权威，使它远非仅超过一小撮摩尼教徒的挑剔，而是超过整个外邦世界既定的敌意；因为，为证实其主张，我们看到它已经按照基督教的规则，使几乎所有的人从他们的偶像崇拜迷信归向唯一天主的崇拜。它不是通过暴力和战争，而是通过真理不可抗拒的力量征服了世界。那么，圣经在哪里赞扬\Person{犹大}呢？除了在他父亲祝福中他超群出众，因为预言基督将出自他的支派降生成人之外，还有什么好话说到他呢？\BibleRef{创 49:8}\par

61. \Person{犹大}，如\Person{福斯图斯}所说，犯了奸淫；此外我们还能指控他将自己的弟弟卖到埃及。这能贬损光明吗？光明在揭示一切时也将丑陋显露。同样，圣经的品格也不受它所记载的恶行影响。毫无疑问，永恒律法要求并禁止违犯自然秩序，夫妻交合应仅为生育子女，且在婚姻庆典之后，以保持和平的纽带。因此，仅为满足邪情私欲而卖淫，为神圣的永恒律法所谴责。购买他人的堕落，使购买者蒙羞；因此，如果犹大明知而与儿媳同寝，罪过会更大（因为，如主所说，\BibleQuote{夫妻不再是两个人，而是一体了}（\BibleRef{玛 19:6}），儿媳如同女儿），但就他自己的意图而言，他与娼妓交合已使自己蒙羞。另一方面，那欺骗公公的女子，并非由于放荡或贪恋不义之财，而是出于渴望从这个家族得子。在两个儿子都失望、且未能得到第三子后，她以诡计从公公那里得子；她所得的抵押物并非作为装饰，而是作为信物。当然，无子胜于不经婚姻而为母亲。然而，她希望公公成为她孩子的父亲，这与对公公的犯罪情感截然不同。当她被犹大下令带出处死时，她出示手杖、项链和戒指，说孩子的父亲就是给她这些信物的人，犹大认出它们，说：\Quote{她比我更有义}——不是赞美她，而是谴责自己。他更责备自己不正当的情欲，而非她求子的愿望，那情欲使他去找一个他以为是娼妓的人。类似地，经上说有些人使\Place{索多玛}成义（\BibleRef{则 16:52}），意思是他们的罪如此之大，以至索多玛相比之下显得义。即使退一步说，这女子并非被认为罪较轻，而是被公公实际赞美，但根据禁止违犯自然秩序——既关乎身体，又首先并主要关乎灵魂——的公义永恒律法，她因不遵守既定的婚姻礼仪而成为罪犯，那么一个罪人赞美另一个罪人，有什么希奇呢？\par

62. \Person{福斯图斯}和\OriginalTerm{摩尼教}{Manichaeism}一般的错误，在于以为这些反对意见能证明什么损害我们，仿佛我们对圣经的崇敬和对圣经权威的承认，迫使我们认同其中所记的一切恶行；然而，我们对圣经的敬意越大，就越应当定断圣经真理本身教导我们定断之事。在圣经中，一切奸淫和通奸都为神圣律法所谴责；因此，当此类行为被叙述而没有明确谴责时，其用意不是要我们赞美它们，而是要我们亲自审判它们。人人都厌恶福音中\Person{黑落德}的残忍，当他听说\Person{基督}诞生而心里不安时，下令屠杀众多婴孩（\BibleRef{玛 2:16}）。但这仅仅是叙述而无谴责。或者，如果摩尼教式的荒谬竟敢否认这个叙述的真实性——因为他们不承认基督的诞生，即令黑落德不安的事——那么让他们读一读犹太人盲目狂怒的记述，这记述毫无责难之词，尽管人人都同样感到憎恶。\par

63. 但有人说，\Person{犹大}与儿媳同寝，却被算作十二族长之一。\Person{犹达斯}出卖了主，难道不也是十二宗徒之一吗？并且这其中一个原是魔鬼，不也被差遣与他们一同传福音吗（\BibleRef{若 6:70}）？对此，会有人说：犹达斯犯罪后上吊自尽，被从宗徒中除去；而犹大在恶行之后，不仅与兄弟们一同受祝福，还从他父亲那里获得特别的尊荣和认可——他父亲在圣经中备受称赞。但从中要学的主要教训是，这个预言并非指犹大，而是指基督，他早已被预言要从犹大支派降生成人；并且提及犹大之罪的真正原因，在于它有助于教导我们在其父的话语中寻找另一层含义，这些话语所表达的赞美，明显不适用于行恶的他。\par

64. 无疑，福斯图斯诽谤的目的在于损害这一断言：基督出生于犹大支派。尤其是在玛窦所记载的家谱中，我们发现了匝辣的名字，他是这妇人塔玛尔为犹大所生的。如果福斯图斯只想指责雅各伯家族，而非基督的诞生，他本可以选取长子勒乌本的事例，他犯了玷污父亲床笫的逆伦之罪，关于这种淫乱，宗徒说，连在外邦人中也没有这样的事。\BibleRef{格前 5:1} 雅各伯也在祝福中提及此事，谴责他儿子的可耻行为。福斯图斯本可以提出这个，因为勒乌本似乎犯了故意乱伦之罪，且此事中并无娼妓的伪装；然而，那是因为塔玛尔的行为——她只想生育子女——在福斯图斯眼中，比出于犯罪情欲的行为更可憎；并且，他难道不是想通过诽谤基督的先祖来贬低降生成人吗？不幸的是，福斯图斯不知道，最真实、最诚实的救主不仅是言语上的导师，也是诞生上的导师。祂肉身的起源给那些将从万民中信从祂的人带来教训：即他们祖先的罪不应成为他们的阻碍。此外，那要召唤好人和坏人参加祂婚宴的新郎，\BibleRef{玛 22:10} 乐意使自己与宾客相似，从好人和坏人中出生。祂由此证实了逾越节祭献的象征——其中命令要取用绵羊或山羊作为逾越节羔羊——即出自义人或恶人。\BibleRef{出 12:3} 祂始终保持着天主性与人性的指示：作为人，祂同意既有好父亲也有坏父亲；作为天主，祂选择了由童贞女而生的奇迹性诞生。\par

65. 因此，福斯图斯对圣经的攻击之亵渎只能损害他自己；因为他所攻击的圣经，如今理应成为普世敬礼的对象。如前所述，圣经如同一面忠实的镜子，其画像绝无奉承；它要么亲自判定人类行为值得赞许或责难，要么让读者自行判断。而且，它不仅区分可责难者和可赞许者，还注意到可责难者应受赞扬、可赞许者应受责难的情况。因此，虽然撒乌耳是可责难的，但他那样仔细地查明在诅咒期间谁吃了食物，并遵从天主的命令宣布严厉判决，这仍然是值得赞扬的。\EditorialNote{撒慕尔纪上十四章}同样，他驱逐境内那些招魂问卜之人也是正确的。\BibleRef{撒上 28:3} 虽然达味是可赞许的，但我们不应赞同或效法他的罪过，天主曾透过先知责备他。同样，般雀·比拉多宣布主无罪，不顾犹太人的控告，这并没有错；\BibleRef{若 19:4} 伯多禄三次否认主也不是值得赞许的；此外，当基督称他为撒殚，因为他不明了天主的事，想劝阻基督接受苦难——即我们的救恩时，他也不是值得赞许的。在这里，伯多禄刚被称为有福的，紧接着就被称为撒殚。从他的宗徒职分和殉道者的荣冠，我们可以看到哪一种性格真正属于他。\par

66. 至于\Person{达味}，我们也读到他有善行与恶行。但是，\Person{达味}的力量何在，他成功的秘诀是什么，这对于瞎眼恶意的\Person{法乌斯托}——他以此攻击圣经与圣人们——并不明显，但对于虔诚的明辨——它顺服于天主的权威，同时正确判断人的行为——则足够清楚。摩尼教徒若读圣经，便会发现天主责备\Person{达味}比\Person{法乌斯托}更严厉。\EditorialNote{撒下 12}但他们还会读到他的忏悔之祭，他对那无情嗜血的敌人超乎寻常的温良——\Person{达味}既虔诚又勇敢，当那敌人屡次落入他手中时，他都释放了他，未加伤害。他们会读到他在天主的惩罚下令人难忘的谦卑——君王的颈项如此屈服于主的轭下，以致他极其忍耐地承受了来自敌人的尖刻嘲讽，尽管他本身武装，并且有武装的随从。当他的同伴为这些对王所说的话而发怒，即将报复那嘲笑者时，他温和地阻止了他，呼求敬畏天主来支持他的王命，并说这灾祸临到他身上，是天主给他的惩罚，是派那人来咒骂他的。\EditorialNote{撒下 16}他们会读到，他如何怀着牧者对托付给他之羊群的爱，愿意为他们而死——当他数点人民之后，天主认为适合藉减少他所夸耀的数目来惩罚他有罪的骄傲。在这场毁灭中，毫无不义的天主，以祂隐秘的审判，既夺去了祂知道不配活命之人的性命，也藉减少数目医治了那以人民数目自夸的虚荣。他们会读到，他因敬畏天主而尊重圣油傅中基督的预象，这使\Person{达味}心里自责，后悔暗中割下\Person{撒乌耳}外衣的一角，为向他证明自己无意杀他，尽管他本可以这样做。他们会读到他对待子女的明智行为，以及他对他们的柔情——当一子患病时，他流了许多泪，以极大的自卑向主恳求；但当那无辜的孩子死后，他却不哀悼；又如，当他那年轻的儿子以悖逆的敌意行可耻的乱伦，玷污父亲的床，并且发起弑父之战时，他愿他活着，又在他被杀时为他哭泣；因为他想到犯下如此罪行的灵魂的永罚，并希望他能活着，藉被带到顺服与悔改来逃避这永罚。这些以及许多其他值得赞扬、堪为模范的事迹，若坦诚查考圣经叙述，都能在这位圣人身上看到，尤其是当我们以谦卑的虔诚与真诚的信德来看天主的判断时——祂知道\Person{达味}心中的秘密，并以无误的察看如此赞许\Person{达味}，以致推崇他为儿子们的典范。\par

67. 必定是因天主圣神如此洞察\Person{达味}内心深处的缘故，当他受先知责斥而说\Quote{我犯了罪}时，在他作了这简短告解之后，便立即被认为堪当被告知他已获得宽恕——也就是说，他被接纳进入永生。因为他未能逃脱父权的惩戒，正如天主在祂的警告中所说的：虽然他藉告解获得了永久的豁免，却仍要受暂时的惩罚所考验。这显著地证明了\Person{达味}信德的坚强以及他温顺服从的精神：当先知告诉他天主已宽恕了他之后，尽管所警告的后果仍然发生了，他既没有指责先知欺骗了他，也没有抱怨天主以宣布宽恕来嘲弄他。这位深具圣德的人，他的灵魂向上主高举，而不是反对天主，他知道，若非主仁慈地接纳了他的告解与悔改，他的罪恶该受永罚。因此，当他反而受暂时的惩戒之苦时，他看出，在宽恕有效的同时，也提供了有益的纪律。\Person{撒乌耳}也曾在受\Person{撒慕尔}责斥时说过\BibleQuote{我犯了罪}。\BibleRef{撒上 15:24}那么，为什么他不像\Person{达味}那样被认为堪当被告知主已赦免了他的罪呢？难道天主偏袒人吗？绝非如此。虽然人耳听到的是同样的话，天主的慧眼却看到了内心的差异。我们从这些事当学的功课是：天国就在我们心里，\BibleRef{路 17:28}我们必须从内心最深处敬拜天主，以便口舌因心里所充满而说话，而不是像古时的人那样用嘴唇尊敬祂，心却远离祂。我们也可以学习判断人——我们看不见他们的心——只能像天主那样判断，祂看见我们所不能看见的，并且祂不能被偏见或误导。我们既以圣经的崇高权威得到了天主对\Person{达味}之看法的最明确宣告，便可将那些持不同意见者的鲁莽视为荒谬可悲。关于这些古人之品格的圣经权威，有它们所含有的预言——现在正在应验——的见证为支持。\par

68. 我们在\WorkTitle{福音}中也看到同样的事：魔鬼用\Person{伯多禄}所说的话承认基督是天主子，但内心却截然不同。因此，尽管话语相同，\Person{伯多禄}因信德受到称赞，而魔鬼的不敬虔却被制止。因为基督并非凭人的感官，而是凭天主的全知，能够审视并无误地辨别话语的源头。此外，有许多人承认基督是永生天主之子，却没有得到像\Person{伯多禄}那样的赞许——不仅包括那些在那日要说\Quote{主啊，主啊}，却受到\Quote{离开我吧}判决的人，也包括那些将站在右边的人。他们可能从未否认过基督，可能从未反对祂为我们的救恩而受难，可能从未强迫外邦人像犹太人那样行事；\BibleRef{迦 2:14}然而，他们不会得到与\Person{伯多禄}同等的尊荣，因为\Person{伯多禄}虽曾做过这一切，却要坐在十二个宝座之一，审判不仅十二支派，而且连天使也要审判。同样，许多从未贪恋他人妻子、从未像\Person{达味}那样谋害丈夫的人，永远不会达到\Person{达味}在天主宠爱中所占据的地位。事物本身如此不可取以致必须完全弃绝，与日后可能出现的丰盛收获之间有天壤之别。因为农夫最喜爱那些经过除草——也许是除去大蓟——之后获得百倍收成的田地，而不是那些从未长过蓟、几乎只结三十倍的田地。\par

69. 同样，\Person{梅瑟}——他在天主全家尽忠，是圣洁、公义、良善法律的仆人；宗徒在引用的经文中论及他的品格；\BibleRef{希 3:5}他也是那些虽不赐予救恩却应许救主的表记的仆人，救主亲自显示这一点说：\Quote{如果你们相信\Person{梅瑟}，你们也会相信我，因为他论及我写了。}——从这段经文我们已经充分回答了摩尼教徒的狂妄诘难。这位\Person{梅瑟}是永生、真实、至高天主的仆人，祂创造天地，不是出于异质的东西，而是从无中创造；不是出于必然的压力，而是出于美善的丰盈；不是通过祂肢体的受苦，而是通过祂言语的权能。这位\Person{梅瑟}谦卑地推辞这项崇高使命，却顺从地接受、忠实地持守、勤勉地完成；他以警醒治理百姓，以严厉责备他们，以热忱爱他们，以忍耐包容他们，在天主面前为他的子民站立，接受祂的训诲，平息祂的义怒。这位伟大而良善的人，不应凭\Person{福斯塔斯}的恶意描述来判断，而应凭天主所说的话来判断，祂的话语真实表达了对这个人的真正看法——祂认识他，因为祂造了他。因为人的罪过也为天主所知，尽管祂不是罪过的创造者；但祂以审判者的身份留意那些拒绝承认罪过的人，并以父亲的身份宽恕那些忏悔的人。我们对祂的仆人\Person{梅瑟}，如上述所描述的，爱慕并钦佩，并尽己所能效法他，尽管我们远不及他的功德，虽然我们未曾杀过一个埃及人，未曾抢夺过任何人，也未曾发动过任何战争；\Person{梅瑟}的这些行为，一种情况是出于这位百姓未来保卫者的热忱，另一种情况是出于天主的命令。\par

70. 可以证明的是，虽然梅瑟在没有受到天主的命令下杀了那个埃及人，但这一行为出于天主的许可，因为从梅瑟的先知性特质来看，它预象了未来的一些事。不过，我现在不使用这个论证，而是将这一行为视为没有象征意义。那么，按照永恒法律，一个没有合法权力的人杀死这人是有罪的，即使他是个坏人，而且是攻击者。但在那些将要成就伟大德性的人心中，常有早期的恶习萌芽，我们仍可在其中看出某种特定德性的禀性，这禀性会在心灵得到适当培育时显现。正如农夫看到一块地长出大量作物，虽然是杂草，却断言它适合种谷物；或看到必须拔除的野生蔓草，仍认为土地适合种有用的葡萄；或看到山上长满野橄榄树，便断定经过培育它会结出好果子：同样，那引导梅瑟自行执法，以防止他住在陌生人中的兄弟所受的冤屈得不到伸张的心性，并非不适合产生德性，但由于缺乏培育，它以不正当的方式显示出多产性。那一位后来借祂的天使召唤梅瑟上了西乃山，交给他神圣使命，将以色列人民从埃及解放出来，并借着荆棘火焰而不被烧毁的神迹显现训练他服从，教导他他的职务；那也是同一位，通过从天上的呼唤对正在迫害教会的扫禄说话，使他谦卑，提拔他，并激励他；或者用象征性的话来说，借着这一击，祂砍下树枝，接上它，并使它结果实。因为保禄的猛烈精力，在他对祖传传统的热忱中迫害教会，以为他是在事奉天主，就像一片杂草显示出丰产的大迹象。同样的情况也发生在伯多禄身上，当他拔剑出鞘来保护主，削去一个攻击者的右耳时，主以一种近乎威胁的口气责备他，说：\BibleQuote{把你的剑放回剑鞘里；因为凡动剑的，必死于剑下。}\BibleRef{玛 26:51} 动剑，就是未经合法权威的许可，使用武器攻击人的性命。主确实曾告诉祂的门徒带剑；但祂没有告诉他们使用它。但是，在这罪之后伯多禄成为教会的牧者，并不比梅瑟在击打埃及人之后成为会众的领袖更不合适。在这两种情况下，过犯并非出于根深蒂固的残忍，而是出于一种可以纠正的仓促热情。两者都是对伤害的愤慨，在一种情况下伴随对兄弟的爱，在另一种情况下伴随对主的爱，虽仍属血肉之爱。这里有需要被制伏或拔除的邪恶；但具有如此潜力的心灵只需要像好土一样被培育，就能在德性上结出果实。\par

71. 那么，至于福斯图斯对夺取埃及人财物一事的反对，他不知自己在说什么。在这件事上，梅瑟不仅没有犯罪，若不这样做反而是罪。这是出于天主的命令，祂知道人的行为和人心，能决定每个人应该遭受什么，以及通过谁的手来执行。那时的人民仍然属血肉，沉迷于世俗的情感；而埃及人公开叛逆天主，因为他们使用天主的受造物黄金来服务偶像，羞辱造物主，并且他们严重压迫外方人，强迫他们无偿劳动。因此埃及人应受惩罚，而以色列人被适宜地用来施加惩罚。或许，这与其说是一个命令，不如说是天主许可希伯来人在此事上按自己的意愿行事；而天主通过梅瑟传达的信息，只是希望他们得知祂的许可。天主在这件事上对人民所说的话，也可能有神秘的缘故。无论如何，天主的命令要顺服地接受，而不是争论。宗徒说：\BibleQuote{谁认识主的心意？或者谁做过祂的参谋？}\BibleRef{罗 11:34} 那么，不论理由是我所说的，还是在天主隐秘的安排中，有一些未知的原因使祂通过梅瑟告诉人民向埃及人借东西并带走，这一点是确定的：这些话是出于某种良善的理由，而梅瑟不能合法地做不同于天主所吩咐的事，将命令的理由留给天主，而仆人的本分是服从。\par

72. 但是，\Person{福斯托}说，不能承认真实而良善的天主曾下达这样的命令。我回答说，这样的命令只能由真实而良善的天主正确地发出，唯有祂知道在每种情况下合适的命令，并且唯有祂不能对任何人施加不应得的痛苦。人心的这种无知虚假的良善同样可以否认基督所说的，并反对由良善天主使恶人受苦，当祂对天使说：\BibleQuote{先将莠子捆成捆，好燃烧它们。}然而，当仆人们想要提前这样做时，他们被阻止了：\BibleQuote{免得你们收集莠子时，连麦子也一起拔出来。}\BibleRef{玛 13:29}因此，唯有真实而良善的天主知道何时、向谁、通过谁命令或容许任何事情。同样，这种人性的良善，或者更确切说是愚昧，可能会反对主容许魔鬼进入猪群，它们怀着恶意的意图请求这样做，\BibleRef{玛 8:31}尤其是\OriginalTerm{摩尼教徒}{Manichaeans}相信不仅猪，甚至最卑微的昆虫都有人类的灵魂。但撇开这些荒谬的观念，这是不可否认的，即我们的主耶稣基督，天主的独生子，因此是真实而良善的天主，应魔鬼的请求容许了他人猪群的毁灭，这暗示了生命的丧失和大量财产的损失。没有人会疯狂到认为基督不能将魔鬼从人身上驱逐出去，而不通过猪的毁灭来满足它们的恶意。因此，如果所有本性的造物主和管理者，在其眷顾中——虽神秘却始终公正——放纵了那些已注定永火的失丧灵魂的暴力和不义倾向，那么为什么埃及人——不义的压迫者——不应该被希伯来人——一个自由的民族——剥夺，他们应为自己被迫且痛苦的劳役索取报酬，尤其是他们失去的世俗财产被埃及人用于不敬的仪式，以辱没造物主？然而，如果\Person{梅瑟}发起了这项命令，或者人民自发地做了，无疑这会是犯罪的；也许人民确实犯了罪，不是在做天主命令或容许的事情上，而是在他们对所取之物的某种欲望上。神圣权威对这一行动所给予的容许，是符合那位的公正而良善的旨意的，祂使用惩罚既是为了约束恶人，也是为了教育祂自己的子民；祂也知道如何给那些能够承担的人更先进的诫命，同时在对待软弱者时从较低的标准开始。至于\Person{梅瑟}，既不能因贪恋财产而责备他，也不能因在任何事例中置疑神圣权威而责备他。\par

73. 按照永恒法律，它要求保存自然秩序，并禁止违反它，有些行为具有中性特征，因此如果人们在未被召唤的情况下做这些事，他们会因冒失而受责备，而当被要求做时，他们则应受称赞。行为、行为者以及行为的权威在自然秩序中都非常重要。\Person{亚巴郎}自愿牺牲他的儿子是令人震惊的疯狂。但他在天主的命令下这样做，证明了他的忠信和服从。这一点被真理的声音如此大声地宣告，以至于\Person{福斯托}急切地寻找某些过错，最终沦落到诽谤指控，却不敢攻击这一行为。他几乎不可能忘记如此著名的事迹，它无需任何研究或反思就自动浮现在脑海中，事实上被那么多口舌传颂，在那么多地方描绘，以至于没有人能假装对此视而不见或充耳不闻。因此，如果\Person{亚巴郎}自愿杀死他的儿子是不自然的，那么他在天主的命令下这样做不仅显示无罪，而且值得称赞的顺从，为什么\Person{福斯托}责备\Person{梅瑟}夺取埃及人的财物？你对纯粹人类行为的不赞成应因对神圣认可的尊重而受到约束。你敢责备天主自己渴望这样的行为吗？那么，\BibleQuote{退到我后面去吧，\Person{撒殚}！因为你所体会的不是天主的事，而是人的事。}但愿这一责备能在你身上产生像在\Person{伯多禄}身上一样的效果，以至于此后你能宣讲关于天主的真理，而你目前却凭软弱的判断指责它！正如\Person{伯多禄}成为热心的使者，向外邦人宣布他起初所反对的，当主提到这是祂的旨意时。\par

74. 现在，如果这一解释足以满足人类的顽固和对正确行为的歪曲误解，即人之情欲放纵与僭越，与服从天主的命令（天主知道该允许或命令什么，也知道时间、人物以及每种情况下的应有作为或忍受）之间的巨大差异，那么\Person{梅瑟}的战争叙述就不会引起惊讶或憎恶，因为在他奉天命进行的战争中，他表现出的不是残暴，而是服从；而天主在命令时，也并非出于残忍，而是出于公义的报应，给每人其所应得，并警告那些需要警告的人。战争的恶是什么？是某些人的死亡——他们迟早要死，好使其他人在和平的臣服中生活吗？这不过是懦弱的厌恶，而非任何宗教情感。战争中真正的恶是暴力之爱、复仇之残忍、激烈而不可和解的敌意、狂野的抵抗、权力之欲，以及诸如此类；而通常正是为了惩罚这些事，当需要以武力实施惩罚时，善人们才在服从天主或某种合法权威之下发动战争，当他们发现自己在人事治理上所处的地位要求他们以正当行为，或使他人以这种方式行事时。否则，\Person{若翰}当那些士兵前来受洗问说：\Quote{我们该作什么？}时，他就会回答：\Quote{扔掉你们的武器；放弃服役；绝不可击打、伤害或致残任何人。}但\Person{若翰}知道，战斗中的这些行为并非凶杀，而是法律所授权，且士兵们并非以此为自己报仇，而是保卫公共安全，于是他回答说：\BibleQuote{不要对人施暴，不要诬告人，对你们的粮饷应当知足。}\BibleRef{路 3:14}但摩尼教徒惯于毁谤\Person{若翰}，让他们听主\Person{耶稣基督}亲自命令将这钱给\Person{凯撒}，正是\Person{若翰}告诉士兵们要知足的。祂说：\BibleQuote{凯撒的，就应归还凯撒。}\BibleRef{玛 22:21}再者，关于那位百夫长，他曾说：\BibleQuote{我也是属人权下的人，我也有兵士在我权下：我对这个说：你去，他就去；对那个说：你来，他就来；对我的奴仆说：你做这个，他就做。}基督曾称赞他的信德；\BibleRef{玛 8:9}祂并未叫他离开军旅。但这里无需深入讨论正义与非正义之道的冗长问题。\par

75. 许多取决于人们发动战争的原因及其所拥有的权威；因为寻求人类和平的自然秩序规定，君王若认为适宜，应有权发动战争，而士兵则应为了共同体的和平与安全履行其军事职责。当战争是出于服从天主的命令而发动，祂要斥责、贬抑或粉碎人类的骄傲时，必须承认这是正义的战争；因为即使那些源于人类情欲的战争，也不能损害天主的永恒福祉，甚至不能伤害祂的圣徒；因为在考验他们的耐心、磨练他们的精神以及承受慈父般的管教中，他们得益而非受损。除非从上赐予，没有人能有什么权力对抗他们。因为\BibleQuote{没有权柄不是从天主来的}\BibleRef{罗 13:1}，祂或是命令或是容许。因此，既然一个义人，或许在一个不虔敬的君王之下服役，可以按照君主的命令战斗，履行其在国家中地位的职责——因为在某些情况下，战斗显然是天主的旨意，而在其他情况不甚明确时，君主的命令可能是不义的，但士兵是无辜的，因为他的地位使服从成为义务——那么，那些奉天主权威发动战争的人，更应无可指责，因为凡服事祂的人都知道，祂绝不会要求错误的事。\par

76. 如果有人以为天主不能命令战争，因为后来\Person{主耶稣基督}曾说：\BibleQuote{我给你们说：不要抵抗恶人；如果有人打你的右脸，你把左脸也转给他。}\BibleRef{玛 5:39} 那么回答是：这里所要求的不是外在行动，而是内心态度。德行的神圣宝座是人心，我们祖先、古时的义人的心也是如此。但秩序要求对事件作如此安排，对时代作如此区分，以便首先显示，甚至现世的福祉（因为人认为世间王国和战胜敌人便是福祉，而普世不敬神者的团体不断向偶像和魔鬼祈求的正是这些东西）也完全受独一无二真天主的控制和支配。因此，在旧约之下，天国的奥秘——将在适当时机启示出来——被隐藏，甚至被遮盖，以尘世应许的外表。但当时期一满，新约——隐藏在旧约预象下——被启示出来时，就要为真理作明确的见证：还有另一生命，为这生命应当轻视今生；还有另一王国，为这王国应当忍耐世上所有王国的反对。因此，殉道者（意思是见证人）这一名称，给了那些按天主的旨意，通过他们的宣认、受苦和死亡，为此作见证的人。这些见证人的数目如此庞大，以至于如果\Person{基督}愿意——祂曾从天高声呼召\Person{扫禄}，将他从狼变为羊，然后派遣他进入狼群中——将他们全部集结为一支军队，并赐予他们战事胜利，如同祂赐予希伯来人一样，那么哪个民族能抵挡他们？哪个王国能不被征服呢？但由于新约的教义是：我们必须事奉天主，不是为今生的暂时福乐，而是为来世的永恒幸福，这一真理通过为那幸福而忍耐常被称为逆境的事，得到了极其有力的证实。因此，在时期圆满时，天主子——由女人所生，生于法律之下，为救赎那在法律之下的人，按肉身说是\Person{达味}的后裔——派遣祂的门徒如同羊进入狼群，并吩咐他们不要害怕那些能杀害肉身却不能杀害灵魂的人，并应许甚至肉身也将完全恢复，连一根头发也不会失落。祂命令\Person{伯多禄}的剑归回鞘内，将祂敌人被割下的耳朵复原如初。祂说祂能求得十二营以上的天使毁灭祂的敌人，但祂必须饮下祂父的旨意所赐给祂的杯。祂以身作则饮了这杯，然后递给祂的门徒，从而在言行上都彰显了忍耐的恩宠。因此，天主从死者中复活了祂，赐给祂一个超越万名的名字，使天上、地上和地下的众生，因\Person{耶稣}的名字都屈膝叩拜，并承认\Person{耶稣基督}是主，以光荣天主圣父。\BibleRef{斐 2:9} 因此，圣祖和先知们在此世拥有王国，为表明这些王国也是由天主赐予和收回的；宗徒和殉道者在此世没有王国，为表明天国的更为优越的可欲性。然而，先知们在那些时代也可以为真理而死，正如主亲自说：\BibleQuote{从\Person{亚伯尔}的血到\Person{匝加利亚}的血。}\BibleRef{玛 23:35} 而在这些日子，自从\Person{基督}的圣咏——以\Person{撒罗满}为预象，撒罗满意为和平使者，因为\Person{基督}是我们的和平\BibleRef{弗 2:14}——中所预言的开始应验：\BibleQuote{世上的众王都要朝拜祂，万民都要事奉祂。}我们已见到基督信仰的皇帝们，他们完全信赖\Person{基督}，战胜了不敬神的敌人——这些敌人寄望于偶像崇拜和魔鬼崇拜的仪式——取得了辉煌胜利。这一事实有公开且不可否认的证据：一方面，魔鬼的预言被证明是虚妄的；另一方面，圣人的预言成了支持；我们现在有记载这些事实的著作。\par

77. 若我们愚昧的反对者因天主在新约恩宠尚未显明时，赐给旧约职员的诫命，与如今旧约的隐晦已除去、赐给新约宣讲者的诫命之间的差异而感到惊讶，他们必会发觉基督自己也在不同时期说不同的话。祂说：\Quote{我派遣你们的时候，不带钱囊、不带口袋、不带鞋，你们缺少什么没有？}他们回答说：\Quote{什么也没有。}祂遂对他们说：\Quote{可是如今，有钱囊的，要带着，口袋也一样；没有剑的，要卖掉自己的外衣，去买一把。}如果摩尼教人在旧约和新约中找到这样不同的话，他们必会宣称这是两约彼此对立的证据。但这里的差异出自同一位的言论。一时祂说：\Quote{我派遣你们不带钱囊、口袋、鞋，你们什么也不缺；}另一时：\Quote{如今有钱囊的，要带着，口袋也一样；有外衣的，要卖掉它买一把剑。}这岂不说明，诫命、劝谕和许可如何能随着不同时期所需的不同安排而改变，且毫无矛盾？若说带钱囊、口袋和买剑的命令有象征意义，那么为何同一上主在古时命令先知们作战，而又禁止宗徒们作战，就不能也有象征意义呢？我们在所引的福音经文中看到，主所说的话由祂的门徒付诸实行。因为，除了起初不带钱囊或口袋而一无所缺（从主的问话和他们的回答可知），如今祂提到买剑，他们说：\Quote{看，这里有两把剑。}祂回答说：\Quote{够了。}因此我们发现伯多禄带着武器，砍掉了来人的耳朵，当时他自发的鲁莽被制止，因为他虽然被告知要带剑，却未被允许使用它。无疑，主要求他们携带武器，却禁止使用，这是奥秘的。但是，祂分施合适的诫命，而他们无保留地服从，这是祂的职权。\par

78. 因此，指控梅瑟作战纯属无根据的诽谤，因为擅自作战的害处，还不及奉命作战而拒绝服从的大。至于竟敢指责天主亲自如此命令，或不信一位公义良善的天主会如此命令，这至少显示人无法考虑到：在天主眷顾（它渗透从最高到最低的一切事物）的眼光中，时间既不能增添什么，也不能减少什么；但一切按各事的本性或功过，或进行、或来临、或停留；而在人中，正确的意志与神圣法律合一，不受约束的情欲被神圣法律的秩序所约束；这样，善人只意愿所命之事，恶人只能做所许之事，同时他因那不义地意愿之事而受惩罚。如此，在一切令人类脆弱感到震惊可怕的事中，真正的恶是不义；其余只是自然属性或道德过失的结果。这义德见于每一情形：人若为自己本身而爱那仅作为手段可欲之物，或为其他事物而寻求那本应为自身所爱之物，便是不义。因为如此，他尽其所能扰乱自身中那永恒律法要求我们遵守的自然秩序。反之，人若只按天主指定的目的使用事物，以天主为万物的终向而享受天主，并在天主内并为天主而享受自己和朋友，便是义人。因为在一个朋友中爱那对天主的爱，便是为天主而爱那朋友。但义与不义，若要成为行为，必是自愿的；否则便没有公义的赏罚；这是任何理智正常的人都不会主张的。那阻止人认识其本分或完全行其所愿的无知与软弱，属于天主隐秘的惩罚安排和祂不可测量的判断，因为在祂毫无不义。如此，我们从天主确实的圣言得知亚当的罪；圣经真实地断言，在亚当内众人都死了，罪恶借着他进入了世界，死亡借罪恶而进入。我们的经验也充分证明，为惩罚这罪，我们的身体腐败，重压灵魂，土帐幕妨碍心智的多种活动；\BibleRef{智 9:15}而且我们知道只有靠恩宠的干预，才能从这惩罚中得救。所以宗徒痛苦地呼喊：\BibleQuote{我这个人真不幸！谁能救我脱离这死亡的身体呢？借着我们的主耶稣基督，天主的恩宠。}\BibleRef{罗 7:24} 我们所知如此；但神圣审判与仁慈分配的缘由——为何这人处于这种状况，那人处于那种——虽属公义，却属未知。但我们确信，这一切或出于天主的仁慈，或出于天主的审判；而造物主安排一切自然产物所用的度量、数目和重量，则隐而不现。因为天主不是罪的主动者，而是罪的掌管者；如此，那因违反本性而为罪的行为，受到审判和掌管，并被安置于其适当的位置和状态，以免给万有的本性带来不和谐和耻辱。事情既然如此，且天主的审判和人的意志运动含有隐藏的理由，为何同样的繁荣（有人善用，有人却因此毁灭）和同样的苦难（有人因此崩溃，有人反受益）——既然人在世上的整个有死生命是一场考验，\BibleRef{约 7:4}——谁能在任何特定情形下，在和平时期，是统治或服务，或安逸或死亡，或在战争时期，是命令或作战，或得胜或被杀，说是好是坏？同时，不变的是：凡是善的，皆因天主的祝福；凡是恶的，皆因天主的审判。\par

79. 因此，谁也不要胆敢鲁莽地控告人，更不用说控告天主了。如果旧约的仆役——他们也是新约的宣告者——的职分在于处死罪人，而新约的仆役——他们也是旧约的解释者——的职分在于被罪人处死，那么两种职分都是事奉同一位天主，祂按照时代不同变换教导方式，教导我们既要向祂祈求现世祝福，也要为祂放弃现世祝福；并教导我们现世的苦难既由祂降下，也应为祂忍受。因此，当\Person{梅瑟}出于对他子民的神圣热忱，希望他们只服从唯一真天主，得知他们堕落去敬拜亲手所造的偶像时，他借着用刀剑惩罚少数人，在当时用有益的恐惧打动他们的心，并给他们未来的警戒，这命令或行动中并没有残忍。这些人所受的公正惩罚，是祂们所犯罪得罪的天主，在祂隐秘审判的深处，指定立即执行的。\Person{梅瑟}如此行，不是出于残忍，而是出于极大的爱，这可以从他为百姓的罪祈祷的话中看出：\BibleQuote{如果祢赦免他们的罪，就赦免；若不然，求祢从祢所写的书上抹去我的名字。}\BibleRef{出 32:32}虔敬的探究者若将屠杀与祈祷相比较，就会在这其中找到最清晰的证据，表明灵魂向魔鬼的偶像卖淫所造成的伤害是多么可怕，以至于这样的爱竟能激起如此的愤怒。我们在\Person{保禄}宗徒身上也看到同样的情形，他不是出于残忍，而是出于爱，将一个人交给撒殚，使肉体败坏，好使灵魂在主耶稣的日子得救。\BibleRef{格前 5:5}他也将其他人交出去，叫他们不再亵渎。\BibleRef{弟前 1:20}在摩尼教人的伪经中，有一部寓言集，是由一些不知名的作者以宗徒的名义出版的。这些书如果在出版时，当时在世的圣贤和有学识的人——他们有能力判断此事——认为内容真实的话，毫无疑问会被教会接纳为正典。其中一个故事说，\Person{多默}宗徒有一次在一个他素不相识的地方参加婚宴，其中一个仆人打了他，他立即用诅咒给这人带来了可怕的惩罚。因为当他到水泉去为客人取水时，一头狮子扑到他身上咬死了他，而那只曾经轻微打了宗徒的手，应验了诅咒，被撕了下来，被一条狗叼到宗徒斜倚着的餐桌前。还有什么比这更残忍的呢？然而，如果我没有弄错的话，故事接着说，宗徒通过使这人在来世获得赦免的恩典来补偿了残忍；这样，这个异国的人民就因惧怕宗徒而敬畏他，认为他是天主所深爱的，而这人则以失去今世生命为代价，获得了永世的福乐。这故事是真是假无关紧要。无论如何，那些将教会正典所拒绝的书视为真实和正典的摩尼教人，必须承认，正如故事所示，主命令的忍耐德行——祂说：\BibleQuote{如果有人打你的右脸，将左脸也转给他}——可能存在于内心的态度中，尽管并未表现在身体行动或言语上。因为当宗徒被打时，他没有把另一边脸转过去，也没有叫那人再打，而是祈祷天主在来世赦免打他的人，但不要放过当时的伤害不受惩罚。他内里保持着善意，外面却希望那人受罚作为榜样。无论摩尼教人对此信得对或错，他们也可能相信，天主的仆人\Person{梅瑟}在砍倒制造和敬拜偶像的人时，也有这样的意图；因为他自己的话表明，他如此恳求赦免他们的拜偶像之罪，以至于请求如果他的祈祷不被垂听，就把他从天主的书上抹去。一个陌生人用手打人，与因离弃天主去崇拜偶像而侮辱天主的行为——尽管天主曾将子民从埃及的奴役中领出，带领他们过海，并用水淹没了追赶他们的敌人——是无法相比的。就惩罚而言，被刀杀与被野兽撕碎也是无法相比的。因为法官执行法律时，判给野兽残食的是比判给刀杀更坏的罪犯。\par

80. 福斯图斯的另一个恶意而亵渎的指控需要回答，即关于主对欧瑟亚先知说的话：\BibleQuote{你去娶淫妇为妻，也收淫妇所生的儿女。}\BibleRef{欧 1:2}关于这段经文，我们对手的不洁心灵如此盲目，竟不明白主在福音中对犹太人说的清楚的话：\BibleQuote{税吏和娼妓比你们先进天国。}\BibleRef{玛 21:31}娼妓离弃淫乱，成为贞洁的妻子，这并不违反真理的慈悲，也不违背基督信仰。事实上，对于一个自称先知的人而言，最不恰当的就是不相信堕落女子在改过自新后，她所有的罪都被赦免了。因此，当先知娶了娼妓为妻时，这既有利于这女子生命的改善，也象征了我们即将谈论的一个真理。但显然，摩尼教人在这件事上反感的是什么；因为他们最大的担忧是使娼妓怀孕。他们宁愿这女子继续做娼妓，以免他们的神被禁锢，也不愿她成为一人的妻子并生儿育女。\par

81. 关于撒罗满，只需说：圣经忠实的叙述对他行为的谴责，远比福斯图斯幼稚激烈的攻击更为严重。圣经忠实而准确地记载了撒罗满起初所有的善，以及他因恶行而丧失起初的善；而福斯图斯在攻击中，如同一个闭着眼睛或根本没有眼睛的人，不从光明中寻求引导，只被狂暴的敌意所驱使。对于虔诚而有辨识力的圣经读者，那些据说有多位妻子的圣人们的贞洁，其证据在于：撒罗满因多妻制而放纵情欲，而非寻求子嗣，被明确记载为贪爱妇女之人。正如不接受人面的真理所告知我们的，这将他引向了偶像崇拜的深渊。\par

82. 现在，我们已经审查了福斯图斯对旧约的所有指摘，并逐一考虑了其功过：或为天主的圣人们辩护，驳斥肉欲的异端者的诽谤；或当这些人有过错时，彰显圣经的卓越与威严。现在让我们再次按福斯图斯指控的顺序，考察这些记载的行动的意义，它们预示什么，又预言什么。这一点我们已在亚巴郎、依撒格和雅各伯身上做了；关于他们，天主说祂是他们的天主，仿佛宇宙的天主只是他们的天主；这不是用无意义的称号荣耀他们，而是因为那唯一能拥有完全而完美知识的那一位，知道这些人真诚而卓越的爱德；也因为这三位圣祖共同构成了未来天主之民的显著预象：不仅由自由妇人生育自由子孙，如撒辣、黎贝加、肋阿和辣黑耳所生，也生育奴婢的子孙，如这位黎贝加生了厄撒乌，对后者说：\Quote{你必服事你的兄弟}\BibleRef{创 27:40}；并且由奴婢妇人生育的不仅有奴婢的子孙，如哈加尔所生，也有自由的子孙，如彼耳哈和齐耳帕所生。同样，在天主之民中，那些精神自由的人不仅生出享受自由的子孙，如同那些被称为\Quote{你们该效法我，如同我效法基督一样}\BibleRef{格前 4:16}的人，他们也生出有罪奴役的子孙，就如西满由斐理伯所生。\BibleRef{宗 8:13}；同样，从肉身的奴婢所生的，不仅有模仿他们的有罪奴役的子孙，也有幸福自由的子孙，如同被称为\Quote{凡他们所吩咐你们的，你们要谨守遵行；但不要效法他们的行为}\BibleRef{玛 23:3}的人。凡正确观察这个预象在天主之民中应验的人，就在和平的联系中保持圣神的合一，恒心至终与一些人联合，并耐心容忍另一些人。关于罗特，我们也已说过，并指出圣经中提及他值得称赞之处和应受责备之处，以及整个叙述的意义。\par

83. 接下来要考虑犹大与儿媳同寝之行动的预言意义。但为使理解力薄弱的人受益，我们先指出：在圣经中，恶行有时并非预示恶，而是预示善。天主眷顾始终保持着其本质上的良善，以至正如上述例子，从通奸行为中生出一个男婴——这是藉自然能力从人的恶中产生的天主善工，而非因父母的过错；同样，在预言性的圣经中，善恶行为都被记载，而这些叙述本身是预言性的，甚至通过记载恶事来预告善事，其功劳归于作者而非作恶者。犹大为了满足邪恶情欲而与塔玛尔同寝时，他并不想通过其放荡行为来预象任何与人的救恩相关的事，正如出卖主的犹达斯也不想通过其背叛产生任何与人的救恩相关的效果。那么，如果主从犹达斯的恶行中藉自己的苦难成就了我们救赎的善工，为什么那位主亲自说过\Quote{他写的是关于我}的先知，为了教导我们，却不让犹大的恶行具有某种善的意义呢？在圣神的引导和默感下，先知编纂了这些行动的叙述，使之成为对他所预言之事的连续预言。在预言善事时，预象行为的善恶无关紧要。如果用红墨水写厄提约丕雅人是黑的，或用黑墨水写高卢人是白的，这并不影响书写所传达的信息。当然，如果是绘画而非书写，颜色错误就是缺陷；同样，当人类行为被呈现为榜样或警告时，其善恶具有重要性。但当行为作为预象被叙述或记载时，行为者的功过并不重要，只要行为与所象征的事物之间存在真实的预象关系。同样，在福音书中关于盖法的事：就他那不义而有害的意图，甚至就他使用那些话的意义而言，他主张一个义人应被不义地处死，这无疑是恶的；然而他的话中却有一种他所不知道的善的意义，他说：\Quote{一个人替百姓死，免得整个民族灭亡，这为你们有益。}关于他，经上记载：\Quote{这话不是由他自己说出的，只因他是大司祭，才预言了耶稣将为百姓而死。}\BibleRef{若 11:50}同样，犹大的行为就他的邪情私欲而言是恶的，但它预表了一种他所不知的伟大善事。他出于自己行了恶，但不是出于自己预表了善。这些前言不仅适用于犹大，也适用于所有其他情况，即恶行的叙述中包含了对善的预言。\par

84. 在塔玛尔——犹大的儿媳身上，我们看到了犹大王国的子民，其君王取自这一支派，相当于塔玛尔的丈夫。塔玛尔意为\Quote{苦涩}；这含义是恰当的，因为这子民给主递上了苦胆的杯。\BibleRef{玛 27:34} 犹大的两个儿子代表两类治理不善的君王：一类行恶，一类不为善。其中一个儿子在\Person{主}面前是邪恶或残忍的；另一个将种子撒在地上，使塔玛尔不能成为母亲。世上只有这两种无用的人：有害的和不肯施予所有之善却将其浪费或撒在地上的。既然行恶比不为善更恶劣，作恶者被称为长子，另一为次子。厄尔，长子的名字，意为\Quote{制皮者}，指我们的原祖父母因被逐出乐园受罚时所穿皮衣的制作者。\BibleRef{创 3:21} 敖难，次子的名字，意为\Quote{他们的悲哀}，即那些他不施善之人因他将所有之善浪费于地上的悲哀。长子之名所隐含的生命丧失，是比次子之名所隐含的缺乏帮助更为严重的恶。两者皆为\Person{天主}所杀，预表了王国从这类人手中被夺去。犹大第三子未与那女子结合，意味着犹大的君王曾有一段时间并非出自该支派。因此，这第三子并未成为塔玛尔的丈夫；塔玛尔代表犹大支派，它继续存在，尽管子民未从它获得君王。故此，这儿子名舍拉，意为\Quote{他的遣散}。这些预像都不适用于那些圣洁公义的人，如\Person{达味}，他们虽生活在那些时代，却属于新约，并以明悟的预言为之服务。再者，在犹大不再有本支派君王的时期，大黑落德不算作由塔玛尔的丈夫所预表的君王之一；因为他是个外邦人，他与子民的结合从未以圣油祝圣。他的权柄是外人的权柄，由罗马人和凯撒赐予。他的儿子们，即分封侯们，也是如此，其中一个与他父亲同名的黑落德，在\Person{主}受难时与比拉多成了朋友。\BibleRef{路 23:12} 这些外邦人是如此明显地区别于犹大的神圣王权，以至于犹太人自己在向基督发怒时公开喊道：\Quote{我们没有君王，只有凯撒。}\BibleRef{若 19:15} 凯撒并非严格意义上的他们的君王，除非在全世界都臣服于罗马的意义上。犹太人如此自我定罪，只是为了表达他们对基督的拒绝，并奉承凯撒。\par

85. 王国从犹大支派被夺去的时候，正是我们\Person{主}、\Person{耶稣基督}、真救主来临的预定时期，他来不是为了损害，而是为了伟大的益处。因此曾预言：\Quote{权杖必不离开犹大，御杖不离他两腿之间，直到那应得者来到；他是万民所期待的。}\BibleRef{创 49:10} 不仅犹太人的王国，连一切政权都已终止，并且如\Person{达尼尔}所预言，那神圣的傅油（基督或受傅者之名即由此而来）也已停止。然后，那应得者来到了，万民所期待的；至圣所被喜乐油傅抹，超越他的同伴。\Person{基督}在大黑落德时期诞生，在分封侯黑落德时期受难。他来到以色列家迷失的羊群中，这由\Person{犹大}前往提默纳剪羊毛所预表：提默纳意为\Quote{衰败}。因为那时权杖已离开犹大，犹太人的一切政权和傅油都已终止，以致那应得者得以到来。据记载，\Person{犹大}带着他的阿杜兰牧人同来，那牧人名依辣；阿杜兰人意为\Quote{水中的证据}。同样，\Person{主}也是带着这证据而来，他确实有比\Person{若翰}更大的证据，\BibleRef{若 5:36} 但为了他软弱的羊群，他使用了水中的证据。依辣这个名字也意为\Quote{我兄弟的异象}。因此，\Person{若翰}看见了兄弟，一位亚巴郎家族中的兄弟，并因\Person{玛利亚}和\Person{依撒伯尔}的亲戚关系；他也认出这同一位是他的主、他的天主，因为正如他自己所说，他领受了祂的圆满。\BibleRef{若 1:6} 因这异象，在妇人所生中，没有兴起比\Person{若翰}更大的；\BibleRef{玛 11:11} 因为一切预言基督的人中，唯有他看见了众多义人和先知渴望看见却没有看见的。他在母腹中就向\Person{基督}致敬；\BibleRef{路 1:44} 他因看见鸽子而更确切地认识祂；因此，作为阿杜兰人，他藉水作证。\Person{主}来剪祂的羊的毛，释放他们脱离痛苦的负担，正如\WorkTitle{雅歌}中赞美教会时所说，她的牙齿像剪过毛的一群母羊。\BibleRef{歌 4:2}\par

86． 接下来，我们有\\Person{塔玛尔}（\\OriginalTerm{塔玛尔}{Tamar}，意为\\Quote{改变}）改变装束；因为\\Person{塔玛尔}也意为\\Quote{改变}。不过，必须保留那苦涩的名称——不是那苦胆给主的苦涩，而是那使\\Person{伯多禄}痛哭的苦涩。\\BibleRef{玛 26:75}因为\\Person{犹大}（\\OriginalTerm{犹大}{Judah}，意为\\Quote{忏悔}）；苦涩与忏悔相混，作为真正悔改的预像。正是这悔改使建立在万民中的教会结出果实。因为\\BibleQuote{基督必须受难，并从死者中复活，并且因祂的名，从耶路撒冷开始，向万民宣讲悔改和罪过的赦免。}\\BibleRef{路 24:46}\\Person{塔玛尔}所穿的衣服中有罪的忏悔；而\\Person{塔玛尔}穿着这衣服坐在\\Place{厄南}（\\OriginalTerm{厄南}{Ænan}）或\\Place{厄纳因}（\\OriginalTerm{厄纳因}{Ænaim}，意为\\Quote{喷泉}）的门口，是蒙召自万民的教会的预像。她像牝鹿奔向水泉，去迎接亚巴郎的苗裔；在那里，她被一个不认识她的人所孕，正如所预言的：\\BibleQuote{我所不认识的人民必事奉我。}塔玛尔在她伪装下接受了一个戒指、一个手镯、一根棍杖；她在蒙召中被标记，在成义中被装饰，在受光荣中被高举。因为\\BibleQuote{祂所预定的人，祂也召叫了他们；祂所召叫的人，祂也使他们成义；祂所成义的人，祂也光荣了他们。}\\BibleRef{罗 8:30}这发生在她仍伪装的时候，正如我所说；在同一种状态下，她怀孕，并在圣德中结出果实。那承诺的山羊羔也被送到她那里，如同送给了妓女。山羊羔代表对罪的谴责，它是由已经提到的\\Person{阿杜兰人}（\\OriginalTerm{阿杜兰人}{Adullamite}）送来的，他仿佛使用了责备的话：\\BibleQuote{毒蛇的种类！}\\BibleRef{玛 3:7}但这种对罪的谴责没有临到她，因为她已被忏悔的苦涩所改变。随后，通过展示戒指、手镯和棍杖的抵押品，她胜过犹太人——他们对她的仓促判断，此时由\\Person{犹大}本人代表；正如今天我们听到犹太人說我们不是基督的子民，也没有亚巴郎的苗裔。但是，当我们展示了我们蒙召、成义和受光荣的确凿标记时，他们将立即感到羞愧，并承认我们比他们更正义。我本应更详细地探讨这一点，仿佛逐个肢解，如主所赐给我的能力，若不是因为需要结束这部作品（它已经比理想的长了）而阻止了这样细致的探究。\par

87． 关于\\Person{达味}（\\OriginalTerm{达味}{David}）罪的预言意义，只言片语即可。叙述中的名字显示了它所预表的内容。\\Person{达味}意为\\Quote{强手}或\\Quote{可羡慕的}；有什么比犹太支派的狮子更强壮呢？祂战胜了世界；或者有什么比先知所说的\\BibleQuote{万民的渴望将要来到}更可羡慕呢？\\BibleRef{盖 2:8}\\Person{巴特舍巴}（\\OriginalTerm{巴特舍巴}{Bersabee}）意为\\Quote{满足之井}或\\Quote{第七井}：这两种解释都符合我们的目的。因此，在\\WorkTitle{雅歌}中，那配偶——即教会——被称为活水之井；\\BibleRef{歌 4:15}或者，数字七代表圣神，正如在五旬节的天数，圣神从天而降。我们从\\WorkTitle{多俾亚传}也得知，五旬节是七周的庆节。\\BibleRef{多 2:1}四十九，即七乘七，加上一以表示合一。与此一致的是宗徒的话：\\BibleQuote{以爱德彼此担待，努力保持圣神的合一，借着和平的联系。}\\BibleRef{弗 4:2}教会因圣神的恩赐而成为满足之井，数字七表示其灵性；因为\\BibleQuote{她里面有活水的泉源涌到永生，凡拥有它的必永不再渴。}\\BibleRef{若 4:13}\\Person{巴特舍巴}的丈夫\\Person{乌黎雅}（\\OriginalTerm{乌黎雅}{Uriah}），从其名字的意义来看，必须被理解为代表魔鬼。所有被天主的恩宠所释放的人，先前都与魔鬼结合，使毫无玷污或皱纹的教会能与她真正的救主结婚。\\Person{乌黎雅}意为\\Quote{我的天主之光}；\\Place{赫特人}（\\OriginalTerm{赫特人}{Hittite}）意谓\\Quote{被剪除的}，这或者指他没有存留在真理中，因他的骄傲而被剪除于他从天主所得的天上之光，或者指他装作光明的天使，因为他跌倒后失去真实的力量，却仍敢说：\\Quote{我的光来自天主。}然后，按字面意义的\\Person{达味}确实犯了大罪，天主通过先知在指责\\Person{达味}时谴责了这罪，\\Person{达味}以他的悔改赎了罪。另一方面，那位万民的渴望，在教会于屋顶上沐浴时爱了她，即当她从世界的污秽中洗净自己，并在灵性默观中超越她的土屋并践踏它时；在开始相识后，祂处死了魔鬼，先将它从她身上完全除去，然后使她自己与祂永远结合。我们憎恨罪恶时，不可忽略预言的意义；我们按祂所应得的爱那以慈悲将我们从魔鬼手中释放的\\Person{达味}时，我们也可以爱那位以谦卑悔改医治其过犯所造成创伤的\\Person{达味}。\par

88. 关于\Person{撒罗满}，无需多言。\WorkTitle{圣经}提到他时，用了最强烈的责备和谴责，却未提及他的悔改和恢复天主的恩宠。在他可悲的堕落中，我也找不到任何与善有关的象征联系。也许他贪恋的那些外邦女子可以被视为代表从外邦人中拣选的教会。如果这些女子为了\Person{撒罗满}而离开她们的神，去敬拜他的天主，这个说法或许还说得通。但事实上，他为了她们而得罪了他的天主，并敬拜了她们的神，因此，要想出任何好的寓意似乎是不可能的。无疑，某些事情被象征了，但那是坏的象征，正如已经解释过的\Person{罗特}妻子和女儿的例子。我们在\Person{撒罗满}身上看到显著的高升和显著的堕落。如今，我们在不同的时期在他身上看到的这种善与恶，先是善，后是恶，在我们这个时代却一同存在于教会内。我认为，\Person{撒罗满}身上善的部分代表教会中善的成员，而他身上恶的部分则代表恶的成员。两者同在一人身上，正如在一个打谷场上，恶和善同在糠秕和麦粒中，或是在一块田里，莠子和麦子一同生长。对\WorkTitle{圣经}中关于\Person{撒罗满}的记载进行更深入的研究，也许会向我或更有学问、更有价值的人揭示出一些更合理的解释。但由于我们此刻正在讨论另一个主题，我们不可让此事打断我们论述的连贯性。\par

89. 关于先知欧瑟亚，我不必解释天主的命令和先知行为的含义，当时天主对他说：\BibleQuote{你去娶一个淫妇为妻，并生出淫乱的子女}，因为圣经本身就告诉我们这一指示的来源和目的。圣经接着这样说：\BibleQuote{因为此地大大行淫，离弃了主。他就去娶了狄布拉因的女儿哥默尔；她怀孕，给他生了一个儿子。主对他说：\InnerQuote{给他起名叫依次勒耳；因为再过不久，我必为依次勒耳的血讨伐犹大家，并终止以色列家的王国。到那日，我必在依次勒耳谷折断以色列的弓。}她又怀孕，生了一个女儿。天主对他说：\InnerQuote{给她起名叫罗鲁哈玛；因为我必不再怜悯以色列家，而要将他们彻底除去。但我要怜悯犹大家，并借他们的天主上主拯救他们，不靠弓、刀、战事、马匹或骑兵。}罗鲁哈玛断奶后，她又怀孕，生了一个儿子。天主说：\InnerQuote{给他起名叫罗阿米；因为你们不是我的人民，我也不作你们的天主。}然而以色列子民的人数必如海沙，不可测量；将来在人们对他们说\InnerQuote{你们不是我的人民}的地方，必有人对他们说\InnerQuote{你们是永生天主的儿子}。那时，以色列子民和犹大子民将聚集在一起，为自己立一个首领，并从那地上来，因为依次勒耳的日子是伟大的。你们要对你们的兄弟说\InnerQuote{我的人民}，对你们的姊妹说\InnerQuote{她获得了怜悯}。}既然在同一本书中，主亲自解释了命令和先知行为的预表意义，并且宗徒的著作宣布了这一预言在新约宣讲中的应验，那么每个人都必须接受对命令和先知行为的这种解释作为真确的解释。因此，保禄宗徒说：\BibleQuote{如此，祂要彰显祂荣耀的丰富在那些蒙怜悯的器皿上，就是祂事先预备好为得荣耀的器皿，就是我们这些蒙召的人，不仅从犹太人中，也从外邦人中。正如祂在欧瑟亚中也说过：\InnerQuote{我要称那不是我人民的人为我的人民，那不蒙爱的为蒙爱的；将来在人们对他们说\InnerQuote{你们不是我的人民}的地方，他们必被称为永生天主的儿子。}}\BibleRef{罗 9:23}保禄在此将预言应用于外邦人。同样，伯多禄写信给外邦人，没有提先知的名字，却借用了他的话语，他说：\BibleQuote{但你们是特选的种族、王家的司祭、圣洁的国民、属于主的子民，为叫你们宣扬那召你们出离黑暗进入祂奇妙光明者的美德；你们从前不是子民，现在却是天主的子民；从前未蒙怜悯，现在却蒙了怜悯。}\BibleRef{伯前 2:9}由此清楚可见，先知的话\BibleQuote{以色列子民的人数必如海沙，不可测量}，以及紧接着的话\BibleQuote{将来在人们对他们说\InnerQuote{你们不是我的人民}的地方，他们必被称为永生天主的儿子}，并非指那属血肉的以色列，而是指那宗徒对外邦人所说的\BibleQuote{你们因此是亚巴郎的子孙，按照恩许是承继人}。\BibleRef{迦 3:29}但是，许多属血肉的以色列中的犹太人已经相信，并且将来还要相信；因为宗徒们就是其中之一，还有耶路撒冷中所有成千上万与宗徒们在一起的人，以及保禄所提到的众教会，当他给迦拉达人写信时说：\BibleQuote{我本人与犹太地区在基督内的教会并不相识。}\BibleRef{迦 1:22}接着，他又解释圣咏中称主为角石的段落，是指祂在自己内将受割礼和未受割礼的两堵墙联合起来，\BibleQuote{为在自己内将两者造成一个新人，如此缔造和平；并借十字架消灭了仇恨，使两者在一个身体内与天主和好；并且祂来向远处的人和近处的人传报了和平}，即向外邦人和犹太人；\BibleQuote{因为祂是我们的和平，祂使两者合为一体}。\BibleRef{弗 2:11}同样，我们发现先知论及犹太人时称他们为犹大子民，论及外邦人时称他们为以色列子民，他说：\BibleQuote{犹大子民和以色列子民将聚集在一起，为自己立一个首领，并从那地上来。}因此，反对一个如此被历史事实所证实的预言，就是反对宗徒和先知的著作；而且不仅仅是反对著作，更是以最鲁莽的方式攻击明如白昼的既定事实的证据。至于犹大叙事，或许不容易从那个名叫塔玛尔的妇人伪装下的妓女身上认出那从外邦迷信的腐败中聚集起来的教会的形象；但在这里，圣经自我解释，且这解释又为宗徒著作所证实，我们不再赘述，而应立即探讨福斯托在天主仆人梅瑟身上所反对的那些事情本身的含义。\par

90. 梅瑟为保护他的一位弟兄而杀死埃及人，自然令我们想起我们的辩护者主基督在这异乡之地杀死攻击我们的魔鬼。正如梅瑟将尸体藏在沙中，同样，魔鬼虽被杀死，却仍隐藏在那些没有坚定扎根的人中。我们知道，主将教会建立于磐石上；那些听了祂的话而实行的人，祂比作一个聪明人，把房子盖在磐石上，在试探面前不屈服不后退；而那些听了却不实行的人，祂比作一个愚昧人，把房子盖在沙上，当房子受考验时，倒塌得很大。\BibleRef{玛 7:24}\par

91. 关于梅瑟奉他的上主天主（祂的命令无不是至公至义）的命令夺取埃及人财物的预言意义，我记得我在题为\WorkTitle{论基督圣道}的书中曾记下当时所想到的：埃及人的金银和衣裳预表了某些学问分支，这些学问可以在外邦人中有益地学习或教导。这可能是真实的解释；或者，我们可以认为金银器皿代表宝贵的灵魂，衣裳代表那些从外邦人中归附天主之民的身体，使他们一同从这世界的埃及得解脱。无论真实的解释如何，圣经的虔敬学者都会确信，在这命令、这行动和这叙述中，有一个目的和一个象征意义。\par

92. 要逐一详述梅瑟的一切战争，未免太长。提起已经说过的话，就足以回答福斯特，说明与阿玛肋克之战的预言和象征性质。还有，圣经的敌人或从未读过任何东西的人对梅瑟发出残忍的指控。福斯特虽未具体指控，但说梅瑟命令并做了许多残忍的事。然而，我了解他们惯常提出并歪曲的那些事，已选取一个具体案例加以辩护，让愿意改正的摩尼教徒，以及所有其他无知而不敬神的人，明白他们的指控毫无根据。现在我们必须探究那命令的预言意义：在梅瑟不在时，许多为自己造偶像的人被杀死，不顾及亲属关系。显然，这些人的杀戮象征了对邪恶原则的争战，这些原则曾引导人民陷入同样的偶像崇拜。我们奉命与这样的邪恶争战，正如\BibleBook{}{圣咏}的话：\BibleQuote{你们应发怒，但不可犯罪。}宗徒也给了类似的命令，他说：\BibleQuote{你们要致死属于地上的肢体，即淫乱、不洁、邪情、恶欲和贪婪，贪婪就是偶像崇拜。}\BibleRef{哥 3:5}\par

93. 要解读梅瑟的第一个行动——用火焚烧牛犊、磨碎成粉、撒在水里让百姓喝——的意义，需要更仔细的考察。他所领受的法版，是天主用手指写的，即圣神的作为；他打碎了法版，或许是因为他判断百姓不配聆听；他焚烧牛犊、磨碎、撒在水里冲走，是为了不让任何残渣留在百姓中间。但他为何要让他们喝下去呢？每个人都必定渴望发现这一行动的典型意义。继续探究，我们可能会发现，牛犊代表了魔鬼的化身，正如万国中那些以魔鬼为首领、行不敬礼的人一样。牛犊是金的，因为偶像崇拜的制度中有一种智慧的假象。对此，宗徒说：\BibleQuote{他们虽然认识天主，却没有以祂为天主而光荣祂，也不感谢祂，反而在思想中成了虚妄，他们愚昧的心就昏暗了。他们自称为智者，反而成了愚昧，将不可朽坏的天主的光荣，变成了可朽坏的人、飞禽、走兽和爬虫的形像。}\BibleRef{罗 1:21} 从这所谓的智慧中出来了金牛犊，它是埃及首领与自称智者的主要偶像崇拜形式之一。因此，牛犊代表外邦偶像崇拜者的全体或团体。主基督用那火焚烧这个不敬神的团体，祂在福音中说：\BibleQuote{我来是把火投在地上。}\BibleRef{路 12:49} 因为，没有什么能躲过祂的热力，当外邦人信祂时，他们就在神圣影响的火中失去了魔鬼的形像。然后整个身体被磨碎，也就是说，在不义的肢体联合解体后，在真理之言下谦卑。然后粉末撒在水里，使以色列人，即福音的宣讲者，在圣洗中接纳那些过去崇拜偶像的人进入他们自己的身体，即基督的身体。对那位以色列人之一的伯多禄，曾论及外邦人说：\BibleQuote{宰了吃。}\BibleRef{宗 10:13} 宰了吃与磨碎喝极为相似。因此，这牛犊借着热忱之火、言语的锐利穿透和圣洗之水，被百姓吞没，而不是百姓被它吞没。\par

94. 因此，当异端者用来反对圣经的那些经文被研究和检验时，它们越是隐晦，我们就越能发现奇妙的奥秘来回答我们的疑问；以至于亵渎者的口被完全堵住，真理的证据使他们窒息，甚至发不出声音。那些不愿在心中接受真理之甘饴的不幸者，必感到真理如同一根塞口物堵在他们的口中。所有那些经文都讲论基督。那现已升天的头，连同仍在世上受苦的身体，是圣经作者全部目的的圆满实现，这圣经被称为\Quote{圣经}。先知书叙述的每一部分都应视为具有比喻意义，除非那些只是作为框架，用于字面或比喻预言这位君王及其子民。因为如同竖琴和其他乐器，乐音并非来自整个乐器，而是来自琴弦；其余部分只是用来固定和绷紧琴弦以便调音，当乐师拨动时发出悦耳的声音；同样，在先知叙述中，由先知之灵选出的细节或预言未来的事件，或本身虽无声音，却连接其他有意义的言论。\par

95. 如果异端者拒绝我们关于那些寓意叙事的解释，甚至坚持只按字面意义理解它们，那么争论这种理解上的差异就如同争论口味上的差异一样徒劳。不过，天主诫命要么具有道德和宗教的性质，要么具有预言的涵义，这一事实必须被相信，无论是否明智。此外，形象化的解释都必须服务于道德和宗教的利益。因此，如果摩尼教徒或任何其他人不同意我们的解释，或在方法上或任何具体意见上与我们有分歧，只要那些天主因其行为和对其诫命的服从而称赞的祖先的品格，在一个除了那些执迷不悟的敌视者之外所有人都承认真的原则上得到辩护；而且圣经的纯洁与尊严，在那些真理的敌人所挑剔的段落中得以维护，其中某些行为要么被赞扬要么被谴责，要么仅仅被叙述以供我们判断它们，这就足够了。\par

96. 事实上，没有什么比这更能教导并有益于圣经的虔诚读者了：除了将值得称赞的品格作为榜样、将应受责备的品格作为警告加以描述外，它还应当叙述善人后退而陷入邪恶的事例，无论他们是重归正途还是继续执迷不悟；也要叙述恶人改变而获得良善的事例，无论他们是坚持良善还是再次堕入邪恶；好让义人不至于在安全的骄傲中自高自大，也不让恶人在无药可救的绝望中刚硬。就连圣经中那些不包含榜样或警告的段落，要么是为连接起见，以便过渡到关键问题，要么因其表面上的多余而暗示着某种隐秘的象征意义。因为在我们所谈论的书籍中，预言性宣告非但不缺乏或稀少，反而是众多而清晰的；如今既然预言实际上已经应验，这些书卷的神圣权威由此获得的见证就强而有力，以至于要认为那些各阶层、各心智的人都表示敬意的书卷中可能存在任何无用或无意义的段落，这简直是疯狂，而这些书卷本身正预言了我们所看到实际发生的事。\par

97. 那么，如果有人读到\Person{达味}的作为——即当主谴责并威胁他时，他为此而悔改——却在叙事中找到鼓励去犯罪，难道圣经该为此受责备吗？难道不是此人自己的罪过，与他为了伤害或毁灭自己而滥用那些为他的康复和释放而写的东西的程度成正比吗？\Person{达味}被树立为悔改的伟大榜样，因为陷入罪中的人要么骄傲地忽视悔改的救治，要么在获得救恩或得到宽恕的绝望中迷失自己。这榜样是为病弱者的益处，而非为健康者的伤害。如果疯子用外科器械自毁，或恶人用它们毁灭他人，那不是外科手术的过错。\par

98. 即使假定我们的祖先——圣祖和先知们，关于他们的虔诚和宗教生活习惯，在\WorkTitle{圣经}中给予了如此良好的记载（每一位认识\WorkTitle{圣经}、且尚未完全丧失理智的人，都必须承认\WorkTitle{圣经}是由天主为人的救恩而赐下的）——正如摩尼教徒所虚假和狂热地指控的那样，是淫荡和残忍的，他们仍然可以被证明不仅优于摩尼教徒所谓的\Quote{选民}，而且也优于他们自己的神。难道男人与女人之间的放荡结合，比无瑕之光与黑暗混合而自贬身价更为恶劣吗？在此，一个人因贪婪和贪欲而将自己的妻子冒充为妹妹，并把她卖给她的情人；但更恶劣、更骇人听闻的是，有人伪装自己的本性以满足罪恶的激情，并无缘无故地甘受玷污和堕落。甚至那些明知与自己女儿同寝的人，也不如那些让自己的肢体参与所有如此粗鄙或更粗鄙的感官玷污的人罪大恶极。而那位摩尼教的神，不也参与了最令人发指的不洁行为的污染吗？再者，如果真如\Person{福斯图斯}所说，\Person{雅各伯}在四个妻子之间来往，不是为了渴望子孙，而是像一只公山羊般放荡，那他仍然没有堕落到你的神的地步：你的神不仅因被囚禁在\Person{雅各伯}及其妻子的身体中，而与他们的所有动作混合，而且，与\Person{福斯图斯}粗俗比喻中的这只公山羊结合，必须承受动物情欲的所有痛苦，每一步都招致新的玷污，因为它参与了雄性的激情、雌性的受孕和幼崽的出生。同样，假定\Person{犹大}不只犯了奸淫，还犯了乱伦，那么你的神也要分担这乱伦情欲的热度和不洁。\Person{达味}悔改了他的罪——爱上别人的妻子并下令杀死她的丈夫；但你的神何时才会悔改，因为他将自己的肢体交给了黑暗种族男女首领的放纵情欲，并且杀死的不是他情妇的丈夫，而是他自己的儿女，并把他们囚禁在正作为他的情人的那些恶魔的肢体中？即使\Person{达味}没有悔改，也没有这样恢复义德，他仍然比你的神好。\Person{达味}可能因这一次行为或一个人所能沾染的这种玷污程度而受玷污；但你的神却在无论何人犯下的一切此类行为中，蒙受其肢体的污染。先知\Person{欧瑟亚}也被\Person{福斯图斯}指控：假定他娶了那妓女是因为对她有邪恶的爱，如果他淫荡而她是个娼妓，按照你们自己的断言，他们的灵魂是你神的本质的一部分和肢体。明白地说，那娼妓本身必定就是你们的神。你们不能假装你的神没有被囚禁在污染的身体中，也不能说他只是临在，同时保持他自己本性的纯洁；你们承认你的神的肢体被如此玷污，以致需要特别的净化。那么，这个你对这位天主的人吹毛求疵的娼妓，即使她没有因为成为贞洁的妻子而变好，她仍然曾是你们的神；至少你们必须承认她的灵魂是你们神的一部分，即使很小。但单个娼妓并不像你们的神那样坏，因为他由于与黑暗种族的混合，参与了每一次卖淫行为；而无论何处发生这种不洁，他都经历相应的遗弃、释放和囚禁体验，而且代代相传，直到这最败坏的部分达到黑暗团块中的最终状态，像一个不可救药的娼妓。你的神无法从他的肢体上避开这些邪恶和这些可耻的憎恶，他被自己无情的敌人不可抗拒地引到这地步；因为只有牺牲他自己的臣民，或者更确切地说，他自己的部分，他才能摧毁那可怕的攻击者。确实，杀死一个埃及人以保全同胞，并没有像这样邪恶。然而\Person{福斯图斯}最荒谬地对此找茬，而他惊人地糊涂，忽略了他自己神的处境。他带走埃及人的金银器皿，难道不是比让他的肢体被黑暗种族带走更好吗？然而，这位不幸之神的崇拜者却挑剔我们天主的仆人所进行的战争，在这些战争中他和他的跟随者总是得胜，以至于在\Person{梅瑟}的领导下，以色列子民俘虏了他们的敌人，无论男女，正如你们的神如果能做到也会做的那样。你们声称控告\Person{梅瑟}做错了，而实际上你们嫉妒他的成功。用刀剑惩罚那些严重得罪了天主的人，这并非残忍。事实上，\Person{梅瑟}为这罪祈求宽恕，甚至自己愿意承担天主的愤怒来代替他们。但即使他是残忍而不是慈悲，他仍然比你们的神好。因为如果他的跟随者中有人被派去冲破敌人的阵线而被俘虏，他在得胜时绝不会在对方没有做错、只是服从命令的情况下判他死刑。然而这正是你们的神要对那被固定在黑暗团块中的自身部分所做的事，因为它服从命令，冒着生命危险为保卫他的王国而前进对抗敌人的身体。但摩尼教徒说，这个部分在漫长的世纪中与邪恶混合和结合后，并没有服从。但为什么？如果服从是自愿的，那么罪责是真实的，惩罚是公正的。但这会导致没有本性是与罪对立的；否则它就不会自愿犯罪；于是整个摩尼教体系立刻倒塌。相反，如果这个部分由于它所对抗的敌人的力量而受苦，并因无法抵抗的力量而被制服，那么惩罚是不公正的，并且是明目张胆的残忍。以必要性为由被辩护的神，正适合于那些拒绝崇拜唯一真天主的人来崇拜。然而，必须承认，无论崇拜这个神是何等卑下，崇拜者仍比他们的神好，因为他们存在，而他不过是虚构的产物。现在让我们继续论述\Person{福斯图斯}的其余荒诞之词。\par

\SourceRecord{Translated by Richard Stothert.；Nicene and Post-Nicene Fathers, First Series；Vol. 4.；Edited by Philip Schaff.；Buffalo, NY: Christian Literature Publishing Co.,；1887.；<http://www.newadvent.org/fathers/140622.htm>.}

% structure: p0001 source=h1 target=section reason=page-title

\section{反对福斯图斯，卷二十三}

\Emph{\Person{福斯图斯}重新回到谱系难题，并坚持认为即使按照\Person{玛窦}福音，耶稣直到受洗时才是天主子。\Person{奥古斯丁}阐述了关于基督位格中神性与人性关系的公教观点。}\par

1. \Person{福斯图斯}说：有一次，当向一大群人讲话时，人群中有人问我：\Quote{你相信耶稣是由玛利亚所生吗？}我回答说：\Quote{你指的是哪个耶稣？因为在希伯来语中，这是好几个人的名字。一个是农的儿子，\Person{梅瑟}的追随者；\BibleRef{出 23:11}另一个是约匝达克大司祭的儿子；\BibleRef{盖 1:1}还有一个被称为达味的儿子；\BibleRef{罗 1:1}另有一个是天主子。\BibleRef{谷 1:1}你问的是哪一个，问我是否相信他是由玛利亚所生？}他回答说：\Quote{当然是天主子。}我说：\Quote{我凭什么证据，口传的还是书面的，来相信这事？}他回答说：\Quote{依据玛窦的权威。}我说：\Quote{玛窦写了什么？}他回答说：\Quote{\BibleQuote{耶稣基督的族谱，达味之子，亚巴郎之子}}。\BibleRef{玛 1:1}然后我说：\Quote{我恐怕你会说：耶稣基督的族谱，天主子；我已准备好纠正你。既然你准确地引用了经文，你仍必须注意这些词语。玛窦并没有自称要记述天主子，而是记述达味之子的族谱。}\par

2. 现在我暂且假设此人说得对，即达味之子是由玛利亚所生。但仍有事实：在这整段谱系中，直到受洗时才提到天主子；因此，你们说这位作者使天主子成为胎中的寄居者，是一种有害的歪曲。这位作者确实似乎在反对这种想法，并在其书的标题中就表明自己免于这样的亵渎，声称他所描述出生的人是达味之子，而非天主子。如果你们注意作者的意思和目的，你们会看到，他希望我们相信的关于天主子耶稣的，与其说是祂由玛利亚所生，不如说是祂在约旦河受洗而成为天主子。他告诉我们，他一开始称为达味之子的人，由若翰施洗，并在此时成为天主子，按路加记载，大约三十岁，当时有声音向祂说：\BibleQuote{祢是我的儿子，我今天生了你。}\BibleRef{路 3:22}由此看来，那据推测三十年前由玛利亚所生的，并不是天主子，而是后来在约旦受洗而成为的，即新人，如同我们在由外邦人的谬误归化并相信天主时一样。这个教义可能与你们所谓的公教信仰相符，也可能不符；无论如何，这是玛窦所说的，如果玛窦是真正的作者。这些话语：\BibleQuote{祢是我的儿子，我今天生了你}，或\BibleQuote{这是我的爱子，我所喜悦的}，并未出现在与玛利亚母性相关的叙述中，而是出现在约旦洗除罪恶之处。这就是所写的；如果你们相信这个教义，你们应被称为玛窦派信徒，因为你们不再是公教徒。公教教义是众所周知的，它与玛窦的表述一样不像真理。按照你们信经的话，你们宣称相信耶稣基督，天主子，由童贞玛利亚所生。因此，按照你们，天主子出自玛利亚；按照玛窦，出自约旦；而我们相信祂出自天主。所以，玛窦的教义，如果我们正确地将作者归于他，既不同于你们的，也不同于我们的；只是我们承认他比你们更谨慎，将出于妇人的出生归于达味之子，而非天主子。至于你们，你们唯一的选择要么是否认这些陈述是由玛窦作出的，要么承认你们已放弃了宗徒们的信仰。\par

3. 至于我们，虽然无人能改变我们的信念，即天主子出于天主，但我们可能会放纵轻信的性格，以至于承认这个虚构：耶稣在约旦成为天主子，但天主子并非由妇人所生。此外，那据称由玛利亚所生的儿子，不能适当称为达味之子，除非确证他是由若瑟所生。你们说他不是，因此你们必须承认他不是达味之子，即使他是玛利亚的儿子。族谱从亚巴郎到达味，再从达味到若瑟，沿希伯来父系传承；既然我们被告知若瑟并非耶稣的真正父亲，耶稣就不能说是达味之子。一开始称耶稣为达味之子，然后又继续讲祂在若瑟婚姻未完成前由玛利亚所生，这完全是疯狂。如果玛利亚的儿子因不是若瑟的儿子而不能称为达味之子，那么天主子这个称号就更不能给予祂了。\par

4. 此外，童贞女本人似乎不属于犹太国王所属的犹大支派——所有人都同意那是达味的支派——而是属于司祭的肋未支派。这从童贞女的父亲约阿金是司祭可知；他的名字并未出现在族谱中。那么，玛利亚如何能与达味拉上关系呢？她的父亲和丈夫都不属于达味支派。因此，玛利亚的儿子不可能成为达味之子，除非你能使母亲与若瑟产生某种联系，要么是他的妻子，要么是他的女儿。\par

5. \Person{奥斯定}答说：公教（亦即宗徒的）道理是：我们的主和救主耶稣基督，按天主性而言是天主子，按血肉而言是达味之子。我们由福音作者和宗徒的著作证明这一点，以致没有人能拒绝我们的证明而不也拒绝这些著作。\Person{福斯图斯}的伎俩是让某人说几句话，却不提出任何证据来反驳福斯图斯的诡辩多端。但尽管他诡计多端，我所提供的证据将使福斯图斯无言以对，只能声称这些段落是神圣记载中的伪造插入——这种答复是用来逃避或试图逃避圣经最明确陈述的力量。我们已在本书中充分揭露了这种批评的荒谬无理和大胆亵渎；为免逾越界限，我们必须避免重复。我们无需汇集散见于圣经的所有段落，这些段落为答复福斯图斯而表明，在最高最神圣权威的书卷中，那被称为天主独生子、甚至与天主同等的天主者，也被称为达味之子，因祂从童贞玛利亚（\Person{若瑟}的妻子）取了仆人的形状。仅以\Person{玛窦}为例，既然福斯图斯的论证涉及这部福音，而此处无法引用整部书，就让愿意者自行阅读，看看\Person{玛窦}如何将那在族谱引言中称为达味之子的叙述延伸至受难与复活。他论及这位达味之子由圣神受孕并由童贞玛利亚诞生。他也引用先知的话：\BibleQuote{看，一位贞女将怀孕生子，人将称祂的名字为\OriginalTerm{厄玛奴耳}{Emmanuel}，即译作\Quote{天主与我们同在}。} 再者，那位甚至在童贞母胎中就被称为\Quote{天主与我们同在}者，据闻当祂受\Person{若翰}洗礼时，听见从天上来的声音说：\BibleQuote{这是我的爱子，我所喜悦的。}\BibleRef{玛 3:17} 福斯图斯难道会说被称为天主比被称为天主子更卑微吗？他似乎这样认为，因为他试图证明，因这声音在洗礼时来自天上，所以按照\Person{玛窦}，祂那时才成为天主子；然而同一位福音作者在前文中引用了先知的神圣宣告，其中称那位由贞女所生的婴孩为\Quote{天主与我们同在}。\par

6. 值得注意的是，这位可怜的好事者在他离题万里的胡言中，不放过任何机会用自己幻想的神话故事来模糊圣经的宣告。例如，他论及\Person{亚巴郎}，说他娶婢女为妻时，不相信天主应许他由\Person{撒辣}得子；而事实上，那时这应许尚未赐下。然后他指责亚巴郎称撒辣为妹是谎言，却不曾阅读圣经权威所记载的关于撒辣家族的事。他也指责亚巴郎之子\Person{依撒格}假称妻子为妹，尽管关于她的家族有明确记述。然后他指责\Person{雅各伯}的四位妻子每天争吵，看谁在他从田间回来时首先占有他，而圣经中对此毫无记载。此人正是假装因虚假而憎恨圣经作者的人，他竟敢如此歪曲甚至福音的记载——尽管福音的权威因具有最充分的证实而被所有人接受——试图使之显得，并非\Person{玛窦}本人（因为那样他就不得不屈服于宗徒的权威），而是某个托名玛窦的人，写了关于基督的事，而这些事他不愿相信，并且企图以辱骂式的机敏加以反驳！\par

7. 在\Place{约旦河}从天上来的声音应与在山上听到的声音相比较。\BibleRef{玛 17:5} 两者中，\BibleQuote{这是我的爱子，我所喜悦的}这句话都不意味祂先前不是天主子；因为祂从童贞母胎中取了\OriginalTerm{仆人的形体}{form of a servant}，却\BibleQuote{具有\OriginalTerm{天主的形体}{form of God}，不以自己与天主同等为僭夺。}\BibleRef{斐 2:6} 同一位宗徒\Person{保禄}在其他地方明确说：\BibleQuote{但是时期一满，天主就派遣了自己的儿子，生于女人，生于法律之下；}\BibleRef{迦 4:4} 这里的女人是按希伯来文的意义，不是妻子，而是女性。天主子按其天主性是达味的主，按血肉出自达味的后裔，也是达味之子。如果我们相信这事不是有益的，同一位宗徒就不会如此强调，他对\Person{弟茂德}说：\BibleQuote{你务要记得：耶稣基督，达味的后裔，从死者中复活了，这合乎我的福音。}\BibleRef{弟后 2:8} 他仔细嘱咐信徒，凡宣讲别的福音与这不同的，应视为可诅咒的。\par

8. 这位攻击神圣福音的人不必因以下事实而感到困惑：基督被称为\Person{达味}之子，尽管祂由童贞女所生，且\Person{若瑟}并非祂的真正父亲；而福音史家\Person{玛窦}所记载的族谱并非追溯至\Person{玛利亚}，而是至\Person{若瑟}。首先，作为丈夫，男子更为尊贵；\Person{若瑟}是\Person{玛利亚}的丈夫，尽管她并未与他同居，因为\Person{玛窦}本人提到，天使称\Person{玛利亚}为\Person{若瑟}的妻子；也是从\Person{玛窦}那里，我们得知\Person{玛利亚}怀孕并非由于\Person{若瑟}，而是由于圣神。但倘若这并非宗徒\Person{玛窦}所写的真实叙述，而是某位冒名者所写的虚假叙述，那么他会在如此紧密相连的段落中如此明显地自相矛盾吗？他先说\Person{达味}之子生于\Person{玛利亚}而无夫妻之实，然后在给出族谱时，却将其追溯至那位与童贞女明确无交合的男士，除非他这样做有其理由？即使假设有两位作者，一位称基督为\Person{达味}之子，并记述了从\Person{达味}至\Person{若瑟}的基督先祖；另一位则不称基督为\Person{达味}之子，并说祂由童贞\Person{玛利亚}所生，未与任何男子交合；这些陈述并非不可调和，以至证明其中一位或两位作者必为虚假。经思考可见，两种叙述可能均为真实；因为\Person{若瑟}可被称为\Person{玛利亚}的丈夫，尽管她只是情感上和精神上的妻子，这种精神的结合比肉体的结合更为亲密。这样，基督童贞之母的丈夫理应在基督祖先的名单中占有一席之地。也可能\Person{玛利亚}本人身上流有\Person{达味}的血统，因此基督的肉身虽出于童贞女，仍源于\Person{达味}的种子。但事实上，这两项陈述均由同一位作者作出，他既告诉我们\Person{若瑟}是\Person{玛利亚}的丈夫，基督之母是童贞女，基督是\Person{达味}种子所出，\Person{若瑟}在\Person{达味}世系的基督祖先名单中；因此，那些偏爱神圣福音权威胜过异端虚构的人必须得出结论：\Person{玛利亚}与\Person{达味}家族并非无关，她之所以被称为\Person{若瑟}的妻子是合适的，因为作为一名女性，她与他有精神上的结合，尽管没有身体上的联系。显然，\Person{若瑟}也不能在族谱中被省略；因为从性别优势来看，这样的省略相当于否认他与那位内在地与他结合的女性的关系；并且基督信徒受教导，不要将肉体的结合视为婚姻中的首要之事，仿佛没有它就不能成为夫妻，而应在基督徒的婚姻中尽可能效法基督的父母，以便与基督的肢体有更亲密的结合。\par

9. 我们相信\Person{玛利亚}以及\Person{若瑟}都出自\Person{达味}家族，因为我们相信圣经，圣经断言基督按肉身是\Person{达味}的种子，并且祂的母亲是童贞\Person{玛利亚}，祂没有人类的父亲。因此，凡否认\Person{玛利亚}与\Person{达味}有血缘关系的人，显然与这些圣经段落的高度权威相抵触；为维护这种反对，他必须从他承认的、由教会接纳为正经和公教的著作中提出证据，而不是随意引用任何著作。在我们现在处理的事情上，只有正经著作对我们有分量；因为只有它们才被普世教会所接受和承认，而教会本身就是这些著作中关于它的预言的应验。因此，我没有义务接纳\Person{福斯图斯}所采用的关于\Person{玛利亚}出生之非正经记载，即她的父亲是名叫\Person{约阿基姆}的肋未支派司祭。但即使我接纳这一记载，我仍会主张\Person{约阿基姆}必以某种方式属于\Person{达味}家族，并曾以某种方式从犹大支派被过继到肋未支派；或者若不是他，则是他的某位祖先；或至少，他虽然生于肋未支派，但仍与\Person{达味}世系有关；正如\Person{福斯图斯}本人承认，\Person{玛利亚}虽属肋未支派，却可许配给犹大支派的丈夫；并且他明确说，若\Person{玛利亚}是\Person{若瑟}的女儿，则基督被称为\Person{达味}之子是适用的。这样，通过\Person{若瑟}的女儿在肋未支派中结婚，她的儿子虽生于肋未支派，仍可恰当地被称为\Person{达味}之子。同样，若那位在\Person{福斯图斯}所引段落中被称为\Person{玛利亚}之父的\Person{约阿基姆}的母亲，在属于犹大支派和\Person{达味}家族时，嫁入了肋未支派，那么就有足够的理由称\Person{约阿基姆}、\Person{玛利亚}和\Person{玛利亚}的儿子属于\Person{达味}的种子。若我有义务重视那部称\Person{约阿基姆}为\Person{玛利亚}之父的伪经，我会采用上述解释，而不会承认福音中有任何虚假，福音中记载了：耶稣基督，天主子，我们的救主，按肉身是\Person{达味}的种子，并且祂由童贞\Person{玛利亚}所生。对我们而言，这些记载这些真理并为我们所相信的圣经的敌人，不能证明它们有任何虚假的指控，这就足够了。\par

10. \Person{福斯图斯}不能假装我无法证明\Person{玛利亚}出于\Person{达味}家族，正如我已证明他无法证明她不是。我从具有权威的圣经中提出最强有力的证据，这些证据宣告基督是\Person{达味}的后裔，并且祂由\Person{童贞玛利亚}无父而生。\Person{福斯图斯}表达了他认为对不当行为最为合适的愤慨，他说，\Quote{使作者谈论天主圣子住在子宫里，是一种有害的曲解。}当然，天主教义教导基督、天主子，由\Person{童贞玛利亚}降生成人，这并不意味着天主圣子只限于子宫内而没有超越的存在，好像祂放弃了治理天地，或好像祂离开了圣父的临在。错误在于摩尼教人，他们的理解力如此局限于物质观念，以至于无法领会天主的圣言，祂是天主的德能和智慧，既安住于自身并与圣父同在，又治理宇宙，以德能达至地极，并温和地安排一切（见：\BibleRef{智 8:1}）。在这令人钦崇的天主眷顾的无瑕过程中，祂为自己指定了一位地上的母亲；为了将祂的仆人从腐朽的束缚中释放出来，祂在这位母亲内取了仆人的形体，即一个可死的身体；祂所取的这身体，祂公开显示出来，当这身体经历苦难乃至死亡之后，祂又使其从死者中复活，并重建了那被拆毁的殿宇。你们这些视此教义为亵渎的人，却将你们神的肢体限制在一切雌性动物的子宫里，从大象到苍蝇，而不仅仅是\Person{童贞玛利亚}的腹中。也许你们轻视真正的基督，因为圣言据说如此在\Person{童贞玛利亚}腹中降生成人，以致在人性中为自己预备了一座殿宇，而祂自己的本性却保持完整不变；另一方面，你们更看重你们的神，因为在血肉束缚的污秽中，在被系于黑暗群众的部分里，他徒劳地寻求帮助，甚至无力寻求帮助。\par

\SourceRecord{Translated by Richard Stothert.；Nicene and Post-Nicene Fathers, First Series；Vol. 4.；Edited by Philip Schaff.；Buffalo, NY: Christian Literature Publishing Co.,；1887.；<http://www.newadvent.org/fathers/140623.htm>.}

% structure: p0001 source=h1 target=section reason=page-title

\section{驳福斯图斯，卷二十四}

\Emph{\Person{福斯图斯}解释摩尼教否认人是天主所造，认为这仅指血肉之人而非属灵之人。\Person{奥斯定}阐明宗徒\Person{保禄}关于血肉与属灵的对比，从而排除摩尼教的观点。}\par

1. \Person{Faustus}说：有人问我们，为何否认人是天主所造。但我们并非全然否认人是天主所造；只是问在何种意义上、何时、如何被造。因为，按宗徒所说，有两个人：一个他有时称为外在的人，通常称为属土的人，有时也称作旧人；另一个他称为内在的人或属天的人或新人。问题是，哪一个是由天主所造？因为我们也有两次出生：一次是自然把我们带到这光明中，用肉体束缚我们；另一次是真理在我们从错误归正、进入信仰时，重生了我们。耶稣在福音中论到这第二次出生时说：\BibleQuote{人若不重生，就不能见天主的国。}\BibleRef{若 3:3}尼苛德摩不明白基督的意思，困惑不解，问这事怎能成就，因为老人不能进入母腹再次出生。耶稣回答说：\BibleQuote{人若不借水和圣神而生，就不能进天主的国。}然后祂接着说：\BibleQuote{由肉生的属于肉，由神生的属于神。}因此，既然我们身体的出生并非唯一的出生，还有另一种我们在圣神内重生的出生，那么由此区分就产生一个重要问题：在这两种出生中，天主是在哪一种里造了我们？出生的方式也是双重的。在普通生殖的羞辱过程中，我们来自肉体情欲的热力；但当我们进入信仰时，是在耶稣基督内，借着圣神，在良好的教导、尊贵和纯洁中形成。为此，在所有宗教，尤其是基督宗教中，孩童被邀请加入教会。这暗示在祂宗徒的话中：\BibleQuote{我的孩子们，我为你们再受产痛，直到基督在你们内形成。}\BibleRef{迦 4:19}因此，问题不在于天主是否造人，而在于祂造的是什么人，以及何时、如何造。因为如果我们是在母腹中成形时，天主就按祂的肖像造了我们——这是外邦人和犹太人共同的信仰，也是你们的信仰——那么天主就造了旧人，并借感官情欲产生我们，这似乎不适合祂的神性。但如果我们是在归正、被引向更好生活时由天主形成——这是基督和祂宗徒的一般教导，也是我们的教导——那么天主就造了我们作新人，在尊贵和纯洁中产生我们，这完全符合祂神圣可敬的尊威。如果你不拒绝保禄的权威，我们将从他那里向你证明天主造的是什么人，以及何时、如何造。他对厄弗所人说：\BibleQuote{你们应脱去照从前生活的旧人，这旧人是因私欲的迷惑而败坏的；并在你们心思的灵里得以更新，并穿上新人，这新人是照着天主在真理的正义和圣洁中创造的。}\BibleRef{弗 4:22}这表明，在按天主的肖像造人这件事上，所说的是另一个人、另一种出生、另一种出生方式。他所说的脱去和穿上，指向接受真理的时刻；他说新人是天主所造，暗示旧人既不是天主所造，也不是按天主所造。当他补充说这新人是在圣洁、正义和真理中造成时，他指出了另一种出生方式，其特征如我所说，与肉体生殖的方式大不相同。既然他宣告只有前者出于天主，那么后者就不是。同样，他写给哥罗森人时用了相同的话：\BibleQuote{脱去旧人及其行为，穿上新人，这新人是在认识天主上得以更新，按造祂者的肖像在你们内形成。}这里他不仅表明天主所造的是新人，还宣告了形成的时间和方式，因为\Quote{在认识天主上}这句话指向信从的时刻。然后他加上\Quote{按造祂者的肖像}，以表明旧人不是天主的肖像，也非由天主形成。此外，接着的话：\BibleQuote{此处没有男和女，没有犹太人和希腊人，没有野蛮人和西提亚人，}\BibleRef{哥 3:9}更清楚地表明，使我们成为男和女、希腊人和犹太人、西提亚人和野蛮人的出生，不是天主在其中形成人的出生；天主所关涉的出生，是我们失去民族、性别和身份的差异，成为与那一位相同的一位，即基督。所以同一宗徒又说：\BibleQuote{凡受洗归于基督的，都穿上了基督：没有犹太人，也没有希腊人，没有奴隶，也没有自由人，没有男人，也没有女人，但众人在基督内都是一位。}\BibleRef{迦 3:27}因此，天主造人，不是在一个人分裂为多人时，而是在多人成为一体时。分裂在于第一次出生，即肉体的出生；合一来自第二次出生，那是非物质的和神圣的。这为我们提供了充分理由，认为肉体的出生应归于本性，而第二次出生应归于至高尊威。所以同一宗徒又对格林多人说：\BibleQuote{我在基督耶稣内借着福音生了你们，}\BibleRef{格前 4:15}并在谈到自己时对迦拉达人写道：\BibleQuote{然而，那从母胎中分别我出来，又借祂的恩宠召我的天主，既然乐意将祂的圣子启示在我里面，叫我在外邦人中传扬祂，我就没有与血肉商量。}\BibleRef{迦 1:15}显然，他处处谈论的第二次即属灵的出生，是我们被天主所造的出生，与第一次出生的不雅相区别——在第一次出生中，我们在尊严和纯洁上与其它动物同等，因为我们是在母腹中受孕、成形并被生下。你可以看到，在此事上，我们之间的争论与其说是关于教义的问题，不如说是关于解释的问题。因为你们认为，那旧人或外在的人或属土的人是被天主所造的；而我们将这一点应用于属天的人，将优越性归于内在的人或新人。我们的意见并非轻率或没有根据，因为我们是从基督和祂的宗徒学来的，他们被证明是世界上首先这样教导的人。\par

2. \Emph{\Person{奥斯定}}回答说：\Person{保禄}宗徒确实用\Quote{内在的人}指心灵的灵性部分，用\Quote{外在的人}指身体和今世的生命；但我们从未见他将这变成两个不同的人，而是一个，由天主所造的整体，包括内在与外在。然而，只有内在部分是按天主的肖像造的，因为它除了非物质外，还是理性的，低级动物不具备这一点。那么，天主并非造了一个按自己肖像的人，又造了一个不按肖像的人；而是造了一个人，包括内在与外在，祂是按自己的肖像造的，不是就拥有身体和今世生命而言，而是就理性心灵能认识天主、以及拥有理性所赋予的超越一切无理性受造物的优越性而言。\Person{福斯图斯}承认内在的人是由天主造的，因为如他所说，它是在认识天主的知识中，按造祂者的肖像更新的。我乐意根据宗徒的权威承认这一点。为什么\Person{福斯图斯}不同样根据那权威承认\BibleQuote{天主随自己的心意，把肢体一一安置在身体上}呢？\BibleRef{格前 12:18} 我们从此同一宗徒得知，天主也是外在的人的创造者。为什么\Person{福斯图斯}只采纳他认为对自己有利的，而省略或拒绝那些推翻摩尼教徒愚昧之言的？此外，在论及属地的和属天的人，区分可朽与不朽、我们在亚当内的和我们在基督内将成之时，宗徒引用法律关于属地的或属肉体身体的宣告，指的正是那本书和那段记载，其中写着天主也造了属地的身体。他在谈论死人如何复活、以及他们以什么身体前来时，用了谷粒的比喻——种子是光秃秃的种下，天主随己意赐予它身体，并给每粒种子自己的身体——这样，顺便推翻了摩尼教徒的错误，他们说谷粒、植物、所有根和芽都是由黑暗种族所造，而非天主，按照他们的说法，天主在产生这些对象时并未施展能力，反倒自己被囚禁其中——在驳斥了摩尼教徒的不虔之后，他继续描述不同种类的肉身。\BibleQuote{凡肉身并不相同}，他说。然后他谈到天上的和地上的身体，然后谈到我们身体的变化，使之成为属神的和属天的。\BibleQuote{它种下时是羞辱的，复活起来是光荣的；种下时是软弱的，复活起来是强壮的；种下的是属血肉的身体，复活起来的是属神的身体。} 然后，为了表明血肉身体的来源，他说：\BibleQuote{有属血肉的身体，也有属神的身体；正如经上所载：第一个人亚当成了有生命的灵魂。} \BibleRef{格前 15:33} 这记载在创世纪中，\BibleRef{创 2:7} 那里叙述了天主如何造人，并赋予祂用泥土所形成之身体生命。宗徒用\Quote{旧人}仅指旧的生命，即罪恶中的生命，是按亚当的方式，关于亚当曾说：\BibleQuote{罪恶借着一人进入了世界，死亡借着罪恶也进入了世界；这样死亡就传给了众人，因为众人都犯了罪。} \BibleRef{罗 5:12} 因此，这整个人，包括内在和外在部分，都因罪恶而成为旧的，并受制于可朽的惩罚。然而，内在的人有恢复，当它按造物主的肖像更新，脱去不义——即旧人，穿上正义——即新人。但当那属血肉的身体被种下，复活起来成为属神的身体时，外在的人也将获得属天品质的尊严；如此，所有受造物得以再造，所有被造物由创造者祂自己重新制造。这在以下话语中简要说明：\BibleQuote{身体因罪恶而死亡；但灵性因正义而生活。然而，若那使耶稣从死者中复活者的圣神住在你们内，那使基督从死者中复活的，也必要借那住在你们内的圣神，使你们可朽的身体活过来。} \BibleRef{罗 8:10} 任何一个受过天主教教义教导的人都知道，男女性别之分在于身体，而不在于心灵的灵性部分，我们正是在那部分按天主的肖像更新。但宗徒在别处教导天主是两者的创造者；因为他说：\BibleQuote{在主内，女也不可无男，男也不可无女；因为女人是由男人而出，男人是由女人而生；但一切都是出于天主。} \BibleRef{格前 11:11} 对于这一切，那些因心地昏昧而远离天主生命的人，由于无知和心硬，所给的唯一答复是：宗徒著作中任何他们喜欢的便是真的，任何不喜欢的便是假的。这就是摩尼教徒的疯狂，他们若停止做摩尼教徒，就变得有智慧。而实际上，若是问他们：祂是重新制造和更新内在的人的那位吗（他们承认内在的人按天主的肖像更新，而且他们自己引用这段话支持这一点；按照\Person{福斯图斯}，当内在的人按天主的形象更新时，天主造人），他们必会回答是。如果我们接着问，天主在什么时候造了祂现在所更新的，他们必须想出一些托词来避免暴露他们的荒谬。因为按照他们，内在的人并非由天主形成、创造或起源，而是天主自己本体的一部分，被派去对抗祂的敌人；它不是因罪恶而变旧，而是被迫地被敌人俘虏和损伤。无需重复他们所说的一切胡话，他们所说的第一个人，不是宗徒所讲的属土的属地的第一个人（\BibleRef{格前 15:47}），而是来自他们自己捏造谎言的仓库。\Person{福斯图斯}虽然选择人作为讨论主题，但对这第一个人只字不提；因为他害怕讨论中的对手可能知道一些关于他的事。\par

\SourceRecord{Translated by Richard Stothert.；Nicene and Post-Nicene Fathers, First Series；Vol. 4.；Edited by Philip Schaff.；Buffalo, NY: Christian Literature Publishing Co.,；1887.；<http://www.newadvent.org/fathers/140624.htm>.}

% structure: p0001 source=h1 target=section reason=page-title

\section{《驳福斯托》第二十五卷}

\Emph{\Person{福斯托}试图通过讽刺旧约中拟人化的描述，来嘲笑正统信仰关于天主无限性的主张。\Person{奥古斯丁}表示，只要摩尼教徒继续认为灵魂和天主是在空间中延伸的，他便无法说服他们接受关于天主无限性的正确观点。}\par

1. \Person{福斯托}说：\Quote{天主是有限的还是无限的？祂必定是有限的，除非你们错称祂为\InnerQuote{亚巴郎、依撒格和雅各伯的天主}；除非你们所这样称呼的那位与你们称为无限的天主不相同。至于\InnerQuote{亚巴郎、依撒格和雅各伯的天主}，割损的印记将这些人与其他民族分开，也划定了天主能力的界限，只延伸到他们。一个能力有限者不能是无限的。此外，在这样的称呼中，你们连亚巴郎之前的古人也不提，比如哈诺客、诺厄、闪以及类似的人，你们承认他们未受割损却为义人；但因他们缺乏这个分辨的印记，你们就不称天主为他们的天主，而只称祂为\InnerQuote{亚巴郎及其后裔的天主}。如今，如果天主是唯一的且无限的，又何必如此小心谨慎地称呼祂，仿佛只提\InnerQuote{天主}还不够，还必须加上祂是谁的天主——即\InnerQuote{亚巴郎、依撒格和雅各伯的天主}；似乎亚巴郎成了你们呼求时的航标，以免在众多神祇中失事？受过割损的犹太人完全有理由这样称呼这位神明，因为他们称天主为\InnerQuote{割损之神}，与\InnerQuote{未受割损之神}相对。但你们为何也这样做，就难以理解了；因为你们并不声称有亚巴郎的印记，却呼求他的天主。若我们理解正确，犹太人和他们的天主似乎互为标记以便相认，以免彼此丢失。所以天主赐给他们可憎的割损印记，为使他们无论在哪一块土地、哪一个民族中，都能因受割损而被认出是属祂的。而他们也用称天主为\InnerQuote{祖先的天主}来标记祂，为使祂无论在哪里，即使在众神之中，一听到\InnerQuote{亚巴郎的天主、依撒格和雅各伯的天主}之名，就立刻知道是在呼求祂。我们常见许多人同名，无人回应，直到被叫到姓氏。同样，牧人或牧者使用烙印以防他人窃取他的财产。你们这样称天主为\InnerQuote{亚巴郎、依撒格和雅各伯的天主}来标记祂，不仅表明祂是有限的，而且表明你们与祂毫无关系，因为你们没有祂用以认出自己子民的割损印记。因此，如果这就是你们敬拜的天主，那么祂无疑是有限的。但若你们说天主是无限的，就必须首先放弃这位有限的神明，并改变你们的呼求方式，为自己过去的错误表示忏悔。我们已从你们自己的立场出发，证明了天主是有限的。但要决定至高真神是否是无限的，我们只需参照善恶之间的对立。如果恶不存在，那么天主无疑是无限的；否则祂只能是有限的。然而，恶无疑存在；因此天主不是无限的。善止之处，恶便起。}\par

2. \Person{奥古斯丁}答：\Quote{认识你的人没有一个会想到问你关于天主无限性的问题，或与你讨论此事。因为，在你的任何观念中有任何属灵程度之前，你首先必须通过单纯的信仰和某些基础知识，使你的心灵摆脱肉体和物质思想的幻象。你的异端却阻止你这样做，因为它总是将身体、灵魂和天主描述为在空间中延伸，无论是有限的还是无限的，而空间的概念只适用于身体。只要这样，你最好别管这件事；因为你不能在此事上教导真理，正如在其他事上一样；而在此事上你也无学习能力，若非你骄傲好斗，你本可在其他事上学习。因为诸如这样的问题：若没有空间能容下天主，祂怎能是有限的？若圣子完全认识祂，祂怎能是无限的？祂怎能是有限的，却又无界限？祂怎能是无限的，却又完善？那位无尺度的怎能是有限的？那位万物的尺度怎能是无限的？——一切肉体的思想都归于虚无；若要除去肉体思想，它必须首先自惭形秽。因此，对于你提出的天主有限还是无限的问题，最好的结束方式是暂时不再谈论，直到你不再远离基督，祂是法律的终向。关于亚巴郎、依撒格和雅各伯的天主，我们已经说得足够多了，以表明那位所有受造物的真天主为何愿意被自己的子民以这个名称亲切地认识。关于割损，我们也已在几处回应了无知的指责。如果摩尼教徒视这个标记为天主所立，作为脱去肉体的适当象征，就不会在它里面找到可笑之处。他们应按基督徒而非异端者的心思考量这仪式；正如经上所载：\BibleQuote{洁净者一切皆洁净。}但是，考虑到下面这句话的真实性：\BibleQuote{为不洁与不信者，无一洁净，连他们的理智与良心也被玷污。}（\BibleRef{铎 1:15}）我们必须提醒我们诙谐的对手，如果割损如他们所说是不雅的，他们应为此哭泣而非嘲笑；因为他们的神在割下的皮肤和流出的血中，既受限制又受玷污。}\par

\SourceRecord{Translated by Richard Stothert.；Nicene and Post-Nicene Fathers, First Series；Vol. 4.；Edited by Philip Schaff.；Buffalo, NY: Christian Literature Publishing Co.,；1887.；<http://www.newadvent.org/fathers/140625.htm>.}

% structure: p0001 source=h1 target=section reason=page-title

\section{驳福斯图斯，第二十六卷}

\Emph{\Person{福斯图斯} 坚持说，耶稣虽非由母而生，却可能藉神圣能力而死，然而他同样拒绝承认耶稣的降生与死亡。\Person{奥古斯定} 坚持说，有些事甚至连天主也不能作，其中之一便是死亡。他驳斥了摩尼教徒的\OriginalTerm{幻象论}{Docetism}。}\par

1. \Person{福斯图斯} 说：\Quote{你问，若耶稣非由母生，他如何能死？嗯，这是一个盖然性，如人在缺乏证据时所使用的那样。我们却要从你们一般所相信的事例来回答这个问题。若真实，它们将证明我们的论点；若虚假，它们对你的帮助不会超过对我们。你说，耶稣如何能死，若他不是人？我要回问你，厄里亚如何不死，尽管他是人？一个必朽者能侵犯不朽的界限，而基督不能将所需的死亡经验加给祂的不朽吗？若厄里亚违反本性永生不死，为何不允许耶稣以同样少的本性违逆，在死亡中停留三天？此外，不仅是厄里亚，还有梅瑟和哈诺客，你们相信他们是不朽的，且带着肉身被提到天上。因此，若耶稣因死亡而为人的论证是好的，那么厄里亚因不死而不是人的论证也是同等的。但厄里亚不是人这一点是假的，尽管他被认为不朽，所以耶稣是人这一点也是假的，尽管他被认为死了。事实是，若你们肯信，希伯来人在耶稣之死和厄里亚之不死两方面都错了。因为耶稣死了和不死者厄里亚不死，是同样不真实的。但你们相信你们所愿的；至于其余，你们诉诸本性。而允许这一诉诸，本性既反对不朽者之死，也反对必朽者之不死。若我们诉诸天主和人有能力实现其目的，那么耶稣能死似乎比厄里亚不能死更为可能；因为耶稣的能力大于厄里亚。但若你们将较弱者提到天上，尽管本性反对，并忘记他的必朽条件，赐他永福，为何我不能承认耶稣若愿意能死，即使我同意祂的死是真实的，而非仅是一个表象？因为，从祂取人性相似之初，祂就在外貌上经历了人性的所有经验，那么祂以表象的死来完成这体系，是完全一致的。}\par

2. 此外，应当记起，这关于本性所允许的可能性的参考，应与耶稣的全部历史相关，而不仅仅是祂的死亡。根据本性，生来盲目的人看见光是不可能的；然而耶稣似乎行了这样一种奇迹，以至于犹太人自己呼喊说，从创世以来从未见过有人开了生来盲人的眼睛。\EditorialNote{若望福音第九章} 同样，治愈枯干的手，使生来哑吧的人有说话和表达的能力，使死人复活，以及身体在腐烂后恢复生机，引起惊异，并且鉴于本性可能和不可能的事，必须显得完全难以置信。然而，作为基督徒，我们相信所有这些事情都由同一个人所作；因为我们不看本性的律，而是看天主的大能运作。还有一个故事，耶稣被从山崖抛下，且安然无恙逃脱。因此，若从高处被抛下时祂没有死，仅仅因为祂选择不死，那么为何祂不能在自己愿意时拥有死的能力？我们以此方式回答你们，因为你们喜欢讨论，且装作用不属于你们的逻辑武器。关于我们自己的信仰，耶稣死了并不比厄里亚不朽更真实。\par

3. \Emph{\Person{奥古斯定}} 回答说：关于哈诺客、厄里亚和梅瑟，我们的信仰并非由福斯图斯的假设决定，而是由圣经的宣告决定，它们建立在最坚实最确定的证据基础之上。处于错误中的人，如你们，不适合决定什么是本性，什么是违反本性。我们承认，违反人类通常经验的事物通常被称为违反本性。因此宗徒使用这些话：\BibleQuote{如果你从野橄榄树上被砍下，且逆着本性被接在好橄榄树上。}\BibleRef{罗 11:24} 这里违反本性，用于违反人类对本性过程的经验；野橄榄接在好橄榄树上却结出橄榄的肥汁而非野果。但是天主，一切本性的作者和创造者，不作任何违反本性的事；因为凡由祂所作，祂安排一切本性的秩序、尺度与比例，在每种情形下必然是本性的。而人本身只有在犯罪时才违反本性；然后藉惩罚被带回本性。正义的本性秩序要求要么罪不应被犯，要么不应不受惩罚。在任何一种情形下，本性秩序都被保存，若不是由灵魂，至少是由天主。因为罪使良心痛苦，并给罪人的心带来悲伤，因失去正义之光，即使没有带来肉体上的痛苦，这些痛苦是为纠正而施加的，或是为不可救药的人保留的。然而，说天主作了一件违反本性的事，并非不当，当这事违反我们所知的本性时。因为我们称通常的普通本性过程为本性；而天主违反此所作的任何事，我们称为异能或奇迹。但反对本性的最高法律，这法律超出不敬虔者和软弱信徒的知识，天主从不行动，如同祂不反对自己一样。至于精神与理性存有，人类灵魂属于此类，他们越分享这不改变的法律与光明，就越清楚看见什么是可能，什么是不可能；反之，他们距离它越远，他们对未来的感知就越少，而对奇特事件的惊讶就越频繁。\par

4. 因此，关于\Person{厄里亚}所遭遇的事，我们虽无知，但仍相信圣经关于他的真实记述。有一件事我们确知：凡天主所愿的，必会发生；除非出于天主的旨意，任何事都不会临到任何人身上。所以，若有人告诉我，某人的肉身可能转变为属天的身体，我承认这种可能性，但我无法知道它是否会发生；我之所以不知，是因为我不明了天主在这事上的旨意。如果这是天主的旨意，那么它会发生，这是完全清楚且不容置疑的。再者，若有人告诉我，某事若非天主阻止它发生，本来会发生，我则确信地回答：将要发生的是天主的行动，而非本来可能发生的其他事件。因为天主预知祂自己未来的行动，因此祂也知道那行动的效果，即阻止本来可能发生的事件；而且，毫无疑义，天主所知的比人所想的更为确定。因此，未来的事不可能不发生，正如过去的事不可能不曾发生；因为天主绝不会愿意同一件事既真又假。所以，凡真正属于未来的事，不可能不发生；凡不发生的事，从未是未来；正如凡真正属于过去的事，确实都发生过。\par

5. 因此，若说\Quote{如果天主是全能的，就让祂使已发生的事变为未发生}，这实际上等于说：\Quote{如果天主是全能的，就让祂使同一件事既真又假。}天主能使任何存在的事物停止存在，只要那事物是现在存在的；例如，祂用死亡终止一个已出生者的存在，因为这里有一个实际存在可以被终止。但若一个事物不存在，则不能使其存在终止。现在，过去的事不再存在；任何可以被终止存在的事物，不可能是过去的事。真正过去的事不再现在；而其过去存在之真理性在于我们的判断，而不在于已不再存在的事物本身。断言某事是过去的命题，当那事物不再存在时，才是真的。天主不能使这样的命题为假，因为祂不能与真理相矛盾。在此情形下的真理，或正确的判断，首先存在于我们的心智中，当我们认知并表达它时。但即便因我们遗忘而从心智中消失，它仍作为真理存留。过去的事物虽不再现在，但它曾存在，这一事实永远为真；而它停止存在后之过去存在的真理性，与它开始存在前之未来存在的真理性相同。这真理不能被天主所否定，因为在祂内寓居着至高不变之真理，一切存在于任何心智或理智中的真理都来自祂的光照。天主并非全能到能够死；但这无能并不妨碍祂的全能。真正的全能属于那真实存在者，祂独自是一切存在的根源，无论是精神的还是形体的存在。造物主按祂的旨意使用祂所有的受造物；祂的旨意合乎真实不变之正义，祂不变地藉此正义，按其本性或行为的功过，促成所有可变事物的变化。因此，没有人会愚昧到否认\Person{厄里亚}作为天主的受造物，可以变坏或变好；或者，全能天主的旨意可以使他以异于常人的方式变化。所以，我们毫无理由怀疑圣经至高权威关于他的记载，除非我们限制天主的权力于我们所熟悉的事物上。\par

6. \Person{福斯图斯}的论点是：如果\Person{厄里亚}是一个人却可以免于死亡，那么\Person{基督}既然不仅是人，为何不能拥有死亡的能力？这等于说：如果人的本性可以变好，为何神性不能变坏？——这是一个软弱的论证，因为人的本性是可变的，而神性是不可变的。这样的推论方法会导致明显的荒谬：如果天主能将永生的荣耀赐予人，祂也必然有能力将自己投入永罚。\Person{福斯图斯}会反驳说，他的论证只涉及天主死亡三天，相对于人的永生。好吧，若你理解死亡三天是指天主所取的属于我们有死本性的肉身的死亡，那么你完全正确；因为福音的真理表明，\Person{基督}死亡三天是为了人的永生。但你若断言神性本身无需取死性也能死亡三天，正如人的本性可得不朽，则显出你既不认识天主，也不认识天主的恩赐。实际上，既然你将自己神的一部分永远束缚于黑暗之团，你怎能避免已经提及的荒谬结论，即天主将自己投入永罚？你还需要证明部分光明是光明，而部分神不是神。简言之，不须论证即给出我们信仰的真实理由：关于\Person{厄里亚}作为一个凡人被提到天上，以及\Person{基督}确实由童贞女所生并确实死在十字架上，这两者我们信靠的基础都是圣经的宣示，相信是虔敬，不信是不虔敬。关于\Person{厄里亚}的记载，你假装否认——因为你会假装任何事。关于\Person{基督}，虽然你甚至不至于说他不能死（尽管他能生），但你仍否认祂由童贞女所生，并断言祂在十字架上的死是虚假的，这等于也否认了死，除非作为迷惑人的戏弄；你仅仅这样承认，是为了让相信这些虚构的人宽容你自己的谎言。\par

7. 浮士德所展现的问题，即\Quote{若耶稣不是由母亲所生，祂怎能死？}这一问题，似乎是一位天主教徒向他提出的。只有那些忽视亚当虽非由母亲所生却已死去这一事实的人，才会提出这样的问题。谁竟敢说圣子若愿意，不能像为亚当造身体那样，为自己造一个真实的人体呢？因为万物是藉着祂造的。\BibleRef{若 1:3} 或者，谁能否认，那位全能者的全能圣子，若愿意，不能从一个属天的物质，或从空气、蒸气中取一个身体，并把它改变成具有人体完全特有的形态，以致祂能作为人生活并死在其中呢？或再进一步，若祂愿意，不从祂所造的任何物质材料中取身体，而是从无中为自己创造真实的血肉，正如万物都是藉着祂从无中被造一样，我们中没有人会反对说祂不能这样做。我们相信祂由童贞玛利亚所生，不是因为祂若非如此就不能以真实的身体显现于人中，而是因为圣经是这样记载的，我们必须相信圣经才能成为基督徒，或得救。因此，我们相信基督由童贞玛利亚所生，因为福音是这样记载的；我们相信祂死在十字架上，因为福音是这样记载的；我们相信祂的诞生和死亡都是真实的，因为福音并非虚构。祂为何选择从女人所取的身体中承受这一切，只有祂自己知道。或许祂以此方式赋予祂所创造的男女两性以重要性和尊荣，取了男人的形态并由女人所生；或者可能有其它原因，我们不得而知。但可以肯定的是，所发生的事正如福音叙述所告诉我们的，而天主的智慧所决定的正应该发生。我们将福音的权威置于一切异端讨论之上；我们景仰天主智慧的谋划，胜过任何受造物的谋划。\par

8. 浮士德呼求我们相信他，并说：\Quote{事实是，你若愿意相信，希伯来人在耶稣的死和厄里亚的不朽两方面都犯了错误。}稍后他又补充说：\Quote{既然祂从起初取了人的肖象，就在表面上经历了人类所有的遭遇，那么祂以看似死亡来结束这一安排，是完全一致的。}这个无耻的撒谎者，竟宣称基督假装死亡，还指望被人相信吗？当基督说：\BibleQuote{人子必须被杀，第三日复活。} 时，祂说了谎言吗？\BibleRef{路 24:7} 而你却告诉我们要相信你所说的话，仿佛你并不说谎？那样的话，伯多禄比基督更真实，他对祂说：\BibleQuote{主，愿这事远离祢，这绝不能临到祢身上。} 为此，祂对他说：\BibleQuote{撒殚，退到我后面去！} \BibleRef{玛 16:22} 这个责备并没有在伯多禄身上白费，因为在他受纠正并完全预备好之后，他甚至宣讲到自己的死，宣讲基督之死的真理。但若伯多禄因认为基督不会死而被称为撒殚，那么当你们不仅否认基督死了，而且断言祂假装死亡时，你们该被称为什么？你们给出的理由是基督看似死亡，因为祂表面上经历了人类所有的遭遇。但祂假装经历了人类所有的遭遇，只是你们反对福音的意见。事实上，当福音记载说耶稣睡了（\BibleRef{玛 8:24}）、祂饿了（\BibleRef{玛 4:2}）、祂渴了（\BibleRef{若 19:28}）、祂忧愁（\BibleRef{玛 26:37}）或喜乐等等时，这些事情都是真实的，并非假装，而是真实的经历；只是祂经历这些，并非出于单纯的自然必然性，而是出于控制的意志和天主能力的运用。在人身上，愤怒、忧愁、睡眠、饥渴常常是不由自主的；在基督身上，这些都是祂自己的意志行为。同样，人是不由自主地出生，并违背己意地受苦；而基督是出于自己的意志而生并受苦。尽管如此，这些事情是真实的；对它们的准确叙述，旨在教导凡相信基督福音的人得着真理，而不是用谎言来欺骗他。\par

\SourceRecord{Translated by Richard Stothert.；Nicene and Post-Nicene Fathers, First Series；Vol. 4.；Edited by Philip Schaff.；Buffalo, NY: Christian Literature Publishing Co.,；1887.；<http://www.newadvent.org/fathers/140626.htm>.}

% structure: p0001 source=h1 target=section reason=page-title

\section{\WorkTitle{驳福斯特 第二十七卷}}

\Emph{\Person{福斯特}警告不要过分推究这一论证，即如果\Person{耶稣}没有出生，祂就不能受苦。\Person{奥斯定}同样接受出生和死亡，依据的是福音叙述的见证，这比\Person{摩尼}的谎言更具权威。}\par

1. \Emph{\Person{福斯特}}说：\Quote{如果\Person{耶稣}没有出生，祂就不能受苦；但既然祂确实受苦了，祂就必定出生了。我劝你不要在这些事情上求助于逻辑推理，否则你的整个信仰将会动摇。因为，即使按照你们，\Person{耶稣}是奇迹般地由童贞女所生；而由后件到前件的推理表明这是假的。因为你的论证可能这样反过来对付你：如果\Person{耶稣}由女人所生，祂就必由男人所生；但祂并非由男人所生，因此祂不是由女人所生。如果，如你所相信的，祂可以不因受孕而出生，为何祂不能不受出生而受苦呢？}\par

2. \Emph{\Person{奥斯定}}回答说：\Quote{你在这里反驳的论证只能被像你这样成功误导无知者的人使用，而不能被那些有足够知识反驳你的人使用。\Person{耶稣}可以不因受孕而出生，也可以不受出生而受苦。祂是否如此取决于祂自己的意志。祂选择不因受孕而出生，而不选择不受出生而受苦。如果你问我如何知道祂出生且受苦，我在忠实的福音叙述中读到这些。如果我问你如何知道你所声称的，你搬出\Person{摩尼}的权威，并指控福音说谎。即使\Person{摩尼}没有将虚假作为基督的一种卓越来陈述，我也不会相信他的陈述。他赞美虚假并非来自祂在基督里所发现的，而是来自他自己的道德品格。}\par

\SourceRecord{Translated by Richard Stothert.；Nicene and Post-Nicene Fathers, First Series；Vol. 4.；Edited by Philip Schaff.；Buffalo, NY: Christian Literature Publishing Co.,；1887.；<http://www.newadvent.org/fathers/140627.htm>.}

% structure: p0001 source=h1 target=section reason=page-title

\section{\WorkTitle{驳福斯图斯，第廿八卷}}

\Emph{福斯图斯再论族谱，并坚持要审查其自身的一致性。奥斯定以圣经权威为立场，坚持玛窦关于基督诞生的叙述必须被接受为定论。}\par

1. \Emph{\Person{福斯图斯}}说：\Quote{你说，基督若没有诞生，就不能死去。我回答说：若他诞生了，他就不能是天主；或者，若他既可以是天主又可以被诞生，为什么他既可以诞生又不能死去？显然，论证和必然推论并不适用于那些涉及对耶稣应当如何描述的问题。答案必须从他自己所说的话，或从他的宗徒们关于他的话中取得。族谱必须就其自身的一致性加以审查，而不是从基督死亡的假设推论到他诞生的事实；因为他可能未曾诞生而受苦，或者他可能诞生了而从未受苦；因为你们自己也承认，在天主没有不可能的事，这与否认基督可能未曾诞生而受苦是矛盾的。}\par

2. \Emph{\Person{奥斯定}}回答说：你总是在回答无人使用的论证，而不是我们真正的论证，这些你无法回答。没有人说基督若没有诞生就不能死；因为亚当虽然未曾诞生，却死了。我们说的是，基督诞生了，因为这不是这个或那个异端者说的，而是记载在圣福音中；他死了，因为这也记载着，不是在某个异端作品中，而是在圣福音中。你搁置了关于对耶稣应当如何描述的真正问题的论证，并引用了他自己所说的以及他的宗徒们关于他的话；然而，当我开始引用他的宗徒玛窦的福音，其中有关于基督诞生的全部叙述时，你立即否认玛窦写了这叙述，尽管整个教会从宗徒统治时期直至我们时代的诸位主教，一直持续不断地见证这一点。你会引用什么权威来反对这一点？也许是摩尼的某本书，其中否认耶稣是由童贞女所生。那么，正如我认定你的书是摩尼的作品，因为它从摩尼生活的时代至今，通过你们领袖的连续传承，在摩尼的门徒中保存并传递下来，同样，我要求你相信我所引用的书是玛窦所写，因为它从玛窦的时代起，在教会中一直传承至今，其间没有中断。那么问题在于，我们是相信一位曾与基督在世时同在的宗徒的话，还是相信一个远在波斯、生于基督之后很久的人的话？但也许你会引用另一本以一位已知被基督拣选的宗徒署名的书；你会在那里发现基督不是由玛利亚所生。那么，既然其中必有一本是假的，问题在于，我们是相信一本从一开始就在基督亲自建立的教会中传承、并通过宗徒及其继任者延续至今、遍及全球而从未中断的书，还是相信一本被这个教会斥为未知的书，而且是由那些通过为基督说谎而证明自己诚实的人提出的？\par

3. 此时你会说：\Quote{查考两部福音中的族谱，看它是否自身一致。}对此的回答已经给出。你的困难在于若瑟如何能有两位父亲。但即使你无法想到一种解释，即一位是他的生父，另一位是收养他的父亲，你也不应如此轻易地使自己反对如此崇高的权威。既然这解释已给了你，我呼吁你承认福音的真理，尤其是停止你对真理的有害而无理的攻击。\par

4. \Person{福斯图斯}似乎很巧妙地引用了\Person{耶稣}论及自己的话。但若不藉祂门徒的叙述，如何能知道这些事？我们若不相信他们告诉我们\Person{基督}由童贞女所生，又怎能相信他们记载\Person{基督}论及自己的话？因为，关于任何自称直接来自\Person{基督}本人的著作，如果真是祂的，怎么会在教会中不被诵读、不被承认、不被视作至高权威？这教会从\Person{基督}本人开始，由祂的宗徒延续，宗徒之后继以主教，一直保持并扩展到今日，且在教会中应验了许多先前的预言，而关于末日的预言也必将实现。至于这样一部著作的出现，需要考虑它从何处发端。假设它出自\Person{基督}本人，那些与祂直接相关的人很可能接受了它，并传给了他人。在这种情况下，该著作的权威将由各团体及其首领的传统充分确立，如我所说过的。那么，谁会在今天如此愚妄，竟相信\Person{摩尼}所发布的基督书信是真的，或者不相信\Person{玛窦}关于基督言行的叙述呢？或者，如果问题是\Person{玛窦}是否为真正的作者，那么，在这方面，谁不会相信他在教会中所找到的——这教会从\Person{玛窦}时代至今有着清晰而不断绝的历史——反而相信一个来自波斯的闯入者，他比\Person{玛窦}晚了两百多年，却要我们相信他对基督言行的叙述而非\Person{玛窦}的？然而，即使是宗徒\Person{保禄}，他在主升天后由天上蒙召，教会若非当时还有宗徒在世，他得以与他们交通并对照自己的福音与他们的福音，从而被认可为属于同一团体，教会也不会相信他。当确认\Person{保禄}所宣讲的与宗徒们所宣讲的相同，且他生活在与他们的共融与和谐中，再加上天主藉\Person{保禄}行如宗徒所行的奇迹赐予见证，他的权威变得如此之大，以致现在教会中接受他的话，就如他自己恰当地说的，基督在他内说话。\BibleRef{格后 13:3}另一方面，\Person{摩尼}却认为基督的教会应该相信他所说的，与圣经相悖，而圣经有如此强而持续的证据支持，教会从中找到郑重的禁令：凡向她宣讲与她所领受的不同者，当受诅咒。\BibleRef{迦 1:8}\par

5. \Person{福斯图斯}告诉我们，他有充分理由断定这些圣经不值得信赖。然而他又说不用论证。但论证也要被驳斥。整个论证的终极目的是要使灵魂相信，它在这世界受苦的原因，是它防止天主失去祂的国度，并且天主的本体与本性如此易变、可朽、受损、受染，以致其中一部分无可救药地被玷污，被祂自己投入黑暗大团中永远受罚，尽管这部分在与祂无害的结合中、未犯任何罪时，祂明知故犯地将它送到那遭受玷污的地方。这就是你们一切论证和虚构的结局；但愿这些论调和谎言也从你们心中和口中消失，好使你们有朝一日不再相信和吐出那些可憎的亵渎！但是，\Person{福斯图斯}说，我从著作本身证明它们在所有点上不可靠，因为它们互相矛盾。那么，如果它们的见证不一致且自相矛盾，为什么不说它们完全不可靠呢？但是，\Person{福斯图斯}说，我说我认为符合真理的话。什么真理？这真理只是你自己的虚构，始于天主的战斗，继而祂的污染，终于祂的定罪。没有人，\Person{福斯图斯}说，会相信自相矛盾的著作。但如果你认为它们自相矛盾，那是因为你不理解它们；因为你所引用支持你意见的段落已显明了你的无知，对你可能引用的任何其他段落也是如此。因此，我们没有理由不相信这些著作，它们有如此重大的见证支持；而这本身就是最好的理由，来宣布那些宣讲与其中所写不同的人当受诅咒。\par

\SourceRecord{Translated by Richard Stothert.；Nicene and Post-Nicene Fathers, First Series；Vol. 4.；Edited by Philip Schaff.；Buffalo, NY: Christian Literature Publishing Co.,；1887.；<http://www.newadvent.org/fathers/140628.htm>.}

% structure: p0001 source=h1 target=section reason=page-title

\section{驳斥福斯图斯，第二十九卷}

\Emph{福斯图斯试图为摩尼教的幻象说辩护。奥斯定坚持认为出生毫无可耻之处。}\par

1. \Person{Faustus}说：如果基督是可见的，且没有经过出生而受苦，这便是巫术。你们这个论据可以反过来对付你们，我们可以回答说，如果祂没有受孕而生下，那便是巫术。童贞女生产不符合自然规律，更不用说生产后仍是童贞女。那么，你们为何拒绝承认基督以超自然方式受苦而不受出生条件限制呢？相信我：实质上，我们双方的信仰都违反自然；但我们的信仰是正派的，而你们的则不是。我们对基督受难给出了至少合理的解释，而你们对祂的出生给出的唯一解释是错误的。总之，我们认为祂只是外表受苦，并非真的死亡；你们相信真实的出生和子宫中的受孕。若非如此，你们只需承认出生也是幻觉，我们所有的争论就结束了。至于你们经常声称的，基督若不经过出生就无法显现或与人交谈，这是荒谬的；因为我们的导师们已经证明，天使们常常显现并与人交谈。\par

2. \Person{Augustine}回复：我们并不说未经出生而死亡是巫术；因为正如我们已说过的，这在亚当身上发生过。但是，即使从未发生过，谁敢说基督若愿意，不能不带从童贞女取的身体而来，却以真实的身体显现，以真实的死亡救赎我们呢？然而，更好的是祂实际所做的那样：由童贞女所生，并以祂的屈尊俯就，为要拯救的两性——男人和女人——赋予尊荣，取了由女人所生的男人的身体。在此，祂强烈地见证反对你们，并驳斥你们的教义，即认为两性是魔鬼的作为。我们在你们教义中所称的巫术，是你们使基督的受难和死亡成为仅仅是外表，以致通过鬼影般的幻觉，祂似乎死了而实际上没有。因此，你们也必须使祂的复活成为鬼影、幻觉和虚假；因为如果没有真实的死亡，就不可能有真实的复活。因此，祂向怀疑的门徒显示的伤痕也必定是虚假的；多默在呼喊\Quote{我的主，我的天主！}时，不是被真理所确定，而是被谎言所欺骗。\BibleRef{若 20:28} 然而你们却要我们相信，你们的舌头说出真理，虽然基督的整个身体是谎言。我们对你们的论点是：你们所造的基督是这样的，除非你们也实行欺骗，否则你们不能成为祂真正的门徒。基督的身体是唯一由童贞女所生的，这并不证明祂的出生中有巫术，正如它是唯一在第三天复活、永不再死的身体一样。你们会说，主的一切奇迹都是巫术，因为它们不寻常吗？它们确实发生了，它们在人们眼中所显现的外表是真实的，而非幻觉；当它们被说成违反自然时，并非它们反对自然，而是它们超越了我们所习惯的自然方式。愿天主保持祂的子民中那些仍在基督里为婴孩的人的心，不受福斯图斯的影响，当他推荐一种义务，即我们应承认基督的出生是幻觉而非真实，以便结束我们的争论时。不，实在说，让我们继续为真理而与他们争辩，而不是与他们同流合污于虚假。\par

3. 但如果我们这样说就能结束争论，为什么我们的对手自己不这么说呢？当他们断言基督的死亡不是真实的而是伪装的，为什么他们又声称祂根本没有出生，甚至连祂死亡的那种出生也没有？如果他们如此尊重福音作者的权威，以至于不得不承认基督至少在外表上受苦，那么同样的权威也见证了祂的出生。确实，两位福音作者记述了出生的事迹；但在所有福音书中，我们都读到耶稣有一位母亲。也许福斯图斯不愿意将出生视为幻觉，因为玛窦和路加所载的族谱差异导致了表面上的不一致。但是，假设一个无知的人，他会认为福音作者在许多关于基督受难的事情上也不一致；假设一个受过教导的人，他会发现完全一致。难道伪装死亡是允许的，而伪装出生就是错误的吗？然而福斯图斯要我们承认出生是伪装的，以便结束争论。我们在接下来回答另一个异议时，会说明我们认为福斯图斯为何不接受任何出生——即使是伪装的出生——的原因。\par

我们否认圣人的身体中有什么可耻之处。诚然，有些肢体被称为不雅观，因为它们不像那些经常显露的肢体那样悦目。但请留意宗徒所说的话，当他从身体的合一与和谐中，向教会劝勉爱德时：\BibleQuote{身上那些似乎软弱的肢体，更是必要的；而我们以为身体上不体面的肢体，我们越加给予光荣；不雅的肢体，反越加显得雅致。因为我们雅致的肢体并不需要什么；但天主将身体搭配好了，把更多的光荣给予了那有缺欠的肢体，免得身体上发生分裂。}\BibleRef{格前 12:22}对这些肢体的放纵和不节制使用是可耻的，而非肢体本身；因为它们不仅被未婚者纯洁地保存，也被圣洁生活的已婚父母保存，在他们身上，自然的欲望不是服务于情欲，而是服务于生育子女的明智目的，因此绝无任何可耻之处。更有甚者，在童贞玛利亚身上，她因信德怀孕了基督的身体，她的肢体毫无可耻之处，因为这些肢体并非用于普通的自然怀孕，而是用于奇迹性的生产。为使我们可以以真诚的心怀怀孕基督，并在宣认中仿佛生出祂，基督的身体理当来自祂母亲的实体，又不损害她身体的纯洁。我们不能认为基督的母亲因祂的出生而有所损失，或生育的恩赐取代了童贞的恩宠。如果这些事件——是真实的而非幻象——是新颖奇异的，且违反自然常道，那么原因在于它们伟大、惊人且神圣；正因如此，它们更真实、稳固和确凿。\Person{福斯图斯}说：天使出现说话，而并非由出生而来。仿佛我们坚持认为基督若非由女人所生，就不能显现或说话！祂能够，但祂不选择如此；而祂所选择的便是最好的。祂选择做祂所做的事是显然的，因为祂行事出于自愿，而非像你们的神那样出于必然。我们知道祂出生了，因为我们并非相信异端者，而是相信基督的福音。\par

\SourceRecord{Translated by Richard Stothert.；Nicene and Post-Nicene Fathers, First Series；Vol. 4.；Edited by Philip Schaff.；Buffalo, NY: Christian Literature Publishing Co.,；1887.；<http://www.newadvent.org/fathers/140629.htm>.}

% structure: p0001 source=h1 target=section reason=page-title

\section{\WorkTitle{驳斥福斯托斯，第三十卷}}

\Emph{\Person{福斯托斯}驳斥了那种暗示，即\Person{保禄}关于那些禁止嫁娶、戒绝肉食等的预言更适用于\OriginalTerm{摩尼教徒}{Manichæans}而非\OriginalTerm{天主教苦修者}{Catholic ascetics}，而后者在教会中受到极高的尊崇。\Person{奥古斯丁}则论证了这预言的适用性，并指出了摩尼教苦修主义与基督教苦修主义之间的区别。}\par

1. \Person{福斯托斯}说：\Quote{你们将\Person{保禄}的话应用在我们身上：\BibleQuote{有些人将离弃信德，听从欺诈的神和魔鬼的教义；他们以伪善说谎，良心如同被烙铁烙过；禁止嫁娶，戒绝食物，这些食物原是天主所造，供信徒以感恩之心领受。}\BibleRef{弟前 4:1}我拒绝承认宗徒说了这话，除非你们先承认\Person{梅瑟}和先知们教导的是魔鬼的教义，他们是说谎的和恶毒的神的解释者；因为他们极力戒绝猪肉和其他肉类，称之为不洁。这个案子必须先解决；你们必须长久仔细地思考应如何看待他们的教导：他们是从天主还是从魔鬼说了这些。关于这些事，要么\Person{梅瑟}和先知们必须与我们一同被定罪，要么我们必须与他们一同被宣告无罪。你们现在谴责我们，如你们所做的，说我们是魔鬼教义的追随者，这是不公正的，因为我们只要求司祭阶层戒绝动物类食物；我们把禁令限于司祭，而你们却认为你们的先知，以及\Person{梅瑟}本人，禁止所有阶层的人吃猪肉、野兔、蹄兔，以及各种乌贼和所有无鳞鱼，这些话不是出于说谎的神，也不是出于魔鬼的教义，而是出于天主，出于圣神。那么，即使\Person{保禄}说了这些话，你们也只能通过谴责\Person{梅瑟}和先知们来说服我；因此，即使你们不会为了理性或真理而这样做，你们也会为了自己的肚子而反驳\Person{梅瑟}。}\par

2. 此外，在你们的\BibleBook{达尼尔}{书}中有关于三青年的记载，你们会发现很难使其与戒绝肉食是魔鬼的教义这一观点相协调。因为我们被告知，他们不仅戒绝了律法所禁止的，甚至戒绝了律法所允许的；\BibleRef{达 1:12}而你们习惯于赞美他们，视他们为殉道者；然而，如果这是宗徒的意见，他们也是遵循了魔鬼的教义。而且\Person{达尼尔}本人宣称，他禁食三周，不吃肉，不喝酒，同时为他的子民祈祷。\BibleRef{达 10:2}他怎能夸耀这魔鬼的教义，以说谎的神的谎言为荣呢？\par

3. 再者，我们该如何看待你们，或你们中较好的基督徒，其中一些人戒绝猪肉，一些人戒绝四足动物的肉，还有一些人戒绝所有动物类食物，而整个教会因此钦佩他们，并以极大的尊敬看待他们，几乎视如神明？你们顽固地拒绝考虑，如果引自宗徒的话是真实可靠的，那么这些人也被魔鬼的教义所误导。还有另一种遵守，没有人敢解释或否认，因为它众所周知，并且在全世界天主教徒的聚会中每年都特别实行——我指的是\OriginalTerm{四旬期}{fast of forty days}的斋戒，在适当的遵守中，一个人必须戒绝所有那些按照这节经文是由天主创造供我们领受的事物，而同时他称这种戒绝是魔鬼的教义。那么，我亲爱的朋友们，我们是否要说，你们在守斋期间，在庆祝\OriginalTerm{基督受难的奥迹}{mysteries of Christ's passion}时，也过着魔鬼般的生活，被诱惑之神所欺骗，说话虚伪，良心被烙铁烙过？如果这不适用于你们，那也不适用于我们。对这节经文，或它的作者，该作何感想？它适用于谁？因为它既不符合旧约的传统，也不符合新约的制度。关于新约，证据来自你们自己的实践；尽管旧约只要求戒绝某些东西，但它仍然要求戒绝。另一方面，你们的这个意见使所有戒绝动物类食物都成为魔鬼的教义。如果这就是你们的信仰，我再说一次，你们必须谴责\Person{梅瑟}，拒绝先知，并对你们自己做出同样的判决；因为正如他们总是戒绝某些种类的食物，你们有时也戒绝所有食物。\par

4. 但若你们以为在食物上加以区别乃是梅瑟和先知们所设立的神圣规诫，而非魔鬼的教义；若\Person{达尼尔}在圣神内遵守了三周的禁食；若青年\Person{阿纳尼雅}、\Person{阿匝黎雅}和\Person{米沙耳}在天主的指引下选择以卷心菜或豆类为生；若你们中那些禁食者并非受魔鬼鼓动；若你们四十天不饮酒不食肉并非出于迷信而是出于天主的命令——那么，我恳求你们想一想，若把这些字句（即\Quote{禁止食物和禁止结婚是魔鬼的教义}）当作\Person{保禄}的话，岂不是完全的疯狂吗？\Person{保禄}不可能说过将童贞女奉献给\Person{基督}是魔鬼的教义。但是你们读到这些话，却如往常一样轻率地应用在我们身上，没有看出这也就把你们的贞女也定为被魔鬼教义所迷惑，而你们自己则成为魔鬼的仆役，不断努力引诱贞女们发此誓愿，以致在你们所有的教会中，贞女的人数几乎超过了已婚妇女。若她们所履行的是魔鬼的意愿而不是\Person{基督}的，你们为何仍坚持这样的做法？为何要陷害可怜的青年女子？但首先，我想知道，使人成为贞女是否在任何情况下都是魔鬼的教义，还是仅仅禁止结婚才是？若是禁止，那就不适用于我们，因为我们同样认为阻止愿意的人乃是愚蠢，而强迫不情愿的人则是罪恶且不敬的。但若你们说，鼓励这提议、不抗拒这样的渴望，就是魔鬼的教义，那么且不说对你们自身的后果，宗徒本人也将陷入危险，因为他必定被认为将魔鬼的教义引入了\Place{依科尼雍}，因为\Person{特克拉}在订婚后，因他的讲论而燃起了保持终身贞洁的渴望。我们该如何说\Person{耶稣}，祂自己是导师、一切圣德的源泉，是发此誓愿的贞女们的未嫁之夫？祂在\WorkTitle{福音}中列举三种阉人：天生的、人为的和自愿的，并将棕榈枝赐给那些\BibleQuote{为天国而自阉的人} \BibleRef{玛 19:12}，意指那些从心中根除了婚姻欲望的男女青年，他们在教会中如同王宫中的阉人。这也是魔鬼的教义吗？那些话也是出于诱惑的精神吗？若\Person{保禄}和\Person{基督}被证明是魔鬼的司祭，那么那在天主内说话的同一个神岂不也是？我不提我们主的其他宗徒：\Person{伯多禄}、\Person{安德肋}、\Person{多默}，以及独身的典范，真福\Person{若望}，他们以各种方式向青年男女推荐这誓愿的高贵，并为我们，也为你们，留下了使人成为贞女的形式。我不提他们，因为你们不接纳他们进入正典，所以你们会毫无顾忌地不虔敬地将魔鬼的教义归给他们。但你们也要对\Person{基督}，或对宗徒\Person{保禄}说同样的话吗？我们知道，\Person{保禄}处处表示对未婚妇女的偏爱胜过已婚者，并在圣\Person{特克拉}的例子中提供了榜样。但若\Person{保禄}传讲给\Person{特克拉}的教义（其他宗徒也传讲）不是魔鬼的教义，我们怎能相信\Person{保禄}留下记录说，劝勉成圣本身就是魔鬼的命令和教义呢？仅仅通过劝勉使人成为贞女而不禁止结婚，这并非你们特有的做法。这也是我们的原则；认为私法可以禁止公法所允许之事的人，不仅是愚人，更是疯子。因此，关于婚姻，我们也鼓励贞女在自愿的情况下保持现状；我们不违背她们的意愿使她们成为贞女。因为我们知道，当公法反对时，意志和自然欲望的力量；何况当法律只是私人的，人人都可自由违背时。因此，若以这种方式使人成为贞女不是罪过，我们与你们一样无罪。若以任何方式使人成为贞女都是错的，你们与我们有罪。所以，你们引用这经文反对我们的用意或意图，是无法说清的。\par

5. \Person{奥斯定}答复说：请听，你们就会知道我们引用这节经文对抗你们的意思和意图，因为你们说自己不知道。并非你们禁食动物食品；因为正如你们所说，我们古代的祖先禁食某些食物，但并非谴责它们，而是具有预表意义，你们不明白，而我在本作品中已经说了所有必要的话。此外，基督徒（不是异端者，而是天主教徒）为了克制身体，使灵魂在祈祷中更谦卑，不但禁食动物食品，也禁食某些植物产品，但并不认为它们是不洁的。少数人总是这样做；在某些季节或日子，如四旬期，几乎所有人或多或少地这样做，根据个人的选择或能力。而你们否认受造物是好的，称其不洁，说动物是由魔鬼在邪恶的实质中最坏的杂质中制造的，因此你们怀着恐惧拒绝它们，认为它们是最残酷、最可憎的监禁你们之神的地方。你们作为一种让步，允许你们的追随者（与司铎不同）食用动物食品；正如宗徒在某种情况下允许的，不是一般意义上的婚姻，而是在婚姻中满足情欲。\BibleRef{格前 7:5} 只有罪才这样被宽容。这就是你们对一切动物食品的感觉；你们从自己的异端中学来了，并教导你们的追随者。你们宽容你们的追随者，因为如我之前所说，他们供应你们必需品；但你们允许他们，却不说不算有罪。至于你们自己，你们避开这种邪恶和不洁的接触；因此我们引用这节经文反对你们的原因，就在宗徒的这些话中，在你们结束引用的那些话之后。也许正是因为这个原因，你们省略了这些话，然后说你们不知道我们引用的意思和意图；因为你们删去对我们的意图的描述，比表述它更合适。因为，在谈到禁食天主所造、为叫信徒以感恩的心领受的肉类之后，宗徒接着说：\BibleQuote{并给那些知道真理的人；因为天主所造的一切都是好的，只要以感恩的心领受，没有一样是可摒弃的：因为天主的话和祈祷使它成圣。}\BibleRef{弟前 4:3} 这是你们否认的；因为你们禁食这类食物的想法、动机和信念是，它们不是预表性的，而是本性邪恶和不洁的。在这方面，你们确实亵渎了造物主；这就是魔鬼的道理。你们不必惊讶，在如此久远之前，圣神就预言了你们的事。\par

6. 同样，如果你们对童贞的劝勉类似于宗徒的教导：\BibleQuote{那结婚的做得好，那不结婚的做得更好}，\BibleRef{格前 7:38} 如果你们教导婚姻是好的，而童贞更好，如同真正属于基督的教会所教导的，你们就不会在圣神的预言中被描述为禁止嫁娶了。一个人禁止某事，就是使其成为邪恶的；但一件好事可以被置于更好的事之后而不被禁止。此外，唯一体面的婚姻，或为正当合理目的而缔结的婚姻，恰恰是你们最憎恨的。所以，尽管你们可能不禁止性交，但你们禁止婚姻；因为婚姻的特点不仅是为满足情欲，而且如婚约所写，是为生育子女。虽然你们允许许多追随者因拒绝或无法服从你们而保留与你们的联系，但你们不能否认你们作出了禁止。这禁令是你们虚假教义的一部分，而容忍只是为了社团的利益。在这里我们看到了原因（我一直推迟到现在才提），即你们使基督的诞生而非死亡成为虚幻假象。死亡是灵魂（即你们之神的本质）与属于其仇敌的身体的分离，因为身体是魔鬼的作为，你们维护并赞同它；因此，按照你们的信经，基督虽然未死，但应该通过看似死去来赞扬死亡。而在诞生方面，你们相信你们的神是被束缚而非释放；所以你们不允许基督甚至以这种虚幻的方式诞生。如果\Person{玛利亚}不作母亲而保持童贞，你们会更赞赏她，超过她作为母亲而不失童贞。所以，你们看，劝勉童贞作为两件好事中更好的，与通过谴责婚姻的真正目的而禁止嫁娶，有着巨大区别；禁食食物作为预表性的遵守或为克制身体，与禁食天主所造的食物因为天主没有创造它，也有巨大区别。一种情况是先知和宗徒的道理；另一种情况是撒谎魔鬼的道理。\par

\SourceRecord{Translated by Richard Stothert.；Nicene and Post-Nicene Fathers, First Series；Vol. 4.；Edited by Philip Schaff.；Buffalo, NY: Christian Literature Publishing Co.,；1887.；<http://www.newadvent.org/fathers/140630.htm>.}

% structure: p0001 source=h1 target=section reason=page-title

\section{驳斥福斯图斯，第三十一卷}

\Emph{经文段落：\BibleQuote{对于纯洁的人，一切都是纯洁的，但对于不洁和污秽的人，没有一样是纯洁的，就连他们的理智和良心也都污秽了}，从\OriginalTerm{摩尼教徒}{Manichaean}和\OriginalTerm{大公的}{Catholic}的角度进行讨论，\Person{福斯图斯}反对将其应用于他的团体，而\Person{奥古斯丁}坚持应用。}\par

1. \Emph{\Person{福斯图斯}}说：\Quote{\BibleQuote{对于纯洁的人，一切都是纯洁的。但对于不洁和污秽的人，没有一样是纯洁的，就连他们的理智和良心也都污秽了。}关于这节经文，也很有疑问，你们是否为了自己的利益，应当相信它是由\Person{保禄}所写。因为由此会推出，\Person{梅瑟}和先知们不仅在其法律中如此区分食物是受魔鬼影响，而且他们自己在理智和良心上也是不洁和污秽的，以至于下面的话也适用于他们：\BibleQuote{他们自称认识天主，但在行为上却否认祂。}\BibleRef{铎 1:16}这对任何人都比对\Person{梅瑟}和先知们更适用，因为人们知道他们的生活方式与认识天主的人极不相称。直到现在，我只想到通奸、欺诈和谋杀玷污了\Person{梅瑟}和先知们的良心；但现在，根据这节经文，显然他们也被玷污了，因为他们把某些东西视为污秽。那么，你们怎能坚持认为，这样的男人曾被赋予了神圣威严的显现？因为经上记载，只有心地纯洁的人才能看见天主。即使他们免于非法罪行，这种对某种食物的迷信禁戒，如果玷污了心智，也足以使他们不能看见神性。\Person{达尼尔}和三位青年的夸耀也永远消失了，直到我们被告知没有不洁之物以前，他们在\OriginalTerm{犹太人}{Jew}中被视为非常纯洁和品德高尚的人，因为他们遵守祖传习俗，小心避免被\OriginalTerm{外邦人}{Gentile}的食物、尤其是祭物所玷污。\BibleRef{达 1:12}现在看来，当他们紧闭嘴唇，不碰血和偶像宴席时，他们在理智和良心上最为污秽。}\par

2. 但也许他们的无知可以原谅他们；因为当时这种\OriginalTerm{基督教}{Christian}的\Quote{万物对纯洁者都是纯洁的}教义尚未出现，他们可能认为有些东西是不洁的。但在\Person{保禄}宣布\Quote{没有一样东西是不纯洁的}，\Quote{禁戒某些食物是魔鬼的教义}，以及\Quote{认为任何东西污秽的人，其心智受玷污、良心被污染}之后，你们就没有借口了，如果你们不仅像我们所说的那样禁戒，而且还将此视为功德，并相信你们越节制，就越蒙\Person{基督}悦纳，或者，按照这个新教义，你们的心智越被玷污、良心越被污染。还应注意到，世上三种宗教——\OriginalTerm{犹太人}{Jew}、\OriginalTerm{外邦人}{Gentile}和\OriginalTerm{基督徒}{Christian}——虽然以非常不同的方式，都将贞洁和禁欲作为净化心智的方法，但\Quote{万物都是纯洁的}这一看法不可能来自三者中的任何一个。它当然不是来自\OriginalTerm{犹太教}{Judaism}，也不是来自\OriginalTerm{异教}{Paganism}（它也区分食物）；唯一的区别是，\OriginalTerm{希伯来}{Hebrew}的动物分类与\OriginalTerm{异教}{Pagan}的不一致。至于\OriginalTerm{基督信仰}{Christian faith}，如果你们认为\Quote{不把任何东西视为污秽}是\OriginalTerm{基督教}{Christianity}特有的，那么你们首先必须承认，你们中间没有\OriginalTerm{基督徒}{Christian}。因为祭偶像之物和自死的动物，暂且不提别的，在你们所有人看来都是极大的污秽。再者，如果这是一种\OriginalTerm{基督徒}{Christian}的做法，那么反对一切禁戒污秽的教义也不能追溯到\OriginalTerm{基督教}{Christianity}。那么，\Person{保禄}怎么能说出与任何宗教都不一致的话呢？事实上，当宗徒从\OriginalTerm{犹太人}{Jew}变成\OriginalTerm{基督徒}{Christian}时，这是习俗的改变多于宗教的改变。至于这节经文的作者，似乎没有宗教支持他的观点。\par

3. 所以，每当你发现圣经中其他可以用来攻击我们信仰的内容时，首先要确保它不反对你，然后才开始攻击我们。例如，你们不断引用的关于\Person{伯多禄}的段落：他曾看见一个器皿从天上降下，里面有各种动物和蛇，当他感到惊讶和惊奇时，有声音对他说：\BibleQuote{伯多禄，杀死并吃掉你在器皿里看到的一切。}他回答说：\BibleQuote{主，凡俗物和不洁之物，我总不吃。}于是声音再次说道：\BibleQuote{天主所洁净的，你不可称为不洁。}\BibleRef{宗 10:11}这确实似乎有寓意，而不是指食物没有区别。但既然你们选择给它这个意义，你们就必须按照\Person{伯多禄}的这个异象，吃掉所有野兽、蝎子、蛇和一般爬行动物。这样，你们就会显示你们确实顺从了据说\Person{伯多禄}所听到的声音。但你们决不可忘记，你们同时谴责了\Person{梅瑟}和先知们，他们曾认为许多东西是污秽的，而照这个声音所说，天主已经圣化了它们。\par

4. \Emph{\Person{奥古斯丁}}回答说：当宗徒说\Quote{对洁净的人一切都是洁净的}时，他指的是天主所创造的本性——正如\Person{梅瑟}在\BibleBook{}{创世纪}中所记：\BibleQuote{天主造了一切，看，都是很好的}\BibleRef{创 1:31}——而不是指预表意义，按照这些意义，天主借着同一位\Person{梅瑟}区分了洁与不洁。关于这一点，我们已经多次详细讨论过，这里不再赘述。很明显，宗徒称那些在新约启示之后仍然主张遵守未来之事的影儿、好像没有这些外邦人就不能获得在基督内的救恩的人为不洁的，因为他们在这一点上是属肉体的思想；并称他们为不信的，因为他们不分辨法律时期与恩宠时期。对他们来说，他说，没有一样是洁净的，因为他们对他们所接受的和所拒绝的都作出了错误而有罪的使用；这对所有不信的人来说都是真的，但特别是对你们摩尼教徒，因为对你们来说，没有任何东西是洁净的。因为，虽然你们非常小心地使你们所用的食物与肉体的玷污分开，但它对你们仍然不是洁净的，因为你们承认它的唯一创造者是魔鬼。你们还认为，通过吃它，你们释放了你们的神，它在这食物中遭受监禁和污染。人们会认为你们可能认为自己洁净，因为你们的胃是净化你们的神的适当地方。但按照你们的意见，你们自己的身体也是黑暗种族的本性和作品；而你们的灵魂仍然受到你们身体的玷污。那么，对你们来说什么是洁净的呢？不是你们所吃的东西；不是你们盛食物的容器；不是你们自己——借着你们，它被洁净。因此你们看到宗徒的话是针对谁的；他表达自己以便包括所有不洁和不信的人，但首先而且主要是谴责你们。因此，对洁净的人来说，一切都是洁净的，按照它们被创造的本性；但对于古代的犹太人民来说，一切都不是按照它们的预表意义洁净的；并且，就身体的健康或社会的习俗而言，并非所有事物都适合我们。但当事物处于适当的位置，并保持自然的秩序时，对洁净的人一切都是洁净的；但对于不洁和不信的人，其中你们是首位的，没有一样是洁净的。如果你们渴望治愈你们麻木的良心，你们可以有益地将宗徒的以下话语应用到自己身上。这些话是：\Quote{他们的心思和良心都被玷污了。}\par

\SourceRecord{Translated by Richard Stothert.；Nicene and Post-Nicene Fathers, First Series；Vol. 4.；Edited by Philip Schaff.；Buffalo, NY: Christian Literature Publishing Co.,；1887.；<http://www.newadvent.org/fathers/140631.htm>.}

% structure: p0001 source=h1 target=section reason=page-title

\section{\WorkTitle{驳斥福斯图斯，第三十二卷}}

\Emph{\Person{福斯图斯}未能理解为何他必须要么全部接受要么全部拒绝新约，而天主教徒却随意接受或拒绝旧约的各部分。\Person{奥古斯丁}否认天主教徒任意对待旧约，并解释了他们对其的态度。}\par

1. \Person{福斯图斯}说：\Quote{你们说，如果我们相信福音，就必须相信其中所写的一切。那么，既然你们相信旧约，为什么不信其中任何部分所写的一切？相反，你们只挑出预言将来有一位犹太人的王的预言，因为你们认为这就是\Person{耶稣}，还有几条共同的道德诫命，例如\BibleQuote{不可杀人}、\BibleQuote{不可奸淫}；其余的一切你们都略过，认为其他事物就如\Person{保禄}所认为的、他视为粪土的事物一样。\BibleRef{斐 3:8} 那么，为什么对我来说，从新约中挑选最纯洁、有助于我救恩的部分，而将你们前辈的篡改（这些篡改损害了新约的尊严和恩宠）置之不理，就显得奇怪或特别呢？}\par

2. 如果圣父的盟约中有些部分我们不必遵守（因为你们将犹太法律归于圣父，并且众所周知其中许多内容使你们震惊和羞愧，以至于在内心你们不再视之为纯洁无瑕，尽管你们相信圣父亲自用他的手指为你们写下了部分，另一部分则由忠实可靠的\Person{梅瑟}书写），那么圣子的盟约也同样可能受到败坏，也同样可能包含可反对的内容；尤其是因为众所周知，它并非由圣子本人或他的宗徒所写，而是很久以后由一些无名之人所写，这些人唯恐被怀疑写了他们一无所知的事情，便将自己的书冠以宗徒或被认为是宗徒追随者的名字，宣称内容是根据这些原典所写。我认为，这样做是对基督门徒的严重不公，因为他们引用这些著作中不一致和矛盾的说法，声称这些福音书是根据这些门徒所写的，而这些福音书充满了错误和矛盾，无论在事实还是观点上，既不能自洽，也不能彼此协调。这无异于诽谤善人，并将不和的罪名加诸门徒的团体。在阅读福音书时，我们内心的清晰意图能看出这些错误；为了避免不义，我们接受那些有助于建立信德、促进主基督及其全能天主父的光荣的有用部分，而拒绝其余与天主和基督的威严不相称、与我们的信仰不一致的部分。\par

3. 回到我之前说的，你们不接受旧约中的一切。你们不接受属肉体的割礼，尽管这是所写的；\BibleRef{创 17:9}也不在安息日停止一切工作，尽管这是命令；\BibleRef{出 31:13}你们不是像\Person{梅瑟}所推荐的，通过奉献和祭献来平息天主，而是将这些视为与基督徒崇拜完全不合，毫无可取之处而抛弃。在某些情况下，你们却做出区分，接受一部分而拒绝另一部分。例如在逾越节（这也是旧约的年度庆节）中，虽然所写的是你们必须宰杀一只羔羊在晚上吃，并且七天之内不可吃发酵的饼，只能吃无酵饼和苦菜，\EditorialNote{出谷纪第十二章}你们接受这个庆节，却不在意庆祝的规则。同样的还有五旬节（或七个星期），以及\Person{梅瑟}命令的某种数量和种类的祭献；\EditorialNote{肋未纪第二十三章}你们遵守庆节，却谴责作为其部分的赎罪礼，因为它们与基督教不合。至于禁止吃外邦人食物的命令，你们是热心相信祭过偶像的物和自死的物是不洁的；但你们却不那么相信禁止吃猪肉、野兔、獾、鲻鱼、乌贼以及所有你们爱吃但\Person{梅瑟}宣布为不洁的鱼。\par

4. 我不认为你们会同意甚至听信这样的事：公公与儿媳同寝，如\Person{犹大}所为；父亲与女儿同寝，如\Person{罗特}所为；先知与妓女同寝，如\Person{欧瑟亚}所为；丈夫为了一夜将妻子卖给情人，如\Person{亚巴郎}所为；一个人娶两姐妹，如\Person{雅各伯}所为；或者你们认为是统治者中最受感动的人豢养成百上千的情妇；或者根据\Person{梅瑟}在\BibleBook{}{申命纪}中关于妻子的规定，如果一个兄弟没有孩子就死了，他的妻子应该嫁给活着的兄弟，他应代替兄弟为她生子；如果那人拒绝这样做，原告应到长老面前，传唤兄弟并劝他履行这种宗教义务；如果他坚持拒绝，则不能不受惩罚，而是那女人要脱下他右脚上的鞋，打他的脸，将他吐上唾沫、咒诅着赶走，使羞辱永存于他的家中。\BibleRef{申 25:5}这些以及类似这些都是旧约的榜样和诫命。如果它们是好的，你们为什么不实行？如果它们是坏的，你们为什么不谴责包含它们的旧约？但如果你们认为这些是伪造的插入，那正是我们对新约的看法。你们无权要求我们承认新约，而你们自己却不承认旧约。\par

5. 既然你坚持旧约与新约均为天主所著，那么，既然你不服从其诫命，更合理、更体面的做法，应当是承认它被不当的添加所败坏，而不是在它真实且未受败坏的情况下如此轻蔑地对待它。因此，我对你忽视旧约要求的一贯解释，至今仍然是：要么你明智到足以将其斥为伪作，要么你胆大妄为、不敬虔到在它们真实的情况下无视它们。无论如何，当你因我接受新约本身而迫使我相信新约文献所包含的一切时，你应当考虑，尽管你声称接受旧约，但你在内心却不信其中的许多事。例如，你不承认旧约的声明为真实或具权威，即\Quote{凡挂在树上的，都是被咒诅的}（\BibleRef{申 21:23}），因为这将适用于耶稣；或者\Quote{凡不在以色列中为弟兄生子立嗣的，都是被咒诅的}（\BibleRef{申 25:5}），因为这将包括所有献身于天主的男女；或者\Quote{凡未受割损的男子，都应从民间铲除}（\BibleRef{创 17:14}），因为这将适用于所有基督徒；或者\Quote{凡违犯安息日的，必须用石头砸死}（\BibleRef{户 15:35}）；或者\Quote{对违犯旧约单一诫命的人，不可施以怜悯}。如果你真的相信这些是神所命令的，那么在基督的时代，你必会首先攻击祂；而现在你也不会与犹太人争吵，因为他们全心全力迫害基督，正是服从了他们自己的神。\par

6. 我知道，你并没有大胆地宣称这些段落是伪作，而是声称这些事在耶稣来临之前是对犹太人的要求；而现在，按照你所说旧约的预言，祂已经来临，祂亲自教导我们应当接受什么，以及应当弃绝什么为过时。至于先知们是否预言了耶稣的来临，我们将在后面查看。同时，我只需说：如果耶稣在旧约中被预言，现在却对其进行如此彻底的批判，并教导我们接受少数事物而抛弃许多事物，那么，新约中所应许的护慰者也同样教导我们接受它的一部分，拒绝另一部分；正如耶稣自己在福音中许诺护慰者时所说：\Quote{祂要引导你们进入一切真理，并要教训你们一切，使你们想起一切。}（\BibleRef{若 16:13}）因此，藉着护慰者的帮助，我们可以对新约采取同样的自由，正如耶稣使你们能对旧约采取的自由一样——除非你认为圣子的新约比圣父的新约更有价值（如果它真是圣父的），以至于新约的许多部分可以被定罪，而旧约却免于一切非难；而且，正如我先前所说，我们知道新约并非由基督或祂的宗徒所写。\par

7. 因此，既然你在旧约中只接受预言和上文所引的普遍实践道德诫命，却废弃了割损、祭献、安息日及其遵守、无酵节，那么，我们为什么不能在新约中只接受那些我们发现在尊崇和赞美圣子威严的话（无论是祂自己说的，还是祂的宗徒说的，且附加条件是：就宗徒而言，这些话是在他们达到成全之后、不再处于不信时说的），而对其他部分——如果是当时说的，便是出自无知或缺乏经验；如果不是，便是由诡诈的对手怀着恶意添加的，或是写作者未经充分考虑而陈述的，并因此被当作真确传承下来的——置之不理呢？试举例如下：耶稣出自妇人的羞耻出生，祂像犹太人一样受割损，像外邦人一样献祭，以屈辱的方式受洗，在旷野中被魔鬼带领，并受到祂最痛苦的诱惑。除了这些例外，再加上任何假借引用旧约之名而插入的内容，我们相信全部，特别是神秘的被钉十字架——象征灵魂在苦难中的创伤；以及耶稣纯正的道德训诲、比喻，和祂全部不朽的教诲——尤其阐述了两个性体的区分，因此无疑必是祂的。因此，你没有理由认为我必须相信福音书的所有内容；因为你自己，正如已被证明的，从旧约中啜饮得如此挑剔，以至于几乎连嘴唇都没有沾湿。\par

8. \Emph{奥斯定}答：我们给予全部旧约圣经应有的赞美，视之为真实而神圣；你却攻击新约圣经，声称它被篡改和败坏。旧约中那些我们不遵守的事，我们视为适合那个时代和那时期子民的适宜规定，并且对我们而言，它们是真理的预象，虽外在的遵守已被废除，但仍具有属灵用途；这一观点被证明是宗徒著作的教导。另一方面，你指责新约中一切你不接受的内容，并声称这些段落不是基督或祂的宗徒所说或所写的。在这些方面，我们之间存在明显的差异。因此，当有人问你为何不接受新约的所有内容，反而在同样书中认可一些事、拒绝许多事为虚假和伪造的插入时，你不可假装效法我们怀着敬畏和信德所做的区分，而是必须为你的冒昧做出交代。\par

9. 如果有人问我们，为什么我们不照旧约的希伯来祖先那样朝拜天主，我们的答复是：新约的祖先们教导我们另一种方式，但这与旧约并无抵触，相反，正如旧约本身所预言的。因为先知如此预言说：\BibleQuote{看，时日将到——上主说——我必要与以色列家和犹大家都订立新约，不像是按照我与他们的祖先所订立的盟约，那时我握住他们的手，领他们出离埃及地。}\BibleRef{耶 31:31} 因此，这预言说明那盟约不会持久，而是将有一个新约。至于有人反对说，我们不属于以色列家或犹大家，我们按照宗徒的教导回答，他称基督为亚巴郎的后裔，并对我们这些属于基督身体的人说：\BibleQuote{因此，你们是亚巴郎的后裔。}\BibleRef{迦 3:29} 此外，若有人问我们，为何我们视那盟约为权威，却不遵守它的规诫，我们在宗徒著作中也找到答复；因为宗徒说：\BibleQuote{不要让任何人在饮食上，或在节期、月朔、安息日上，判断你们，这些原是将来之事的影子。}\BibleRef{哥 2:16} 从这里我们学到：我们应当阅读这些规诫，承认它们是神圣的制度，以保存预言之记忆，因为它们是将来之事的影子；同时，我们无需在意那些因我们不继续外在遵守而判断我们的人；正如宗徒在其他地方为同一目的说：\BibleQuote{这些事发生在他们身上，是为了作鉴戒，并且写下来，是为警戒我们这些生活在末世的人。}\BibleRef{格前 10:11} 因此，当我们在旧约书卷中读到任何我们不需要在新约中遵守、甚至被禁止的事时，我们不应指责它，而应询问其含义；因为这种遵守的中断本身表明它没有被谴责，而是被应验了。关于这一点，我们已经多次说过。\par

10. 以\Person{福斯图斯}无知地用来指控旧约的这条要求为例：一个人应娶他兄弟的妻子，为兄弟立嗣，使其冠以兄弟之名。这不正是预表每位福音宣讲者在教会中应如此劳苦，为他已故的兄弟——即为我们死去的基督——立嗣，并使这嗣子冠以祂的名吗？此外，宗徒现在不是以预表形式遵守这条要求，而是以属灵实况：他责备那些他说用福音在耶稣基督里所生的人，\BibleRef{格前 4:15} 并指出他们想属于保禄的错误。他说：\BibleQuote{保禄为你们被钉死了吗？或者你们受洗是归于保禄的名下吗？}\BibleRef{格前 1:13} 好像他在说：我为我的已故兄弟生了你们；你们的名字是基督徒，不是保禄派。再者，任何被教会拣选却拒绝福音职务的人，理应遭到教会的藐视。所以我们看见，吐唾沫在脸上伴随着一种羞辱的记号：放下一只脚的鞋，以将那人排除在宗徒所说\BibleQuote{用和平的福音当作预备，穿在脚上}\BibleRef{弗 6:15} 的那些人之外，以及先知所说的：\BibleQuote{那传报福音、宣传和平、传报佳美的，他们的脚步是多么美丽啊！}\BibleRef{依 52:7} 那持守福音信德，既使自己受益，又随时准备应召服务教会的人，被恰当地描绘为双脚穿鞋。但那些认为只要相信就足以保全自己，而逃避造福他人之责的人，则蒙受赤足的羞辱——不是在预表意义上，而是在实际上。\par

11. \Person{福斯图斯}无端反对我们遵守逾越节，嘲笑我们与犹太人的遵守方式不同：因为在福音中，我们拥有真正的羔羊，不再是影子，而是实体；我们不是预表死亡，而是每日纪念它，尤其在年度庆节中。因此，我们的逾越节日期也与犹太人的遵守不同，因为我们加入了主日，即基督复活的那一天。至于无酵节，所有信仰纯正的基督徒都遵守它，不是用旧生活的酵母，即邪恶，而是用信德的真诚和真实；\BibleRef{格前 5:8} 不是为了七天，而是常常，正如数字七所预表的，因为日子总是以七计数。如果这种遵守在今世有些困难，因为通往生命的道路是狭窄的，\BibleRef{玛 7:13} 将来的赏报却是确定的；这种困难预表在苦菜中，它们有点苦味。\par

12. 我们也遵守五旬节，即从主受难与复活起的第五十天，因为在这一天祂将所应许的圣神护慰者赐给了我们；正如犹太人的逾越节所预表的：在羔羊被宰杀后的第五十天，梅瑟在西乃山用天主的手指接受律法。\EditorialNote{出谷纪第十九至三十一章} 如果你们读福音，会发现圣神在那里被称为天主的手指。\BibleRef{路 11:20} 发生在某些日子上的重要事件每年在教会中予以纪念，以便此节期的再次来临保存对这些如此重要而有益之事的记忆。因此，若你们问我们为何遵守逾越节，那是因为基督那时为我们被祭献。若你们问我们为何不保留犹太人的仪式，那是因为它们预表了未来的事实，而我们将这些事实作为过去的事来纪念；未来与过去之间的区别体现在我们所用的不同语词上。关于这一点，我们已经说得够多了。\par

13. 再者，若你问，为何在先前预表性的律法所禁止的各种食物中，我们只禁戒祭偶像的祭品和自死的动物，你且听——只要你愿意选择真理而非无谓的毁谤。基督徒不宜吃祭偶像之物的理由，宗徒已给出：他说：\BibleQuote{我不愿意你们与魔鬼相交。}并非他指责作为祭献本身，因为先祖们以此预表基督用来救赎我们的祭献之血。他先说：\BibleQuote{外邦人所祭献的，是祭献给魔鬼，而不是献给天主；}然后加上这话：\BibleQuote{我不愿意你们与魔鬼相交。}\BibleRef{格前 10:20} 若祭肉的本质不洁，则即使无知而吃也必受玷污。但明知而禁戒的原因不在于食物的本性，而是为了良心的缘故，免得显得与魔鬼相交。至于自死的动物，我以为禁食这类肉的理由是：自死动物的肉已腐败，不利于健康，而健康是食物的首要考虑。古代在洪水后向诺厄本人所吩咐的倒血之律\BibleRef{创 9:6}，其意义我们已解释过；许多人认为《宗徒大事录》中所指的正是这点，即我们读到外邦人应禁戒奸淫、祭偶像之物和血\BibleRef{宗 15:29}——也就是未放血的肉。其他人对这经文有不同解释，认为\Quote{禁戒血}是指不沾染杀人的罪行。要解决这问题会耗费太多时间，且非必要。因为即使承认宗徒当时要求基督徒禁戒动物的血、不吃勒死的物，在我看来他们是为时宜而选择了一项容易遵守、对任何人都非重担的诫命，为使外邦人与以色列人因那使两者在祂内合一的角石而共有此规；\BibleRef{弗 2:11} 同时，他们也被提醒：全教会在诺厄方舟中如何被预表——那时天主颁布了这条命令；这预表在宗徒时代因外邦人加入信仰而开始应验。然而，自那割礼与未割礼两道墙虽已在角石中合一却仍保留某些区别特征的时期结束以来，如今教会已全然变为外邦人，以至于其中找不到外表上的以色列人；因此，没有基督徒感到有义务禁戒画眉或小鸟因其血未放出，或禁戒兔子因其被击颈而死未流血。若有谁仍怕碰这些东西，必被其余人嘲笑：普遍深信这真理：\BibleQuote{不是入口的污秽人，而是出口的。}\BibleRef{玛 15:11} 即邪恶在于犯罪，而不在于任何日常食物的本性。\par

14. 至于古人的事迹——那些在愚昧无知者看来似乎有罪而实则不然的事，以及那些真正有罪的事——我们已解释过为何它们被记载下来，以及这如何增进了而非损害圣经的尊严。同样，关于\BibleQuote{凡挂在树上的，都是被咒骂的}以及\BibleQuote{在以色列中不兴起后裔}的咒诅，我们在回应福斯特的反对时于适当地方已经回答。至于所有无论何种反对——无论是我们已分别回答过的，还是包含在我们正考虑着的福斯特言论中的——我们诉诸我们既定的原则，据此我们维护神圣圣经的权威。这原则是：旧约各书所写的一切，都应被接纳并赞美，视为极其真实且对永生极其有益；那些不再外在遵守的诫命，应理解为在当时最为适宜，并视为未来事物的预像，其应验如今我们可看见。因此，凡在那些时代忽略遵守这些象征性诫命的人，按神圣法令被正义地判处应受刑罚，正如现今若有谁不敬地亵渎新约的圣事——这些圣事与旧礼规的差异仅在此时与彼时不同——也要受罚一样。因为正如旧约圣事的义人宁死不从，当受称赞；新约的殉道者也当同样受称赞。又如病人不应挑剔医疗方法，因为今天开一种药、明天开另一种，起初要求的后来禁止，因为疗法正赖于此；同样，从亚当直到世界终结、朽坏的身体沉重压着灵魂\BibleRef{智 9:15}的人类，患病而伤痛，不应挑剔天主的药方，即便有时同样的规条被命令，有时旧规条被另一种替换；尤其曾有过律法将改变的应许。\par

15. 因此，\Person{福斯特}所提出的类比是没有力量的：他类比基督在旧约中指出我们该相信什么和拒绝什么（旧约中预告了基督自己），与护慰者在你们中间同样在新约中行事（新约中类似地预告了他）。假如在旧约中有任何东西我们斥为错误、或非天主权威、或不真实，那么这种类比或许有一定道理。但我们并无此类行为：我们接受一切，无论是我们遵守为行为准则的，还是我们不再遵守但仍承认曾是先知性规条的——它们曾受命，如今已实现。此外，护慰者的应许见于那些书卷中（你们并不接受其全部内容），他的派遣则记录在你们甚至不愿提名的书中。因为，如上所述且反复说过，在\WorkTitle{宗徒大事录}中明确叙述了圣神于五旬节那天降临，发生的效果表明他是谁。因为所有初领圣神的人都说了各种方言；\EditorialNote{宗徒大事录 第二章}在这标记中有一个应许：在万民中，后来的教会将忠实地宣讲关于圣神以及圣父和圣子的道理。\par

16. 那么，你们为什么不接受新约中的一切？是因为这些书卷没有基督宗徒的权威，还是因为宗徒教导了错误的东西？你回答说，这些书卷没有宗徒的权威。宗徒在他们的教导中犯了错误——这是外教人说的话。但是，你们能拿出什么证据来证明这些书卷的流传不能追溯到宗徒？你们回答说，在许多事情上，它们自相矛盾，并且彼此矛盾。再没有比这更不真实的了；事实上，是你们不理解。在每一个情形中，当\Person{福斯特}提出你们认为矛盾的地方时，我们已经指出并不存在矛盾；而且我们会在每一个其他情形中同样做。读者或学生胆敢将过错归咎于具有如此权威的圣经，而不承认自己的愚笨，这是不能容忍的。护慰者曾教导你们，这些作品不是宗徒所著，而是别人冒用他们的名字写的吗？但是，你们从哪里证明你们从护慰者学到了这一点？如果你们说护慰者是由基督所应许并派遣的，我们回答说，你们的护慰者既非基督所应许，也非所派遣；并且我们也向你们指出他何时派遣了他所应许的护慰者。你们有什么证据证明基督派遣了你们的护慰者？你们从哪里得到证据支持你们的报告者，或者更确切地说，误导者？你们回答说，你们在福音中找到证据。在哪部福音中？你们并不接受所有的福音，你们说福音被篡改过。你们要先控告你们的见证人腐败，然后要求他的证词？你们愿意时相信他，不愿意时又不相信他，这等于只相信你们自己。如果我们准备相信你们，就根本不需要见证人。此外，在关于圣神作为护慰者的应许中，经上说：\BibleQuote{他要引导你们进入一切真理。}\BibleRef{若 16:13}但是，一个教导你们基督是欺骗者的人，怎么能引导你们进入一切真理呢？再者，如果你们证明福音中关于护慰者应许的所有言论，除了\Person{摩尼}之外不能适用于任何人，正如先知们的预言适用于基督一样；并且如果你们从那些你们说真实的抄本中引用段落，我们也许可以说，正是在这一点上，为了证明\Person{摩尼}是唯一所指的人，这些段落已经被你们教派篡改了。你们对此唯一的回答是，你们不可能篡改已经为所有基督徒所拥有的文件；因为在这种尝试一开始，就会因比较更古老的抄本而被揭露。但是，如果这证明你们不可能败坏这些书卷，那也证明任何人都不可能败坏它们。第一个胆敢如此行的人，会被更古老抄本的比较而判明有罪；尤其因为圣经不仅以一种语言存在，而是以多种语言存在。事实上，虚假的读本有时可以通过比较更古老的抄本或原文来纠正。因此，你们要么承认这些文件是真实的，那么你们的异端就无法站立片刻；要么它们是伪造的，那么你们就不能使用它们的权威来支持你们关于护慰者的教义，这样你们就反驳了自己。\par

17. 再者，在护慰者的应许中所说的话，表明它绝不可能指多年以后才来的摩尼。因为若望清楚地说，圣神要在主复活升天后立即降临：\BibleQuote{圣神还没有赐下，因为耶稣还没有受到光荣。}\BibleRef{若 7:39} 现在，如果圣神没有赐下的原因是耶稣还没有受到光荣，那么祂必然在耶稣受光荣之后立即赐下。同样，卡塔弗里吉亚派声称他们领受了所应许的护慰者；因而他们背离了公教信仰，禁止保禄所允许的事，并谴责他视为合法的第二次婚姻。他们曲解了关于圣神的话：\BibleQuote{祂要引导你们进入一切真理}，仿佛保禄和其他宗徒没有教导全部的真理，而是为卡塔弗里吉亚派的护慰者留下了余地。他们也强解保禄的话：\BibleQuote{我们现在所知道的，只是局部的；我们所预言的，也只是局部的；但当那圆满的一来到，那局部的就要被废除。}\BibleRef{格前 13:9} 他们声称宗徒知道和预言是局部的，当他说：\BibleQuote{他愿怎样做，可以随意；他若结婚，并不犯罪}\BibleRef{格前 7:36}，而这话被弗里吉亚护慰者的圆满所废除。如果他们被告知，他们被教会的权威所谴责——这教会是古代应许的对象，并遍布全世界——他们就会回答说，这正应验了关于护慰者的话：\BibleQuote{世界不能接受祂}。\BibleRef{若 14:17} 那么，那些经文：\BibleQuote{祂要引导你们进入一切真理}、\BibleQuote{当那圆满的一来到，那局部的就要被废除}、\BibleQuote{世界不能接受祂}，岂不正是你们发现其中预言了摩尼的经文吗？因此，任何以护慰者之名出现的异端都胆敢对这些经文作出同样看似合理的应用。因为没有一种异端不称自己是真理；它越是骄傲，就越可能称自己为圆满的真理；因而它声称要引导进入一切真理；并且既然由它而来的那圆满者已经来到，它便试图废除宗徒的教导，而它自己的错误是与这教导对立的。至于教会，因着宗徒的恳切告诫：\BibleQuote{谁若给你们宣讲福音，与你们所接受的不同，当受诅咒}，\BibleRef{迦 1:9} 当异端的宣讲者开始被全世界宣布为受诅咒时，他岂不立即喊道：这正是经上所写的\BibleQuote{世界不能接受祂}吗？\par

18. 那么，你们在哪里能找到所需的证据，证明你们是从护慰者那里得知福音不是宗徒写的呢？相反，我们有证据表明圣神、护慰者，在耶稣受光荣之后立即降临。因为\BibleQuote{祂还没有赐下，因为耶稣还没有受到光荣}。我们也有证据表明祂引导进入一切真理，因为通往真理的唯一道路是爱，而宗徒说：\BibleQuote{天主的爱，借着所赐给我们的圣神，已倾注在我们心中了。}\BibleRef{罗 5:5} 我们也指出，在\BibleQuote{当那圆满的一来到}这句话中，保禄说的是在享受永生中的圆满。因为他在同一处说：\BibleQuote{我们现在是借着镜子观看，模糊不清，到那时，就要面对面了。}\BibleRef{格前 13:12} 你们不能合理地坚持说我们在此世面对面看见天主。因此那圆满的并没有临到你们。由此可见宗徒在这主题上的想法。这圆满要等到若望所说的成就时才临到圣徒：\BibleQuote{现在我们是天主的子女，将来如何，还没有显明；但我们知道，当祂显现时，我们必要相似祂，因为我们要看见祂实在怎样。}\BibleRef{若一 3:2} 那时，我们将借着圣神被引导进入一切真理，我们现在已领受了这真理的抵押。再者，\BibleQuote{世界不能接受祂}这句话清楚地指向那些在圣经中通常被称为世界的人——世界的爱好者、恶人、或属血气的人；宗徒论到他们说：\BibleQuote{属血气的人不能领受天主圣神的事。}\BibleRef{格前 2:14} 那些被称为属这世界的人，他们除了物质的事物——这世界中的感官对象——之外，什么也不能了解；正如你们的情况，当你们钦羡太阳和月亮时，就认为一切神圣之事都与它们相似。你们这些欺骗人而又被人欺骗的人，称这个愚蠢理论的作者为护慰者。但既然你们没有证据证明他是护慰者，你们就没有可靠的基础来断言福音著作——你们只接受一部分——不是宗徒所写的。因此，你们仅剩的论据是：这些著作包含有损基督荣耀的事，例如：祂由童贞女所生、祂受割损、为祂献上惯常的祭品、祂受洗、祂受魔鬼的试探。\par

19. 除上述例外，包括旧约的见证在内，你（依\Person{福斯图斯}的话）声称接受其余一切，尤其是神秘钉十字架，象征灵魂在受难中的创伤；以及耶稣的道德训诫，和他整个不朽的教诲，其中特别阐明了二本性的区分，因此必定是他的。你显然意图剥夺\WorkTitle{圣经}的一切权威，并让每个人的心灵成为判断\WorkTitle{圣经}段落之认可与否的法官。这并非在信仰的事上服从\WorkTitle{圣经}，而是使\WorkTitle{圣经}服从于你。每个人并非以\WorkTitle{圣经}的崇高权威作为认可的理由，而是以个人的认可作为认为段落正确的理由。因此，如果你抛弃权威，那么，你这可怜的、微弱的灵魂，被肉性云雾所遮蔽，我恳求你，你将投靠什么呢？撇开权威，让我们听听你信仰的理由。是否通过逻辑过程，你关于天主本性的冗长叙述必定以这一惊人宣告告终：这一本性容易受到伤害和败坏？你又如何知道有八个大陆和十重天，以及\Person{阿特拉斯}支撑着世界，世界悬挂在伟大的宇宙持有者之下，以及无数同类事物？谁是你们的权威？当然是\Person{摩尼}，你会说。但是，不幸的人，这不是看见，而是信仰。因此，如果你服从一个晦涩且无理性之权威的命令，接受一堆无穷的虚构，以至于你相信所有这些事，因为它们写在那些被你误判为可靠的书籍中，尽管没有其真理的证据，那么为什么不服从\WorkTitle{福音}的权威？\WorkTitle{福音}如此根基稳固，如此确认，如此普遍承认和钦佩，并且从宗徒时代直到今天拥有不间断的一系列见证，以便你能有一种明智的信仰，并且能认识到你所有的反驳都是愚蠢和乖张的果实；并且，比起相信天主本性容易受到伤害和败坏，而且在遭受污染和囚禁后不能完全释放和净化，反而被至高的神性必然性判决永远惩罚在黑暗的团块中，这一观点更真实：即天主的不变本性应取有死性的一部分，以便在与这有死性结合时不被损伤，而是真实地做和受凡有死性为拯救从其取来的人类所应做和应受的一切？\par

20. 你回答说，你相信\Person{摩尼}没有证明的事，因为他如此清楚地证明了世界上善与恶二本性的存在。但这是你不幸错觉的根源；因为正如在福音书中，在世界上，你对邪恶的观念完全来自对蛇、火、毒药等令人不愉快事物在你感官上的效果；而你所知道的唯一善就是对感官有愉悦效果的事物，如令人愉快的味道、甜美的气味、阳光，以及任何强烈吸引你眼睛、鼻孔、味觉或其他感觉器官的东西。但是，如果你一开始就把\WorkTitle{自然之书}看作万有创造者的作品，并且相信你自己有限的理解可能在似乎有不当之处时出错，而不是敢于挑剔天主的作品，你就不会陷入这些不敬的愚蠢和亵渎的幻想之中，由于你对邪恶的真正本质一无所知，你把所有邪恶都归咎于天主。\par

21. 我们现在可以回答这个问题：我们如何知道这些书是由宗徒们写的。简而言之，我们知道的途径与你（误以为你偏爱的书籍权威是\Person{摩尼}写的）知道的途径相同。因为，假如有人对此提出质疑，并与你辩论，声称你归于\Person{摩尼}的书籍并非他的作品；你唯一的回答只能是嘲笑这种无端质疑由如此广泛连续见证所确认的事实的荒谬。因此，正如确定这些书是\Person{摩尼}的作品，并且对于一个出生于多年之后的人来说，凭空提出自己的反对意见并展开讨论是可笑的；同样，我们可以断言，对于\Person{摩尼}或其追随者来说，对最初确凿可靠、从宗徒时代直到今天通过不断的保管者系列精心传承的著作提出这样的反对意见，是荒谬的，或者更可悲的。\par

22. 现在我们只需比较摩尼与宗徒们的权威。两者的著作真实性同样确定。但除非一个人不再追随派遣宗徒的基督，否则无人会将摩尼与宗徒相提并论。谁若没有误解基督的话，又岂能在其中发现两种彼此对立、各有其本原的性体的道理？再者，宗徒们，作为真理的门徒，宣告基督的诞生与受难是真实的事件；而摩尼，自夸引人进入一切真理，却引领我们走向一位其受难被他宣告为幻觉的基督。宗徒们说，基督在祂从亚巴郎后裔所取的血肉中受了割损；摩尼却说，天主按其本性被黑暗种族肢解。宗徒们说，按当时的制度，为婴孩时期的基督按我们的本性献了祭；摩尼却说，不是人性的一部分，而是天主实体本身的一部分，必须被引入敌对种族的本性中，以祭献给全体恶魔。宗徒们说，基督为给我们树立榜样，在约但河受了洗；摩尼却说，天主将自己浸入黑暗的污秽中，并且永远不能完全出来，而那不能净化的部分将被罚入永罚。宗徒们说，基督按我们的本性受了恶魔首领的试探；摩尼却说，天主的一部分被恶魔种族俘虏。在基督受试探时，祂抵抗了试探者；而在天主被俘时，被俘的部分即使在得胜后也无法恢复其本源。总之，摩尼以改善的面貌宣讲另一个福音，这是魔鬼的道理；而宗徒们，按照基督的道理，命令若有人宣讲另一个福音，当受诅咒。\BibleRef{迦 1:8}\par

\SourceRecord{Translated by Richard Stothert.；Nicene and Post-Nicene Fathers, First Series；Vol. 4.；Edited by Philip Schaff.；Buffalo, NY: Christian Literature Publishing Co.,；1887.；<http://www.newadvent.org/fathers/140632.htm>.}

% structure: p0001 source=h1 target=section reason=page-title

\section{\WorkTitle{驳斥福斯图斯，第三十三卷}}

\Emph{福斯图斯认为与亚巴郎、依撒格和雅各伯同席并非什么大荣耀，因为他厌恶旧约所描述的他们的道德品格。他为自己对圣经的主观批评辩护。奥古斯丁总结论点，声称胜利，并劝勉摩尼教人尽管旧约存在困难，仍要放弃对旧约的反对，并承认天主教会的权威。}\par

1. \Emph{福斯图斯}说：你引用福音中的话：\BibleQuote{有许多人要从东方和西方来，与亚巴郎、依撒格和雅各伯一同坐席于天国}\BibleRef{玛 8:11}，并问我们为何不承认这些祖先。现在，我们决不会嫉妒任何人，只要天主怜悯他，将他从毁灭中带出，获得救恩。然而，在这样的情况下，我们应承认这怜悯行为中彰显的仁慈，而非那人的功德，因为他的生活无疑应受责备。因此，关于犹太人的祖先——亚巴郎、依撒格和雅各伯，如果这段经文是真实的话，基督在这节经文中提到了他们，尽管他们过着邪恶的生活，正如我们从他们的后裔梅瑟或者那部被称为 \WorkTitle{创世纪} 的历史的作者那里得知的，这部历史描述了他们的行为极其骇人、可憎；我们愿意承认，他们毕竟可能在天国里，那个他们既不信仰也不盼望的地方，这一点从他们的书卷中看得很清楚。但是，必须记住，正如你们自己所承认的，如果他们确实达到了这节经文所说的地方，那与他们的恶行应得的深渊地狱是截然不同的，他们的得救是我主基督的工作，是他神秘苦难的结果。谁会嫉妒同一位主赐予那个十字架上的强盗得救，并且基督说就在那一天他将与基督同在圣父的乐园里？\BibleRef{路 23:43} 有谁如此心硬，不认可这仁慈的行为？然而，这并不因为耶稣赦免了一个盗贼，我们就必须认可盗贼的习惯和行为；同样，也不因耶稣赦免了税吏和娼妓的罪，并声明他们将比那些高傲行事的人先进入天国。\BibleRef{玛 21:31} 因为，当祂赦免了那个被犹太人指控有罪、并且正在行淫时被捉住的妇人时，祂告诉她说：\BibleQuote{不要再犯罪了。}\BibleRef{若 8:3} 所以，如果祂对亚巴郎、依撒格和雅各伯也做了类似的事，一切赞美都属于祂；因为这样的行为对于灵魂是合适的，\BibleQuote{祂使太阳照耀善人，也照耀恶人，降雨给义人，也给不义的人。}\BibleRef{玛 5:45} 有一件事在你的教义中令我困惑：为什么你们只局限于犹太人的祖先，而不认为外邦人的族长也分享了赎主的恩宠？特别是因为基督教会由他们的子女构成，多于由亚巴郎、依撒格和雅各伯的后裔构成。你将会说，外邦人敬拜偶像，而犹太人敬拜全能的天主，因此耶稣只关心犹太人。这样看来，敬拜全能的天主是通往地狱的必然之路，而圣子必须来帮助敬拜圣父的人。你愿意怎么说都行。就我而言，我准备与你们一同相信，祖先们进入天堂，并非凭借他们自己的功德，而是凭借那比罪恶更强大的神圣慈悲。\par

2. 然而，判断这节经文是否确实对基督说的也有困难，因为叙述中存在不一致。因为两位福音作者，\Person{玛窦}和\Person{路加}，都同样叙述了那位仆人患病、这些\Person{耶稣}的话被认为是对他说的百夫长，说\Person{耶稣}在\Place{以色列}中没有见过这么大的信德，甚至在这位外邦人和外教人身上也没有，因为他说他不配让\Person{耶稣}进他的屋檐下，只愿\Person{耶稣}说一句话，他的仆人就会痊愈；只有\Person{玛窦}补充说\Person{耶稣}继续说道：\BibleQuote{我实在告诉你们：将有许多人从东方和西方来，与\Person{亚巴郎}、\Person{依撒格}和\Person{雅各伯}一同坐席在天国里；但天国之子要被赶到外边的黑暗里去。}所谓将有许多人来，是指外教人，因为在那百夫长身上，虽然他是外邦人，却找到了如此大的信德；而天国之子是指犹太人，在他们身上找不到信德。\Person{路加}虽然也在他的福音中提到了这件事，作为基督奇迹叙述的一部分，但没有提到\Person{亚巴郎}、\Person{依撒格}和\Person{雅各伯}。如果说他省略了，因为\Person{玛窦}已经说了，那为什么他还要讲百夫长和他的仆人的故事呢，既然\Person{玛窦}巧妙叙述中已经详细记载了？但这段经文是有问题的。因为，在描述百夫长向\Person{耶稣}的请求时，\Person{玛窦}说他亲自来求医治；而\Person{路加}说他没有，而是派了犹太人的长老，他们在\Person{耶稣}可能轻视身为外邦人的百夫长（因为他们认为\Person{耶稣}是一个彻底的犹太人）的情况下，开始说服他，说他是值得\Person{耶稣}为他做此事的人，因为他爱他们的民族，还为他们建造了会堂；这里再次想当然地认为天主之子关心一个外邦百夫长认为为犹太人建造会堂是适宜的。事实上，这些有问题的词也在\Person{路加}福音中出现，也许是因为他反思后认为它们可能是真实的；但它们出现在另一个地方，且上下文完全不同。这段经文是\Person{耶稣}对他的门徒说：\BibleQuote{你们要努力进窄门；因为将来有许多人来想进去，却不能进去。等到家主起来关了门，你们开始站在外面敲门，说：\InnerQuote{主啊，给我们开门。}他要回答你们说：\InnerQuote{我不认识你们。}那时你们要说：\InnerQuote{我们在你面前吃过喝过，你在我们的街市上教导过我们；}但他要对你们说：\InnerQuote{我不知道你们从哪里来；离开我吧，所有作恶的人。}在那里会有哀哭和切齿，当你们看见\Person{亚巴郎}、\Person{依撒格}、\Person{雅各伯}和众先知都在天主的国里，而你们却被赶出去。还有人从东方、西方、北方、南方来，要在天主的国里坐席。}\BibleRef{路 13:24}至于说许多人将被关在天主国之外、他们只背负基督的名却不做基督工作的地方，\Person{玛窦}并没有省略；但他这里没有提到\Person{亚巴郎}、\Person{依撒格}和\Person{雅各伯}。同样地，\Person{路加}提到了百夫长和他的仆人，但在那个上下文中没有提及\Person{亚巴郎}、\Person{依撒格}和\Person{雅各伯}。既然这些话是在何时说的不确定，我们可以自由地怀疑它们是否被说过。\par

3. 我们带着批判的眼光研究存在如此不一致和矛盾的圣经，不是没有理由的。通过这样审视一切，将一段经文与另一段比较，我们确定哪些包含\Person{基督}的实际话语，哪些可能是或不是真实的。因为你们的前辈在吾主的话语中做了许多篡改，这些话语以他的名义出现，却与他的教义不一致。此外，正如我们一再证明的，这些著作不是\Person{基督}或其宗徒的作品，而是在他们离世很久之后，由一些默默无闻的半犹太人编造的谣言和信仰的集合，甚至彼此之间也不和谐，并以宗徒或被认为是宗徒追随者的名义出版，以便给所有这些错误和谎言披上宗徒权威的外衣。但无论你如何看待这一点，关于这节经文，我再说一遍，我不坚持拒绝它。对我的立场来说，足够的是，正如我之前所说，且你不得不承认，在吾主来临之前，\Place{以色列}的所有圣祖和先知都因他们的罪躺在地狱的黑暗中。即使他们可能被\Person{基督}恢复到光明和自由，这与他们可憎生活的特质无关。我们憎恨和避开的不是他们的人，而是他们的特质；不是他们现在被净化后的样子，而是他们过去不洁时的样子。所以，不管你对这节经文怎么看，它都不影响我们：因为如果它是真实的，它只是彰显了\Person{基督}的良善和怜悯；如果是伪造的，那些写它的人应该受到责备。我们的立场一如既往地安全。\par

\Person{奥古斯丁}答道：这真是可怜的安全感！当你因憎恨先祖不洁而自相矛盾时，你又为你那不洁的神而悲伤。你承认，自从救主降临以来，先祖们已恢复纯洁，并享有真福者的安息；而你的神即使在救主降临之后，仍躺在黑暗里，仍沉溺于不义的海洋，仍在一切污秽的泥沼中打滚。因此，这些人不仅生前比你的神更好，而且死时也比你的神更幸福。关于在基督降生成人之前离开此生的义人的居所，以及他们的状况是否因基督的受难——就是他们所相信将要来临、受苦、复活，并借着预言之神的引导以适当言语预言的那一位——而改善，若有可能从圣经中清楚发现，便当从圣经中探求；我们不需要采纳所有人粗浅的看法，更不需要接受那些误入如此严重错误的异端者的意见。\Person{浮士德}在此企图从侧门引入一种观念，即我们在今生之后可以得到某种东西，除了我们今世行为应得的报应之外。你最好在活着的时候放弃你的错误，并接受和持守天主教信仰的真理。否则，当天主开始对不义之人执行祂的威胁时，不义之人的期望将悲惨地落空。\par

关于\Person{浮士德}对先祖生活的诽谤，我已给予我认为充分的答复。他们死时受到惩罚，或是在主的受难后得以成义，这并非我们从主对他们的称赞所学的；主曾告诫犹太人，如果他们是\Person{亚巴郎}的子女，就该做\Person{亚巴郎}的作为，并说\Person{亚巴郎}渴望看见祂的日子，看见了就欢喜快乐（\BibleRef{若 8:39}）；又说他将那个被骄傲的富人藐视的受苦穷人，由天使带到他的怀抱中，即某个深藏之福乐安息的角落（\BibleRef{路 16:23}）。至于宗徒\Person{保禄}，又当如何理解？当\Person{保禄}称赞\Person{亚巴郎}时，他是否认为死后可以成义？他说\Person{亚巴郎}在未受割损之前就相信了天主，这就算为他的正义（\BibleRef{罗 4:3}）。\Person{保禄}如此重视这一点，以至于他指定我们成为\Person{亚巴郎}的子女的唯一根据——即使我们并非从他肉身所生——乃是跟随他信德的足迹。\par

6. 你们在错误中如此顽固，不顾圣经的见证，以致无法从你们身上得到任何成果；因为每当引用任何东西反对你们时，你们竟敢说那不是宗徒写的，而是某个假托他名的人写的。你们所宣讲的魔鬼的道理与基督信仰的道理如此对立，以致你们若否认宗徒著作的真实性，就无法作为自称为基督徒的人继续坚持它。你们怎能这样损害自己的灵魂呢？如果不在福音和宗徒著作中，你们还能从哪里找到任何权威呢？如果我们怀疑那些由宗徒们自己建立的、在各地如此显赫的教会所归于宗徒名下的著作的宗徒来源，同时又承认那些由异端者提出的、反对教会的著作无疑是宗徒的作品（而这些异端者的作者，即他们得名的创始人，生活在宗徒之后很久），那么，我们如何能确定任何书的作者呢？难道我们没有看到，在世俗文学中，有许多著名作者的作品在他们去世后被假托出版，这些作品因其与确定无疑的作品不一致，或因其在作者生前不为人知，无法得到作者或其朋友确认的证明而被拒绝吗？举一个例子：最近不是有一些书以著名的医生\Person{希波克拉底}的名字出版，但未被医生们承认为权威吗？尽管在风格和内容上有一定相似，这个决定仍然未改变；因为与\Person{希波克拉底}的真本相比，这些书被发现较逊色；此外，在确定他的真本作者身份时，这些书并未被确认是他的。再次，那些用来与可疑作品比较并据此拒绝可疑作品的书，可以确凿地归于\Person{希波克拉底}；任何否认这些书作者的人只会被嘲笑，这仅仅是因为从\Person{希波克拉底}时代至今有一连串的见证，使得现在或将来对此有任何怀疑都是不合理的。我们如何知道\Person{柏拉图}、\Person{亚里士多德}、\Person{西塞罗}、\Person{瓦罗}及其他类似作家的作品作者，难道不是通过连续不断的证据吗？同样，许多对教会书籍的注释，虽无正典权威，却显露出有益的渴望和探究精神。在每种情况下，作者身份如何确定，难道不是通过作者尽可能在其生前将作品公之于众，并通过信息的连续传递，随着信仰变得更加普遍而变得更加确定，直至我们今日；以致当我们被问及某书的作者时，我们可以毫不困难地回答吗？但为何谈论古书呢？就拿我们眼前的书来说：如果若干年后，有人否认这本书是我写的，或者否认\Person{浮士德}的那本书是他写的，那么除了某些当代人掌握并一代代传递到遥远未来的信息之外，还能从哪里找到事实证据呢？由此看来，凡没有屈服于说谎魔鬼的恶意和欺骗性暗示的人，都不会被激情蒙蔽到否认宗徒教会——一个既众多又忠信的弟兄团体——能够将他们的著作原封不动地传给后代，正如宗徒们最初的驻地一直由连续的主教继承至今，尤其是我们习惯看到这种情况发生在教会内外的普通著作中。\par

7. 然而福斯托斯在福音中发现了矛盾。不如说，福斯托斯以错误的精神读福音，他愚昧到不能理解，盲目到不能看见。倘若你被虔诚所推动，而非被派系精神所误导，你便很容易通过查考这些段落，发现作者们之间一种奇妙而最具教诲的和谐。谁在阅读同一事件的两篇叙述时，会因为一篇省略了另一篇所提及的内容，或一篇虽简洁但却与另一篇详述的内容实质一致而仅指出所作之事以及如何作为，就指控其中一位或两位作者有错误或谎言？这正是福斯托斯在试图指控福音真理时所为，仿佛路加省略了玛窦记载的基督某句话，就意味着路加否认玛窦的陈述。此事中并无真正的困难；制造困难表明缺乏思考或思考能力。确实，在百夫长的叙述中有一点在信徒之间被讨论，而学问不大的不信者也提出反对，他们在受教后仍不放弃错误，证明他们的好争论。这一点是：玛窦说百夫长来到耶稣跟前\Quote{恳求祂，说}；而路加说他托犹太人的长老去见耶稣，同样请求祂医治他那患病的仆人；当耶稣快到房子时，他又派了别人去，通过他们说他不配让耶稣进他的房子，也不配亲自到耶稣面前。那么，我们如何读到玛窦的话：\Quote{他来到祂跟前，恳求祂，说：\InnerQuote{我的仆人瘫痪在家，痛苦不堪？}}解释是：玛窦的叙述正确但简短，提到百夫长来到耶稣跟前，却没有说他亲自来还是由别人代劳，也未说明关于他仆人的话是他自己说的还是通过别人说的。但人们不是常说某人接近某物，即使并未到达吗？甚至\Quote{到达}这一最强烈的表达形式，也常用在某人通过他人行动的情况，比如我们说他把案件带到法庭，他到达了法官面前；或者他到达了某个有权势者面前，虽然可能是通过朋友，而且此人可能并未见到他所称的\Quote{到达其面前}的那个人。从\Quote{到达}一词，我们给那些通过野心手段，亲自或通过朋友接近那些几乎不可接近的权贵之心的人起名为\Quote{钻营者}（\OriginalTerm{钻营者}{Perventors}）。那么我们在阅读时就要忘记通常的言语习惯吗？还是圣经必须有它自己的语言？好争辩的批评者的挑剔就这样被引用通常的言语形式所驳斥。\par

8. 那些不是以好辩而是以平静相信的精神审视此事的人，受邀来到耶稣面前，不是外表地而是内心地，不是身体临在而是藉信德的能力，正如百夫长所做的那样；这样他们就会更好地理解玛窦的叙述。对他们，圣咏中说：\Quote{你们来接近祂，并获得光明；你们的脸面不会羞愧。}由此我们得知，这位百夫长——他的信德受到高度赞扬——比那些传达他口信的人更真实地来到了基督面前。我们在患血漏的妇人身上找到一个类似案例，她因触摸基督衣边而被治愈，当时基督说：\Quote{有人触摸了我。}门徒们对基督说\Quote{有人触摸了我}感到惊奇，因为人群拥挤着他。事实上，他们这样回答：\Quote{人群拥挤着你，你却说：\InnerQuote{有人触摸了我？}}\BibleRef{路 8:43} 现在，正如人群拥挤基督而妇人触摸了他一样，使者被派到基督那里，但百夫长真正来到了他面前。在玛窦那里，我们遇到一种并非罕见的表达方式，同时也具有象征意义；而在路加那里，则是对整个事件的简单叙述，这样引起我们注意玛窦记录的方式。我希望那些以如此琐碎的困难为基础对福音提出愚蠢反对的人，能够自己讲述一个故事两次，诚实地真实叙述所发生之事，然后他的话语被写下并读给他听。那么我们将看到，他是否不会在一次讲述中比另一次多说或少说；是否不仅词语顺序而且事物顺序都会改变；他是否不会把自己的一些意见放入别人口中，因为虽然从未听他说过，却完全知道他的想法；他是否不会有时用几句话总结他之前详细叙述的内容。在这些以及其他或许可以归纳为规律的方式中，两人对同一事物的叙述，或同一人对同一事物的两次叙述，可能在许多方面有差异而不对立，不同而不矛盾。这样，可怜的摩尼教徒（\OriginalTerm{摩尼教徒}{Manichaeans}）用以窒息自己以保持错误精神并拒斥一切可能引导他们走向救恩之事的全部束缚就被解开了。\par

9. 如今，\Person{福斯特}的所有诽谤均已被驳斥，至少那些在本书中详细论述并依主所赐予我的能力充分解释的内容，我已尽数处理。我要对你们这些卷入那些惊人而该受罚的谬误中的人提出一句忠告：如果你们承认圣经的最高权威，那么你们就应当承认那自从基督亲自建立以来，通过祂宗徒的服务，以及通过宗徒座位上的主教们的合法继承，至今仍保留在普世、并为众人所知的权威。在那里，\WorkTitle{旧约}的疑难也得到了解决，其预言也得以应验。如果你们要求证明，那么首先考虑你们自己是什么，你们是多么不适宜理解自己灵魂的本性，更不用说理解天主了；我指的是理性的理解，正是你们自称渴望或曾经渴望的那种理解，而不是轻信幻想的观念。你们既然承认这一无能（只要你们保持现状，这种无能就会持续），那么至少可以诉诸每个人理性的自然确信——只要未被谬误败坏——即天主的性体和本质是完美不变与不朽的。承认这个，或相信这个，你们就不再是摩尼教徒，这样假以时日，你们就可以成为天主教徒。\par

\SourceRecord{Translated by Richard Stothert.；Nicene and Post-Nicene Fathers, First Series；Vol. 4.；Edited by Philip Schaff.；Buffalo, NY: Christian Literature Publishing Co.,；1887.；<http://www.newadvent.org/fathers/140633.htm>.}
